% Generated mechanically from the frozen descent.tex target.
% Do not edit this generated file; edit the bound locale source instead.

% OK, start here.
%

\setcounter{chapter}{34}
\chapter{下降}


\phantomsection
\label{descent-section-phantom}


\section{引言}
\label{descent-section-introduction}

\noindent
在概形上的拓扑一章中（见 Topologies, Section
\ref{topologies-section-introduction}），我们引入了概形的 Zariski、
\'etale、fppf、光滑、syntomic 和 fpqc 覆盖。本章讨论何种概形上结构
可以沿这类覆盖下降。例如可参阅 \cite{Gr-I}、\cite{Gr-II}、
\cite{Gr-III}、\cite{Gr-IV}、\cite{Gr-V} 和 \cite{Gr-VI}。
本章还旨在拟凝聚层、概形等对象的下降问题这一较少形式化的背景中，
引入下降、下降数据和有效下降数据等概念。形式化的概念，即位点上的叠，
将在叠一章中讨论（见 Stacks, Section
\ref{stacks-section-introduction}）。

\section{拟凝聚层的下降数据}
\label{descent-section-equivalence}

\noindent
本章采用如下约定：投影态射
$\text{pr}_i : X \times \ldots \times X \to X$ 的编号从 $i = 0$ 开始。
因此有 $\text{pr}_0, \text{pr}_1 : X \times X  \to X$、
$\text{pr}_0, \text{pr}_1, \text{pr}_2 : X \times X \times X  \to X$
等等。


\begin{definition}
\label{descent-definition-descent-datum-quasi-coherent}
设 $S$ 为概形，并设 $\{f_i : S_i \to S\}_{i \in I}$ 是以 $S$ 为靶的
态射族。
\begin{enumerate}
\item 相对于给定族的一个{\it 拟凝聚层下降数据
$(\mathcal{F}_i, \varphi_{ij})$}由如下资料组成：对每个 $i \in I$，
$S_i$ 上的一个拟凝聚层 $\mathcal{F}_i$；以及对每一对
$(i, j) \in I^2$，拟凝聚 $\mathcal{O}_{S_i \times_S S_j}$-模
$\varphi_{ij} : \text{pr}_0^*\mathcal{F}_i \to \text{pr}_1^*\mathcal{F}_j$
的一个同构，并要求对每个指标三元组 $(i, j, k) \in I^3$，图
$$
\xymatrix{
\text{pr}_0^*\mathcal{F}_i \ar[rd]_{\text{pr}_{01}^*\varphi_{ij}}
\ar[rr]_{\text{pr}_{02}^*\varphi_{ik}} & &
\text{pr}_2^*\mathcal{F}_k \\
& \text{pr}_1^*\mathcal{F}_j \ar[ru]_{\text{pr}_{12}^*\varphi_{jk}} &
}
$$
作为 $\mathcal{O}_{S_i \times_S S_j \times_S S_k}$-模的图可交换。
这称为{\it 余圈条件}。
\item 一个{\it 下降数据态射
$\psi : (\mathcal{F}_i, \varphi_{ij}) \to
(\mathcal{F}'_i, \varphi'_{ij})$}由一族 $\mathcal{O}_{S_i}$-模态射
$\psi_i : \mathcal{F}_i \to \mathcal{F}'_i$，即
$\psi = (\psi_i)_{i\in I}$，给出，并要求所有图
$$
\xymatrix{
\text{pr}_0^*\mathcal{F}_i \ar[r]_{\varphi_{ij}} \ar[d]_{\text{pr}_0^*\psi_i}
& \text{pr}_1^*\mathcal{F}_j \ar[d]^{\text{pr}_1^*\psi_j} \\
\text{pr}_0^*\mathcal{F}'_i \ar[r]^{\varphi'_{ij}} &
\text{pr}_1^*\mathcal{F}'_j \\
}
$$
均可交换。
\end{enumerate}
\end{definition}

\noindent
应当牢记的一个典型例子如下。设 $S = \bigcup S_i$ 是开覆盖。
在这种情形下，我们已经在 Sheaves, Section
\ref{sheaves-section-glueing-sheaves} 中见过集合层的下降数据，
在那里称之为“相对于给定覆盖的集合层粘合数据”。此外，我们证明了
粘合数据范畴等价于 $S$ 上的层范畴。我们将证明：当
$\{S_i \to S\}_{i\in I}$ 是 fpqc 覆盖时，上述背景中也有相应结论。

\medskip\noindent
在极端情形中，覆盖 $\{S \to S\}$ 由 $\text{id}_S$ 给出，此时下降数据
必为 $(\mathcal{F}, \text{id}_\mathcal{F})$ 的形式。余圈条件保证，
在这种情形下，$\mathcal{F}$ 上唯一允许的态射就是恒等态射。特别地，
下述引理表明：对每个拟凝聚 $\mathcal{O}_S$-模层和任意给定的态射族，
都相应地存在唯一的下降数据。

\begin{lemma}
\label{descent-lemma-refine-descent-datum}
设 $\mathcal{U} = \{U_i \to U\}_{i \in I}$ 和
$\mathcal{V} = \{V_j \to V\}_{j \in J}$ 是具有固定靶的概形态射族。
设 $(g, \alpha : I \to J, (g_i)) : \mathcal{U} \to \mathcal{V}$
是具有固定靶的态射族之间的态射；见 Sites, Definition
\ref{sites-definition-morphism-coverings}。设
$(\mathcal{F}_j, \varphi_{jj'})$ 是拟凝聚层相对于态射族
$\{V_j \to V\}_{j \in J}$ 的一个下降数据。则
\begin{enumerate}
\item 系统
$$
\left(g_i^*\mathcal{F}_{\alpha(i)},
(g_i \times g_{i'})^*\varphi_{\alpha(i)\alpha(i')}\right)
$$
是相对于态射族 $\{U_i \to U\}_{i \in I}$ 的下降数据。
\item 此构造关于下降数据 $(\mathcal{F}_j, \varphi_{jj'})$ 具有函子性。
\item 给定具有固定靶的态射族之间的第二个态射
$(g', \alpha' : I \to J, (g'_i))$，且 $g = g'$，则存在下降数据之间的
函子性同构
$$
(g_i^*\mathcal{F}_{\alpha(i)},
(g_i \times g_{i'})^*\varphi_{\alpha(i)\alpha(i')})
\cong
((g'_i)^*\mathcal{F}_{\alpha'(i)},
(g'_i \times g'_{i'})^*\varphi_{\alpha'(i)\alpha'(i')}).
$$
\end{enumerate}
\end{lemma}

\begin{proof}
略。提示：给出第 (3) 项中下降数据同构的态射
$g_i^*\mathcal{F}_{\alpha(i)} \to (g'_i)^*\mathcal{F}_{\alpha'(i)}$
是态射 $\varphi_{\alpha(i)\alpha'(i)}$ 沿
$(g_i, g'_i) : U_i \to V_{\alpha(i)} \times_V V_{\alpha'(i)}$ 的拉回。
\end{proof}

\noindent
任意态射族 $\mathcal{U} = \{S_i \to S\}_{i \in I}$ 都以唯一方式成为
平凡覆盖 $\{S \to S\}$ 的加细。若 $\mathcal{F}$ 是 $S$ 上的拟凝聚层，
则将上述步骤所得的相对于 $\mathcal{U}$ 的下降数据简记为
$(\mathcal{F}|_{S_i}, can)$。

\begin{definition}
\label{descent-definition-descent-datum-effective-quasi-coherent}
设 $S$ 为概形，并设 $\{S_i \to S\}_{i \in I}$ 是以 $S$ 为靶的态射族。
\begin{enumerate}
\item 设 $\mathcal{F}$ 为拟凝聚 $\mathcal{O}_S$-模。称 $\mathcal{F}$
相对于覆盖 $\{S \to S\}$ 的唯一下降数据为{\it 平凡下降数据}。
\item 平凡下降数据在 $\{S_i \to S\}$ 上的拉回称为{\it 典范下降数据}，
记作 $(\mathcal{F}|_{S_i}, can)$。
\item 若相对于给定覆盖的拟凝聚层下降数据
$(\mathcal{F}_i, \varphi_{ij})$ 满足：存在 $S$ 上的拟凝聚层
$\mathcal{F}$，使得 $(\mathcal{F}_i, \varphi_{ij})$ 同构于
$(\mathcal{F}|_{S_i}, can)$，则称它是{\it 有效的}。
\end{enumerate}
\end{definition}

\begin{lemma}
\label{descent-lemma-zariski-descent-effective}
设 $S$ 为概形，并设 $S = \bigcup U_i$ 是开覆盖。拟凝聚层相对于态射族
$\mathcal{U} = \{U_i \to S\}$ 的任意下降数据都是有效的。此外，从拟凝聚
$\mathcal{O}_S$-模范畴到相对于 $\mathcal{U}$ 的下降数据范畴的函子是全忠实的。
\end{lemma}

\begin{proof}
这立即由 Sheaves, Section \ref{sheaves-section-glueing-sheaves}
以及拟凝聚性是局部性质这一事实得出；见 Modules, Definition
\ref{modules-definition-quasi-coherent}。
\end{proof}

\noindent
为进一步推进，我们首先需要研究环上模的情形。










\section{模的下降}
\label{descent-section-descent-modules}

\noindent
设 $R \to A$ 为环同态。由 Simplicial, Example
\ref{simplicial-example-push-outs-simplicial-object}，它给出余单纯 $R$-代数
$$
\xymatrix{
A
\ar@<1ex>[r]
\ar@<-1ex>[r]
&
A \otimes_R A
\ar@<0ex>[l]
\ar@<2ex>[r]
\ar@<0ex>[r]
\ar@<-2ex>[r]
&
A \otimes_R A \otimes_R A
\ar@<1ex>[l]
\ar@<-1ex>[l]
}
$$
将其记为 $(A/R)_\bullet$，于是 $(A/R)_n$ 是 $A$ 在 $R$ 上的
$(n + 1)$ 重张量积。给定映射 $\varphi : [n] \to [m]$，$R$-代数同态
$(A/R)_\bullet(\varphi)$ 为
$$
a_0 \otimes \ldots \otimes a_n
\longmapsto
\prod\nolimits_{\varphi(i) = 0} a_i
\otimes
\prod\nolimits_{\varphi(i) = 1} a_i
\otimes \ldots \otimes
\prod\nolimits_{\varphi(i) = m} a_i
$$
其中约定空乘积为 $1$。因此，采用 Simplicial, Section
\ref{simplicial-section-cosimplicial-object} 中的记号，最初几个态射为
$$
\begin{matrix}
\delta^1_0 & : & a_0 & \mapsto & 1 \otimes a_0 \\
\delta^1_1 & : & a_0 & \mapsto & a_0 \otimes 1 \\
\sigma^0_0 & : & a_0 \otimes a_1 & \mapsto & a_0a_1 \\
\delta^2_0 & : & a_0 \otimes a_1 & \mapsto & 1 \otimes a_0 \otimes a_1 \\
\delta^2_1 & : & a_0 \otimes a_1 & \mapsto & a_0 \otimes 1 \otimes a_1 \\
\delta^2_2 & : & a_0 \otimes a_1 & \mapsto & a_0 \otimes a_1 \otimes 1 \\
\sigma^1_0 & : & a_0 \otimes a_1 \otimes a_2 & \mapsto & a_0a_1 \otimes a_2 \\
\sigma^1_1 & : & a_0 \otimes a_1 \otimes a_2 & \mapsto & a_0 \otimes a_1a_2
\end{matrix}
$$
等等。

\medskip\noindent
一个 $R$-模 $M$ 给出余单纯 $(A/R)_\bullet$-模
$(A/R)_\bullet \otimes_R M$。换言之，$M_n = (A/R)_n \otimes_R M$，
并利用 $R$-代数同态 $(A/R)_n \to (A/R)_m$ 定义
$M \otimes_R (A/R)_\bullet$ 上相应的态射。

\medskip\noindent
拟凝聚层下降数据在模的背景中的对应概念如下。

\begin{definition}
\label{descent-definition-descent-datum-modules}
设 $R \to A$ 为环同态。
\begin{enumerate}
\item 一个{\it 模相对于 $R \to A$ 的下降数据 $(N, \varphi)$}
由一个 $A$-模 $N$ 和一个 $A \otimes_R A$-模同构
$$
\varphi : N \otimes_R A \to A \otimes_R N
$$
给出，并要求满足{\it 余圈条件}：$A \otimes_R A \otimes_R A$-模态射图
$$
\xymatrix{
N \otimes_R A \otimes_R A \ar[rr]_{\varphi_{02}}
\ar[rd]_{\varphi_{01}}
& &
A \otimes_R A \otimes_R N \\
& A \otimes_R N \otimes_R A \ar[ru]_{\varphi_{12}} &
}
$$
可交换（记号见下文）。
\item 一个{\it 下降数据态射 $(N, \varphi) \to (N', \varphi')$}
是一个 $A$-模态射 $\psi : N \to N'$，并要求图
$$
\xymatrix{
N \otimes_R A \ar[r]_\varphi \ar[d]_{\psi \otimes \text{id}_A} &
A \otimes_R N \ar[d]^{\text{id}_A \otimes \psi} \\
N' \otimes_R A \ar[r]^{\varphi'} &
A \otimes_R N'
}
$$
可交换。
\end{enumerate}
\end{definition}

\noindent
在此定义中采用如下记号：$\varphi_{01} = \varphi \otimes \text{id}_A$，
$\varphi_{12} = \text{id}_A \otimes \varphi$；若
$\varphi(n \otimes 1) = \sum a_i \otimes n_i$，则
$\varphi_{02}(n \otimes 1 \otimes 1) = \sum a_i \otimes 1 \otimes n_i$。
三者都是 $A \otimes_R A \otimes_R A$-模同态。等价地，
$$
\varphi_{ij}
=
\varphi \otimes_{(A/R)_1, \ (A/R)_\bullet(\tau^2_{ij})} (A/R)_2
$$
其中 $\tau^2_{ij} : [1] \to [2]$ 是由 $0 \mapsto i$、$1 \mapsto j$
给出的映射。具体而言，
$(A/R)_{\bullet}(\tau^2_{02})(a_0 \otimes a_1) =
a_0 \otimes 1 \otimes a_1$,
其余情形类似\footnote{注意，若
$\{i, j, k\} = [2] = \{0, 1, 2\}$，则
$\tau^2_{ij} = \delta^2_k$；见 Simplicial, Definition
\ref{simplicial-definition-face-degeneracy}。}。

\medskip\noindent
为陈述下一个引理，还需要一些记号。设 $(N, \varphi)$ 是相对于环同态
$R \to A$ 的下降数据。对 $n \geq 0$ 和 $i \in [n]$，令
$$
N_{n, i} =
A \otimes_R
\ldots
\otimes_R A \otimes_R N \otimes_R A \otimes_R
\ldots
\otimes_R A
$$
其中因子 $N$ 位于第 $i$ 个位置。这是一个 $(A/R)_n$-模。引入映射
$\tau^n_i : [0] \to [n]$，$0 \mapsto i$，便可看出
$$
N_{n, i} = N \otimes_{(A/R)_0, \ (A/R)_\bullet(\tau^n_i)} (A/R)_n
$$
对 $0 \leq i \leq j \leq n$，令 $\tau^n_{ij} : [1] \to [n]$ 为把
$0$ 映至 $i$、把 $1$ 映至 $j$ 的映射。与上文类似，同态 $\varphi$
诱导同构
$$
\varphi^n_{ij}
=
\varphi \otimes_{(A/R)_1, \ (A/R)_\bullet(\tau^n_{ij})} (A/R)_n :
N_{n, i} \longrightarrow N_{n, j}
$$
（当 $i < j$ 时为 $(A/R)_n$-模同构）。若 $i = j$，则令
$\varphi^n_{ij} = \text{id}$。由于这些都是同构，我们可以把因子 $N$
移至任意所需位置。余圈条件恰好说明，移动方式并不影响结果（例如，
复合其中两个同构或一次完成移动所得结果相同）。最后，对任意
$\beta : [n] \to [m]$，将态射
$$
N_{\beta, i} : N_{n, i} \to N_{m, \beta(i)}
$$
定义为唯一满足下式的 $(A/R)_\bullet(\beta)$-半线性映射：
$$
N_{\beta, i}(1 \otimes \ldots \otimes n \otimes \ldots \otimes 1)
=
1 \otimes \ldots \otimes n \otimes \ldots \otimes 1
$$
对所有 $n \in N$ 均成立。这提示了下述引理。

\begin{lemma}
\label{descent-lemma-descent-datum-cosimplicial}
设 $R \to A$ 为环同态。给定下降数据 $(N, \varphi)$，可以按如下规则
为其配备一个余单纯 $(A/R)_\bullet$-模
$N_\bullet$\footnote{严格说来，应写成 $(N, \varphi)_\bullet$。}：
令 $N_n = N_{n, n}$，并对给定的 $\beta : [n] \to [m]$ 定义
$$
N_\bullet(\beta) = (\varphi^m_{\beta(n)m}) \circ N_{\beta, n} :
N_{n, n} \longrightarrow N_{m, m}.
$$
此步骤关于下降数据具有函子性。
\end{lemma}

\begin{proof}
若 $\varphi(n \otimes 1) = \sum \alpha_i \otimes x_i$，则最初几个态射为
$$
\begin{matrix}
\delta^1_0 & : & N & \to & A \otimes N & n & \mapsto & 1 \otimes n \\
\delta^1_1 & : & N & \to & A \otimes N & n & \mapsto &
\sum \alpha_i \otimes x_i\\
\sigma^0_0 & : & A \otimes N & \to & N & a_0 \otimes n & \mapsto & a_0n \\
\delta^2_0 & : & A \otimes N & \to & A \otimes A \otimes N &
a_0 \otimes n & \mapsto & 1 \otimes a_0 \otimes n \\
\delta^2_1 & : & A \otimes N & \to & A \otimes A \otimes N &
a_0 \otimes n & \mapsto & a_0 \otimes 1 \otimes n \\
\delta^2_2 & : & A \otimes N & \to & A \otimes A \otimes N &
a_0 \otimes n & \mapsto & \sum a_0 \otimes \alpha_i \otimes x_i \\
\sigma^1_0 & : & A \otimes A \otimes N & \to & A \otimes N &
a_0 \otimes a_1 \otimes n & \mapsto & a_0a_1 \otimes n \\
\sigma^1_1 & : & A \otimes A \otimes N & \to & A \otimes N &
a_0 \otimes a_1 \otimes n & \mapsto & a_0 \otimes a_1n
\end{matrix}
$$
其中采用 Simplicial, Section
\ref{simplicial-section-cosimplicial-object} 的记号。我们先验证两个关系
$\sigma^0_0 \circ \delta^1_0 = \text{id}$ 和
$\sigma^0_0 \circ \delta^1_1 = \text{id}$。第一个关系
$\sigma^0_0 \circ \delta^1_0 = \text{id}$ 由上述态射的显式描述立即可见。
证明第二个关系时必须使用余圈条件（因为任意同构
$\varphi : N \otimes_R A \to A \otimes_R N$ 未必满足该关系）。记
$p = \sigma^0_0 \circ \delta^1_1 : N \to N$。由上述态射的描述可知，
$p$ 也等于
$$
p = \varphi \otimes \text{id} :
N = (N \otimes_R A) \otimes_{(A \otimes_R A)} A
\longrightarrow
(A \otimes_R N) \otimes_{(A \otimes_R A)} A = N
$$
由于 $\varphi$ 是同构，$p$ 也是同构。对某些 $\alpha_i \in A$ 和
$x_i \in N$，写作
$\varphi(n \otimes 1) = \sum \alpha_i \otimes x_i$。于是
$p(n) = \sum \alpha_ix_i$。再对某些 $\alpha_{ij} \in A$ 和
$y_j \in N$，写作
$\varphi(x_i \otimes 1) = \sum \alpha_{ij} \otimes y_j$。余圈条件给出
$$
\sum \alpha_i \otimes \alpha_{ij} \otimes y_j
=
\sum \alpha_i \otimes 1 \otimes x_i.
$$
这意味着 $p(n) = \sum \alpha_ix_i = \sum \alpha_i\alpha_{ij}y_j =
\sum \alpha_i p(x_i) = p(p(n))$。因此 $p$ 是幂等算子；又因它是同构，
故它必为恒等态射。

\medskip\noindent
要完整证明 $N_\bullet$ 是余单纯模，需要验证 Simplicial, Remark
\ref{simplicial-remark-relations-cosimplicial} 中的全部五类关系。
$\sigma$ 之间的复合关系是显然的；$\delta$ 之间的复合关系归结为
$\varphi$ 的余圈条件。完全仿照上文即可验证关系
$\sigma_j \circ \delta_j = \sigma_j \circ \delta_{j + 1} = \text{id}$。
最后，$\delta$ 与 $\sigma$ 的其余复合关系对任意 $\varphi$ 都成立。
\end{proof}

\noindent
注意，可以为 $R$-模 $M$ 配备一个典范下降数据，即
$(M \otimes_R A, can)$，其中
$can : (M \otimes_R A) \otimes_R A \to A \otimes_R (M \otimes_R A)$
是显然的态射：
$(m \otimes a) \otimes a' \mapsto a \otimes (m \otimes a')$.

\begin{lemma}
\label{descent-lemma-canonical-descent-datum-cosimplicial}
设 $R \to A$ 为环同态，$M$ 为 $R$-模。与该典范下降数据相联系的余单纯
$(A/R)_\bullet$-模同构于余单纯模 $(A/R)_\bullet \otimes_R M$。
\end{lemma}

\begin{proof}
略。
\end{proof}

\begin{definition}
\label{descent-definition-descent-datum-effective-module}
设 $R \to A$ 为环同态。若存在一个 $R$-模 $M$ 以及从
$(M \otimes_R A, can)$ 到 $(N, \varphi)$ 的下降数据同构，
则称下降数据 $(N, \varphi)$ 是{\it 有效的}。
\end{definition}

\noindent
设 $R \to A$ 为环同态，$(N, \varphi)$ 为下降数据。可以取与
$N_\bullet$ 相联系的上链复形 $s(N_\bullet)$（见 Simplicial, Section
\ref{simplicial-section-dold-kan-cosimplicial}）。它具有如下形状：
$$
N \to A \otimes_R N \to A \otimes_R A \otimes_R N \to \ldots
$$
这些态射可描述如下。第一个态射是
$$
n \longmapsto 1 \otimes n - \varphi(n \otimes 1).
$$
第二个态射在纯张量上的取值为
$$
a \otimes n \longmapsto 1 \otimes a \otimes n
- a \otimes 1 \otimes n + a \otimes \varphi(n \otimes 1).
$$
此后的规律是清楚的。

\medskip\noindent
在特殊情形 $N = A \otimes_R M$ 中，对任意 $m \in M$，元素
$1 \otimes m$ 都属于与余单纯模 $(A/R)_\bullet \otimes_R M$ 相联系的
上链复形的第一个态射之核。因此得到增广上链复形
\begin{equation}
\label{descent-equation-extended-complex}
0 \to M \to A \otimes_R M \to A \otimes_R A \otimes_R M \to \ldots
\end{equation}
这里把 $0$ 置于次数 $-2$，把模 $M$ 置于次数 $-1$，把模
$A \otimes_R M$ 置于次数 $0$，依此类推。注意，当 $M = R$ 时，
此复形的形状为
$$
0 \to R \to A \to A \otimes_R A \to A \otimes_R A \otimes_R A \to \ldots
$$

\begin{lemma}
\label{descent-lemma-with-section-exact}
设 $R \to A$ 有截面。则对任意 $R$-模 $M$，增广上链复形
(\ref{descent-equation-extended-complex}) 正合。
\end{lemma}

\begin{proof}
由 Simplicial, Lemma
\ref{simplicial-lemma-push-outs-simplicial-object-w-section}，态射
$R \to (A/R)_\bullet$ 是余单纯 $R$-代数之间的同伦等价（这里 $R$
表示常值余单纯 $R$-代数）。由于 $\otimes_R M$ 是从 $R$-代数范畴到
$R$-模范畴的函子，$M \to (A/R)_\bullet \otimes_R M$ 因而是余单纯
$R$-模范畴中的同伦等价；见 Simplicial, Lemma
\ref{simplicial-lemma-functorial-homotopy}。由此，相联系复形之间的诱导态射
也是同伦等价；见 Simplicial, Lemma
\ref{simplicial-lemma-homotopy-s-Q}。与常值余单纯 $R$-模 $M$ 相联系的
复形为
$$
\xymatrix{
M \ar[r]^0 & M \ar[r]^1 & M \ar[r]^0 & M \ar[r]^1 & M \ldots
}
$$
故结论成立（增广版本只是在开头再添一个 $M$）。
\end{proof}

\begin{lemma}
\label{descent-lemma-ff-exact}
设 $R \to A$ 忠实平坦；见 Algebra, Definition
\ref{algebra-definition-flat}。则对任意 $R$-模 $M$，增广上链复形
(\ref{descent-equation-extended-complex}) 正合。
\end{lemma}

\begin{proof}
假设能证明存在忠实平坦环同态 $R \to R'$，使结论对环同态
$R' \to A' = R' \otimes_R A$ 成立，则结论也对 $R \to A$ 成立。
事实上，对任意 $R$-模 $M$，余单纯模
$(M \otimes_R R') \otimes_{R'} (A'/R')_\bullet$ 正是余单纯模
$R' \otimes_R (M \otimes_R (A/R)_\bullet)$。因此，复形
$(M \otimes_R R') \otimes_{R'} (A'/R')_\bullet$ 的上同调消失，
再结合 $R \to R'$ 的忠实平坦性，即推出与
$M \otimes_R (A/R)_\bullet$ 相联系的复形之上同调消失。增广复形在次数
$-1$ 和 $0$ 的上同调群消失也可同理证明（证明从略）。

\medskip\noindent
这样的忠实平坦扩张确实存在。取 $R' = A$ 即可，因为环同态
$R' = A \to A' = A \otimes_R A$ 有截面 $a \otimes a' \mapsto aa'$，
从而可应用 Lemma \ref{descent-lemma-with-section-exact}。
\end{proof}

\noindent
下面说明此复形与有效性问题的关系。

\begin{lemma}
\label{descent-lemma-recognize-effective}
设 $R \to A$ 为忠实平坦环同态，$(N, \varphi)$ 为下降数据。则
$(N, \varphi)$ 有效，当且仅当典范态射
$$
A \otimes_R H^0(s(N_\bullet)) \longrightarrow N
$$
是同构。
\end{lemma}

\begin{proof}
若 $(N, \varphi)$ 有效，则可写成 $N = A \otimes_R M$ 且
$\varphi = can$。由 Lemmas
\ref{descent-lemma-canonical-descent-datum-cosimplicial} 和
\ref{descent-lemma-ff-exact}，可得 $H^0(s(N_\bullet)) = M$。反之，设引理中的
态射为同构。在此情形令 $M = H^0(s(N_\bullet))$。这是 $N$ 的一个
$R$-子模，即 $M = \{n \in N \mid 1 \otimes n = \varphi(n \otimes 1)\}$。
只需验证：通过同构 $A \otimes_R M \to N$，典范下降数据与
$\varphi$ 相符。验证从略。
\end{proof}

\begin{lemma}
\label{descent-lemma-descent-descends}
设 $R \to A$ 为忠实平坦环同态，且 $R \to R'$ 忠实平坦。令
$A' = R' \otimes_R A$。若相对于 $R' \to A'$ 的全部下降数据均有效，
则相对于 $R \to A$ 的全部下降数据也均有效。
\end{lemma}

\begin{proof}
设 $(N, \varphi)$ 是相对于 $R \to A$ 的下降数据。令
$N' = R' \otimes_R N = A' \otimes_A N$，并以
$\varphi' = \text{id}_{R'} \otimes \varphi$ 表示下降数据 $\varphi$
的基变换。于是 $(N', \varphi')$ 是相对于 $R' \to A'$ 的下降数据，且
$H^0(s(N'_\bullet)) = R' \otimes_R H^0(s(N_\bullet))$。此外，态射
$A' \otimes_{R'} H^0(s(N'_\bullet)) \to N'$ 可与 $A$-模态射
% P05-ERR-0018
$A \otimes_R H^0(s(N)) \to N$ 沿忠实平坦态射 $A \to A'$ 的基变换
等同。故由 Lemma \ref{descent-lemma-recognize-effective} 得出结论。
\end{proof}

\noindent
下面是本节的主要结果。它的证明可能显得有些笨拙；更高观点的处理方式见
下文 Remark \ref{descent-remark-homotopy-equivalent-cosimplicial-algebras}。

\begin{proposition}
\label{descent-proposition-descent-module}
\begin{slogan}
模沿忠实平坦环同态的有效下降。
\end{slogan}
设 $R \to A$ 为忠实平坦环同态。则
\begin{enumerate}
\item 模相对于 $R \to A$ 的任意下降数据都是有效的；
\item 从 $R$-模到下降数据范畴的函子
$M \mapsto (A \otimes_R M, can)$ 是等价；
\item 其逆函子由 $(N, \varphi) \mapsto H^0(s(N_\bullet))$ 给出。
\end{enumerate}
\end{proposition}

\begin{proof}
我们只证明 (1)，省略 (2) 和 (3) 的证明。由于 $R \to A$ 忠实平坦，
存在忠实平坦基变换 $R \to R'$，使
$R' \to A' = R' \otimes_R A$ 有截面（如 Lemma
\ref{descent-lemma-ff-exact} 的证明中那样取 $R' = A$ 即可）。因此，利用
Lemma \ref{descent-lemma-descent-descends}，可以假设 $R \to A$ 有一个截面，
记为 $\sigma : A \to R$。设 $(N, \varphi)$ 是相对于 $R \to A$ 的
下降数据。令
$$
M = H^0(s(N_\bullet)) = \{n \in N \mid 1 \otimes n = \varphi(n \otimes 1)\}
\subset
N
$$
由 Lemma \ref{descent-lemma-recognize-effective}，只需证明
$A \otimes_R M \to N$ 是同构。

\medskip\noindent
取元素 $n \in N$。对某些 $a_i \in A$ 和 $x_i \in N$，写作
$\varphi(n \otimes 1) = \sum a_i \otimes x_i$。由 Lemma
\ref{descent-lemma-descent-datum-cosimplicial}，在 $N$ 中有
$n = \sum a_i x_i$（因为任意余单纯对象中都有
$\sigma^0_0 \circ \delta^1_1 = \text{id}$）。接着，对某些
$a_{ij} \in A$ 和 $y_j \in N$，写作
$\varphi(x_i \otimes 1) = \sum a_{ij} \otimes y_j$。余圈条件意味着
$$
\sum a_i \otimes a_{ij} \otimes y_j = \sum a_i \otimes 1 \otimes x_i
$$
在 $A \otimes_R A \otimes_R N$ 中成立。由此得到两点：
\begin{enumerate}
\item 对第一个 $A$ 应用 $\sigma$，得到
$\sum \sigma(a_i) \varphi(x_i \otimes 1) = \sum \sigma(a_i) \otimes x_i$；
\item 对中间的 $A$ 应用 $\sigma$，得到
$\sum_i a_i \otimes \sum_j \sigma(a_{ij}) y_j = \sum a_i \otimes x_i$。
\end{enumerate}
由 (1) 可知 $\sum \sigma(a_i) x_i \in M$。将这一点应用于 $x_i$，
便知对每个 $i$ 都有
% P05-ERR-0019
$\sum \sigma(a_{ij})y_i \in M$。把 (2) 中的等式相乘，得到
$\sum_i a_i (\sum_j \sigma(a_{ij}) y_j) = \sum a_i x_i = n$。
故 $A \otimes_R M \to N$ 是满射。最后，设 $m_i \in M$ 且
$\sum a_i m_i = 0$。对 $\sum a_im_i \otimes 1$ 应用 $\varphi$，可见
$\sum a_i \otimes m_i = 0$。换言之，$A \otimes_R M \to N$ 是单射，
结论成立。
\end{proof}

\begin{remark}
\label{descent-remark-standard-covering}
设 $R$ 为环，且 $f_1, \ldots, f_n\in R$ 生成单位理想。环
$A = \prod_i R_{f_i}$ 是忠实平坦 $R$-代数。注意，余单纯环
$(A/R)_\bullet$ 在次数 $n$ 的环为
$$
\prod\nolimits_{i_0, \ldots, i_n} R_{f_{i_0}\ldots f_{i_n}}
$$
因此上述结果重新得到 Algebra, Lemmas
\ref{algebra-lemma-standard-covering}、\ref{algebra-lemma-cover-module}
和 \ref{algebra-lemma-glue-modules}。不过，由于高次数中的正合性，
上述结果实际给出了更强的结论：它蕴含仿射概形上拟凝聚层的
{\v C}ech 上同调平凡。由此得到 Cohomology of Schemes, Lemma
\ref{coherent-lemma-cech-cohomology-quasi-coherent-trivial} 的另一个证明。
\end{remark}

\begin{remark}
\label{descent-remark-homotopy-equivalent-cosimplicial-algebras}
设 $R$ 为环，$A_\bullet$ 为余单纯 $R$-代数。在此背景中，下降数据
对应于满足如下性质的余单纯 $A_\bullet$-模 $M_\bullet$：对每个
$n, m \geq 0$ 和每个 $\varphi : [n] \to [m]$，态射
$M(\varphi) : M_n \to M_m$ 诱导同构
$$
M_n \otimes_{A_n, A(\varphi)} A_m \longrightarrow M_m.
$$
称这样的余单纯模为{\it 笛卡尔模}。在此背景中，Proposition
\ref{descent-proposition-descent-module} 的证明可分为以下步骤：
\begin{enumerate}
\item 若 $R \to R'$ 和 $R \to A$ 均忠实平坦，则当相对于
$(R' \otimes_R A)/R'$ 的下降数据有效时，相对于 $A/R$ 的下降数据也有效。
\item 设 $A$ 为 $R$-代数。相对于 $A/R$ 的下降数据对应于笛卡尔
$(A/R)_\bullet$-模。
\item 若 $R \to A$ 有截面，则 $(A/R)_\bullet$ 与 $R$ 同伦等价；
这里后者是取值为 $R$ 的常值余单纯 $R$-代数。
\item 若 $A_\bullet \to B_\bullet$ 是余单纯 $R$-代数之间的同伦等价，
则函子 $M_\bullet \mapsto M_\bullet \otimes_{A_\bullet} B_\bullet$
诱导笛卡尔 $A_\bullet$-模范畴与笛卡尔 $B_\bullet$-模范畴之间的等价。
\end{enumerate}
关于 (1)，见 Lemma \ref{descent-lemma-descent-descends}；(2) 使用 Lemma
\ref{descent-lemma-descent-datum-cosimplicial}；(3) 已在 Lemma
\ref{descent-lemma-with-section-exact} 的证明中见过（它依赖 Simplicial,
Lemma \ref{simplicial-lemma-push-outs-simplicial-object-w-section}）。
此外，只要从正确角度思考，(4) 是显然的！
\end{remark}








\section{普遍单射态射的下降}
\label{descent-section-descent-universally-injective}

\noindent
代数几何中的许多构造都使用{\it 下降}技术。例如，为了在给定空间上
构造对象，可以先在投影到该空间的某个稍大空间上工作；也可以沿覆盖态射
拉回，以验证空间或态射的某种性质。这类技术是否有用，当然取决于能否
识别出足够广泛的一类{\it 有效下降态射}。现代代数几何的 Grothendieck
式发展初期，人们已经证明：{\it 拟紧}且{\it 忠实平坦}的态射类对于
对象、态射及其许多性质的下降都是有效的。

\medskip\noindent
与通常一样，这一陈述归结为环与模的性质。Grothendieck 证明：要使环同态
$f: R \to S$ 成为模的有效下降态射，$f$ 忠实平坦便足够。然而，此条件
排除了许多自然例子；例如，任意可裂环同态都是有效下降态射。甚至在
忠实平坦下降的证明中就自然出现了一个这样的例子：对任意环同态
$f: R \to S$，无论它是否平坦，$1_S \otimes f: S \to S \otimes_R S$
都由乘法态射给出分裂。

\medskip\noindent
于是可以追问：是否存在一个自然的环论条件，既涵盖忠实平坦态射，
又涵盖可裂环同态，并且蕴含模的有效下降？令人意外的是（至少本节作者
对此感到意外），这个问题的完整答案早在约 1970 年就已为人所知！
不难验证，$f: R \to S$ 成为模的有效下降态射的一个必要条件是：$f$
在 $R$-模范畴中{\it 普遍单射}；也就是说，对任意 $R$-模 $M$，态射
$1_M  \otimes f: M \to M \otimes_R S$ 必须为单射。事实证明，此条件
也充分。例如，若 $f$ 在 $R$-模范畴中可裂（但在环范畴中未必可裂），
则 $f$ 是模的有效下降态射。
 
\medskip\noindent
这一结果的历史略显曲折。Olivier \cite{olivier} 最先断言了它，并把
普遍单射态射称为{\it 纯}态射，但没有清楚说明证明。从 Joyal 与 Tierney
的工作 \cite{joyal-tierney} 中可以提取出此结果；不过据我们所知，
文献中最早出现的独立证明是 Mesablishvili 的证明
\cite{mesablishvili1}。本节的第一个目的就是阐述 Mesablishvili 的证明。
除改正若干笔误外，他的原始叙述几乎不需修改；唯一的例外是，我们将明确
说明代数几何中下降数据的惯常定义与 \cite{mesablishvili1} 所用定义之间
的关系。该证明完全是范畴论的，因而可以用单子的语言表述（并应用于其他
背景）；见 \cite{janelidze-tholen}。

\medskip\noindent
本节的第二个目的是汇集一些信息，说明模、代数和态射的哪些性质可以沿
普遍单射环同态下降。有限模与平坦模的情形由 Mesablishvili
\cite{mesablishvili2} 处理。


\subsection{范畴论预备知识}
\label{descent-subsection-category-prelims}

\noindent
先回顾几个可简化叙述的范畴论基本概念。本小节固定一个环境范畴。

\medskip\noindent
回忆：对两个态射 $g_1, g_2: B \to C$，$g_1$ 和 $g_2$ 的
{\it 等化子}是满足
$g_1 \circ f = g_2 \circ f$ 且对该性质具有泛性质的态射
$f: A \to B$。后一个陈述意味着，任何不含虚线箭头的可交换图
$$
\xymatrix{A' \ar[rd]^e \ar@/^1.5pc/[rrd] \ar@{-->}[d] & & \\
A \ar[r]^f & B \ar@<1ex>[r]^{g_1} \ar@<-1ex>[r]_{g_2} & 
C
}
$$
都能以唯一方式补全。在此情形下，也称图
\begin{equation}
\label{descent-equation-equalizer}
\xymatrix{
A \ar[r]^f  & B \ar@<1ex>[r]^{g_1} \ar@<-1ex>[r]_{g_2} & C
}
\end{equation}
为等化子。反转箭头即得到{\it 余等化子}的定义。见 Categories,
Sections \ref{categories-section-equalizers} 和
\ref{categories-section-coequalizers}。

\medskip\noindent
由于等化子涉及泛性质，应用协变函子通常不会保持等化子性质。与单态射和
满态射的情形一样，在某些情况下可以通过给出分裂来绕过这一问题。

\begin{definition}
\label{descent-definition-split-equalizer}
一个{\it 可裂等化子}是满足 $g_1 \circ f = g_2 \circ f$ 的图
(\ref{descent-equation-equalizer})，并且存在辅助态射 $h : B \to A$ 和
$i : C \to B$，使得
\begin{equation}
\label{descent-equation-split-equalizer-conditions}
h \circ f = 1_A, \quad f \circ h = i \circ g_1, \quad i \circ g_2 = 1_B.
\end{equation}
\end{definition}

\noindent
关键在于，箭头之间的这些等式迫使 (\ref{descent-equation-equalizer}) 成为等化子：
写成 $e = f \circ (h \circ e)$，态射 $e$ 便唯一地通过 $f$ 分解。因此，
对可裂等化子应用协变函子仍得到可裂等化子；应用反变函子则得到
{\it 可裂余等化子}，其定义是显然的。

\subsection{普遍单射态射}
\label{descent-subsection-universally-injective}

\noindent
回忆，$\textit{Rings}$ 表示含 $1$ 交换环的范畴。对 $\textit{Rings}$ 的对象
$R$，以 $\text{Mod}_R$ 表示 $R$-模范畴。

\begin{remark}
\label{descent-remark-reflects}
设 $F : \mathcal{A} \to \mathcal{B}$ 是 Abel 范畴之间的正合函子，
且把非零对象映为非零对象，则它反映单射和满射。事实上，正合性意味着
$F$ 保持核与余核（比较 Homology, Section
\ref{homology-section-functors}）。例如，若 $f : R \to S$ 是忠实平坦
环同态，则 $\bullet \otimes_R S: \text{Mod}_R \to \text{Mod}_S$
具有这些性质。
\end{remark}

\noindent
设 $R$ 为环。回忆，若对所有 $P \in \text{Mod}_R$，态射
$f \otimes 1_P: M \otimes_R P \to N \otimes_R P$ 都是单射，
则称 $\text{Mod}_R$ 中的态射 $f : M \to N$ 是{\it 普遍单射}的。
见 Algebra, Definition \ref{algebra-definition-universally-injective}。

\begin{definition}
\label{descent-definition-universally-injective}
若环同态 $f: R \to S$ 作为 $\text{Mod}_R$ 中的态射是普遍单射的，
则称它是{\it 普遍单射}的。
\end{definition}

\begin{example}
\label{descent-example-split-injection-universally-injective}
$\text{Mod}_R$ 中任意可裂单射都是普遍单射的。特别地，
$\textit{Rings}$ 中任意可裂单射都是普遍单射的。
\end{example}

\begin{example}
\label{descent-example-cover-universally-injective}
设 $R$ 为环，且 $f_1, \ldots, f_n \in R$ 生成单位理想，则态射
$R \to R_{f_1} \oplus \ldots \oplus R_{f_n}$ 是普遍单射的。尽管这立即
由 Lemma \ref{descent-lemma-faithfully-flat-universally-injective} 得出，直接验证
仍有启发性：立刻可化归到 $R$ 为局部环的情形。此时某个 $f_i$ 必为单位，
故态射 $R \to R_{f_i}$ 是同构。
\end{example}

\begin{lemma}
\label{descent-lemma-faithfully-flat-universally-injective}
任意忠实平坦环同态都是普遍单射的。
\end{lemma}

\begin{proof}
这只是 Algebra, Lemma
\ref{algebra-lemma-faithfully-flat-universally-injective} 的改述。
\end{proof}

\noindent
\cite{mesablishvili1} 中的关键观察是：利用 $\text{Mod}_R$ 中内射余生成元
的通常构造，可以借助分裂有效地改述普遍单射性。

\begin{definition}
\label{descent-definition-C}
设 $R$ 为环。定义反变函子
{\it $C$} $ : \text{Mod}_R \to \text{Mod}_R$ 如下：
$$
C(M) = \Hom_{\textit{Ab}}(M, \mathbf{Q}/\mathbf{Z}),
$$
其中 $C(M)$ 上的 $R$-作用由 $rf(s) = f(rs)$ 给出。
\end{definition}

\noindent
在 More on Algebra, Section
\ref{more-algebra-section-injectives-modules} 中，此函子记为
$M \mapsto M^\vee$。

\begin{lemma}
\label{descent-lemma-C-is-faithful}
对任意环 $R$，函子 $C : \text{Mod}_R \to \text{Mod}_R$ 正合，
并且反映单射和满射。
\end{lemma}

\begin{proof}
正合性见 More on Algebra, Lemma \ref{more-algebra-lemma-vee-exact}；
其余性质由此得出，见 Remark \ref{descent-remark-reflects}。
\end{proof}

\begin{remark}
\label{descent-remark-adjunction}
我们将经常使用 $\Hom$ 与张量积之间的标准伴随，其形式是反变函子的
自然同构
\begin{equation}
\label{descent-equation-adjunction}
C(\bullet_1 \otimes_R \bullet_2) \cong \Hom_R(\bullet_1, C(\bullet_2)): 
\text{Mod}_R \times \text{Mod}_R \to \text{Mod}_R
\end{equation}
它把 $f: M_1 \otimes_R M_2 \to \mathbf{Q}/\mathbf{Z}$ 映为
$m_1 \mapsto 
(m_2 \mapsto f(m_1 \otimes m_2))$。见 Algebra, Lemma
\ref{algebra-lemma-hom-from-tensor-product-variant}。由此可得如下推论：若
$$
\xymatrix@C=9pc{
C(M) \ar@<1ex>[r] \ar@<-1ex>[r] & C(N) \ar[r] & C(P)
}
$$
是 $\text{Mod}_R$ 中的可裂余等化子图，则对任意
$Q \in \text{Mod}_R$，图
$$
\xymatrix@C=9pc{
C(M \otimes_R Q) \ar@<1ex>[r] \ar@<-1ex>[r] & C(N \otimes_R Q) \ar[r] & C(P 
\otimes_R Q)
}
$$
也是可裂余等化子图。
\end{remark}

\begin{lemma}
\label{descent-lemma-split-surjection}
设 $R$ 为环。$\text{Mod}_R$ 中的态射 $f: M \to N$ 普遍单射，当且仅当
$C(f): C(N) \to C(M)$ 是可裂满射。
\end{lemma}

\begin{proof}
由 (\ref{descent-equation-adjunction})，对任意 $P \in \text{Mod}_R$ 都有可交换图
$$
\xymatrix@C=9pc{
\Hom_R( P, C(N)) \ar[r]_{\Hom_R(P,C(f))} \ar[d]^{\cong} &
\Hom_R(P,C(M)) \ar[d]^{\cong} \\
C(P \otimes_R N ) \ar[r]^{C(1_{P} \otimes f)} & C(P \otimes_R M ).
}
$$
若 $f$ 普遍单射，则
$1_{C(M)} \otimes f: C(M) \otimes_R M \to 
C(M) \otimes_R N$ 是单射，
故取 $P = C(M)$ 时，上图的两行都是满射。因此，可以把
$1_{C(M)} \in \Hom_R(C(M), C(M))$ 提升为某个使 $C(f)$ 分裂的
% P05-ERR-0020
$g \in \Hom_R(C(N), C(M))$。反之，若 $C(f)$ 是可裂满射，则上图的
两行都是满射，故由 Lemma \ref{descent-lemma-C-is-faithful}，
$1_{P} \otimes f$ 是单射。
\end{proof}

\begin{remark}
\label{descent-remark-functorial-splitting}
设 $f: M \to N$ 是 $\text{Mod}_R$ 中的普遍单射态射。选取 $C(f)$ 的一个
分裂 $g$ 后，可对每个 $P \in \text{Mod}_R$ 构造
$C(1_P \otimes f)$ 的函子性分裂。事实上，由
(\ref{descent-equation-adjunction})，这等价于关于 $P$ 函子性地分裂
$\Hom_R(P, 
C(f))$，而态射 $g \circ \bullet$ 即实现了这一点。
\end{remark}


\subsection{模及其态射的下降}
\label{descent-subsection-descent-modules-morphisms}

\noindent
本小节始终固定一个环同态 $f: R \to S$。如 Section
\ref{descent-section-descent-modules} 所见，可以用余单纯代数的语言讨论模的
下降数据；但本小节倾向于使用更具体的术语。

\medskip\noindent
对 $i = 1, 2, 3$，令 $S_i$ 为 $S$ 在 $R$ 上的 $i$ 重张量积。按下列公式
定义环同态 $\delta_0^1, \delta_1^1: S_1 \to S_2$、
$\delta_{01}^1, \delta_{02}^1, \delta_{12}^1: S_1 \to S_3$ 以及
$\delta_0^2, \delta_1^2, \delta_2^2: S_2 \to S_3$：
\begin{align*}
\delta^1_0  (a_0) & =  1 \otimes a_0 \\
\delta^1_1  (a_0) & = a_0 \otimes 1 \\
\delta^2_0  (a_0 \otimes a_1) & =  1 \otimes a_0 \otimes a_1 \\
\delta^2_1  (a_0 \otimes a_1) & =  a_0 \otimes 1 \otimes a_1 \\
\delta^2_2  (a_0 \otimes a_1) & =  a_0 \otimes a_1 \otimes 1 \\
\delta_{01}^1(a_0) & = 1 \otimes 1 \otimes a_0 \\
\delta_{02}^1(a_0) & = 1 \otimes a_0 \otimes 1 \\
\delta_{12}^1(a_0) & = a_0 \otimes 1 \otimes 1.
\end{align*}
换言之，上指标表示源环，下指标表示插入因子 1 的位置。（此记号与
Section \ref{descent-section-descent-modules} 中引入的记号相容。）

\medskip\noindent
由 Definition \ref{descent-definition-descent-datum-modules} 回忆\footnote{准确地说，
这里的 $\theta$ 是 Definition \ref{descent-definition-descent-datum-modules} 中
$\varphi$ 的逆。}：对 $M \in \text{Mod}_S$，$M$ 相对于 $f$ 的一个
{\it 下降数据}是一个 $S_2$-模同构
$$
\theta :
M \otimes_{S,\delta^1_0} S_2
\longrightarrow
M \otimes_{S,\delta^1_1} S_2
$$
并且满足{\it 余圈条件}
\begin{equation}
\label{descent-equation-cocycle-condition}
(\theta \otimes \delta_2^2) \circ (\theta \otimes \delta_2^0) = (\theta \otimes 
\delta_2^1):
M \otimes_{S, \delta^1_{01}} S_3 \to M \otimes_{S,\delta^1_{12}} S_3.
\end{equation}
以 $DD_{S/R}$ 表示配备相对于 $f$ 的下降数据的 $S$-模所成的范畴。

\medskip\noindent
例如，给定 $M_0 \in \text{Mod}_R$ 并选取同构
$M \cong M_0 \otimes_R S$，把 $M \otimes_{S,\delta^1_0} S_2$ 和
$M \otimes_{S,\delta^1_1} S_2$ 自然地等同于 $M_0 \otimes_R S_2$，
便得到一个下降数据。特别地，此构造定义函子
$f^*: \text{Mod}_R \to DD_{S/R}$。

\begin{definition}
\label{descent-definition-effective-descent}
函子 $f^*: \text{Mod}_R \to DD_{S/R}$ 称为{\it 沿 $f$ 的底扩张}。
若 $f^*$ 全忠实，则称 $f$ 是{\it 模的下降态射}；若 $f^*$ 是范畴等价，
则称 $f$ 是{\it 模的有效下降态射}。
\end{definition}

\noindent
我们的目标是证明：当 $f$ 普遍单射时，可以利用 $\theta$ 在 $M$ 中定位
$M_0$。这一过程关键地使用若干等化子图。

\begin{lemma}
\label{descent-lemma-equalizer-M}
对 $(M,\theta) \in DD_{S/R}$，图
\begin{equation}
\label{descent-equation-equalizer-M}
\xymatrix@C=8pc{
M \ar[r]^{\theta \circ (1_M \otimes \delta_0^1)} &
M \otimes_{S, \delta_1^1} S_2
\ar@<1ex>[r]^{(\theta \otimes \delta_2^2) \circ (1_M \otimes \delta^2_0)}
\ar@<-1ex>[r]_{1_{M \otimes S_2} \otimes \delta^2_1} & 
M \otimes_{S, \delta_{12}^1} S_3
}
\end{equation}
是可裂等化子。
\end{lemma}

\begin{proof}
按下列公式定义环同态 $\sigma^0_0: S_2 \to S_1$ 以及 $\sigma_0^1, 
\sigma_1^1: S_3 \to S_2$：
\begin{align*}
\sigma^0_0 (a_0 \otimes a_1) & = a_0a_1 \\
\sigma^1_0 (a_0 \otimes a_1 \otimes a_2) & = a_0a_1 \otimes a_2 \\
\sigma^1_1 (a_0 \otimes a_1 \otimes a_2) & = a_0 \otimes a_1a_2.
\end{align*}
取辅助态射为
$1_M \otimes \sigma_0^0: M \otimes_{S, \delta_1^1} S_2 \to M$
和 $1_M \otimes \sigma_0^1: M \otimes_{S,\delta_{12}^1} S_3 \to M \otimes_{S, 
\delta_1^1} S_2$。(\ref{descent-equation-split-equalizer-conditions}) 所要求的
相容性中，第一个由余圈条件 (\ref{descent-equation-cocycle-condition}) 与
$\sigma_1^1$ 作张量积得出，其余则是直接的。
\end{proof}

\begin{lemma}
\label{descent-lemma-equalizer-CM}
对 $(M, \theta) \in DD_{S/R}$，图
\begin{equation}
\label{descent-equation-coequalizer-CM}
\xymatrix@C=8pc{
C(M \otimes_{S, \delta_{12}^1} S_3)
\ar@<1ex>[r]^{C((\theta \otimes \delta_2^2) \circ (1_M \otimes \delta^2_0))}
\ar@<-1ex>[r]_{C(1_{M \otimes S_2} \otimes \delta^2_1)} &
C(M \otimes_{S, \delta_1^1} S_2 )
\ar[r]^{C(\theta \circ (1_M \otimes \delta_0^1))} & C(M).
}
\end{equation}
由对 (\ref{descent-equation-equalizer-M}) 应用 $C$ 得到，并且是可裂余等化子。
\end{lemma}

\begin{proof}
略。
\end{proof}

\begin{lemma}
\label{descent-lemma-equalizer-S}
图
\begin{equation}
\label{descent-equation-equalizer-S}
\xymatrix@C=8pc{
S_1 \ar[r]^{\delta^1_1} &
S_2 \ar@<1ex>[r]^{\delta^2_2} \ar@<-1ex>[r]_{\delta^2_1} & 
S_3
}
\end{equation}
是可裂等化子。
\end{lemma}

\begin{proof}
在 Lemma \ref{descent-lemma-equalizer-M} 中取 $(M, \theta) = f^*(S)$。
\end{proof}

\noindent
这提示了 $f^*$ 的一个候选拟逆函子的定义。

\begin{definition}
\label{descent-definition-pushforward}
定义函子 {\it $f_*$} $: DD_{S/R} \to \text{Mod}_R$：令
$f_*(M, \theta)$ 为 $M$ 的 $R$-子模，使得图
\begin{equation}
\label{descent-equation-equalizer-f}
\xymatrix@C=8pc{f_*(M,\theta) \ar[r] & M \ar@<1ex>^{\theta \circ (1_M \otimes 
\delta_0^1)}[r] \ar@<-1ex>_{1_M \otimes \delta_1^1}[r] & 
M \otimes_{S, \delta_1^1} S_2 
}
\end{equation}
是等化子。
\end{definition}

\noindent
利用 Lemma \ref{descent-lemma-equalizer-M}，并注意限制函子
$\text{Mod}_S \to \text{Mod}_R$ 是底扩张函子
$\bullet \otimes_R S: \text{Mod}_R \to \text{Mod}_S$ 的右伴随，
可知 $f_*$ 是 $f^*$ 的右伴随。

\medskip\noindent
现在可以证明关键引理。在忠实平坦的情形，结论是显然的（见 Remark
\ref{descent-remark-descent-lemma}）；一般情形则需要一些论证。

\begin{lemma}
\label{descent-lemma-descent-lemma}
若 $f$ 普遍单射，则图
\begin{equation}
\label{descent-equation-equalizer-f2}
\xymatrix@C=8pc{
f_*(M, \theta) \otimes_R S
\ar[r]^{\theta \circ (1_M \otimes \delta_0^1)} &
M \otimes_{S, \delta_1^1} S_2 
\ar@<1ex>[r]^{(\theta \otimes \delta_2^2) \circ (1_M \otimes \delta^2_0)}
\ar@<-1ex>[r]_{1_{M \otimes S_2} \otimes \delta^2_1} &
M \otimes_{S, \delta_{12}^1} S_3
}
\end{equation}
由 (\ref{descent-equation-equalizer-f}) 在 $R$ 上与 $S$ 作张量积得到，
并且是等化子。
\end{lemma}

\begin{proof}
由 Lemma \ref{descent-lemma-split-surjection} 和 Remark
\ref{descent-remark-functorial-splitting}，态射
$C(1_N \otimes f): C(N \otimes_R S) \to C(N)$ 可关于 $N$ 函子性地分裂。
这给出下列可交换图中上方的竖直箭头：
$$
\xymatrix@C=8pc{
C(M \otimes_{S, \delta_1^1} S_2)
\ar@<1ex>^{C(\theta \circ (1_M \otimes \delta_0^1))}[r]
\ar@<-1ex>_{C(1_M \otimes \delta_1^1)}[r] \ar[d] &
C(M) \ar[r]\ar[d] & C(f_*(M,\theta)) \ar@{-->}[d] \\
C(M \otimes_{S,\delta_{12}^1} S_3)
\ar@<1ex>^{C((\theta \otimes \delta_2^2) \circ (1_M \otimes \delta^2_0))}[r]
\ar@<-1ex>_{C(1_{M \otimes S_2} \otimes \delta^2_1)}[r] \ar[d] &
C(M \otimes_{S, \delta_1^1} S_2 )
\ar[r]^{C(\theta \circ (1_M \otimes \delta_0^1))}
\ar[d]^{C(1_M \otimes \delta_1^1)} &
C(M) \ar[d] \ar@{=}[dl] \\
C(M \otimes_{S, \delta_1^1} S_2)
\ar@<1ex>[r]^{C(\theta \circ (1_M \otimes \delta_0^1))}
\ar@<-1ex>[r]_{C(1_M \otimes \delta_1^1)} &
C(M) \ar[r] &
C(f_*(M,\theta))
}
$$
其中沿各列的复合都是恒等态射。第二行是余等化子图
(\ref{descent-equation-coequalizer-CM})，因而产生虚线箭头。由右上方的方块得到
辅助态射 $C(f_*(M,\theta)) \to 
C(M)$ 和 $C(M) \to C(M\otimes_{S,\delta_1^1} S_2)$，它们蕴含第一行
是可裂余等化子图。由 Remark \ref{descent-remark-adjunction}，可在 $C$ 内部与
$S$ 作张量积，得到可裂余等化子图
$$
\xymatrix@C=8pc{
C(M \otimes_{S,\delta_2^2 \circ \delta_1^1} S_3)
\ar@<1ex>^{C((\theta \otimes \delta_2^2) \circ (1_M \otimes \delta^2_0))}[r] 
\ar@<-1ex>_{C(1_{M \otimes S_2} \otimes \delta^2_1)}[r] &
C(M \otimes_{S, \delta_1^1} S_2 )
\ar[r]^{C(\theta \circ (1_M \otimes \delta_0^1))} &
C(f_*(M,\theta) \otimes_R S).
}
$$
由 Lemma \ref{descent-lemma-C-is-faithful}，可知
(\ref{descent-equation-equalizer-f2}) 也必为等化子。
\end{proof}

\begin{remark}
\label{descent-remark-descent-lemma}
若 $f$ 是 $\text{Mod}_R$ 中的可裂单射，则可直接分裂 $f$，无须使用
$C$，从而简化论证。若 $f$ 忠实平坦，情况更为简单：此时 Lemma
\ref{descent-lemma-descent-lemma} 的结论立即成立，因为在 $R$ 上与 $S$ 作张量积
保持所有等化子。
\end{remark}

\begin{theorem}
\label{descent-theorem-descent}
下列条件等价。
\begin{enumerate}
\item[(a)] 态射 $f$ 是模的下降态射。
\item[(b)] 态射 $f$ 是模的有效下降态射。
\item[(c)] 态射 $f$ 普遍单射。
\end{enumerate}
\end{theorem}

\begin{proof}
(b) 蕴含 (a) 是显然的。下面验证 (a) 蕴含 (c)。若 $f$ 不是普遍单射，
则可找到 $M \in \text{Mod}_R$，使态射
$1_M \otimes f: M \to M \otimes_R S$ 有非零的核 $N$。自然投影
$M \to M/N$ 不是同构，但它在 $DD_{S/R}$ 中的像是同构。因此 $f^*$
并非全忠实。

\medskip\noindent
最后验证 (c) 蕴含 (b)。由 Lemmas \ref{descent-lemma-equalizer-M} 和
\ref{descent-lemma-descent-lemma}，对 $(M, \theta) \in DD_{S/R}$，自然态射
$f^* f_*(M,\theta) \to M$ 是 $S$-模同构。对
$M_0 \in \text{Mod}_R$，需要证明 $M_0 \to f_*f^*M_0$ 是同构。
由伴随性，复合 $f^*M_0 \to f^*f_*f^*M_0 \to f^*M_0$ 是恒等态射，
而上文表明第二个箭头是同构。因此
$f^*M_0 \to f^*f_*f^*M_0$ 是同构。由 Algebra, Lemma
\ref{algebra-lemma-base-change-along-universally-injective-ring-map}，
可知 $M_0 \to f_* f^* M_0$ 是同构。因此，$f_*$ 与 $f^*$ 互为拟逆函子，
结论得证。
\end{proof}

\subsection{模的性质的下降}
\label{descent-subsection-descent-properties-modules}

\noindent
本小节始终固定一个普遍单射环同态 $f : R \to S$、一个对象
$M \in \text{Mod}_R$ 以及一个环同态 $R \to A$。现在研究 $M$ 或 $A$
的哪些性质可以在沿 $f$ 作底扩张后检验。先从 \cite{mesablishvili2}
中的一些结果开始。

\begin{lemma}
\label{descent-lemma-flat-to-injective}
若 $M \in \text{Mod}_R$ 平坦，则 $C(M)$ 是内射 $R$-模。
\end{lemma}

\begin{proof}
设 $0 \to N \to P \to Q \to 0$ 是 $\text{Mod}_R$ 中的正合列。
由于 $M$ 平坦，
$$
0 \to N \otimes_R M \to P \otimes_R M \to Q \otimes_R M \to 0
$$
正合。由 Lemma \ref{descent-lemma-C-is-faithful}，
$$
0 \to C(Q \otimes_R M) \to C(P \otimes_R M) \to C(N \otimes_R M) \to 0
$$
正合。由 (\ref{descent-equation-adjunction})，后一个序列可改写为
$$
0 \to \Hom_R(Q, C(M)) \to \Hom_R(P, C(M)) \to \Hom_R(N, C(M)) \to 0.
$$
因此 $C(M)$ 是 $\text{Mod}_R$ 中的内射对象。
\end{proof}

\begin{theorem}
\label{descent-theorem-descend-module-properties}
若 $S$-模 $M \otimes_R S$ 具有下列性质之一：
\begin{enumerate}
\item[(a)]
有限生成；
\item[(b)]
有限表现；
\item[(c)]
平坦；
\item[(d)]
忠实平坦；
\item[(e)]
有限投射；
\end{enumerate}
则 $R$-模 $M$ 也具有相应性质（反之亦然）。
\end{theorem}

\begin{proof}
为证明 (a)，在 $\text{Mod}_S$ 中选取 $M \otimes_R S$ 的有限生成集
$\{n_i\}$。将每个 $n_i$ 写成 $\sum_j m_{ij} \otimes s_{ij}$，其中
$m_{ij} \in M$ 且 $s_{ij} \in S$。设 $F$ 是以 $e_{ij}$ 为基的有限自由
$R$-模，并令 $R$-模态射 $F \to M$ 把 $e_{ij}$ 映至 $m_{ij}$。于是
$F \otimes_R S\to M \otimes_R S$ 是满射，故
$\Coker(F \to M) \otimes_R S$ 为零，进而 $\Coker(F \to M)$ 为零。
这证明了 (a)。

\medskip\noindent
为证明 (b)，设 $M \otimes_R S$ 有限表现。由 (a)，$M$ 有限生成。选取核为
$K$ 的满射 $R^{\oplus n} \to M$。于是
% P05-ERR-0021
$K \otimes_R S \to S^{\oplus r} \to M \otimes_R S \to 0$ 正合。
由 Algebra, Lemma \ref{algebra-lemma-extension}，
$S^{\oplus r} \to M \otimes_R S$ 的核是有限 $S$-模。因此可找到有限多个
元素 $k_1, \ldots, k_t \in K$，使 $k_i \otimes 1$ 在
$S^{\oplus r}$ 中的像生成 $S^{\oplus r} \to M \otimes_R S$ 的核。
设 $K' \subset K$ 是由 $k_1, \ldots, k_t$ 生成的子模。则
$M' = R^{\oplus r}/K'$ 是有限表现 $R$-模，并有态射 $M' \to M$，
使 $M' \otimes_R S \to M \otimes_R S$ 为同构。故如所需，
$M' \cong M$。

\medskip\noindent
为证明 (c)，设 $0 \to M' \to M'' \to M \to 0$ 是 $\text{Mod}_R$ 中的短正合列。
由于 $\bullet \otimes_R S$ 是右正合函子，
$M'' \otimes_R S \to M \otimes_R S$ 为满射。故由 Lemma
\ref{descent-lemma-C-is-faithful}，态射
$C(M \otimes_R S) \to C(M'' \otimes_R S)$ 为单射。若
$M \otimes_R S$ 平坦，则 Lemma \ref{descent-lemma-flat-to-injective} 表明
$C(M \otimes_R S)$ 是 $\text{Mod}_S$ 中的内射对象。因此单射
$C(M \otimes_R S) \to C(M'' \otimes_R S)$ 在 $\text{Mod}_S$ 中可裂，
从而在 $\text{Mod}_R$ 中也可裂。由 Lemma
\ref{descent-lemma-split-surjection}，$C(M \otimes_R S) \to C(M)$ 是可裂满射，
所以 $C(M) \to C(M'')$ 是 $\text{Mod}_R$ 中的可裂单射。也就是说，序列
$$
0 \to C(M) \to C(M'') \to C(M') \to 0
$$
可裂正合。对 $N \in \text{Mod}_R$，由 (\ref{descent-equation-adjunction}) 可知
$$
0 \to C(M \otimes_R N) \to C(M'' \otimes_R N) \to C(M' \otimes_R N) \to 0
$$
可裂正合。由 Lemma \ref{descent-lemma-C-is-faithful}，
$$
0 \to M' \otimes_R N \to M'' \otimes_R N \to M \otimes_R N \to 0
$$
正合。这蕴含 $M$ 在 $R$ 上平坦。事实上，
% P05-ERR-0022
取 $M'$ 为满射到 $M$ 的自由模，可知对所有模 $N$ 都有
$\text{Tor}_1^R(M, N) = 0$，于是可应用 Algebra, Lemma
\ref{algebra-lemma-characterize-flat}。这证明了 (c)。

\medskip\noindent
为由 (c) 推出 (d)，注意：若 $N \in \text{Mod}_R$ 且 $M \otimes_R N$
为零，则 $M \otimes_R S \otimes_S (N \otimes_R S) \cong (M \otimes_R N) \otimes_R 
S$ 为零。因此 $N \otimes_R S$ 为零，进而 $N$ 为零。

\medskip\noindent
此时，为推出 (e)，只需回忆：$M$ 有限生成且投射，当且仅当它有限表现且
平坦。见 Algebra, Lemma \ref{algebra-lemma-finite-projective}。
\end{proof}

\noindent
对 $R$-代数也有如下版本。

\begin{theorem}
\label{descent-theorem-descend-algebra-properties}
若 $S$-代数 $A \otimes_R S$ 具有下列性质之一：
\begin{enumerate}
\item[(a)]
有限型；
\item[(b)]
有限表现；
\item[(c)]
形式非分歧；
\item[(d)]
非分歧；
\item[(e)]
\'etale；
\end{enumerate}
则 $R$-代数 $A$ 也具有相应性质（反之当然也成立）。
\end{theorem}

\begin{proof}
为证明 (a)，选取 $A \otimes_R S$ 在 $S$ 上的有限生成集 $\{x_i\}$。
将每个 $x_i$ 写成 $\sum_j y_{ij} \otimes s_{ij}$，其中
$y_{ij} \in A$ 且 $s_{ij} \in S$。令 $F$ 为以变量 $e_{ij}$ 生成的多项式
$R$-代数，并令把 $e_{ij}$ 映至 $y_{ij}$ 的 $R$-代数同态为
% P05-ERR-0023
$F \to M$。于是 $F \otimes_R S\to A \otimes_R S$ 是满射，故
$\Coker(F \to A) \otimes_R S$ 为零，进而 $\Coker(F \to A)$ 为零。
这证明了 (a)。

\medskip\noindent
为证明 (b)，设 $A \otimes_R S$ 是有限表现 $S$-代数。由 (a)，$A$ 在
$R$ 上有限型。选取核为 $I$ 的满射 $R[x_1, \ldots, x_n] \to A$。
于是 $I \otimes_R S \to S[x_1, \ldots, x_n] \to A \otimes_R S \to 0$
正合。由 Algebra, Lemma
\ref{algebra-lemma-finite-presentation-independent}，
$S[x_1, \ldots, x_n] \to A \otimes_R S$ 的核是有限生成理想。因此可找到
有限多个元素 $y_1, \ldots, y_t \in I$，使 $y_i \otimes 1$ 在
$S[x_1, \ldots, x_n]$ 中的像生成
$S[x_1, \ldots, x_n] \to A \otimes_R S$ 的核。令 $I' \subset I$ 为
$y_1, \ldots, y_t$ 生成的理想。则
$A' = R[x_1, \ldots, x_n]/I'$ 是有限表现 $R$-代数，并有态射
$A' \to A$，使 $A' \otimes_R S \to A \otimes_R S$ 为同构。
故如所需，$A' \cong A$。

\medskip\noindent
为证明 (c)，回忆：$A$ 在 $R$ 上形式非分歧，当且仅当相对微分模
$\Omega_{A/R}$ 消失；见 Algebra, Lemma
\ref{algebra-lemma-characterize-formally-unramified} 或
\cite[Proposition~17.2.1]{EGA4}。由于
$\Omega_{(A \otimes_R S)/S} = \Omega_{A/R} \otimes_R S$，
此消失性由 Theorem \ref{descent-theorem-descent} 下降。

\medskip\noindent
为由前述情形推出 (d)，回忆：$A$ 在 $R$ 上非分歧，当且仅当 $A$
在 $R$ 上形式非分歧且有限型；见 Algebra, Lemma
\ref{algebra-lemma-formally-unramified-unramified}。

\medskip\noindent
为证明 (e)，回忆：由 Algebra, Lemma
\ref{algebra-lemma-etale-flat-unramified-finite-presentation} 或
\cite[Th\'eor\`eme~17.6.1]{EGA4}，代数 $A$ 在 $R$ 上 \'etale，
当且仅当 $A$ 在 $R$ 上平坦、非分歧且有限表现。
\end{proof}

\begin{remark}
\label{descent-remark-when-locally-split}
若能找到忠实平坦环同态 $g: R \to T$，使 $T \to S \otimes_R T$ 具有
某种额外结构，事情就会更容易。例如，若能保证
$T \to S \otimes_R T$ 在 $\textit{Rings}$ 中可裂，则可以推出：
凡是在底扩张下稳定且能沿忠实平坦态射下降的模或代数性质，也能沿
普遍单射态射下降。一个显然的猜想是寻找 $g$，使 $T$ 不仅忠实平坦，
而且是 $\text{Mod}_R$ 中的内射对象；但即使对 $R = \mathbf{Z}$，
这样的同态也不可能存在。
\end{remark}
















\section{拟凝聚层的 fpqc 下降}
\label{descent-section-fpqc-descent-quasi-coherent}

\noindent
模的平坦下降的主要应用，是拟凝聚层相对于 fpqc 覆盖的相应下降命题。

\begin{lemma}
\label{descent-lemma-standard-fpqc-covering}
设 $S$ 为仿射概形，并设
$\mathcal{U} = \{f_i : U_i \to S\}_{i = 1, \ldots, n}$ 是 $S$ 的标准
fpqc 覆盖；见 Topologies, Definition
\ref{topologies-definition-standard-fpqc}。拟凝聚层相对于
$\mathcal{U} = \{U_i \to S\}$ 的任意下降数据都是有效的。此外，
从拟凝聚 $\mathcal{O}_S$-模范畴到相对于 $\mathcal{U}$ 的下降数据范畴的
函子是全忠实的。
\end{lemma}

\begin{proof}
这只是用概形语言改述 Proposition \ref{descent-proposition-descent-module}。
首先注意，拟凝聚层相对于 $\mathcal{U}$ 的下降数据 $\xi$，恰好等同于
拟凝聚层相对于覆盖
$\mathcal{U}' = \{\coprod_{i = 1, \ldots, n} U_i \to S\}$ 的下降数据
$\xi'$。此外，$\xi$ 的有效性等同于 $\xi'$ 的有效性。因此可以假设
$n = 1$，即 $\mathcal{U} = \{U \to S\}$，其中 $U$ 与 $S$ 均为仿射概形。
在此情形，下降数据对应于模相对于环同态
$$
\Gamma(S, \mathcal{O})
\longrightarrow
\Gamma(U, \mathcal{O}).
$$
的下降数据。由于 $U \to S$ 满且平坦，此环同态忠实平坦。换言之，
Proposition \ref{descent-proposition-descent-module} 适用，结论成立。
\end{proof}

\begin{proposition}
\label{descent-proposition-fpqc-descent-quasi-coherent}
设 $S$ 为概形，并设 $\mathcal{U} = \{\varphi_i : U_i \to S\}$ 是 fpqc
覆盖；见 Topologies, Definition \ref{topologies-definition-fpqc-covering}。
拟凝聚层相对于 $\mathcal{U} = \{U_i \to S\}$ 的任意下降数据都是有效的。
此外，从拟凝聚 $\mathcal{O}_S$-模范畴到相对于 $\mathcal{U}$ 的下降数据
范畴的函子是全忠实的。
\end{proposition}

\begin{proof}
设 $S = \bigcup_{j \in J} V_j$ 是仿射开覆盖。对 $j, j' \in J$，
以 $V_{jj'} = V_j \cap V_{j'}$ 表示交（它未必仿射）。对开子集
$V \subset S$，记
$\mathcal{U}_V = \{V \times_S U_i \to V\}_{i \in I}$；这是 fpqc 覆盖
（Topologies, Lemma \ref{topologies-lemma-fpqc}）。由 fpqc 覆盖的定义，
对每个 $j \in J$，可找到有限集 $K_j$、映射
$\underline{i} : K_j \to I$，以及仿射开集
$U_{\underline{i}(k), k} \subset U_{\underline{i}(k)}$（$k \in K_j$），
使 $\mathcal{V}_j = \{U_{\underline{i}(k), k} \to V_j\}_{k \in K_j}$
是 $V_j$ 的标准 fpqc 覆盖。当然，$\mathcal{V}_j$ 是
$\mathcal{U}_{V_j}$ 的加细。图示为
$$
\xymatrix{
\mathcal{V}_j \ar[r] \ar@{~>}[d] &
\mathcal{U}_{V_j} \ar[r] \ar@{~>}[d] &
\mathcal{U} \ar@{~>}[d] \\
V_j \ar@{=}[r] & V_j \ar[r] & S
}
$$
其中上方的水平箭头是具有固定靶的态射族之间的态射（见 Sites,
Definition \ref{sites-definition-morphism-coverings}）。

\medskip\noindent
为证明命题，依次证明从拟凝聚层到下降数据的函子忠实、全且本质满。

\medskip\noindent
忠实性。设 $\mathcal{F}$、$\mathcal{G}$ 是 $S$ 上的拟凝聚层，且
$a, b : \mathcal{F} \to \mathcal{G}$ 是 $\mathcal{O}_S$-模同态。
假设对所有 $i$ 都有 $\varphi_i^*(a) = \varphi_i^*(b)$。取 $s \in S$。
则对某个 $i \in I$ 和 $u \in U_i$，有 $s = \varphi_i(u)$。由于
$\mathcal{O}_{S, s} \to \mathcal{O}_{U_i, u}$ 平坦，因而忠实平坦
（Algebra, Lemma \ref{algebra-lemma-local-flat-ff}），可知
$a_s = b_s : \mathcal{F}_s \to \mathcal{G}_s$。故 $a = b$。

\medskip\noindent
全性。设 $\mathcal{F}$、$\mathcal{G}$ 是 $S$ 上的拟凝聚层，并设
$a_i : \varphi_i^*\mathcal{F} \to \varphi_i^*\mathcal{G}$ 是
$\mathcal{O}_{U_i}$-模同态，满足在
% P05-ERR-0024
$U_i \times_U U_j$ 上 $\text{pr}_0^*a_i = \text{pr}_1^*a_j$。
拉回这些态射，得到态射
$$
% P05-ERR-0025
a_k :
\varphi_{i(k)}^*\mathcal{F}|_{U_{\underline{i}(k), k}}
\longrightarrow
\varphi_{i(k)}^*\mathcal{G}|_{U_{\underline{i}(k), k}}
$$
其中 $k \in K_j$，记号同上。此外，Lemma
\ref{descent-lemma-refine-descent-datum} 保证这些态射定义了
$\mathcal{V}_j$ 上两个（典范）下降数据之间的态射。因此，由 Lemma
\ref{descent-lemma-standard-fpqc-covering}，相应地得到唯一态射
$a_j : \mathcal{F}|_{V_j} \to \mathcal{G}|_{V_j}$。为证明
$a_j|_{V_{jj'}} = a_{j'}|_{V_{jj'}}$，利用已经证明的忠实性，并注意：
$a_j$ 与 $a_{j'}$ 都等同于（典范）下降数据态射 $(a_i)_{i \in I}$
在任一同时加细 $\mathcal{V}_{j, V_{jj'}}$ 和
$\mathcal{V}_{j', V_{jj'}}$ 的覆盖上的拉回。例如，覆盖
% P05-ERR-0026
$\mathcal{V}_{jj'} =
\{V_k \times_S V_{k'} \to V_{jj'}\}_{k \in K_j, k' \in K_{j'}}$
即可。

\medskip\noindent
本质满性。设 $\xi = (\mathcal{F}_i, \varphi_{ii'})$ 是拟凝聚层相对于覆盖
$\mathcal{U}$ 的下降数据。拉回此下降数据，得到拟凝聚层相对于 $V_j$
的覆盖 $\mathcal{V}_j$ 的下降数据 $\xi_j$。再次由 Lemma
\ref{descent-lemma-standard-fpqc-covering}，存在 $V_j$ 上的拟凝聚层
$\mathcal{F}_j$，使其相联系的典范下降数据同构于 $\xi_j$。由上文证明的
全性，存在同构
$$
\phi_{jj'} :
\mathcal{F}_j|_{V_{jj'}}
\longrightarrow
\mathcal{F}_{j'}|_{V_{jj'}}
$$
它对应于 $\xi_j$ 与 $\xi_{j'}$ 在 $\mathcal{V}_{jj'}$ 上的拉回之间的
下降数据同构。为证明这些态射 $\phi_{jj'}$ 满足余圈条件，在三重交
$V_{jj'j''}$ 上使用上文证明的忠实性。于是，由 Lemma
\ref{descent-lemma-zariski-descent-effective}，各层 $\mathcal{F}_j$ 如所需地粘合成
拟凝聚层 $\mathcal{F}$。还需验证：与 $\mathcal{F}$ 相联系的、相对于
$\mathcal{U}$ 的典范下降数据同构于最初的下降数据。此验证从略。
\end{proof}









\section{拟凝聚层的 Galois 下降}
\label{descent-section-galois-descent}

\noindent
拟凝聚层的 Galois 下降只是拟凝聚层 fpqc 下降的一个特殊情形。本节说明
如何把 Galois 下降转化为 fpqc 下降，再应用先前的结果得出结论。

\medskip\noindent
设 $k'/k$ 为域扩张。则 $\{\Spec(k') \to \Spec(k)\}$ 是 fpqc 覆盖。
设 $X$ 为 $k$ 上的概形。对 $k$-代数 $A$，令
$X_A = X \times_{\Spec(k)} \Spec(A)$。由 Topologies, Lemma
\ref{topologies-lemma-fpqc}，$\{X_{k'} \to X\}$ 是 fpqc 覆盖。注意
$$
X_{k'} \times_X X_{k'} = X_{k' \otimes_k k'}
\quad\text{且}\quad
X_{k'} \times_X X_{k'} \times_X X_{k'} = X_{k' \otimes_k k' \otimes_k k'}
$$
因此，拟凝聚层相对于 $\{X_{k'} \to X\}$ 的下降数据由如下资料给出：
$X_{k'}$ 上的一个拟凝聚层 $\mathcal{F}$，以及 $X_{k' \otimes_k k'}$
上的同构 $\varphi : \text{pr}_0^*\mathcal{F} \to \text{pr}_1^*\mathcal{F}$，
它在 $X_{k' \otimes_k k' \otimes_k k'}$ 上满足显然的余圈条件。
下面具体说明 Galois 扩张情形中这意味着什么。

\medskip\noindent
设 $k'/k$ 是 Galois 群为 $G = \text{Gal}(k'/k)$ 的有限 Galois 扩张。
则有 $k$-代数同构
$$
k' \otimes_k k' \longrightarrow \prod\nolimits_{\sigma \in G} k',\quad
a \otimes b \longrightarrow \prod a\sigma(b)
$$
并且
$$
k' \otimes_k k' \otimes_k k' \longrightarrow
\prod\nolimits_{(\sigma, \tau) \in G \times G} k',\quad
a \otimes b \otimes c \longrightarrow \prod a\sigma(b)\sigma(\tau(c))
$$
这里选取 $a\sigma(b)\sigma(\tau(c))$ 而非 $a\sigma(b)\tau(c)$，是因为
前者使下述公式更简洁，但这并非必需。给定 $\sigma \in G$，记
$$
f_\sigma = \text{id}_X \times \Spec(\sigma) :
X_{k'} \longrightarrow X_{k'}
$$
请注意，由于 $\Spec(-)$ 是反变函子，有
$f_{\sigma \tau} = f_\tau \circ f_\sigma$，而不是相反的次序。利用上面的
第一个同构，得到等同
$$
X_{k' \otimes_k k'} = \coprod\nolimits_{\sigma \in G} X_{k'}
$$
使得 $\text{pr}_0$ 对应于态射
$$
\coprod\nolimits_{\sigma \in G} X_{k'}
\xrightarrow{\coprod \text{id}}
X_{k'}
$$
而 $\text{pr}_1$ 对应于态射
$$
\coprod\nolimits_{\sigma \in G} X_{k'}
\xrightarrow{\coprod f_\sigma}
X_{k'}
$$
由此可见，$X_{k'}$ 上 $\mathcal{F}$ 的下降数据 $\varphi$ 对应于一族同构
$\varphi_\sigma : \mathcal{F} \to f_\sigma^*\mathcal{F}$。为写出余圈条件，
使用等同
$$
X_{k' \otimes_k k' \otimes_k k'} =
\coprod\nolimits_{(\sigma, \tau) \in G \times G} X_{k'}.
$$
它由上述代数同构得到。在此等同下，态射 $\text{pr}_{01}$ 对应于
$$
\coprod\nolimits_{(\sigma, \tau) \in G \times G} X_{k'}
\longrightarrow
\coprod\nolimits_{\sigma \in G} X_{k'}
$$
该态射通过恒等态射把指标为 $(\sigma, \tau)$ 的分支映到指标为
$\sigma$ 的分支。态射 $\text{pr}_{12}$ 对应于
$$
\coprod\nolimits_{(\sigma, \tau) \in G \times G} X_{k'}
\longrightarrow
\coprod\nolimits_{\sigma \in G} X_{k'}
$$
该态射通过 $f_\sigma$ 把指标为 $(\sigma, \tau)$ 的分支映到指标为
$\tau$ 的分支。最后，态射 $\text{pr}_{02}$ 对应于
$$
\coprod\nolimits_{(\sigma, \tau) \in G \times G} X_{k'}
\longrightarrow
\coprod\nolimits_{\sigma \in G} X_{k'}
$$
该态射通过恒等态射把指标为 $(\sigma, \tau)$ 的分支映到指标为
$\sigma\tau$ 的分支。因此余圈条件
$$
\text{pr}_{02}^*\varphi = \text{pr}_{12}^*\varphi \circ \text{pr}_{01}^*\varphi
$$
转化为每一对 $(\sigma, \tau)$ 所对应的条件
$$
\varphi_{\sigma\tau} = f_\sigma^*\varphi_\tau \circ \varphi_\sigma
$$
这是从 $\mathcal{F} \to f_{\sigma\tau}^*\mathcal{F}$ 的态射等式。
（一切恰好相容；例如，$\varphi_\sigma$ 的靶是
$f_\sigma^*\mathcal{F}$，而 $f_\sigma^*\varphi_\tau$ 的源也正是
$f_\sigma^*\mathcal{F}$。）

\begin{lemma}
\label{descent-lemma-galois-descent}
设 $k'/k$ 是 Galois 群为 $G$ 的（有限）Galois 扩张，$X$ 为 $k$ 上的
概形。拟凝聚 $\mathcal{O}_X$-模范畴等价于系统
$(\mathcal{F}, (\varphi_\sigma)_{\sigma \in G})$ 所成的范畴，其中
\begin{enumerate}
\item $\mathcal{F}$ 是 $X_{k'}$ 上的拟凝聚模；
\item $\varphi_\sigma : \mathcal{F} \to f_\sigma^*\mathcal{F}$ 是模同构；
\item 对所有 $\sigma, \tau \in G$，均有
$\varphi_{\sigma\tau} = f_\sigma^*\varphi_\tau \circ \varphi_\sigma$。
\end{enumerate}
这里 $f_\sigma = \text{id}_X \times \Spec(\sigma) : X_{k'} \to X_{k'}$。
\end{lemma}

\begin{proof}
如上所见，引理中的资料 $(\mathcal{F}, (\varphi_\sigma)_{\sigma \in G})$
等同于 fpqc 覆盖 $\{X_{k'} \to X\}$ 的下降数据。因此引理由 Proposition
\ref{descent-proposition-fpqc-descent-quasi-coherent} 得出。
\end{proof}

\noindent
上述结论有如下稍一般的版本。设有满的有限 \'etale 态射 $X \to Y$、
有限群 $G$ 以及群同态
$G^{opp} \to \text{Aut}_Y(X), \sigma \mapsto f_\sigma$，使态射
$$
G \times X \longrightarrow X \times_Y X,\quad
(\sigma, x) \longmapsto (x, f_\sigma(x))
$$
为同构。则上述相同结论成立。

\begin{lemma}
\label{descent-lemma-galois-descent-more-general}
设 $X \to Y$、$G$ 和 $f_\sigma : X \to X$ 如上。拟凝聚
$\mathcal{O}_Y$-模范畴等价于系统
$(\mathcal{F}, (\varphi_\sigma)_{\sigma \in G})$ 所成的范畴，其中
\begin{enumerate}
\item $\mathcal{F}$ 是拟凝聚 $\mathcal{O}_X$-模；
\item $\varphi_\sigma : \mathcal{F} \to f_\sigma^*\mathcal{F}$ 是模同构；
\item 对所有 $\sigma, \tau \in G$，均有
$\varphi_{\sigma\tau} = f_\sigma^*\varphi_\tau \circ \varphi_\sigma$。
\end{enumerate}
\end{lemma}

\begin{proof}
由于 $X \to Y$ 是满的有限 \'etale 态射，$\{X \to Y\}$ 是 fpqc 覆盖。
又因 $G \times X \to X \times_Y X$，$(\sigma, x) \mapsto (x, f_\sigma(x))$
为同构，可知 $G \times G \times X \to X \times_Y X \times_Y X$，
$(\sigma, \tau, x) \mapsto (x, f_\sigma(x), f_{\sigma\tau}(x))$ 也是同构。
利用这些等同，引理中资料的范畴等同于拟凝聚层相对于覆盖
% P05-ERR-0027
$\{x \to Y\}$ 的下降数据范畴。因此引理由 Proposition
\ref{descent-proposition-fpqc-descent-quasi-coherent} 得出。
\end{proof}









\section{模的有限性性质的下降}
\label{descent-section-descent-finiteness}

\noindent
本节证明：要检验拟凝聚模是否满足某些有限性条件，可以在覆盖的各成员上
进行检验。

\begin{lemma}
\label{descent-lemma-finite-type-descends}
设 $X$ 为概形，$\mathcal{F}$ 为拟凝聚 $\mathcal{O}_X$-模，并设
$\{f_i : X_i \to X\}_{i \in I}$ 是 fpqc 覆盖，使每个
$f_i^*\mathcal{F}$ 都是有限型 $\mathcal{O}_{X_i}$-模。则
$\mathcal{F}$ 是有限型 $\mathcal{O}_X$-模。
\end{lemma}

\begin{proof}
略。仿射情形见 Algebra, Lemma
\ref{algebra-lemma-descend-properties-modules}。
\end{proof}

\begin{lemma}
\label{descent-lemma-finite-type-descends-fppf}
设 $f : (X, \mathcal{O}_X) \to (Y, \mathcal{O}_Y)$ 是局部赋环空间态射，
$\mathcal{F}$ 是 $\mathcal{O}_Y$-模层。若
\begin{enumerate}
\item $f$ 作为拓扑空间映射是开映射；
\item $f$ 满且平坦；
\item $f^*\mathcal{F}$ 有限型，
\end{enumerate}
则 $\mathcal{F}$ 有限型。
\end{lemma}

\begin{proof}
取一点 $y \in Y$，并选取映至 $y$ 的点 $x \in X$。选取开集
$x \in U \subset X$ 以及 $f^*\mathcal{F}(U)$ 中在 $U$ 上生成
$f^*\mathcal{F}$ 的元素 $s_1, \ldots, s_n$。由于
$f^*\mathcal{F} =
f^{-1}\mathcal{F} \otimes_{f^{-1}\mathcal{O}_Y} \mathcal{O}_X$，
缩小 $U$ 后可假设 $s_i = \sum t_{ij} \otimes a_{ij}$，其中
$t_{ij} \in f^{-1}\mathcal{F}(U)$ 且 $a_{ij} \in \mathcal{O}_X(U)$。
进一步缩小 $U$，可假设 $t_{ij}$ 来自某个截面
$s_{ij} \in \mathcal{F}(V)$，其中 $V \subset Y$ 为满足
$f(U) \subset V$ 的开集。令 $N$ 为截面 $s_{ij}$ 的个数，并考虑态射
$$
\sigma = (s_{ij}) : \mathcal{O}_V^{\oplus N} \to \mathcal{F}|_V
$$
由这些截面的选取可知 $f^*\sigma|_U$ 为满射。因此对每个 $u \in U$，
态射
$$
\sigma_{f(u)} \otimes_{\mathcal{O}_{Y, f(u)}} \mathcal{O}_{X, u} :
\mathcal{O}_{X, u}^{\oplus N}
\longrightarrow
\mathcal{F}_{f(u)} \otimes_{\mathcal{O}_{Y, f(u)}} \mathcal{O}_{X, u}
$$
为满射。由于 $f$ 平坦，局部环同态
$\mathcal{O}_{Y, f(u)} \to \mathcal{O}_{X, u}$ 平坦，因而忠实平坦
（Algebra, Lemma \ref{algebra-lemma-local-flat-ff}）。故
$\sigma_{f(u)}$ 为满射。又因 $f$ 是开映射，$f(U)$ 是 $y$ 的开邻域，
证明完成。
\end{proof}

\begin{lemma}
\label{descent-lemma-finite-presentation-descends}
设 $X$ 为概形，$\mathcal{F}$ 为拟凝聚 $\mathcal{O}_X$-模，并设
$\{f_i : X_i \to X\}_{i \in I}$ 是 fpqc 覆盖，使每个
$f_i^*\mathcal{F}$ 都是有限表现 $\mathcal{O}_{X_i}$-模。则
$\mathcal{F}$ 是有限表现 $\mathcal{O}_X$-模。
\end{lemma}

\begin{proof}
略。仿射情形见 Algebra, Lemma
\ref{algebra-lemma-descend-properties-modules}。
\end{proof}

\begin{lemma}
\label{descent-lemma-locally-generated-by-r-sections-descends}
设 $X$ 为概形，$\mathcal{F}$ 为拟凝聚 $\mathcal{O}_X$-模，并设
$\{f_i : X_i \to X\}_{i \in I}$ 是 fpqc 覆盖，使每个
$f_i^*\mathcal{F}$ 作为 $\mathcal{O}_{X_i}$-模都局部由 $r$ 个截面生成。
则 $\mathcal{F}$ 作为 $\mathcal{O}_X$-模局部由 $r$ 个截面生成。
\end{lemma}

\begin{proof}
由 Lemma \ref{descent-lemma-finite-type-descends}，$\mathcal{F}$ 有限型。因此
Nakayama 引理（Algebra, Lemma \ref{algebra-lemma-NAK}）表明：
$\mathcal{F}$ 在一点 $x \in X$ 的某个邻域上由 $r$ 个截面生成，当且仅当
$\dim_{\kappa(x)} \mathcal{F}_x \otimes \kappa(x) \leq r$。选取一个 $i$
以及映至 $x$ 的点 $x_i \in X_i$。则
$\dim_{\kappa(x)} \mathcal{F}_x \otimes \kappa(x) = 
\dim_{\kappa(x_i)} (f_i^*\mathcal{F})_{x_i} \otimes \kappa(x_i)$
由于 $f_i^*\mathcal{F}$ 局部由 $r$ 个截面生成，右端为 $\leq r$。
\end{proof}

\begin{lemma}
\label{descent-lemma-flat-descends}
设 $X$ 为概形，$\mathcal{F}$ 为拟凝聚 $\mathcal{O}_X$-模，并设
$\{f_i : X_i \to X\}_{i \in I}$ 是 fpqc 覆盖，使每个
$f_i^*\mathcal{F}$ 都是平坦 $\mathcal{O}_{X_i}$-模。则 $\mathcal{F}$
是平坦 $\mathcal{O}_X$-模。
\end{lemma}

\begin{proof}
略。仿射情形见 Algebra, Lemma
\ref{algebra-lemma-descend-properties-modules}。
\end{proof}

\begin{lemma}
\label{descent-lemma-finite-locally-free-descends}
设 $X$ 为概形，$\mathcal{F}$ 为拟凝聚 $\mathcal{O}_X$-模，并设
$\{f_i : X_i \to X\}_{i \in I}$ 是 fpqc 覆盖，使每个
$f_i^*\mathcal{F}$ 都是有限局部自由 $\mathcal{O}_{X_i}$-模。则
$\mathcal{F}$ 是有限局部自由 $\mathcal{O}_X$-模。
\end{lemma}

\begin{proof}
这由如下事实得出：拟凝聚层有限局部自由，当且仅当它有限表现且平坦；
见 Algebra, Lemma \ref{algebra-lemma-finite-projective}。具体而言，若每个
$f_i^*\mathcal{F}$ 都平坦且有限表现，则由 Lemmas
\ref{descent-lemma-flat-descends} 和 \ref{descent-lemma-finite-presentation-descends}，
$\mathcal{F}$ 也如此。
\end{proof}

\noindent
局部投射拟凝聚层的定义见 Properties, Section
\ref{properties-section-locally-projective}。

\begin{lemma}
\label{descent-lemma-locally-projective-descends}
设 $X$ 为概形，$\mathcal{F}$ 为拟凝聚 $\mathcal{O}_X$-模，并设
$\{f_i : X_i \to X\}_{i \in I}$ 是 fpqc 覆盖，使每个
$f_i^*\mathcal{F}$ 都是局部投射 $\mathcal{O}_{X_i}$-模。则
$\mathcal{F}$ 是局部投射 $\mathcal{O}_X$-模。
\end{lemma}

\begin{proof}
略。对 Zariski 覆盖，这是 Properties, Lemma
\ref{properties-lemma-locally-projective}；仿射情形则是 Algebra, Theorem
\ref{algebra-theorem-ffdescent-projectivity}。
\end{proof}

\begin{remark}
\label{descent-remark-locally-free-descends}
局部自由性是拟凝聚模的一个不在 fpqc 拓扑中下降的性质。事实上，设
$R$ 为环，$M$ 为投射 $R$-模，它是秩 1 局部自由模的可数直和
$M = \bigoplus L_n$，但自身并非局部自由；见 Examples, Lemma
\ref{examples-lemma-projective-not-locally-free}。作忠实平坦基变换
$$
R \longrightarrow
\bigoplus\nolimits_{m \geq 1}
\bigoplus\nolimits_{(i_1, \ldots, i_m) \in \mathbf{Z}^{\oplus m}}
L_1^{\otimes i_1} \otimes_R \ldots \otimes_R L_m^{\otimes i_m}
$$
后，$M$ 成为自由模。但我们不知道 fppf 覆盖的情形如何。换言之，
我们不知道下述问题的答案：设 $A \to B$ 是有限表现的忠实平坦环同态，
$M$ 是满足 $M \otimes_A B$ 自由的 $A$-模。$M$ 是否为局部自由
$A$-模？事实证明，若 $A$ 是 Noether 环，则答案肯定；这由
\cite{Bass} 的结果得出。但一般情形的答案仍未知。若你知道答案或有参考
文献，请发送邮件至
\href{mailto:stacks.project@gmail.com}{stacks.project@gmail.com}.
\end{remark}

\noindent
这里还补充两个与上述结果相关、但性质略有不同的结果。

\begin{lemma}
\label{descent-lemma-finite-over-finite-module}
设 $f : X \to Y$ 为概形态射，$\mathcal{F}$ 为拟凝聚
$\mathcal{O}_X$-模，并假设 $f$ 是有限态射。则 $\mathcal{F}$ 是有限型
$\mathcal{O}_X$-模，当且仅当 $f_*\mathcal{F}$ 是有限型
$\mathcal{O}_Y$-模。
\end{lemma}

\begin{proof}
由于 $f$ 有限，它是仿射的。这把问题化归到如下情形：$f$ 是由有限环同态
$A \to B$ 给出的态射 $\Spec(B) \to \Spec(A)$。此外，此时
$\mathcal{F} = \widetilde{M}$ 是与 $B$-模 $M$ 相联系的模层。注意，
$M$ 作为 $B$-模有限，当且仅当 $M$ 作为 $A$-模有限；见 Algebra,
Lemma \ref{algebra-lemma-finite-module-over-finite-extension}。结合
Properties, Lemma \ref{properties-lemma-finite-type-module} 即得引理。
\end{proof}

\begin{lemma}
\label{descent-lemma-finite-finitely-presented-module}
设 $f : X \to Y$ 为概形态射，$\mathcal{F}$ 为拟凝聚
$\mathcal{O}_X$-模，并假设 $f$ 有限且有限表现。则 $\mathcal{F}$ 是有限表现
$\mathcal{O}_X$-模，当且仅当 $f_*\mathcal{F}$ 是有限表现
$\mathcal{O}_Y$-模。
\end{lemma}

\begin{proof}
由于 $f$ 有限，它是仿射的。这把问题化归到如下情形：$f$ 是由有限且
有限表现的环同态 $A \to B$ 给出的态射 $\Spec(B) \to \Spec(A)$。
此外，此时 $\mathcal{F} = \widetilde{M}$ 是与 $B$-模 $M$ 相联系的模层。
注意，$M$ 作为 $B$-模有限表现，当且仅当 $M$ 作为 $A$-模有限表现；
见 Algebra, Lemma
\ref{algebra-lemma-finite-finitely-presented-extension}。结合 Properties,
Lemma \ref{properties-lemma-finite-presentation-module} 即得引理。
\end{proof}
















\section{拟凝聚层与拓扑（一）}
\label{descent-section-quasi-coherent-sheaves}

\noindent
本节的结果表明：概形 $S$ 上的拟凝聚模范畴，与拓扑一章中和 $S$
相联系的许多位点上的拟凝聚模范畴之间存在自然等价。

\medskip\noindent
设 $S$ 为概形，$\mathcal{F}$ 为拟凝聚 $\mathcal{O}_S$-模。考虑函子
\begin{equation}
\label{descent-equation-quasi-coherent-presheaf}
(\Sch/S)^{opp} \longrightarrow \textit{Ab},
\quad
(f : T \to S) \longmapsto \Gamma(T, f^*\mathcal{F}).
\end{equation}

\begin{lemma}
\label{descent-lemma-sheaf-condition-holds}
设 $S$ 为概形，$\mathcal{F}$ 为拟凝聚 $\mathcal{O}_S$-模，并设
$\tau \in \{Zariski, \linebreak[0] \etale, \linebreak[0] smooth,
\linebreak[0] syntomic, \linebreak[0] fppf, \linebreak[0] fpqc\}$。
(\ref{descent-equation-quasi-coherent-presheaf}) 中定义的函子，对 $S$ 上任意概形
$T$ 的任意 $\tau$-覆盖 $\{T_i \to T\}_{i \in I}$ 都满足层条件。
\end{lemma}

\begin{proof}
对 $\tau \in \{Zariski, \linebreak[0] \etale, \linebreak[0] smooth,
\linebreak[0] syntomic, \linebreak[0] fppf\}$，任意 $\tau$-覆盖也是 fpqc
覆盖；见 Topologies, Lemmas
\ref{topologies-lemma-zariski-etale},
\ref{topologies-lemma-zariski-etale-smooth},
\ref{topologies-lemma-zariski-etale-smooth-syntomic},
\ref{topologies-lemma-zariski-etale-smooth-syntomic-fppf}，以及
\ref{topologies-lemma-zariski-etale-smooth-syntomic-fppf-fpqc}.
因此只需对 fpqc 覆盖证明结论。设
$\{f_i : T_i \to T\}_{i \in I}$ 是 fpqc 覆盖，且给定
$f : T \to S$。假设有一族截面
$s_i \in \Gamma(T_i , f_i^*f^*\mathcal{F})$，满足
$s_i|_{T_i \times_T T_j} = s_j|_{T_i \times_T T_j}$。需要找出相应截面
$s \in \Gamma(T, f^*\mathcal{F})$。可以把 $s_i$ 重新解释为一族态射
$\varphi_i : f_i^*\mathcal{O}_T = \mathcal{O}_{T_i} \to f_i^*f^*\mathcal{F}$，
它与 $T$ 上拟凝聚层 $\mathcal{O}_T$ 和 $f^*\mathcal{F}$ 相联系的典范
下降数据相容。因此，由 Proposition
\ref{descent-proposition-fpqc-descent-quasi-coherent}，可将它们（唯一地）下降为
态射 $\mathcal{O}_T \to f^*\mathcal{F}$，后者给出所需截面 $s$。
\end{proof}

\noindent
特别地，可以作如下定义。

\begin{definition}
\label{descent-definition-structure-sheaf}
设 $\tau \in \{Zariski, \linebreak[0] \etale, \linebreak[0]
smooth, \linebreak[0] syntomic, \linebreak[0] fppf\}$，$S$ 为概形，
$\Sch_\tau$ 为包含 $S$ 的大位点，且 $\mathcal{F}$ 为拟凝聚
$\mathcal{O}_S$-模。
\begin{enumerate}
\item 大位点 $(\Sch/S)_\tau$ 的{\it 结构层}是环层
$T/S \mapsto \Gamma(T, \mathcal{O}_T)$，记为 $\mathcal{O}$ 或
$\mathcal{O}_S$。
\item 若 $\tau = Zariski$ 或 $\tau = \etale$，则小位点 $S_{Zar}$ 或
$S_\etale$ 的{\it 结构层}是环层 $T/S \mapsto \Gamma(T, \mathcal{O}_T)$，
记为 $\mathcal{O}$ 或 $\mathcal{O}_S$。
\item 大位点 $(\Sch/S)_\tau$ 上{\it 与 $\mathcal{F}$ 相联系的
$\mathcal{O}$-模层}是 $\mathcal{O}$-模层
$(f : T \to S) \mapsto \Gamma(T, f^*\mathcal{F})$，记为
$\mathcal{F}^a$（也常简记为 $\mathcal{F}$）。
\item 若 $\tau = Zariski$ 或 $\tau = \etale$，则小位点 $S_{Zar}$ 或
$S_\etale$ 上{\it 与 $\mathcal{F}$ 相联系的 $\mathcal{O}$-模层}是
$\mathcal{O}$-模层 $(f : T \to S) \mapsto \Gamma(T, f^*\mathcal{F})$，
记为 $\mathcal{F}^a$（也常简记为 $\mathcal{F}$）。
\end{enumerate}
\end{definition}

\noindent
注意，各种情形中都使用同一记号 $\mathcal{F}^a$。这不会真正造成混淆，
因为按定义，给出层 $\mathcal{F}^a$ 的规则不依赖于所考察的位点。

\begin{remark}
\label{descent-remark-Zariski-site-space}
在 Topologies, Lemma \ref{topologies-lemma-Zariski-usual} 中已经看到：
概形 $S$ 的小 Zariski 位点与作为拓扑空间的 $S$ 等价，其意义是相应的
层范畴自然等价。现在 $S_{Zar}$ 还配备了结构层 $\mathcal{O}$，由此可见，
赋环位点 $(S_{Zar}, \mathcal{O})$ 上的模层等同于赋环空间
$(S, \mathcal{O}_S)$ 上的模层。
\end{remark}

\begin{remark}
\label{descent-remark-change-topologies-ringed}
设 $f : T \to S$ 为概形态射。Topologies, Section
\ref{topologies-section-change-topologies} 所列的每个位点态射
$f_{sites}$ 都成为赋环位点态射。具体而言，每个此类位点态射
$f_{sites} : (\Sch/T)_\tau \to (\Sch/S)_{\tau'}$ 或
$f_{sites} : (\Sch/S)_\tau \to S_{\tau'}$ 均由连续函子
$S'/S \mapsto T \times_S S'/S$ 给出。因此，给定 $S'/S$，令
$$
f_{sites}^\sharp :
\mathcal{O}(S'/S)
\longrightarrow
f_{sites, *}\mathcal{O}(S'/S) =
\mathcal{O}(T \times_S S'/T)
$$
为通常的态射
$\text{pr}_{S'}^\sharp : \mathcal{O}(S') \to \mathcal{O}(T \times_S S')$。
类似地，态射
$i_f : \Sh(T_\tau) \to \Sh((\Sch/S)_\tau)$
（其中 $\tau \in \{Zar, \etale\}$；见 Topologies, Lemmas
\ref{topologies-lemma-put-in-T} 和 \ref{topologies-lemma-put-in-T-etale}）
成为赋环拓扑斯态射，因为 $i_f^{-1}\mathcal{O} = \mathcal{O}$。下面列出
若干特殊情形：
\begin{enumerate}
\item 大位点态射 $f_{big} : (\Sch/X)_{fppf} \to (\Sch/Y)_{fppf}$
成为赋环位点态射
$$
(f_{big}, f_{big}^\sharp) :
((\Sch/X)_{fppf}, \mathcal{O}_X)
\longrightarrow
((\Sch/Y)_{fppf}, \mathcal{O}_Y)
$$
如 Modules on Sites, Definition \ref{sites-modules-definition-ringed-site}
所述。大 syntomic、光滑、\'etale 和 Zariski 位点的情形类似。
\item 小位点态射 $f_{small} : X_\etale \to Y_\etale$ 成为赋环位点态射
$$
(f_{small}, f_{small}^\sharp) :
(X_\etale, \mathcal{O}_X)
\longrightarrow
(Y_\etale, \mathcal{O}_Y)
$$
如 Modules on Sites, Definition \ref{sites-modules-definition-ringed-site}
所述。小 Zariski 位点的情形类似。
\end{enumerate}
\end{remark}

\noindent
设 $S$ 为概形。显然，给定（例如）$(\Sch/S)_{Zar}$ 上的
$\mathcal{O}$-模，其到（例如）$(\Sch/S)_{fppf}$ 的拉回正是 fppf 层化。
为说明比较大位点与小位点时的情形，有如下引理。

\begin{lemma}
\label{descent-lemma-compare-sites}
设 $S$ 为概形。以
$$
\begin{matrix}
\text{id}_{\tau, Zar} & : & (\Sch/S)_\tau \to S_{Zar}, &
\tau \in \{Zar, \etale, smooth, syntomic, fppf\} \\
\text{id}_{\tau, \etale} & : &
(\Sch/S)_\tau \to S_\etale, &
\tau \in \{\etale, smooth, syntomic, fppf\} \\
\text{id}_{small, \etale, Zar} & : & S_\etale \to S_{Zar},
\end{matrix}
$$
表示 Remark \ref{descent-remark-change-topologies-ringed} 中的赋环位点态射。
设 $\mathcal{F}$ 为 $\mathcal{O}_S$-模层，并把它视为 $S_{Zar}$ 上的
$\mathcal{O}$-模层。则
\begin{enumerate}
\item $(\text{id}_{\tau, Zar})^*\mathcal{F}$ 是 $(\Sch/S)_\tau$ 上
Zariski 层
$$
(f : T \to S) \longmapsto \Gamma(T, f^*\mathcal{F})
$$
的 $\tau$-层化；
\item $(\text{id}_{small, \etale, Zar})^*\mathcal{F}$ 是 $S_\etale$ 上
Zariski 层
$$
(f : T \to S) \longmapsto \Gamma(T, f^*\mathcal{F})
$$
的 \'etale 层化。
\end{enumerate}
设 $\mathcal{G}$ 为 $S_\etale$ 上的 $\mathcal{O}$-模层。则
\begin{enumerate}
\item[(3)] $(\text{id}_{\tau, \etale})^*\mathcal{G}$ 是 \'etale 层
$$
(f : T \to S) \longmapsto \Gamma(T, f_{small}^*\mathcal{G})
$$
的 $\tau$-层化，其中 $f_{small} : T_\etale \to S_\etale$ 是 Remark
\ref{descent-remark-change-topologies-ringed} 中赋环小 \'etale 位点的态射。
\end{enumerate}
\end{lemma}

\begin{proof}
证明 (1)。首先注意，当 $\tau = Zar$ 时结论成立，因为此时有拓扑斯态射
$i_f : \Sh(T_{Zar}) \to \Sh((\Sch/S)_{Zar})$，并且作为
$T_{Zar} \to S_{Zar}$ 的态射，有
$\text{id}_{\tau, Zar} \circ i_f = f_{small}$；见 Topologies, Lemmas
\ref{topologies-lemma-put-in-T} 和
\ref{topologies-lemma-morphism-big-small}。由于拉回具有传递性（见
Modules on Sites, Lemma
\ref{sites-modules-lemma-push-pull-composition-modules}），可知如所需，
$i_f^*(\text{id}_{\tau, Zar})^*\mathcal{F} = f_{small}^*\mathcal{F}$。
因此，由本引理前的评注，$(\text{id}_{\tau, Zar})^*\mathcal{F}$ 是预层
$T \mapsto \Gamma(T, f^*\mathcal{F})$ 的 $\tau$-层化。

\medskip\noindent
(3) 的证明与 (1) 的证明完全相同，只是使用 Topologies, Lemmas
\ref{topologies-lemma-put-in-T-etale} 和
\ref{topologies-lemma-morphism-big-small-etale}。(2) 的证明从略。
\end{proof}

\begin{remark}
\label{descent-remark-change-topologies-ringed-sites}
Remark \ref{descent-remark-change-topologies-ringed} 和
Lemma \ref{descent-lemma-compare-sites} 有如下应用：
\begin{enumerate}
\item 设 $S$ 为概形。构造 $\mathcal{F} \mapsto \mathcal{F}^a$ 是沿赋环位点态射
$\text{id}_{\tau, Zar} : ((\Sch/S)_\tau, \mathcal{O})
\to (S_{Zar}, \mathcal{O})$
或态射
$\text{id}_{small, \etale, Zar} :
(S_\etale, \mathcal{O}) \to (S_{Zar}, \mathcal{O})$ 的拉回。
\item 设 $f : X \to Y$ 为概形态射。对 Remark
\ref{descent-remark-change-topologies-ringed} 中任一赋环位点态射 $f_{sites}$，都有
$$
(f^*\mathcal{F})^a = f_{sites}^*\mathcal{F}^a.
$$
这由 (1) 以及拉回与赋环位点态射的复合相容这一事实得出；见
Modules on Sites, Lemma
\ref{sites-modules-lemma-push-pull-composition-modules}。
\end{enumerate}
\end{remark}

\begin{lemma}
\label{descent-lemma-quasi-coherent-gives-quasi-coherent}
设 $S$ 为概形，$\mathcal{F}$ 为拟凝聚 $\mathcal{O}_S$-模，并设
$\tau \in \{Zariski, \linebreak[0] \etale, \linebreak[0]
smooth, \linebreak[0] syntomic, \linebreak[0] fppf\}$。
\begin{enumerate}
\item 层 $\mathcal{F}^a$ 是 $(\Sch/S)_\tau$ 上的拟凝聚 $\mathcal{O}$-模，
其定义见 Modules on Sites, Definition
\ref{sites-modules-definition-site-local}。
\item 若 $\tau = Zariski$ 或 $\tau = \etale$，则层 $\mathcal{F}^a$ 是
$S_{Zar}$ 或 $S_\etale$ 上的拟凝聚 $\mathcal{O}$-模，其定义见
Modules on Sites, Definition \ref{sites-modules-definition-site-local}。
\end{enumerate}
\end{lemma}

\begin{proof}
选取 Zariski 覆盖 $\{S_i \to S\}$，使得有正合列
$$
\bigoplus\nolimits_{k \in K_i} \mathcal{O}_{S_i} \longrightarrow
\bigoplus\nolimits_{j \in J_i} \mathcal{O}_{S_i} \longrightarrow
\mathcal{F}|_{S_i} \longrightarrow 0
$$
其中 $K_i$ 和 $J_i$ 是某些指标集。由赋环空间上拟凝聚层的定义，
这样的覆盖存在（见 Modules, Definition
\ref{modules-definition-quasi-coherent}）。

\medskip\noindent
证明 (1)。设 $\tau \in \{Zariski, \linebreak[0] fppf, \linebreak[0]
\etale, \linebreak[0] smooth, \linebreak[0] syntomic\}$。显然，
$\mathcal{F}^a|_{(\Sch/S_i)_\tau}$ 也位于正合列
$$
\bigoplus\nolimits_{k \in K_i} \mathcal{O}|_{(\Sch/S_i)_\tau}
\longrightarrow
\bigoplus\nolimits_{j \in J_i} \mathcal{O}|_{(\Sch/S_i)_\tau}
\longrightarrow
\mathcal{F}^a|_{(\Sch/S_i)_\tau} \longrightarrow 0
$$
中。因此，由 Modules on Sites, Lemma
\ref{sites-modules-lemma-local-final-object}，$\mathcal{F}^a$ 拟凝聚。

\medskip\noindent
证明 (2)。令 $\tau = \etale$。显然，$\mathcal{F}^a|_{(S_i)_\etale}$
也位于正合列
$$
\bigoplus\nolimits_{k \in K_i} \mathcal{O}|_{(S_i)_\etale}
\longrightarrow
\bigoplus\nolimits_{j \in J_i} \mathcal{O}|_{(S_i)_\etale}
\longrightarrow
\mathcal{F}^a|_{(S_i)_\etale} \longrightarrow 0
$$
中。因此，由 Modules on Sites, Lemma
\ref{sites-modules-lemma-local-final-object}，$\mathcal{F}^a$ 拟凝聚。
$\tau = Zariski$ 的情形类似（事实上，由于相应的赋环拓扑斯相同，
这只是同义反复）。
\end{proof}

\begin{lemma}
\label{descent-lemma-fully-faithful-associated}
设 $S$ 为概形，并设
$\tau \in \{Zariski, \linebreak[0] \etale, \linebreak[0]
smooth, \linebreak[0] syntomic, \linebreak[0] fppf\}$。Definition
\ref{descent-definition-structure-sheaf} 中每个函子
$\mathcal{F} \mapsto \mathcal{F}^a$，即
$$
\QCoh(\mathcal{O}_S) \to \QCoh((\Sch/S)_\tau, \mathcal{O})
\quad\text{或}\quad
\QCoh(\mathcal{O}_S) \to \QCoh(S_\tau, \mathcal{O})
$$
都是全忠实的。
\end{lemma}

\begin{proof}
（由 Lemma \ref{descent-lemma-quasi-coherent-gives-quasi-coherent}，确实得到所示的
函子。）我们把 $S$ 上的 $\mathcal{O}_S$-模与
$(S_{Zar}, \mathcal{O}_S)$ 上的模等同。对拟凝聚模 $\mathcal{F}$，函子
$\mathcal{F} \mapsto \mathcal{F}^a$ 由沿某个赋环位点态射 $f$ 的拉回给出；
见 Remark \ref{descent-remark-change-topologies-ringed-sites}。在每种情形中，
函子 $f_*$ 均由沿嵌入函子 $S_{Zar} \to S_\tau$ 或
$S_{Zar} \to (\Sch/S)_\tau$ 的限制给出（这些位点态射的定义方式见
Topologies, Section \ref{topologies-section-change-topologies}）。结合
$f^*\mathcal{F} = \mathcal{F}^a$ 的描述可知，只要 $\mathcal{F}$ 拟凝聚，
就有 $f_*f^*\mathcal{F} = \mathcal{F}$。于是
$$
\Hom_\mathcal{O}(\mathcal{F}^a, \mathcal{G}^a) =
\Hom_\mathcal{O}(f^*\mathcal{F}, f^*\mathcal{G}) =
\Hom_{\mathcal{O}_S}(\mathcal{F}, f_*f^*\mathcal{G}) =
\Hom_{\mathcal{O}_S}(\mathcal{F}, \mathcal{G})
$$
如所需成立。
\end{proof}

\begin{proposition}
\label{descent-proposition-equivalence-quasi-coherent}
设 $S$ 为概形，并设
$\tau \in \{Zariski, \linebreak[0] \etale, \linebreak[0]
smooth, \linebreak[0] syntomic, \linebreak[0] fppf\}$。
\begin{enumerate}
\item 函子 $\mathcal{F} \mapsto \mathcal{F}^a$ 定义范畴等价
$$
\QCoh(\mathcal{O}_S)
\longrightarrow
\QCoh((\Sch/S)_\tau, \mathcal{O})
$$
它连接 $S$ 上的拟凝聚层范畴与 $S$ 的大 $\tau$-位点上的拟凝聚
$\mathcal{O}$-模范畴。
\item 令 $\tau = Zariski$ 或 $\tau = \etale$。函子
$\mathcal{F} \mapsto \mathcal{F}^a$ 定义范畴等价
$$
\QCoh(\mathcal{O}_S)
\longrightarrow
\QCoh(S_\tau, \mathcal{O})
$$
它连接 $S$ 上的拟凝聚层范畴与 $S$ 的小 $\tau$-位点上的拟凝聚
$\mathcal{O}$-模范畴。
\end{enumerate}
\end{proposition}

\begin{proof}
Lemma \ref{descent-lemma-quasi-coherent-gives-quasi-coherent} 已说明该函子有良好定义；
Lemma \ref{descent-lemma-fully-faithful-associated} 已说明它全忠实。为完成证明，
将证明 $(\Sch/S)_\tau$ 上的拟凝聚 $\mathcal{O}$-模给出拟凝聚层相对于
$S$ 的某个 $\tau$-覆盖的下降数据。构造出此下降数据后，应用 Proposition
\ref{descent-proposition-fpqc-descent-quasi-coherent}，得到 $S$ 上相应的拟凝聚层。

\medskip\noindent
设 $\mathcal{G}$ 是 $S$ 的大 $\tau$-位点上的拟凝聚 $\mathcal{O}$-模。
由 Modules on Sites, Definition
\ref{sites-modules-definition-site-local}，存在 $S$ 的 $\tau$-覆盖
$\{S_i \to S\}_{i \in I}$，使每个限制
$\mathcal{G}|_{(\Sch/S_i)_\tau}$ 都有整体表现
$$
\bigoplus\nolimits_{k \in K_i} \mathcal{O}|_{(\Sch/S_i)_\tau}
\longrightarrow
\bigoplus\nolimits_{j \in J_i} \mathcal{O}|_{(\Sch/S_i)_\tau}
\longrightarrow
\mathcal{G}|_{(\Sch/S_i)_\tau} \longrightarrow 0
$$
其中 $J_i$ 和 $K_i$ 是某些指标集。我们断言，这蕴含
$\mathcal{G}|_{(\Sch/S_i)_\tau}$ 等于 $S_i$ 上某个拟凝聚层
$\mathcal{F}_i$ 所对应的 $\mathcal{F}_i^a$。事实上，这对直和
$\bigoplus\nolimits_{k \in K_i} \mathcal{O}|_{(\Sch/S_i)_\tau}$
以及
$\bigoplus\nolimits_{j \in J_i} \mathcal{O}|_{(\Sch/S_i)_\tau}$.
是显然的。因此，$\mathcal{G}|_{(\Sch/S_i)_\tau}$ 是某个态射
$\varphi : \mathcal{K}_i^a \to \mathcal{L}_i^a$ 的余核，其中
$\mathcal{K}_i$、$\mathcal{L}_i$ 是 $S_i$ 上的拟凝聚层。由
$(\ )^a$ 的全忠实性，$\varphi = \phi^a$，其中
$\phi : \mathcal{K}_i \to \mathcal{L}_i$ 是 $S_i$ 上的拟凝聚层态射。
于是显然
$\mathcal{G}|_{(\Sch/S_i)_\tau} \cong \Coker(\phi)^a$
如所断言。

\medskip\noindent
由于 $\mathcal{G}$ 定义在整个范畴 $(\Sch/S)_\tau$ 上，可知
$$
% P05-ERR-0028
(\text{pr}_0^*\mathcal{F}_i)^a
\cong
\mathcal{G}|_{(\Sch/(S_i \times_S S_j))_\tau}
\cong
(\text{pr}_1^*\mathcal{F})^a
$$
是 $(\Sch/(S_i \times_S S_j))_\tau$ 上的 $\mathcal{O}$-模同构。
因此，再次利用全忠实性，得到典范同构
$$
\phi_{ij} :
\text{pr}_0^*\mathcal{F}_i
\longrightarrow
\text{pr}_1^*\mathcal{F}_j
$$
它们是 $S_i \times_S S_j$ 上的拟凝聚模同构。省略它们满足余圈条件的验证。
既然它们确实满足该条件，由拟凝聚层相对于覆盖 $\{S_i \to S\}$ 的下降
有效性（Proposition \ref{descent-proposition-fpqc-descent-quasi-coherent}），
存在 $S$ 上的拟凝聚层 $\mathcal{F}$，使
$\mathcal{F}|_{S_i} \cong \mathcal{F}_i$ 与给定下降数据相容。换言之，
给定 $\mathcal{O}$-模同构
$$
\phi_i :
\mathcal{F}^a|_{(\Sch/S_i)_\tau}
\longrightarrow
\mathcal{G}|_{(\Sch/S_i)_\tau}
$$
它们在 $S_i \times_S S_j$ 上相符。由于
$\SheafHom_\mathcal{O}(\mathcal{F}^a, \mathcal{G})$ 是层（Modules on Sites,
Lemma \ref{sites-modules-lemma-internal-hom}），可知存在恢复上述同构
$\phi_i$ 的 $\mathcal{O}$-模态射 $\mathcal{F}^a \to \mathcal{G}$。
因此该态射是同构，结论成立。

\medskip\noindent
位点 $S_\etale$ 和 $S_{Zar}$ 的情形以完全相同的方式证明。
\end{proof}

\begin{lemma}
\label{descent-lemma-equivalence-quasi-coherent-properties}
设 $S$ 为概形，并设
$\tau \in \{Zariski, \linebreak[0] \etale, \linebreak[0]
smooth, \linebreak[0] syntomic, \linebreak[0] fppf\}$。令 $\mathcal{P}$
为 Modules on Sites, Definitions \ref{sites-modules-definition-global}、
\ref{sites-modules-definition-site-local} 和
\ref{sites-modules-definition-flat} 中定义的模的性质之一\footnote{列表如下：
自由、有限自由、由整体截面生成、由 $r$ 个整体截面生成、由有限多个整体
截面生成、具有整体表现、具有整体有限表现、局部自由、有限局部自由、
局部由截面生成、局部由 $r$ 个截面生成、有限型、有限表现、凝聚或平坦。}。
Proposition \ref{descent-proposition-equivalence-quasi-coherent} 中由规则
$\mathcal{F} \mapsto \mathcal{F}^a$ 定义的范畴等价
$$
\QCoh(\mathcal{O}_S)
\longrightarrow
\QCoh((\Sch/S)_\tau, \mathcal{O})
\quad\text{且}\quad
\QCoh(\mathcal{O}_S)
\longrightarrow
\QCoh(S_\tau, \mathcal{O})
$$
具有性质
$$
\mathcal{F}\text{ 具有 }\mathcal{P}
\Leftrightarrow
\mathcal{F}^a\text{ 具有 }\mathcal{P}\text{ 作为一个 }\mathcal{O}\text{-模}
$$
但 $\mathcal{P}$ 为“局部自由”或“凝聚”时可能例外。若
$\mathcal{P}=$“凝聚”，且 $S$ 局部 Noether、$\tau$ 为 Zariski 或
\'etale，则等价 $\QCoh(\mathcal{O}_S) \to \QCoh(S_\tau, \mathcal{O})$
满足上述性质。
\end{lemma}

\begin{proof}
对整体性质，即 Modules on Sites, Definition
\ref{sites-modules-definition-global} 中定义的性质，结论是直接的。
对局部性质，可以利用 Modules on Sites, Lemma
\ref{sites-modules-lemma-local-final-object}，把“$\mathcal{F}^a$ 具有
$\mathcal{P}$”转化为 $X$ 的某个覆盖的各成员上的性质。因此，结论由
Lemmas \ref{descent-lemma-finite-type-descends}、
\ref{descent-lemma-finite-presentation-descends}、
\ref{descent-lemma-locally-generated-by-r-sections-descends}、
\ref{descent-lemma-flat-descends} 和 \ref{descent-lemma-finite-locally-free-descends} 得出。
在局部 Noether 概形上，拟凝聚模凝聚等价于它有限型（见 Cohomology of
Schemes, Lemma \ref{coherent-lemma-coherent-Noetherian}），故这化归到
有限型模的情形（细节从略）。
\end{proof}










\section{拟凝聚模的上同调与拓扑}
\label{descent-section-quasi-coherent-cohomology}

\noindent
本节证明拟凝聚模的上同调不依赖于拓扑的选择。

\begin{lemma}
\label{descent-lemma-standard-covering-Cech}
设 $S$ 为概形。设
\begin{enumerate}
\item[(a)] $\tau \in \{Zariski, \linebreak[0] fppf, \linebreak[0]
\etale, \linebreak[0] smooth, \linebreak[0] syntomic\}$
且 $\mathcal{C} = (\Sch/S)_\tau$；或
\item[(b)] $\tau = \etale$ 且 $\mathcal{C} = S_\etale$；或
\item[(c)] $\tau = Zariski$ 且 $\mathcal{C} = S_{Zar}$。
\end{enumerate}
设 $\mathcal{F}$ 为 $\mathcal{C}$ 上的 Abel 层，
$U \in \Ob(\mathcal{C})$ 为仿射对象，并设
$\mathcal{U} = \{U_i \to U\}_{i = 1, \ldots, n}$ 是 $\mathcal{C}$ 中的
标准仿射 $\tau$-覆盖。则
\begin{enumerate}
\item $\mathcal{V} = \{\coprod_{i = 1, \ldots, n} U_i \to U\}$ 是 $U$ 的
$\tau$-覆盖；
\item $\mathcal{U}$ 是 $\mathcal{V}$ 的加细；
\item 诱导的 {\v C}ech 复形态射（Cohomology on Sites,
Equation (\ref{sites-cohomology-equation-map-cech-complexes})）
$$
\check{\mathcal{C}}^\bullet(\mathcal{V}, \mathcal{F})
\longrightarrow
\check{\mathcal{C}}^\bullet(\mathcal{U}, \mathcal{F})
$$
是复形同构。
\end{enumerate}
\end{lemma}

\begin{proof}
这是因为
$$
(\coprod\nolimits_{i_0 = 1, \ldots, n} U_{i_0}) \times_U
\ldots \times_U
(\coprod\nolimits_{i_p = 1, \ldots, n} U_{i_p})
=
\coprod\nolimits_{i_0, \ldots, i_p \in \{1, \ldots, n\}}
U_{i_0} \times_U \ldots \times_U U_{i_p}
$$
并且由于不交并是 $\tau$-覆盖，有
$\mathcal{F}(\coprod_a V_a) = \prod_a \mathcal{F}(V_a)$。
\end{proof}

\begin{lemma}
\label{descent-lemma-standard-covering-Cech-quasi-coherent}
设 $S$ 为概形，$\mathcal{F}$ 为 $S$ 上的拟凝聚层，并设
$\tau$、$\mathcal{C}$、$U$、$\mathcal{U}$ 如 Lemma
\ref{descent-lemma-standard-covering-Cech}。则有复形同构
$$
\check{\mathcal{C}}^\bullet(\mathcal{U}, \mathcal{F}^a)
\cong
s((A/R)_\bullet \otimes_R M)
$$
（见 Section \ref{descent-section-descent-modules}），其中
$R = \Gamma(U, \mathcal{O}_U)$、$M = \Gamma(U, \mathcal{F}^a)$，且
$R \to A$ 是忠实平坦环同态。特别地，
$$
\check{H}^p(\mathcal{U}, \mathcal{F}^a) = 0
$$
对所有 $p \geq 1$ 均成立。
\end{lemma}

\begin{proof}
由 Lemma \ref{descent-lemma-standard-covering-Cech}，
$\check{\mathcal{C}}^\bullet(\mathcal{U}, \mathcal{F}^a)$ 同构于
$\check{\mathcal{C}}^\bullet(\mathcal{V}, \mathcal{F}^a)$，其中
$\mathcal{V} = \{V \to U\}$ 且
$V = \coprod_{i = 1, \ldots n} U_i$ 也仿射。令
$A = \Gamma(V, \mathcal{O}_V)$。由于 $\{V \to U\}$ 是
$\tau$-覆盖，$R \to A$ 忠实平坦。另一方面，由 $\mathcal{F}^a$ 的定义，
次数 $p$ 项 $\check{\mathcal{C}}^p(\mathcal{V}, \mathcal{F}^a)$ 为
$$
\Gamma(V \times_U \ldots \times_U V, \mathcal{F}^a)
=
\Gamma(\Spec(A \otimes_R \ldots \otimes_R A), \mathcal{F}^a)
=
A \otimes_R \ldots \otimes_R A \otimes_R M
$$
省略验证 {\v C}ech 复形中的态射与复形
$s((A/R)_\bullet \otimes_R M)$ 中的态射一致。上同调的消失由
Lemma \ref{descent-lemma-ff-exact} 给出。
\end{proof}

\begin{proposition}
\label{descent-proposition-same-cohomology-quasi-coherent}
\begin{slogan}
无论使用哪种拓扑，拟凝聚层的上同调都相同。
\end{slogan}
设 $S$ 为概形，$\mathcal{F}$ 为 $S$ 上的拟凝聚层，并设
$\tau \in \{Zariski, \linebreak[0] \etale, \linebreak[0]
smooth, \linebreak[0] syntomic, \linebreak[0] fppf\}$。
\begin{enumerate}
\item 有典范同构
$$
H^q(S, \mathcal{F}) = H^q((\Sch/S)_\tau, \mathcal{F}^a).
$$
\item 有典范同构
$$
H^q(S, \mathcal{F}) =
H^q(S_{Zar}, \mathcal{F}^a) =
H^q(S_\etale, \mathcal{F}^a).
$$
\end{enumerate}
\end{proposition}

\begin{proof}
当 $q = 0$ 时，结论由 $\mathcal{F}^a$ 的定义立即得出。令
$\mathcal{C} = (\Sch/S)_\tau$，或 $\mathcal{C} = S_\etale$，
或 $\mathcal{C} = S_{Zar}$。

\medskip\noindent
应用 Cohomology on Sites, Lemma
\ref{sites-cohomology-lemma-cech-vanish-collection}，其中取
$\mathcal{F} = \mathcal{F}^a$，取
$\mathcal{B} \subset \Ob(\mathcal{C})$ 为 $\mathcal{C}$ 中仿射概形的集合，
并取 $\text{Cov} \subset \text{Cov}_\mathcal{C}$ 为标准仿射
$\tau$-覆盖的集合。该引理的假设 (3) 由 Lemma
\ref{descent-lemma-standard-covering-Cech-quasi-coherent} 满足。因此，对
$\mathcal{C}$ 的每个仿射对象 $U$，有 $H^p(U, \mathcal{F}^a) = 0$。

\medskip\noindent
接着，令 $U \in \Ob(\mathcal{C})$ 为任意分离对象，以
$f : U \to S$ 表示结构态射，并设 $U = \bigcup U_i$ 是仿射开覆盖。
也可把它视为 $\mathcal{C}$ 中 $U$ 的 $\tau$-覆盖
$\mathcal{U} = \{U_i \to U\}$。注意
$U_{i_0} \times_U \ldots \times_U U_{i_p} =
U_{i_0} \cap \ldots \cap U_{i_p}$ 是仿射的，因为假设了 $U$ 分离。
由 Cohomology on Sites, Lemma
\ref{sites-cohomology-lemma-cech-spectral-sequence-application} 及上述结果，
可知
$$
H^p(U, \mathcal{F}^a) = \check{H}^p(\mathcal{U}, \mathcal{F}^a)
= H^p(U, f^*\mathcal{F})
$$
其中最后一个等号由 Cohomology of Schemes, Lemma
\ref{coherent-lemma-cech-cohomology-quasi-coherent} 给出。特别地，若
$S$ 分离，可取 $U = S$ 且 $f = \text{id}_S$，命题即得证。建议读者
跳过证明的其余部分（或者重写它以给出更清楚的论述）。

\medskip\noindent
在 $S$ 上选取内射分解 $\mathcal{F} \to \mathcal{I}^\bullet$，并在
$\mathcal{C}$ 上选取内射分解 $\mathcal{F}^a \to \mathcal{J}^\bullet$。
记 $\mathcal{J}^n|_S$ 为 $\mathcal{J}^n$ 在 $S$ 的开集上的限制；由于开覆盖
也是 $\tau$-覆盖，它是拓扑空间 $S$ 上的一个层。由此得到复形
$$
0 \to \mathcal{F} \to \mathcal{J}^0|_S \to \mathcal{J}^1|_S \to \ldots
$$
它是正合的，因为它在任意仿射开集 $U \subset S$ 上的截面复形都是正合的
（这是由上面得到的 $H^p(U, \mathcal{F}^a)$ 在 $p > 0$ 时消失而知）。
因此，由 Derived Categories, Lemma
\ref{derived-lemma-morphisms-lift}，存在复形态射
$\mathcal{J}^\bullet|_S \to \mathcal{I}^\bullet$；特别地，它诱导态射
$$
R\Gamma(\mathcal{C}, \mathcal{F}^a)
=
\Gamma(S, \mathcal{J}^\bullet)
\longrightarrow
\Gamma(S, \mathcal{I}^\bullet)
=
R\Gamma(S, \mathcal{F}).
$$
取上同调便得到态射
$H^n(\mathcal{C}, \mathcal{F}^a) \to H^n(S, \mathcal{F})$，我们须证明
它是同构。令 $\mathcal{U} : S = \bigcup U_i$ 为一个仿射开覆盖；也可将其
视为一个 $\tau$-覆盖。由上述构造，得到双复形的态射
$$
\check{\mathcal{C}}^\bullet(\mathcal{U}, \mathcal{J})
=
\check{\mathcal{C}}^\bullet(\mathcal{U}, \mathcal{J}|_S)
\longrightarrow
\check{\mathcal{C}}^\bullet(\mathcal{U}, \mathcal{I}).
$$
它诱导谱序列之间的态射
$$
{}^\tau\! E_2^{p, q} = \check{H}^p(\mathcal{U}, \underline{H}^q(\mathcal{F}^a))
\longrightarrow
E_2^{p, q} = \check{H}^p(\mathcal{U}, \underline{H}^q(\mathcal{F}))
$$
% P05-ERR-0029
第一个谱序列收敛到 $H^{p + q}(\mathcal{C}, \mathcal{F})$，第二个收敛到
$H^{p + q}(S, \mathcal{F})$。另一方面，我们已经看到，所诱导的态射
${}^\tau\! E_2^{p, q} \to E_2^{p, q}$ 都是双射（因为所有交集都是仿射
概形中的开集，因而是分离的）。所以态射
$H^n(\mathcal{C}, \mathcal{F}^a) \to H^n(S, \mathcal{F})$ 也都是同构，
命题得证。
\end{proof}

\begin{proposition}
\label{descent-proposition-equivalence-quasi-coherent-functorial}
设 $f : T \to S$ 为概形态射。
\begin{enumerate}
\item Proposition \ref{descent-proposition-equivalence-quasi-coherent} 中的范畴等价
与拉回相容。更确切地说，对 $S$ 上的任意拟凝聚层 $\mathcal{G}$，有
$f^*(\mathcal{G}^a) = (f^*\mathcal{G})^a$。
\item Proposition \ref{descent-proposition-equivalence-quasi-coherent} 的第 (1) 部分
中的范畴等价一般{\bf 不}与前推相容。
\item 若 $f$ 拟紧且拟分离，并且 $\tau \in \{Zariski, \etale\}$，则
$f_*$ 和 $f_{small, *}$ 保持拟凝聚层，而且图表
$$
\xymatrix{
\QCoh(\mathcal{O}_T)
\ar[rr]_{f_*} \ar[d]_{\mathcal{F} \mapsto \mathcal{F}^a} & &
\QCoh(\mathcal{O}_S)
\ar[d]^{\mathcal{G} \mapsto \mathcal{G}^a} \\
\QCoh(T_\tau, \mathcal{O}) \ar[rr]^{f_{small, *}} & &
\QCoh(S_\tau, \mathcal{O})
}
$$
交换；也就是说，$f_{small, *}(\mathcal{F}^a) = (f_*\mathcal{F})^a$。
\end{enumerate}
\end{proposition}

\begin{proof}
第 (1) 部分由 Remark \ref{descent-remark-change-topologies-ringed-sites} 中的讨论
得出。第 (2) 部分只是一个警告，可以解释如下：首先，这个陈述无法作出
精确表述，因为一般而言，$f_*$ 不把拟凝聚层变为拟凝聚层。即使对 $f$
（以及 $f$ 的任意基变换）确实如此，大位点上的相容性也会意味着
$f_*\mathcal{F}$ 的形成与任意基变换交换，而这一般并不成立。一个具体
例子是拟紧开浸入
$j : X = \mathbf{A}^2_k \setminus \{0\} \to \mathbf{A}^2_k = Y$
，其中 $k$ 是一个域。我们有 $j_*\mathcal{O}_X = \mathcal{O}_Y$，但沿着
$0$ 态射基变换到 $\Spec(k)$ 后，可见此前推为零。

\medskip\noindent
下面在 $\tau = \etale$ 的情形证明 (3)。注意，$f$ 以及 $f$ 的任意基变换
都把拟凝聚层变为拟凝聚层；见 Schemes, Lemma
\ref{schemes-lemma-push-forward-quasi-coherent}。等式
$f_{small, *}(\mathcal{F}^a) = (f_*\mathcal{F})^a$ 意味着，对任意 \'etale
态射 $g : U \to S$，都有
$\Gamma(U, g^*f_*\mathcal{F}) = \Gamma(U \times_S T, (g')^*\mathcal{F})$
，其中 $g' : U \times_S T \to T$ 为投影。此式由 Cohomology of Schemes,
Lemma \ref{coherent-lemma-flat-base-change-cohomology} 成立。
\end{proof}

\begin{lemma}
\label{descent-lemma-higher-direct-images-small-etale}
设 $f : T \to S$ 为拟紧且拟分离的概形态射，并设 $\mathcal{F}$ 为 $T$ 上
的拟凝聚层。对于 \'etale 拓扑或 Zariski 拓扑，都有典范同构
$R^if_{small, *}(\mathcal{F}^a) = (R^if_*\mathcal{F})^a$.
\end{lemma}

\begin{proof}
我们对 \'etale 拓扑证明此结论；省略 Zariski 拓扑情形的证明。由
Cohomology of Schemes, Lemma
\ref{coherent-lemma-quasi-coherence-higher-direct-images}，层
$R^if_*\mathcal{F}$ 是拟凝聚的，故该断言有意义。层
$R^if_{small, *}\mathcal{F}^a$ 是与下列预层相关联的层：
$$
U \longmapsto H^i(U \times_S T, \mathcal{F}^a)
$$
其中 $g : U \to S$ 是 $S_\etale$ 的一个对象；见 Cohomology on Sites,
Lemma \ref{sites-cohomology-lemma-higher-direct-images}。按我们的约定，
右端是 $\mathcal{F}^a$ 在局部化 $T_\etale/U \times_S T$ 上的限制的
\'etale 上同调，而该局部化等于 $(U \times_S T)_\etale$。由
Proposition \ref{descent-proposition-same-cohomology-quasi-coherent}，这个预层与
下列预层相同：
$$
U \longmapsto
H^i(U \times_S T, (g')^*\mathcal{F}),
$$
其中 $g' : U \times_S T \to T$ 为投影。若 $U$ 仿射，则由 Cohomology of
Schemes, Lemma
\ref{coherent-lemma-quasi-coherence-higher-direct-images-application}，
它与 $H^0(U, R^if'_*(g')^*\mathcal{F})$ 相同。再由 Cohomology of
Schemes, Lemma \ref{coherent-lemma-flat-base-change-cohomology}，这等于
$H^0(U, g^*R^if_*\mathcal{F})$，即 $(R^if_*\mathcal{F})^a$ 在 $U$ 上的值。
因此，模层 $R^if_{small, *}(\mathcal{F}^a)$ 与
$(R^if_*\mathcal{F})^a$ 在 $S_\etale$ 的每个仿射对象上的值典范同构，
从而这两个层典范同构。
\end{proof}




\section{拟凝聚层与拓扑，II}
\label{descent-section-quasi-coherent-sheaves-bis}

\noindent
我们继续拓扑章中的讨论，比较概形 $S$ 上的拟凝聚模与关联于 $S$ 的任一
位点上的拟凝聚模。


\begin{lemma}
\label{descent-lemma-compare-etale-zariski-flat}
Lemma \ref{descent-lemma-compare-sites} 中的赋环位点态射
$\text{id}_{small, \etale, Zar} : S_\etale \to S_{Zar}$ 是平坦的。
\end{lemma}

\begin{proof}
记 $\epsilon = \text{id}_{small, \etale, Zar}$，并以
$\mathcal{O}_\etale$ 和 $\mathcal{O}_{Zar}$ 分别表示 $S_\etale$ 与
$S_{Zar}$ 上的结构层。我们须证明 $\mathcal{O}_\etale$ 是平坦的
$\epsilon^{-1}\mathcal{O}_{Zar}$-模。回忆 \'etale 态射是开映射；见
Morphisms, Lemma \ref{morphisms-lemma-etale-open}。由此（以及层的拉回
构造）可知，$\epsilon^{-1}\mathcal{O}_{Zar}$ 是 $S_\etale$ 上预层
$\mathcal{O}'$ 的层化；该预层把 \'etale 态射 $f : V \to S$ 映到
$\mathcal{O}_S(f(V))$。若 $V$ 与 $U = f(V) \subset S$ 都是仿射的，
则 $V \to U$ 是仿射概形之间的 \'etale 态射，因而对应于一个 \'etale
环同态。由于 \'etale 环同态平坦，可见
$\mathcal{O}_S(U) = \mathcal{O}'(V) \to
\mathcal{O}_\etale(V) = \mathcal{O}_V(V)$ 是平坦的。最后，对于每个
\'etale 态射 $f : V \to S$，即 $S_\etale$ 的每个对象，都存在仿射开覆盖
$V = \bigcup V_i$，使得对所有 $i$，$f(V_i)$ 都是 $S$ 中的仿射开集
\footnote{事实上，对 $y \in V$，选取仿射开集 $y \in V' \subset V$，
使 $f(V')$ 包含于某个仿射开集 $U \subset S$。再选取仿射开集
$f(y) \in U' \subset f(V')$。于是 $V'' = f^{-1}(U') \subset V'$ 是仿射
的，因为它等于 $U' \times_U V'$；而 $f(V'') = U'$ 也仿射。}。
结论遂由 Modules on Sites, Lemma
\ref{sites-modules-lemma-flatness-sheafification-refined} 得出。
\end{proof}

\begin{lemma}
\label{descent-lemma-equivalence-quasi-coherent-limits}
设 $S$ 为概形，并设 $\tau \in \{Zariski, \linebreak[0] \etale,
\linebreak[0] smooth, \linebreak[0] syntomic, \linebreak[0] fppf\}$。函子
$$
\QCoh(\mathcal{O}_S)
\longrightarrow
\textit{Mod}((\Sch/S)_\tau, \mathcal{O})
\quad\text{且}\quad
\QCoh(\mathcal{O}_S)
\longrightarrow
\textit{Mod}(S_\tau, \mathcal{O})
$$
由 Proposition \ref{descent-proposition-equivalence-quasi-coherent} 中的规则
$\mathcal{F} \mapsto \mathcal{F}^a$ 定义，并且具有下列性质：
\begin{enumerate}
\item 全忠实；
\item 与直和交换；
\item 与余极限交换；
\item 右正合；
\item 作为函子
$\QCoh(\mathcal{O}_S) \to \textit{Mod}(S_\etale, \mathcal{O})$,
是正合的；
\item 作为函子
$\QCoh(\mathcal{O}_S) \to
\textit{Mod}((\Sch/S)_\tau, \mathcal{O})$
一般{\bf 不}正合；
\item 给定两个拟凝聚 $\mathcal{O}_S$-模
$\mathcal{F}$、$\mathcal{G}$，有
$(\mathcal{F} \otimes_{\mathcal{O}_S} \mathcal{G})^a =
\mathcal{F}^a \otimes_\mathcal{O} \mathcal{G}^a$,
\item 若 $\tau = \etale$ 或 $\tau = Zariski$，并给定两个拟凝聚
$\mathcal{O}_S$-模 $\mathcal{F}$、$\mathcal{G}$，其中 $\mathcal{F}$
是有限表示的，则在 $\textit{Mod}(S_\tau, \mathcal{O})$ 中有
$(\SheafHom_{\mathcal{O}_S}(\mathcal{F}, \mathcal{G}))^a =
\SheafHom_\mathcal{O}(\mathcal{F}^a, \mathcal{G}^a)$；
\item 给定两个拟凝聚 $\mathcal{O}_S$-模 $\mathcal{F}$、$\mathcal{G}$，
一般在 $\textit{Mod}((\Sch/S)_\tau, \mathcal{O})$ 中{\bf 没有}
$(\SheafHom_{\mathcal{O}_S}(\mathcal{F}, \mathcal{G}))^a =
\SheafHom_\mathcal{O}(\mathcal{F}^a, \mathcal{G}^a)$
，即使 $\mathcal{F}$ 是有限表示的也如此；
\item 给定 $\mathcal{O}$-模的短正合列
$0 \to \mathcal{F}_1^a \to \mathcal{E} \to \mathcal{F}_2^a \to 0$
，则 $\mathcal{E}$ 是拟凝聚的\footnote{警告：这个说法容易引起误解。
见第 (6) 部分。}；也就是说，$\mathcal{E}$ 属于该函子的本质像。
\end{enumerate}
\end{lemma}

\begin{proof}
第 (1) 部分已见于 Proposition
\ref{descent-proposition-equivalence-quasi-coherent}。

\medskip\noindent
我们在 Schemes, Section \ref{schemes-section-quasi-coherent} 中已经看到，
概形上拟凝聚层的余极限仍为拟凝聚层。此外，Remark
\ref{descent-remark-change-topologies-ringed-sites} 表明
$\mathcal{F} \mapsto \mathcal{F}^a$ 是某个赋环位点态射的拉回函子，因而
与所有余极限交换；见 Modules on Sites, Lemma
\ref{sites-modules-lemma-exactness-pushforward-pullback}。所以 (3) 成立，
其特殊情形 (2) 也成立。

\medskip\noindent
这也表明该函子右正合（即与有限余极限交换），从而得到 (4)。

\medskip\noindent
函子 $\QCoh(\mathcal{O}_S) \to
\textit{Mod}(S_\etale, \mathcal{O})$，
$\mathcal{F} \mapsto \mathcal{F}^a$，是左正合的，因为 \'etale 态射
平坦；见 Morphisms, Lemma \ref{morphisms-lemma-etale-flat}。这证明了 (5)。

\medskip\noindent
为说明 (6)，设 $S = \Spec(\mathbf{Z})$。此时
$2 : \mathcal{O}_S \to \mathcal{O}_S$ 是单射，但在 $(\Sch/S)_\tau$ 上
所关联的 $\mathcal{O}$-模态射并不是单射，因为
$2 : \mathbf{F}_2 \to \mathbf{F}_2$ 不是单射，而
$\Spec(\mathbf{F}_2)$ 是 $(\Sch/S)_\tau$ 的一个对象。

\medskip\noindent
第 (7) 部分成立，因为如上所述，函子
$\mathcal{F} \mapsto \mathcal{F}^a$ 是赋环位点态射的拉回函子，而这类
函子由 Modules on Sites, Lemma
\ref{sites-modules-lemma-tensor-product-pullback} 与张量积交换。

\medskip\noindent
若 $\tau = Zariski$，则 (8) 显然成立，因为 $S_{Zar}$ 上
$\mathcal{O}$-模的范畴与拓扑空间 $S$ 上 $\mathcal{O}_S$-模的范畴相同。
若 $\tau = \etale$，则 (8) 成立，这是因为如上所述，函子
$\mathcal{F} \mapsto \mathcal{F}^a$ 是平坦赋环位点态射
$(S_\etale, \mathcal{O}) \to (S_{Zar}, \mathcal{O}_S)$ 的拉回函子；见
Lemma \ref{descent-lemma-compare-etale-zariski-flat}。由 Modules on Sites, Lemma
\ref{sites-modules-lemma-pullback-internal-hom}，沿平坦赋环位点态射的拉回
与从有限表示模出发取内部 Hom 交换。

\medskip\noindent
为说明 (9)，设 $S = \Spec(\mathbf{Z})$。令
$\mathcal{F} = \Coker(2 : \mathcal{O}_S \to \mathcal{O}_S)$
且 $\mathcal{G} = \mathcal{O}_S$。于是
$\mathcal{F}^a = \Coker(2 : \mathcal{O} \to \mathcal{O})$ 且
$\mathcal{G}^a = \mathcal{O}$。所以
$\SheafHom_\mathcal{O}(\mathcal{F}^a, \mathcal{G}^a) = \mathcal{O}[2]$
等于 $\mathcal{O}$ 中的 $2$-挠部分，它并非为零；见 (6) 的证明。另一
方面，模 $\SheafHom_{\mathcal{O}_S}(\mathcal{F}, \mathcal{G})$ 为零。

\medskip\noindent
证明 (10)。设
$0 \to \mathcal{F}_1^a \to \mathcal{E} \to \mathcal{F}_2^a \to 0$
为 $\mathcal{O}$-模的短正合列，其中 $\mathcal{F}_1$ 与
$\mathcal{F}_2$ 是 $S$ 上的拟凝聚层。考虑它在 $S_{Zar}$ 上的限制
$$
0 \to \mathcal{F}_1 \to \mathcal{E}|_{S_{Zar}} \to \mathcal{F}_2
$$
。由 Proposition \ref{descent-proposition-same-cohomology-quasi-coherent}，对任意
仿射开集 $U \subset S$，都有
$H^1(U, \mathcal{F}_1^a) = H^1(U, \mathcal{F}_1) = 0$。故上述序列在右端
也是正合的。由 Schemes, Section \ref{schemes-section-quasi-coherent}，
可知 $\mathcal{F} = \mathcal{E}|_{S_{Zar}}$ 是拟凝聚的。因此得到交换图
$$
\xymatrix{
& \mathcal{F}_1^a \ar[r] \ar[d] &
\mathcal{F}^a \ar[r] \ar[d] &
\mathcal{F}_2^a \ar[r] \ar[d] & 0 \\
0 \ar[r] &
\mathcal{F}_1^a \ar[r] &
\mathcal{E} \ar[r] &
\mathcal{F}_2^a \ar[r] & 0
}
$$
为完成证明，只须说明上行也右正合。为此，再次以
$U = \Spec(A) \subset S$ 表示 $S$ 的一个仿射开集。上面已经看到
$0 \to \mathcal{F}_1(U) \to \mathcal{E}(U) \to \mathcal{F}_2(U) \to 0$
是正合的。对任意仿射概形 $V/U$，若 $V = \Spec(B)$，则态射
$\mathcal{F}_1^a(V) \to \mathcal{E}(V)$ 为单射。按定义，有
$\mathcal{F}_1^a(V) = \mathcal{F}_1(U) \otimes_A B$。单射
$\mathcal{F}_1^a(V) \to \mathcal{E}(V)$ 分解为
$$
\mathcal{F}_1(U) \otimes_A B \to
\mathcal{E}(U) \otimes_A B \to \mathcal{E}(V)
$$
考虑形如 $B = A \oplus M$ 的 $A$-代数 $B$，可见
$\mathcal{F}_1(U) \to \mathcal{E}(U)$ 普遍单射（见 Algebra, Definition
\ref{algebra-definition-universally-injective}）。由于
$\mathcal{E}(U) = \mathcal{F}(U)$，可知 $\mathcal{F}_1 \to \mathcal{F}$
在任意基变换后仍为单射；等价地，
$\mathcal{F}_1^a \to \mathcal{F}^a$ 是单射。
\end{proof}

\begin{lemma}
\label{descent-lemma-properties-quasi-coherent}
设 $S$ 为概形。$S_\etale$ 上拟凝聚模的范畴
$\QCoh(S_\etale, \mathcal{O})$ 具有下列性质：
\begin{enumerate}
\item 拟凝聚层的任意直和都是拟凝聚的。
\item 拟凝聚层的任意余极限都是拟凝聚的。
\item 拟凝聚层态射的核与余核都是拟凝聚的。
\item 给定 $\mathcal{O}$-模的短正合列
$0 \to \mathcal{F}_1 \to \mathcal{F}_2 \to \mathcal{F}_3 \to 0$
，若三者中任意两个是拟凝聚的，则第三个也是拟凝聚的。
\item 给定两个拟凝聚 $\mathcal{O}$-模，其张量积是拟凝聚的。
\item 给定两个拟凝聚 $\mathcal{O}$-模 $\mathcal{F}$、$\mathcal{G}$，
其中 $\mathcal{F}$ 是有限表示的，则内部 Hom
$\SheafHom_\mathcal{O}(\mathcal{F}, \mathcal{G})$
是拟凝聚的。
\end{enumerate}
\end{lemma}

\begin{proof}
相应事实对于概形 $S$ 上的拟凝聚模成立；见 Schemes, Section
\ref{schemes-section-quasi-coherent}。我们利用 Lemma
\ref{descent-lemma-equivalence-quasi-coherent-limits} 把这些事实转移到
$S_\etale$ 上。

\medskip\noindent
证明 (1)。设 $\mathcal{F}_i$，$i \in I$，为
$\QCoh(S_\etale, \mathcal{O})$ 中的一族对象。对 $S$ 上某些拟凝聚模
$\mathcal{G}_i$，写成 $\mathcal{F}_i = \mathcal{G}_i^a$。由所引引理，
$\bigoplus \mathcal{F}_i = (\bigoplus \mathcal{G}_i)^a$，结论得证。

\medskip\noindent
证明 (2)。设 $\mathcal{I} \to \QCoh(S_\etale, \mathcal{O})$，
$i \mapsto \mathcal{F}_i$，为一个图。写成
$\mathcal{F}_i = \mathcal{G}_i^a$，便得到图
$\mathcal{I} \to \QCoh(\mathcal{O}_S)$。由所引引理，
$\colim \mathcal{F}_i = (\colim \mathcal{G}_i)^a$，结论得证。

\medskip\noindent
证明 (3)。设 $a : \mathcal{F} \to \mathcal{F}'$ 为
$\QCoh(S_\etale, \mathcal{O})$ 中的一个箭头。对 $S$ 上拟凝聚模之间的
某个态射 $b : \mathcal{G} \to \mathcal{G}'$，写成 $a = b^a$。由所引
引理，有 $\Ker(a) = \Ker(b)^a$ 且
$\Coker(a) = \Coker(b)^a$，结论得证。

\medskip\noindent
证明 (4)。除已知 $\mathcal{F}_1$ 与 $\mathcal{F}_3$ 拟凝聚的情形外，
结论都由 (3) 得出。在这一例外情形，写成
$\mathcal{F}_1 = \mathcal{G}_1^a$ 与
$\mathcal{F}_3 = \mathcal{G}_3^a$，其中 $\mathcal{G}_i$ 是 $S$ 上的
拟凝聚层。结论由 Lemma
\ref{descent-lemma-equivalence-quasi-coherent-limits} 的第 (10) 部分得出。

\medskip\noindent
证明 (5)。设 $\mathcal{F}$ 与 $\mathcal{F}'$ 属于
$\QCoh(S_\etale, \mathcal{O})$。写成
$\mathcal{F} = \mathcal{G}^a$ 与
$\mathcal{F}' = (\mathcal{G}')^a$，其中 $\mathcal{G}$ 与
$\mathcal{G}'$ 是 $S$ 上的拟凝聚层。由所引引理，有
$\mathcal{F} \otimes_\mathcal{O} \mathcal{F}' =
(\mathcal{G} \otimes_{\mathcal{O}_S} \mathcal{G}')^a$
，结论得证。

\medskip\noindent
证明 (6)。设 $\mathcal{F}$ 与 $\mathcal{G}$ 属于
$\QCoh(S_\etale, \mathcal{O})$，且 $\mathcal{F}$ 是有限表示的。写成
$\mathcal{F} = \mathcal{H}^a$ 与 $\mathcal{G} = (\mathcal{I})^a$，其中
$\mathcal{H}$ 与 $\mathcal{I}$ 是 $S$ 上的拟凝聚层。由 Lemma
\ref{descent-lemma-equivalence-quasi-coherent-properties}，可知 $\mathcal{H}$ 是
有限表示的。由 Lemma \ref{descent-lemma-equivalence-quasi-coherent-limits} 的
第 (8) 部分，有
$\SheafHom_\mathcal{O}(\mathcal{F}, \mathcal{G}) =
(\SheafHom_{\mathcal{O}_S}(\mathcal{H}, \mathcal{I}))^a$
，结论得证。
\end{proof}

\begin{lemma}
\label{descent-lemma-properties-quasi-coherent-on-big}
设 $S$ 为概形，并设 $\tau \in \{Zariski, \linebreak[0] \etale, \linebreak[0]
smooth, \linebreak[0] syntomic, \linebreak[0] fppf\}$。
$(\Sch/S)_\tau$ 上拟凝聚模的范畴
$\QCoh((\Sch/S)_\tau, \mathcal{O})$ 具有下列性质：
\begin{enumerate}
\item 拟凝聚层的任意直和都是拟凝聚的。
\item 拟凝聚层的任意余极限都是拟凝聚的。
\item 拟凝聚层态射的余核是拟凝聚的。
\item 给定 $\mathcal{O}$-模的短正合列
$0 \to \mathcal{F}_1 \to \mathcal{F}_2 \to \mathcal{F}_3 \to 0$
，若 $\mathcal{F}_1$ 与 $\mathcal{F}_3$ 拟凝聚，则
$\mathcal{F}_2$ 也拟凝聚。
\item 给定两个拟凝聚 $\mathcal{O}$-模，其张量积是拟凝聚的。
\item 给定两个拟凝聚 $\mathcal{O}$-模 $\mathcal{F}$、$\mathcal{G}$，
其中 $\mathcal{F}$ 有限局部自由，则内部 Hom
$\SheafHom_\mathcal{O}(\mathcal{F}, \mathcal{G})$
是拟凝聚的。
\end{enumerate}
\end{lemma}

\begin{proof}
相应事实对于概形 $S$ 上的拟凝聚模成立；见 Schemes, Section
\ref{schemes-section-quasi-coherent}。我们利用 Lemma
\ref{descent-lemma-equivalence-quasi-coherent-limits} 把这些事实转移到
$(\Sch/S)_\tau$ 上。

\medskip\noindent
证明 (1)。设 $\mathcal{F}_i$，$i \in I$，为
$\QCoh((\Sch/S)_\tau, \mathcal{O})$ 中的一族对象。对 $S$ 上某些拟凝聚模
$\mathcal{G}_i$，写成 $\mathcal{F}_i = \mathcal{G}_i^a$。由所引引理，
$\bigoplus \mathcal{F}_i = (\bigoplus \mathcal{G}_i)^a$，结论得证。

\medskip\noindent
证明 (2)。设 $\mathcal{I} \to \QCoh((\Sch/S)_\tau, \mathcal{O})$，
$i \mapsto \mathcal{F}_i$，为一个图。写成
$\mathcal{F}_i = \mathcal{G}_i^a$，便得到图
$\mathcal{I} \to \QCoh(\mathcal{O}_S)$。由所引引理，
$\colim \mathcal{F}_i = (\colim \mathcal{G}_i)^a$，结论得证。

\medskip\noindent
证明 (3)。设 $a : \mathcal{F} \to \mathcal{F}'$ 为
$\QCoh((\Sch/S)_\tau, \mathcal{O})$ 中的一个箭头。对 $S$ 上拟凝聚模之间
的某个态射 $b : \mathcal{G} \to \mathcal{G}'$，写成 $a = b^a$。由所引
引理，有 $\Coker(a) = \Coker(b)^a$（因为余核是余极限），结论得证。

\medskip\noindent
证明 (4)。写成 $\mathcal{F}_1 = \mathcal{G}_1^a$ 与
$\mathcal{F}_3 = \mathcal{G}_3^a$，其中 $\mathcal{G}_i$ 是 $S$ 上的
拟凝聚层。结论由 Lemma
\ref{descent-lemma-equivalence-quasi-coherent-limits} 的第 (10) 部分得出。

\medskip\noindent
证明 (5)。设 $\mathcal{F}$ 与 $\mathcal{F}'$ 属于
$\QCoh((\Sch/S)_\tau, \mathcal{O})$。写成
$\mathcal{F} = \mathcal{G}^a$ 与 $\mathcal{F}' = (\mathcal{G}')^a$，其中
$\mathcal{G}$ 与 $\mathcal{G}'$ 是 $S$ 上的拟凝聚层。由所引引理，有
$\mathcal{F} \otimes_\mathcal{O} \mathcal{F}' =
(\mathcal{G} \otimes_{\mathcal{O}_S} \mathcal{G}')^a$
，结论得证。

\medskip\noindent
证明 (6)。对某个拟凝聚 $\mathcal{O}_S$-模，写成
$\mathcal{F} = \mathcal{H}^a$。由 Lemma
\ref{descent-lemma-equivalence-quasi-coherent-properties}，可知 $\mathcal{H}$ 有限
局部自由。该问题在 $S$ 上是 Zariski 局部的（略去细节），所以可设
$\mathcal{H} = \mathcal{O}_S^{\oplus n}$ 有限自由。于是
$\mathcal{F} = \mathcal{O}^{\oplus n}$，而
$\SheafHom_\mathcal{O}(\mathcal{F}, \mathcal{G}) = \mathcal{G}^{\oplus n}$
是拟凝聚的。
\end{proof}

\begin{example}
\label{descent-example-internal-hom-not-qcoh}
设 $S$ 为概形。对 Lemma \ref{descent-lemma-properties-quasi-coherent-on-big} 中
考虑的某个拓扑 $\tau$，设 $\mathcal{F}$ 与 $\mathcal{G}$ 为
$(\Sch/S)_\tau$ 上的拟凝聚模。一般而言，
$\SheafHom_\mathcal{O}(\mathcal{F}, \mathcal{G})$
并非拟凝聚的，即使 $\mathcal{F}$ 是有限表示的也如此。事实上，取
$S = \Spec(\mathbf{Z})$，
$\mathcal{F} = \Coker(2 : \mathcal{O} \to \mathcal{O})$,
且 $\mathcal{G} = \mathcal{O}$。此时
$\SheafHom_\mathcal{O}(\mathcal{F}, \mathcal{G}) = \mathcal{O}[2]$
等于 $\mathcal{O}$ 中的 $2$-挠部分，而它不是拟凝聚的。
\end{example}

\begin{lemma}
\label{descent-lemma-qc-colimits}
设 $S$ 为概形。
\begin{enumerate}
\item 范畴 $\QCoh((\Sch/S)_{fppf}, \mathcal{O})$ 有余极限，并且这些
余极限与下列范畴中的余极限一致：
$\textit{Mod}((\Sch/S)_{Zar}, \mathcal{O})$,
$\textit{Mod}((\Sch/S)_\etale, \mathcal{O})$，以及
$\textit{Mod}((\Sch/S)_{fppf}, \mathcal{O})$.
\item 给定 $\QCoh((\Sch/S)_{fppf}, \mathcal{O})$ 中的
$\mathcal{F}, \mathcal{G}$，在 $\textit{Mod}((\Sch/S)_{Zar}, \mathcal{O})$、
$\textit{Mod}((\Sch/S)_\etale, \mathcal{O})$ 或
$\textit{Mod}((\Sch/S)_{fppf}, \mathcal{O})$ 中计算的张量积
$\mathcal{F} \otimes_\mathcal{O} \mathcal{G}$ 相互一致，其共同值是
$\QCoh((\Sch/S)_{fppf}, \mathcal{O})$ 的对象。
\item 给定 $\QCoh((\Sch/S)_{fppf}, \mathcal{O})$ 中的
$\mathcal{F}, \mathcal{G}$，且 $\mathcal{F}$ 有限局部自由（在 fppf 拓扑
中，等价地在 \'etale 拓扑中，或等价地在 Zariski 拓扑中），则在
$\textit{Mod}((\Sch/S)_{Zar}, \mathcal{O})$、
$\textit{Mod}((\Sch/S)_\etale, \mathcal{O})$ 或
$\textit{Mod}((\Sch/S)_{fppf}, \mathcal{O})$ 中计算的内部 Hom
$\SheafHom_\mathcal{O}(\mathcal{F}, \mathcal{G})$
相互一致，其共同值是 $\QCoh((\Sch/S)_{fppf}, \mathcal{O})$ 的对象。
\end{enumerate}
\end{lemma}

\begin{proof}
本引理以略有不同的方式汇集上面证明的结果。首先，由 Lemma
\ref{descent-lemma-properties-quasi-coherent-on-big}，我们已经知道 (1)、(2) 或
(3) 中构造的输出属于 $\QCoh((\Sch/S)_\tau, \mathcal{O})$。还须在每种
情形证明结果与所用拓扑无关。关键在于 Proposition
\ref{descent-proposition-equivalence-quasi-coherent} 中的等价
$\QCoh(\mathcal{O}_S) \to \QCoh((\Sch/S)_\tau, \mathcal{O})$,
$\mathcal{F} \mapsto \mathcal{F}^a$
由同一个公式给出，而与拓扑 $\tau \in \{Zariski, \etale, fppf\}$ 的选择
无关。

\medskip\noindent
证明 (1)。设 $\mathcal{I} \to \QCoh((\Sch/S)_{fppf}, \mathcal{O})$，
$i \mapsto \mathcal{F}_i$，为一个图。写成
$\mathcal{F}_i = \mathcal{G}_i^a$，便得到图
$\mathcal{I} \to \QCoh(\mathcal{O}_S)$。由 Lemma
\ref{descent-lemma-equivalence-quasi-coherent-limits}，对
$\tau \in \{Zariski, \etale, fppf\}$，在
$\textit{Mod}((\Sch/S)_\tau, \mathcal{O})$ 中有
$\colim \mathcal{F}_i = (\colim \mathcal{G}_i)^a$。这证明了 (1)。

\medskip\noindent
证明 (2)。写成 $\mathcal{F} = \mathcal{H}^a$ 与
$\mathcal{G} = (\mathcal{I})^a$，其中 $\mathcal{H}$ 与 $\mathcal{I}$
是 $S$ 上的拟凝聚层。于是
% P05-ERR-0030
$\mathcal{F} \otimes_\mathcal{O} \mathcal{G} =
(\mathcal{H} \otimes_\mathcal{O} \mathcal{I})^a$
，并且由 Lemma \ref{descent-lemma-equivalence-quasi-coherent-limits}，对
$\tau \in \{Zariski, \etale, fppf\}$，此等式在
$\textit{Mod}((\Sch/S)_\tau, \mathcal{O})$ 中成立。这证明了 (2)。

\medskip\noindent
证明 (3)。设 $\mathcal{F}$ 与 $\mathcal{G}$ 属于
$\QCoh((\Sch/S)_{fppf}, \mathcal{O})$。写成
$\mathcal{F} = \mathcal{H}^a$，其中 $\mathcal{H}$ 是 $S$ 上的拟凝聚层。
由 Lemma \ref{descent-lemma-equivalence-quasi-coherent-properties}，有
\begin{align*}
\mathcal{F}\text{ 在 fppf 拓扑中有限局部自由}
& \Leftrightarrow
\mathcal{H}\text{ 在下列对象上有限局部自由： }S \\
& \Leftrightarrow
\mathcal{F}\text{ 在平展拓扑中有限局部自由} \\
& \Leftrightarrow
\mathcal{H}\text{ 在下列对象上有限局部自由： }S \\
& \Leftrightarrow
\mathcal{F}\text{ 在 Zariski 拓扑中有限局部自由}
\end{align*}
这解释了第 (3) 部分的括号内陈述。现在，若这些等价条件成立，则
$\mathcal{H}$ 有限局部自由。Modules on Sites, Section
\ref{sites-modules-section-internal-hom} 中
$\SheafHom_\mathcal{O}(\mathcal{F}, \mathcal{G})$ in
的构造，仅依赖于作为模预层的 $\mathcal{F}$ 与 $\mathcal{G}$（只有
输出 $\SheafHom$ 是否为层，才依赖于 $\mathcal{F}$ 与 $\mathcal{G}$
是否为层）。
\end{proof}










\section{拟凝聚模与仿射概形}
\label{descent-section-alternative-quasi-coherent}

\noindent
设 $S$ 为概形\footnote{在本节中，正如 Topologies, Section
\ref{topologies-section-change-topologies} 中一样，我们选择各位点
$(\Sch/S)_\tau$，使得当
$\tau \in \{Zariski, \etale, smooth, syntomic, fppf\}$ 时，它们具有相同
的底范畴。于是各位点 $(\textit{Aff}/S)_\tau$ 也具有相同的底范畴。}。
设 $\tau \in \{Zariski, \etale, smooth, syntomic, fppf\}$。回忆，
$(\textit{Aff}/S)_\tau$ 是 $(\Sch/S)_\tau$ 中对象为仿射概形的全子范畴；
将覆盖规定为标准 $\tau$-覆盖，便把它变为一个位点。由 Topologies,
Lemmas
\ref{topologies-lemma-affine-big-site-Zariski},
\ref{topologies-lemma-affine-big-site-etale},
\ref{topologies-lemma-affine-big-site-smooth},
\ref{topologies-lemma-affine-big-site-syntomic}，以及
\ref{topologies-lemma-affine-big-site-fppf}
，有拓扑斯等价
$g : \Sh((\textit{Aff}/S)_\tau) \to \Sh((\Sch/S)_\tau)$
，其拉回函子由限制给出。回忆 $\mathcal{O}$ 表示 $(\Sch/S)_\tau$ 上的
结构层；为明确区分记号，暂以 $\mathcal{O}_{\textit{Aff}}$ 表示
$\mathcal{O}$ 在 $(\textit{Aff}/S)_\tau$ 上的限制。于是得到赋环拓扑斯的
等价
\begin{equation}
\label{descent-equation-alternative-ringed}
(\Sh((\textit{Aff}/S)_\tau), \mathcal{O}_{\textit{Aff}})
\longrightarrow
(\Sh((\Sch/S)_\tau), \mathcal{O})
\end{equation}
。由此也可以比较拟凝聚模。

\begin{lemma}
\label{descent-lemma-quasi-coherent-alternative}
设 $S$ 为概形，并设 $\mathcal{F}$ 为 $(\textit{Aff}/S)_{fppf}$ 上的
$\mathcal{O}_{\textit{Aff}}$-模预层。下列条件等价：
\begin{enumerate}
\item 对 $(\textit{Aff}/S)_{fppf}$ 的每个态射 $U \to U'$，态射
$\mathcal{F}(U') \otimes_{\mathcal{O}(U')} \mathcal{O}(U) \to \mathcal{F}(U)$
是同构；
\item $\mathcal{F}$ 是 $(\textit{Aff}/S)_{Zar}$ 上的层，并且按 Modules on
Sites, Definition \ref{sites-modules-definition-site-local} 的意义，是赋环
位点 $((\textit{Aff}/S)_{Zar}, \mathcal{O}_{\textit{Aff}})$ 上的拟凝聚模；
\item 对 \'etale 拓扑，条件与 (2) 相同；
\item 对 smooth 拓扑，条件与 (2) 相同；
\item 对 syntomic 拓扑，条件与 (2) 相同；
\item 对 fppf 拓扑，条件与 (2) 相同；
\item 经等价 (\ref{descent-equation-alternative-ringed})，$\mathcal{F}$ 对应于
$(\Sch/S)_{Zar}$,
$(\Sch/S)_\etale$,
$(\Sch/S)_{smooth}$,
$(\Sch/S)_{syntomic}$，或
$(\Sch/S)_{fppf}$
之一上的拟凝聚模；
\item $\mathcal{F}$ 由 Section \ref{descent-section-quasi-coherent-sheaves} 中
所述过程从唯一的拟凝聚 $\mathcal{O}_S$-模 $\mathcal{G}$ 得到。
\end{enumerate}
\end{lemma}

\begin{proof}
由于拟凝聚模的概念是内蕴的（Modules on Sites, Lemma
\ref{sites-modules-lemma-special-locally-free}），可知等价
(\ref{descent-equation-alternative-ringed}) 诱导拟凝聚模范畴之间的等价。
Proposition \ref{descent-proposition-equivalence-quasi-coherent} 说明，在
$\Sch/S$ 上研究拟凝聚模时使用哪种拓扑并不重要，并且还说明 (8) 与
(7) 相同。因此 (2) -- (8) 全都等价。

\medskip\noindent
假设等价条件 (2) -- (8) 成立，并令 $\mathcal{G}$ 如 (8) 中所述。设
$h : U \to U' \to S$ 为 $\textit{Aff}/S$ 的一个态射。以
$f : U \to S$ 与 $f' : U' \to S$ 表示结构态射，于是
$f = f' \circ h$。我们有
$\mathcal{F}(U') = \Gamma(U', (f')^*\mathcal{G})$ 以及
$\mathcal{F}(U) = \Gamma(U, f^*\mathcal{G}) = \Gamma(U, h^*(f')^*\mathcal{G})$。
故 (1) 由 Schemes, Lemma \ref{schemes-lemma-widetilde-pullback} 成立。

\medskip\noindent
假设 (1) 成立。为完成证明，只须证明 (2)。设 $U$ 为
$(\textit{Aff}/S)_{Zar}$ 的对象，记 $U = \Spec(R)$。标准开覆盖
$U = U_1 \cup \ldots \cup U_n$ 由 $U_i = D(f_i)$ 给出，其中某些元素
$f_1, \ldots, f_n \in R$ 生成 $R$ 的单位理想。由性质 (1)，有
$$
\mathcal{F}(U_i) =
\mathcal{F}(U) \otimes_R R_{f_i} =
\mathcal{F}(U)_{f_i}
$$
and
$$
\mathcal{F}(U_i \cap U_j) =
\mathcal{F}(U) \otimes_R R_{f_if_j} =
\mathcal{F}(U)_{f_if_j}
$$
所以由 Algebra, Lemma \ref{algebra-lemma-cover-module}，可知
$\mathcal{F}$ 是 $(\textit{Aff}/S)_{Zar}$ 上的层。选取自由 $R$-模表示
$$
\bigoplus\nolimits_{k \in K} R
\longrightarrow
\bigoplus\nolimits_{l \in L} R
\longrightarrow
\mathcal{F}(U)
\longrightarrow 0
$$
。由性质 (1) 以及张量积的右正合性，可知对
$(\textit{Aff}/S)_{Zar}$ 中的每个态射 $U' \to U$，都得到表示
$$
% P05-ERR-0031
\bigoplus\nolimits_{k \in K} \mathcal{O}_{Aff}(U')
\longrightarrow
\bigoplus\nolimits_{l \in L} \mathcal{O}_{Aff}(U')
\longrightarrow
\mathcal{F}(U')
\longrightarrow 0
$$
换言之，可见 $\mathcal{F}$ 在局部化范畴
$(\textit{Aff}/S)_{Zar}/U$ 上的限制有表示
$$
\bigoplus\nolimits_{k \in K} \mathcal{O}_{Aff}|_{(\textit{Aff}/S)_{Zar}/U}
\longrightarrow
\bigoplus\nolimits_{l \in L} \mathcal{O}_{Aff}|_{(\textit{Aff}/S)_{Zar}/U}
\longrightarrow
\mathcal{F}|_{(\textit{Aff}/S)_{Zar}/U}
\longrightarrow 0
$$
尽管记号十分繁复，这便完成了证明。
\end{proof}

\noindent
我们继续本节引言中开始的讨论。设 $\tau \in \{Zariski, \etale\}$。
回忆，$S_{affine, \tau}$ 是 $S_\tau$ 中对象为仿射概形的全子范畴；将覆盖
规定为标准 $\tau$-覆盖，便把它变为一个位点。见 Topologies,
Definitions \ref{topologies-definition-big-small-Zariski} 与
\ref{topologies-definition-big-small-etale}。分别由 Topologies, Lemmas
\ref{topologies-lemma-alternative-zariski} 与
\ref{topologies-lemma-alternative}，有拓扑斯等价
$g : \Sh(S_{affine, \tau}) \to \Sh(S_\tau)$，其拉回函子由限制给出。
回忆 $\mathcal{O}$ 表示 $S_\tau$ 上的结构层；为明确区分记号，暂以
$\mathcal{O}_{affine}$ 表示 $\mathcal{O}$ 在 $S_{affine, \tau}$ 上的限制。
于是得到赋环拓扑斯的等价
\begin{equation}
\label{descent-equation-alternative-small-ringed}
(\Sh(S_{affine, \tau}), \mathcal{O}_{affine})
\longrightarrow
(\Sh(S_\tau), \mathcal{O})
\end{equation}
。由此也可以比较拟凝聚模。

\begin{lemma}
\label{descent-lemma-quasi-coherent-alternative-small}
设 $S$ 为概形，设 $\tau \in \{Zariski, \etale\}$，并设 $\mathcal{F}$ 为
$S_{affine, \tau}$ 上的 $\mathcal{O}_{affine}$-模预层。下列条件等价：
\begin{enumerate}
\item 对 $S_{affine, \tau}$ 的每个态射 $U \to U'$，态射
$\mathcal{F}(U') \otimes_{\mathcal{O}(U')} \mathcal{O}(U) \to \mathcal{F}(U)$
是同构；
\item $\mathcal{F}$ 是 $S_{affine, \tau}$ 上的层，并且按 Modules on
Sites, Definition \ref{sites-modules-definition-site-local} 的意义，是赋环
位点 $(S_{affine, \tau}, \mathcal{O}_{affine})$ 上的拟凝聚模；
\item 经等价 (\ref{descent-equation-alternative-small-ringed})，$\mathcal{F}$ 对应
于 $S_\tau$ 上的拟凝聚模；
\item $\mathcal{F}$ 由 Section \ref{descent-section-quasi-coherent-sheaves} 中
所述过程从唯一的拟凝聚 $\mathcal{O}_S$-模 $\mathcal{G}$ 得到。
\end{enumerate}
\end{lemma}

\begin{proof}
我们对 \'etale 拓扑的情形证明此结论。

\medskip\noindent
假设 (1) 成立。为证明 $\mathcal{F}$ 是层，设
$\mathcal{U} = \{U_i \to U\}_{i = 1, \ldots, n}$ 为
$S_{affine, \etale}$ 的一个覆盖。由我们对 $\mathcal{F}$ 的假设，
% P05-ERR-0032
$\mathcal{F}$ 关于 $\mathcal{U}$ 的层条件归结为证明
$$
0 \to \mathcal{F}(U) \to
\prod \mathcal{F}(U) \otimes_{\mathcal{O}(U)} \mathcal{O}(U_i) \to
\prod \mathcal{F}(U) \otimes_{\mathcal{O}(U)} \mathcal{O}(U_i \times_U U_j)
$$
正合。这是因为 $\mathcal{O}(U) \to \prod \mathcal{O}(U_i)$ 忠实平坦
（由 Lemma \ref{descent-lemma-standard-covering-Cech} 以及
$S_{affine, \etale}$ 中的覆盖是标准 \'etale 覆盖这一事实），因而可以
应用 Lemma \ref{descent-lemma-ff-exact}。接着证明 $\mathcal{F}$ 在
$S_{affine, \etale}$ 上拟凝聚。具体而言，对 $S_{affine, \etale}$ 中的
$U$，置 $R = \mathcal{O}(U)$，并选取自由 $R$-模表示
$$
\bigoplus\nolimits_{k \in K} R
\longrightarrow
\bigoplus\nolimits_{l \in L} R
\longrightarrow
\mathcal{F}(U)
\longrightarrow 0
$$
。由性质 (1) 以及张量积的右正合性，可知对 $S_{affine, \etale}$ 中的
每个态射 $U' \to U$，都得到表示
$$
\bigoplus\nolimits_{k \in K} \mathcal{O}(U')
\longrightarrow
\bigoplus\nolimits_{l \in L} \mathcal{O}(U')
\longrightarrow
\mathcal{F}(U')
\longrightarrow 0
$$
换言之，可见 $\mathcal{F}$ 在局部化范畴 $S_{affine, etale}/U$ 上的限制
有表示
$$
% P05-ERR-0033
\bigoplus\nolimits_{k \in K} \mathcal{O}_{affine}|_{S_{affine, \etale}/U}
\longrightarrow
\bigoplus\nolimits_{l \in L} \mathcal{O}_{affine}|_{S_{affine, \etale}/U}
\longrightarrow
\mathcal{F}|_{S_{affine, \etale}/U}
\longrightarrow 0
$$
，这正是证明 $\mathcal{F}$ 拟凝聚所需的表示。尽管记号十分繁复，
这便完成了 (1) 蕴含 (2) 的证明。

\medskip\noindent
由于拟凝聚模的概念是内蕴的（Modules on Sites, Lemma
\ref{sites-modules-lemma-special-locally-free}），可知等价
(\ref{descent-equation-alternative-small-ringed}) 诱导拟凝聚模范畴之间的等价。
所以 (2) 与 (3) 等价。

\medskip\noindent
(3) 与 (4) 的等价性由 Proposition
\ref{descent-proposition-equivalence-quasi-coherent} 得出。

\medskip\noindent
假设 (4) 成立，下面证明 (1)。令 $\mathcal{G}$ 如 (4) 中所述。设
$h : U \to U' \to S$ 为 $S_{affine, \etale}$ 的一个态射。以
$f : U \to S$ 与 $f' : U' \to S$ 表示结构态射，于是
$f = f' \circ h$。我们有
$\mathcal{F}(U') = \Gamma(U', (f')^*\mathcal{G})$ 以及
$\mathcal{F}(U) = \Gamma(U, f^*\mathcal{G}) = \Gamma(U, h^*(f')^*\mathcal{G})$。
故 (1) 由 Schemes, Lemma \ref{schemes-lemma-widetilde-pullback} 成立。

\medskip\noindent
省略 Zariski 拓扑情形的证明。
\end{proof}





\section{寄生模}
\label{descent-section-parasitic}

\noindent
寄生模是限制到在基概形上平坦的概形后为零的模。其正式定义如下。

\begin{definition}
\label{descent-definition-parasitic}
设 $S$ 为概形，设 $\tau \in \{Zar, \etale,
smooth, syntomic, fppf\}$，并设 $\mathcal{F}$ 为 $(\Sch/S)_\tau$ 上的
$\mathcal{O}$-模预层。
\begin{enumerate}
\item 若对每个平坦态射 $U \to S$ 都有 $\mathcal{F}(U) = 0$，则称
$\mathcal{F}$ 是{\it 寄生的}\footnote{这可能不是标准术语。}。
\item 若对每个 $\tau$-覆盖 $\{U_i \to S\}_{i \in I}$ 以及所有 $i$，
都有 $\mathcal{F}(U_i) = 0$，则称 $\mathcal{F}$ 是
{\it 关于 $\tau$-拓扑寄生的}。
\end{enumerate}
\end{definition}

\noindent
若 $\tau = fppf$，这意味着只要 $U \to S$ 平坦且局部有限表示，就有
$\mathcal{F}|_{U_{Zar}} = 0$；其他情形类似。

\begin{lemma}
\label{descent-lemma-cohomology-parasitic}
设 $S$ 为概形，设 $\tau \in \{Zar, \etale, smooth,
syntomic, fppf\}$，并设 $\mathcal{G}$ 为 $(\Sch/S)_\tau$ 上的
$\mathcal{O}$-模预层。
\begin{enumerate}
\item 若 $\mathcal{G}$ 关于 $\tau$-拓扑寄生，则对每个在 $S$ 中开、
或在 $S$ 上 \'etale、或在 $S$ 上 smooth、或在 $S$ 上 syntomic、或在
$S$ 上平坦且局部有限表示的相应 $U$，都有
$H^p_\tau(U, \mathcal{G}) = 0$。
\item 若 $\mathcal{G}$ 寄生，则对每个在 $S$ 上平坦的 $U$，都有
$H^p_\tau(U, \mathcal{G}) = 0$。
\end{enumerate}
\end{lemma}

\begin{proof}
证明 $\tau = fppf$ 的情形；其他情形的证明完全相同。假设意味着，对任意
平坦且局部有限表示的 $U \to S$，都有 $\mathcal{G}(U) = 0$。把
Cohomology on Sites, Lemma
\ref{sites-cohomology-lemma-cech-vanish-collection} 应用于由平坦且局部
有限表示的 $U \to S$ 组成的子集
$\mathcal{B} \subset \Ob((\Sch/S)_{fppf})$，以及由 $\mathcal{B}$ 中对象的
所有 fppf 覆盖组成的集合 $\text{Cov}$。
\end{proof}

\begin{lemma}
\label{descent-lemma-direct-image-parasitic}
设 $f : T \to S$ 为概形态射。对 $(\Sch/T)_\tau$ 上的任意寄生
$\mathcal{O}$-模，前推 $f_*\mathcal{F}$ 与高阶直像
$R^if_*\mathcal{F}$ 都是 $(\Sch/S)_\tau$ 上的寄生 $\mathcal{O}$-模。
\end{lemma}

\begin{proof}
回忆，$R^if_*\mathcal{F}$ 是与下列预层相关联的层：
$$
U \mapsto H^i((\Sch/U \times_S T)_\tau, \mathcal{F})
$$
；见 Cohomology on Sites, Lemma
\ref{sites-cohomology-lemma-higher-direct-images}。若 $U \to S$ 平坦，则
$U \times_S T \to T$ 作为其基变换也是平坦的。故由 Lemma
\ref{descent-lemma-cohomology-parasitic}，上示群为零。若
$\{U_i \to U\}$ 是 $\tau$-覆盖，则 $U_i \times_S T \to T$ 也平坦。
因此很清楚，上示预层的层化在所有平坦于 $S$ 的概形 $U$ 上为零。
\end{proof}

\begin{lemma}
\label{descent-lemma-quasi-coherent-and-flat-base-change}
设 $S$ 为概形，设 $\tau \in \{Zar, \etale\}$，并设 $\mathcal{G}$ 为
$(\Sch/S)_{fppf}$ 上的 $\mathcal{O}$-模层，满足：
\begin{enumerate}
\item $\mathcal{G}|_{S_\tau}$ 拟凝聚；
\item 对每个平坦且局部有限表示的态射 $g : U \to S$，典范态射
$g_{\tau, small}^*(\mathcal{G}|_{S_\tau}) \to \mathcal{G}|_{U_\tau}$
是同构。
\end{enumerate}
那么，对每个在 $S$ 上平坦且局部有限表示的 $U$，都有
$H^p(U, \mathcal{G}) = H^p(U, \mathcal{G}|_{U_\tau})$。
\end{lemma}

\begin{proof}
令 $\mathcal{F}$ 为 $\mathcal{G}|_{S_\tau}$ 到大 fppf 位点
$(\Sch/S)_{fppf}$ 的拉回。注意 $\mathcal{F}$ 拟凝聚。存在典范比较态射
$\varphi : \mathcal{F} \to \mathcal{G}$；由假设 (1) 与 (2)，它对所有
平坦且局部有限表示的 $g : U \to S$ 诱导同构
$\mathcal{F}|_{U_\tau} \to \mathcal{G}|_{U_\tau}$
。所以在短正合列
$$
0 \to \Ker(\varphi) \to \mathcal{F} \to \Im(\varphi) \to 0
$$
and
$$
0 \to \Im(\varphi) \to \mathcal{G} \to \Coker(\varphi) \to 0
$$
中，层 $\Ker(\varphi)$ 与 $\Coker(\varphi)$ 关于 fppf 拓扑寄生。由
Lemma \ref{descent-lemma-cohomology-parasitic}，可知对平坦且局部有限表示的
$g : U \to S$，态射
$H^p(U, \mathcal{F}) \to H^p(U, \mathcal{G})$ 是同构。对于
$\mathcal{F}$，结论由 Proposition
\ref{descent-proposition-same-cohomology-quasi-coherent} 成立，故命题得证。
\end{proof}













\section{Fpqc 覆盖是普遍有效满态射}
\label{descent-section-fpqc-universal-effective-epimorphisms}

\noindent
我们应用上述材料证明一个有趣的结果，即 Lemma
\ref{descent-lemma-fpqc-universal-effective-epimorphisms}。由 Sites, Section
\ref{sites-section-representable-sheaves}，该引理蕴含：当
$\tau \in \{Zariski, fppf, \etale, smooth, syntomic\}$ 时，任一位点
$(\Sch/S)_\tau$ 上的可表预层都是层。先证明一个辅助引理。

\begin{lemma}
\label{descent-lemma-equiv-fibre-product}
对概形 $X$，以 $|X|$ 表示其底集合。设 $f : X \to S$ 为概形态射。则
$$
|X \times_S X| \to |X| \times_{|S|} |X|
$$
为满射。
\end{lemma}

\begin{proof}
这立即由 Schemes, Lemma \ref{schemes-lemma-points-fibre-product} 中对
纤维积上点的描述得出。
\end{proof}


\begin{lemma}
\label{descent-lemma-universal-effective-epimorphism-affine}
设 $\{f_i : X_i \to X\}_{i \in I}$ 为一族仿射概形态射。下列条件等价：
\begin{enumerate}
\item 对任意拟凝聚 $\mathcal{O}_X$-模 $\mathcal{F}$，有
$$
\Gamma(X, \mathcal{F}) =
\text{Equalizer}\left(
\xymatrix{
\prod\nolimits_{i \in I} \Gamma(X_i, f_i^*\mathcal{F})
\ar@<1ex>[r] \ar@<-1ex>[r] &
\prod\nolimits_{i, j \in I}
\Gamma(X_i \times_X X_j, (f_i \times f_j)^*\mathcal{F})
}
\right)
$$
\item $\{f_i : X_i \to X\}_{i \in I}$ 在仿射概形范畴中是普遍有效满态射
（Sites, Definition \ref{sites-definition-universal-effective-epimorphisms}）。
\end{enumerate}
\end{lemma}

\begin{proof}
假设 (2) 成立，并设 $\mathcal{F}$ 为拟凝聚 $\mathcal{O}_X$-模。考虑概形
（Constructions, Section \ref{constructions-section-spec}）
$$
X' = \underline{\Spec}_X(\mathcal{O}_X \oplus \mathcal{F})
$$
，其中 $\mathcal{O}_X \oplus \mathcal{F}$ 是 $\mathcal{O}_X$-代数，其乘法为
$(f, s)(f', s') = (ff', fs' + f's)$.
若 $s_i \in \Gamma(X_i, f_i^*\mathcal{F})$ 是截面，则 $s_i$ 确定下式中的
唯一元素：
$$
\Gamma(X' \times_X X_i, \mathcal{O}_{X' \times_X X_i}) =
\Gamma(X_i, \mathcal{O}_{X_i}) \oplus \Gamma(X_i, f_i^*\mathcal{F})
$$
省略该等式的证明。若 $(s_i)_{i \in I}$ 属于 (1) 中的等化子，则利用对
任意概形 $T$ 都成立的等式
$$
\Mor(T, \mathbf{A}^1_\mathbf{Z}) = \Gamma(T, \mathcal{O}_T)
$$
，可知这些截面定义一族态射
$h_i : X' \times_X X_i \to \mathbf{A}^1_\mathbf{Z}$，并且作为态射
$(X' \times_X X_i) \times_{X'} (X' \times_X X_j) \to \mathbf{A}^1_\mathbf{Z}$，
有 $h_i \circ \text{pr}_1 = h_j \circ \text{pr}_2$。由假设 (2)，得到与
各态射 $h_i$ 相容的态射 $h : X' \to \mathbf{A}^1_\mathbf{Z}$，继而确定
元素 $s \in \Gamma(X, \mathcal{F})$。省略验证 $s$ 在
$\Gamma(X_i, f_i^*\mathcal{F})$ 中映到 $s_i$。

\medskip\noindent
假设 (1)。设 $T$ 为仿射概形，并设 $h_i : X_i \to T$ 为一族态射，使得
对所有 $i, j \in I$，在 $X_i \times_X X_j$ 上有
$h_i \circ \text{pr}_1 = h_j \circ \text{pr}_2$。于是
$$
\prod h_i^\sharp :
\Gamma(T, \mathcal{O}_T)
\to
\prod \Gamma(X_i, \mathcal{O}_{X_i})
$$
映入等化子；对 $\mathcal{F} = \mathcal{O}_X$ 应用本引理的假设，便得到
环同态
$\Gamma(T, \mathcal{O}_T) \to \Gamma(X, \mathcal{O}_X)$
。该环同态对应于态射 $h : X \to T$，且 $h_i = h \circ f_i$。所以这族
态射是有效满态射。

\medskip\noindent
设 $p : Y \to X$ 为仿射概形态射。应用上一段的结果，我们来证明 $f_i$
的基变换 $g_i : Y_i \to Y$ 构成有效满态射。事实上，若 $\mathcal{G}$
为拟凝聚 $\mathcal{O}_Y$-模，则
$$
\Gamma(Y, \mathcal{G}) = \Gamma(X, p_*\mathcal{G}),\quad
\Gamma(Y_i, g_i^*\mathcal{G}) = \Gamma(X, f_i^*p_*\mathcal{G}),
$$
and
$$
\Gamma(Y_i \times_Y Y_j, (g_i \times g_j)^*\mathcal{G}) =
\Gamma(X, (f_i \times f_j)^*p_*\mathcal{G})
$$
，这是由平凡基变换公式（Cohomology of Schemes, Lemma
\ref{coherent-lemma-affine-base-change}）得到的。因此可见，这族 $g_i$
满足本引理的性质 (1)。
\end{proof}

\begin{lemma}
\label{descent-lemma-universal-effective-epimorphism-surjective}
设 $\{f_i : X_i \to X\}_{i \in I}$ 为一族概形态射。
\begin{enumerate}
\item 若这族态射在概形范畴中是普遍有效满态射，则 $\coprod f_i$ 为满射。
\item 若 $X$ 与 $X_i$ 仿射，且这族态射在仿射概形范畴中是普遍有效
满态射，则 $\coprod f_i$ 为满射。
\end{enumerate}
\end{lemma}

\begin{proof}
省略。提示：沿 $\Spec(\kappa(x)) \to X$ 作基变换，即可看出任意
$x \in X$ 都必须在像中。
\end{proof}

\begin{lemma}
\label{descent-lemma-check-universal-effective-epimorphism-affine}
设 $\{f_i : X_i \to X\}_{i \in I}$ 为一族概形态射。若对每个以 $Y$
为仿射概形的态射 $Y \to X$，基变换族 $g_i : Y_i \to Y$ 都构成有效
满态射，则 $f_i$ 这一族态射在概形范畴中构成普遍有效满态射。
\end{lemma}

\begin{proof}
设 $Y \to X$ 为概形态射。我们须证明基变换
$g_i : Y_i \to Y$ 构成有效满态射。为此，设给定概形 $T$ 以及态射
$h_i : Y_i \to T$，并且在 $Y_i \times_Y Y_j$ 上有
$h_i \circ \text{pr}_1 = h_j \circ \text{pr}_2$。选取仿射开覆盖
$Y = \bigcup V_\alpha$。令 $V_{\alpha, i}$ 为 $V_\alpha$ 在 $Y_i$ 中的
逆像。于是 $V_{\alpha, i}  \to V_\alpha$ 是 $f_i$ 沿
$V_\alpha \to X$ 的基变换。因此由假设，各限制
$h_i|_{V_{\alpha, i}}$ 来自概形态射 $h_\alpha : V_\alpha \to T$。
留给读者验证这些态射在交集上一致，并定义所需态射 $Y \to T$。见
Schemes, Section \ref{schemes-section-glueing-schemes} 中的讨论。
\end{proof}

\begin{lemma}
\label{descent-lemma-universal-effective-epimorphism}
设 $\{f_i : X_i \to X\}_{i \in I}$ 为一族仿射概形态射。假设 Lemma
\ref{descent-lemma-universal-effective-epimorphism-affine} 中的等价条件成立，并且
还假设对任意仿射概形态射 $Y \to X$，映射
$$
\coprod X_i \times_X Y \longrightarrow Y
$$
是拓扑空间的商映射（Topology, Definition
\ref{topology-definition-submersive}）。那么这族态射在概形范畴中是
普遍有效满态射。
\end{lemma}

\begin{proof}
由 Lemma \ref{descent-lemma-check-universal-effective-epimorphism-affine}，只须沿以
$Y$ 为仿射概形的 $Y \to X$ 对这族态射作基变换。置
$Y_i = X_i \times_X Y$。设 $T$ 为概形，并设 $h_i : Y_i \to T$ 为一族
态射，使得在 $Y_i \times_Y Y_j$ 上有
$h_i \circ \text{pr}_1 = h_j \circ \text{pr}_2$。注意，作为集合，$Y$ 是
从 $\coprod Y_i \times_Y Y_j$ 到 $\coprod Y_i$ 的两个映射的余等化子。
事实上，满射性由 Lemma
\ref{descent-lemma-universal-effective-epimorphism-surjective} 的仿射情形给出，
单射性由 Lemma \ref{descent-lemma-equiv-fibre-product} 给出。因此存在与各
$h_i$ 相容的底集合映射 $h : Y \to T$。由本引理的第二个条件，可知
$h$ 连续！所以，若 $y \in Y$，且 $U \subset T$ 是 $h(y)$ 的仿射开邻域，
则可找到仿射开集 $V \subset Y$，使得 $h(V) \subset U$。置
$V_i = Y_i \times_Y V = X_i \times_X V$，即可应用 Lemma
\ref{descent-lemma-universal-effective-epimorphism-affine} 中证明的结果，看出
$h|_V : V \to U \subset T$ 来自唯一的仿射概形态射
$h_V : V \to U$，且对所有 $i$，它作为概形态射与
$h_i|_{V_i}$ 一致。粘合这些 $h_V$（见 Schemes, Section
\ref{schemes-section-glueing-schemes}），便得到所需态射 $Y \to T$。
\end{proof}

\begin{lemma}
\label{descent-lemma-open-fpqc-covering}
设 $\{f_i : T_i \to T\}_{i \in I}$ 为 fpqc 覆盖。假设对每个 $i$ 都有开
子集 $W_i \subset T_i$，并且对所有 $i, j \in I$，作为
$T_i \times_T T_j$ 的开子集，有
$\text{pr}_0^{-1}(W_i) = \text{pr}_1^{-1}(W_j)$。则存在唯一的开子集
$W \subset T$，使得对每个 $i$ 都有 $W_i = f_i^{-1}(W)$。
\end{lemma}

\begin{proof}
把 Lemma \ref{descent-lemma-equiv-fibre-product} 应用于映射
$\coprod_{i \in I} T_i \to T$。它蕴含存在子集 $W \subset T$，使得对每个
$i$ 都有 $W_i = f_i^{-1}(W)$；具体地，$W = \bigcup f_i(W_i)$。为证明
$W$ 开，可以在 $T$ 上 Zariski 局部地工作。因此可设 $T$ 仿射。利用
Topologies, Definition \ref{topologies-definition-fpqc-covering}，可选取
一个加细 $\{T_i \to T\}_{i \in I}$ 的标准 fpqc 覆盖
$\{g_j : V_j \to T\}_{j \in J}$。令 $\alpha : J \to I$ 与
$h_j : V_j \to T_{\alpha(j)}$ 如 Sites, Definition
\ref{sites-definition-morphism-coverings} 中所述。于是
$g_j^{-1}(W) = h_j^{-1}(W_{\alpha(j)})$。所以可设
$\{f_i : T_i \to T\}$ 是标准 fpqc 覆盖。在这种情形，对态射
$\coprod T_i \to T$ 应用 Morphisms, Lemma
\ref{morphisms-lemma-fpqc-quotient-topology}，即可断定 $W$ 开。
\end{proof}

\begin{lemma}
\label{descent-lemma-fpqc-universal-effective-epimorphisms}
设 $\{T_i \to T\}$ 为 fpqc 覆盖；见 Topologies, Definition
\ref{topologies-definition-fpqc-covering}。则 $\{T_i \to T\}$ 在概形范畴
中是普遍有效满态射；见 Sites, Definition
\ref{sites-definition-universal-effective-epimorphisms}。换言之，概形范畴
上的每个可表函子都满足 fpqc 拓扑的层条件；见 Topologies, Definition
\ref{topologies-definition-sheaf-property-fpqc}。
\end{lemma}

\begin{proof}
设 $S$ 为概形。我们须证明如下陈述：给定态射
$\varphi_i : T_i \to S$，满足
$\varphi_i|_{T_i \times_T T_j} = \varphi_j|_{T_i \times_T T_j}$，则存在
唯一的态射 $T \to S$，其在每个 $T_i$ 上的限制为 $\varphi_i$。换言之，
我们须证明函子 $h_S = \Mor_{\Sch}( - , S)$ 满足 fpqc 拓扑的层性质。

\medskip\noindent
若 $\{T_i \to T\}$ 是 Zariski 覆盖，则结论由 Schemes, Lemma
\ref{schemes-lemma-glue} 得出。因此 Topologies, Lemma
\ref{topologies-lemma-sheaf-property-fpqc} 把问题归约到覆盖
$\{X \to Y\}$ 的情形，其中该覆盖由仿射概形之间的单个满平坦态射给出。

\medskip\noindent
第一种证明。由 Lemma \ref{descent-lemma-sheaf-condition-holds}，拟凝聚模关于
$\{X \to Y\}$ 满足层条件。由 Lemma \ref{descent-lemma-open-fpqc-covering}，态射
$X \to Y$ 是普遍商映射。因此可应用 Lemma
\ref{descent-lemma-universal-effective-epimorphism}，看出 $\{X \to Y\}$ 是普遍
有效满态射。

\medskip\noindent
第二种证明。令 $R \to A$ 为与满平坦态射 $\pi : X \to Y$ 对应的忠实
平坦环同态。设 $f : X \to S$ 为态射，并且作为态射
$X \times_Y X = \Spec(A \otimes_R A) \to S$，有
$f \circ \text{pr}_1 = f \circ \text{pr}_2$。由 Lemma
\ref{descent-lemma-equiv-fibre-product}，可知在底集合映射的意义下，对某个
（集合论）映射 $g : \Spec(R) \to S$，$f$ 具有形式
$f = g \circ \pi$。由 Morphisms, Lemma
\ref{morphisms-lemma-fpqc-quotient-topology} 以及 $f$ 连续这一事实，可知
$g$ 连续。

\medskip\noindent
取 $y \in Y = \Spec(R)$。选取包含 $g(y)$ 的仿射开集 $U \subset S$，
记 $U = \Spec(B)$。由上述结果，可选取 $r \in R$，使得
$y \in D(r) \subset g^{-1}(U)$。$f$ 在 $\pi^{-1}(D(r))$ 上取值于 $U$
的限制对应于环同态 $B \to A_r$。由对 $f$ 的假设，所诱导的两个环同态
$B \to A_r \otimes_{R_r} A_r = (A \otimes_R A)_r$ 相等。注意
$R_r \to A_r$ 忠实平坦。由 Lemma \ref{descent-lemma-ff-exact}，两个箭头
$A_r \to A_r \otimes_{R_r} A_r$ 的等化子是 $R_r$。所以
$B \to A_r$ 唯一地通过一个映射 $B \to R_r$ 分解。该映射继而给出概形
态射 $D(r) \to U \to S$；见 Schemes, Lemma
\ref{schemes-lemma-morphism-into-affine}。

\medskip\noindent
到目前为止我们证明了什么？我们已经证明，对任意素理想
$\mathfrak p \subset R$，都存在标准仿射开集 $D(r) \subset \Spec(R)$，
使得态射 $f|_{\pi^{-1}(D(r))} : \pi^{-1}(D(r)) \to S$ 唯一地通过某个
概形态射 $D(r) \to S$ 分解。省略验证这些态射粘合为所需态射
$\Spec(R) \to S$。
\end{proof}

\begin{lemma}
\label{descent-lemma-coequalizer-fpqc-local}
考虑概形 $X, Y, Z$、态射 $a, b : X \to Y$ 以及态射
$c : Y \to Z$，其中 $c \circ a = c \circ b$。置
$d = c \circ a = c \circ b$。若存在 fpqc 覆盖 $\{Z_i \to Z\}$，使得
\begin{enumerate}
\item 对所有 $i$，态射 $Y \times_{c, Z} Z_i \to Z_i$ 是
$(a, 1) : X \times_{d, Z} Z_i \to Y \times_{c, Z} Z_i$ 与
$(b, 1) : X \times_{d, Z} Z_i \to Y \times_{c, Z} Z_i$ 的余等化子；
\item 对所有 $i$ 与 $i'$，态射
$Y \times_{c, Z} (Z_i \times_Z Z_{i'}) \to (Z_i \times_Z Z_{i'})$
是
$(a, 1) : X \times_{d, Z} (Z_i \times_Z Z_{i'}) \to
Y \times_{c, Z} (Z_i \times_Z Z_{i'})$ 与
$(b, 1) : X \times_{d, Z} (Z_i \times_Z Z_{i'}) \to
Y \times_{c, Z} (Z_i \times_Z Z_{i'})$
的余等化子。
\end{enumerate}
则 $c$ 是 $a$ 与 $b$ 的余等化子。
\end{lemma}

\begin{proof}
事实上，由 Lemma \ref{descent-lemma-fpqc-universal-effective-epimorphisms}，
对概形 $T$，态射 $Z \to T$ 等同于一族在交叠上一致的态射
$Z_i \to T$。
\end{proof}















\section{态射有限性与光滑性性质的下降}
\label{descent-section-descent-finiteness-morphisms}

\noindent
本节证明态射的若干性质（光滑、局部有限表示等）沿忠实平坦态射下降。
先从代数版本开始。（关注 ``Noetherian'' 情形的读者应参看 Lemma
\ref{descent-lemma-finite-type-local-source-fppf-algebra}，而不是下一个引理。）

\begin{lemma}
\label{descent-lemma-flat-finitely-presented-permanence-algebra}
设 $R \to A \to B$ 为环同态。假设 $R \to B$ 有限表示，而
$A \to B$ 忠实平坦且有限表示。则 $R \to A$ 有限表示。
\end{lemma}

\begin{proof}
考虑代数 $C = B \otimes_A B$，以及由 $p(b) = b \otimes 1$ 与
$q(b) = 1 \otimes b$ 给出的两个映射 $p, q : B \to C$。两个复合
$A \to B \to C$ 当然相同。注意，$p : B \to C$ 是 $A \to B$ 的基变换，
因而平坦且有限表示；所以环同态 $R \to C$ 作为
$R \to B \to C$ 的复合也是有限表示的。

\medskip\noindent
我们将利用 Algebra, Lemma
\ref{algebra-lemma-characterize-finite-presentation} 的判别准则，证明
$R \to A$ 有限表示。设 $S$ 为任意 $R$-代数，并设
$S = \colim_{\lambda \in \Lambda} S_\lambda$ 写成 $R$-代数的有向余极限。
设 $A \to S$ 为 $R$-代数同态。我们须证明 $A \to S$ 通过某个
$S_\lambda$ 分解。考虑环 $B' = S \otimes_A B$ 与
$C' = S \otimes_A C = B' \otimes_S B'$。由于 $B$ 在 $A$ 上忠实平坦且
有限表示，$B'$ 在 $S$ 上也忠实平坦且有限表示。把 Algebra, Lemma
\ref{algebra-lemma-flat-finite-presentation-limit-flat} 的第 (2) 部分应用于
偶对 $(S \to B', B')$ 和系统 $(S_\lambda)$，可知存在
$\lambda_0 \in \Lambda$ 以及平坦有限表示的 $S_{\lambda_0}$-代数
$B_{\lambda_0}$，使得
$B' = S \otimes_{S_{\lambda_0}} B_{\lambda_0}$。对
$\lambda \geq \lambda_0$，置
$B_\lambda = S_\lambda \otimes_{S_{\lambda_0}} B_{\lambda_0}$ 且
$C_\lambda = B_\lambda \otimes_{S_\lambda} B_\lambda$。

\medskip\noindent
暂时中断论证，以证明当 $\lambda$ 足够大时，
$S_\lambda \to B_\lambda$ 忠实平坦。（这其实应该另立为其他地方的一个
引理，或许放在极限章中。）由于
$\Spec(B_{\lambda_0}) \to \Spec(S_{\lambda_0})$ 平坦且有限表示，它是开
映射（见 Morphisms, Lemma \ref{morphisms-lemma-fppf-open}）。令
$I \subset S_{\lambda_0}$ 为理想，使得 $V(I) \subset \Spec(S_{\lambda_0})$
是该像的补集。注意，像的形成与基变换交换。因此，由于
$\Spec(B') \to \Spec(S)$ 为满射，且
$B' = B_{\lambda_0} \otimes_{S_{\lambda_0}} S$，可知 $IS = S$。所以对
某个 $\lambda \geq \lambda_0$，有 $IS_{\lambda} = S_\lambda$。对这个指标
以及所有更大的 $\lambda$，态射
$\Spec(B_\lambda) \to \Spec(S_\lambda)$ 都是满射。

\medskip\noindent
仿照证明第一段的记号，以
$p_\lambda, q_\lambda : B_\lambda \to C_\lambda$ 表示两个典范映射。
于是 $B' = \colim_{\lambda \geq \lambda_0} B_\lambda$ 且
$C' = \colim_{\lambda \geq \lambda_0} C_\lambda$。由于 $B$ 与 $C$ 在
$R$ 上有限表示，反复应用 Algebra, Lemma
\ref{algebra-lemma-characterize-finite-presentation}，可知存在
$\lambda \geq \lambda_0$ 以及 $R$-代数同态
$B \to B_\lambda$、$C \to C_\lambda$，使得图表
$$
\xymatrix{
C \ar[rr] & &
C_\lambda \\
B \ar[rr]
\ar@<1ex>[u]^-p
\ar@<-1ex>[u]_-q
& &
B_\lambda
\ar@<1ex>[u]^-{p_\lambda}
\ar@<-1ex>[u]_-{q_\lambda}
}
$$
交换。这意味着 $A \to B \to B_\lambda$ 映入 $p_\lambda$ 与
$q_\lambda$ 的等化子。由 Lemma \ref{descent-lemma-ff-exact}，可知
$S_\lambda$ 是 $p_\lambda$ 与 $q_\lambda$ 的等化子。因此得到所需环同态
$A \to S_\lambda$，命题得证。
\end{proof}

\noindent
下面给出关于有限型性质的较容易版本。

\begin{lemma}
\label{descent-lemma-finite-type-local-source-fppf-algebra}
设 $R \to A \to B$ 为环同态。假设 $R \to B$ 有限型，而
$A \to B$ 忠实平坦且有限表示。则 $R \to A$ 有限型。
\end{lemma}

\begin{proof}
由 Algebra, Lemma
\ref{algebra-lemma-descend-faithfully-flat-finite-presentation}，存在交换图
$$
\xymatrix{
R \ar[r] \ar@{=}[d] &
A_0 \ar[d] \ar[r] &
B_0 \ar[d] \\
R \ar[r] & A \ar[r] & B
}
$$
，其中 $R \to A_0$ 有限表示，$A_0 \to B_0$ 忠实平坦且有限表示，并且
$B = A \otimes_{A_0} B_0$。由假设，$R \to B$ 有限型，所以可以向
$A_0$ 添加一些元素，并假设映射 $B_0 \to B$ 为满射！此时，由于
$A_0 \to B_0$ 忠实平坦，而
$$
(A_0 \to A) \otimes_{A_0} B_0 \cong (B_0 \to B)
$$
为满射，所以 $A_0 \to A$ 也为满射。命题得证。
\end{proof}

\begin{lemma}
\label{descent-lemma-flat-finitely-presented-permanence}
\begin{reference}
\cite[IV, 17.7.5 (i) and (ii)]{EGA}.
\end{reference}
设
$$
\xymatrix{
X \ar[rr]_f \ar[rd]_p & &
Y \ar[dl]^q \\
& S
}
$$
为概形态射的交换图。假设 $f$ 满、平坦且局部有限表示，并假设 $p$
局部有限表示（相应地，局部有限型）。则 $q$ 局部有限表示（相应地，
局部有限型）。
\end{lemma}

\begin{proof}
该问题在 $S$ 与 $Y$ 上是局部的。因此可设 $S$ 与 $Y$ 仿射。由于 $f$
平坦且局部有限表示，$f$ 是开映射（Morphisms, Lemma
\ref{morphisms-lemma-fppf-open}）。又因为 $Y$ 拟紧，存在有限多个仿射开集
$X_i \subset X$，使得 $Y = \bigcup f(X_i)$。显然可以用
$\coprod X_i$ 替换 $X$，从而也可设 $X$ 仿射。在这种情形，本引理等价
于上面的 Lemma \ref{descent-lemma-flat-finitely-presented-permanence-algebra}
（相应地，Lemma \ref{descent-lemma-finite-type-local-source-fppf-algebra}）。
\end{proof}

\noindent
我们用它改进前面得到的一些关于态射的结果。

\begin{lemma}
\label{descent-lemma-syntomic-smooth-etale-permanence}
设
$$
\xymatrix{
X \ar[rr]_f \ar[rd]_p & &
Y \ar[dl]^q \\
& S
}
$$
为概形态射的交换图。假设
\begin{enumerate}
\item $f$ 满且 syntomic（相应地，smooth；相应地，\'etale）；
\item $p$ syntomic（相应地，smooth；相应地，\'etale）。
\end{enumerate}
则 $q$ syntomic（相应地，smooth；相应地，\'etale）。
\end{lemma}

\begin{proof}
把 Morphisms, Lemmas
\ref{morphisms-lemma-syntomic-permanence},
\ref{morphisms-lemma-smooth-permanence}，以及
\ref{morphisms-lemma-etale-permanence-two}
与上面的 Lemma \ref{descent-lemma-flat-finitely-presented-permanence} 结合即可。
\end{proof}

\noindent
事实上，可以把这一结果加强如下。

\begin{lemma}
\label{descent-lemma-smooth-permanence}
设
$$
\xymatrix{
X \ar[rr]_f \ar[rd]_p & &
Y \ar[dl]^q \\
& S
}
$$
为概形态射的交换图。假设
\begin{enumerate}
\item $f$ 满、平坦且局部有限表示；
\item $p$ smooth（相应地，\'etale）。
\end{enumerate}
则 $q$ smooth（相应地，\'etale）。
\end{lemma}

\begin{proof}
假设 (1) 成立且 $p$ smooth。由 Lemma
\ref{descent-lemma-flat-finitely-presented-permanence}，可知 $q$ 局部有限表示。
由 Morphisms, Lemma \ref{morphisms-lemma-flat-permanence}，可知 $q$ 平坦。
所以现在只须证明 $q$ 的纤维光滑；见 Morphisms, Lemma
\ref{morphisms-lemma-smooth-flat-smooth-fibres}。对 $s \in S$，把 Varieties,
Lemma \ref{varieties-lemma-flat-under-smooth} 应用于平坦满态射
$X_s \to Y_s$，即得结论。省略 \'etale 情形的证明。
\end{proof}

\begin{remark}
\label{descent-remark-smooth-permanence}
在 Lemma \ref{descent-lemma-smooth-permanence} 的假设 (1) 以及 $p$ smooth 的条件
下，$X \to Y$ 不一定 smooth。一个反例是 $S = \Spec(k)$、
$X = \Spec(k[s])$、$Y = \Spec(k[t])$，且 $f$ 由 $t \mapsto s^2$ 给出。
不过，关于 $f$ 的结构的一些信息，另见 Lemma
\ref{descent-lemma-syntomic-permanence}。
\end{remark}

\begin{lemma}
\label{descent-lemma-syntomic-permanence}
设
$$
\xymatrix{
X \ar[rr]_f \ar[rd]_p & &
Y \ar[dl]^q \\
& S
}
$$
为概形态射的交换图。假设
\begin{enumerate}
\item $f$ 满、平坦且局部有限表示；
\item $p$ syntomic。
\end{enumerate}
则 $q$ 与 $f$ 都 syntomic。
\end{lemma}

\begin{proof}
% P05-ERR-0034
由 Lemma \ref{descent-lemma-flat-finitely-presented-permanence}，可知 $q$ 有限表示。
由 Morphisms, Lemma \ref{morphisms-lemma-flat-permanence}，可知 $q$ 平坦。
由 Morphisms, Lemma
\ref{morphisms-lemma-syntomic-locally-standard-syntomic}，现在只须证明
$Y \to S$ 的纤维以及 $X \to Y$ 的纤维的局部环都是局部完全交环。为此，
可取 $X \to Y \to S$ 在一点 $s \in S$ 上的纤维；也就是说，可以假设
$S$ 是某个域的谱。选取一点 $x \in X$，其像为 $y \in Y$，并考虑环同态
$$
\mathcal{O}_{Y, y} \longrightarrow \mathcal{O}_{X, x}
$$
。这是局部 Noether 环之间的平坦局部同态。局部环
$\mathcal{O}_{X, x}$ 是完全交。于是可应用 Avramov 的结果；见 Divided
Power Algebra, Lemma \ref{dpa-lemma-avramov}；从而断定
$\mathcal{O}_{Y, y}$ 与
$\mathcal{O}_{X, x}/\mathfrak m_y\mathcal{O}_{X, x}$ 都是完全交环。
\end{proof}

\noindent
下面这类引理有时很有用。

\begin{lemma}
\label{descent-lemma-curiosity}
设 $X \to Y \to Z$ 为概形态射。令 $P$ 为概形态射的下列性质之一：
平坦、局部有限型、局部有限表示。假设 $X \to Z$ 具有 $P$，并且
$\{X \to Y\}$ 可由 $Y$ 的一个 fppf 覆盖加细。则 $Y \to Z$ 具有 $P$。
\end{lemma}

\begin{proof}
设 $\Spec(C) \subset Z$ 为仿射开集，并设 $\Spec(B) \subset Y$ 为映入
$\Spec(C)$ 的仿射开集。关于 $X \to Y$ 的假设蕴含：可以找到标准仿射
fppf 覆盖 $\{\Spec(B_j) \to \Spec(B)\}$ 以及提升
$x_j : \Spec(B_j) \to X$。由于 $\Spec(B_j)$ 拟紧，可以找到位于
$\Spec(B)$ 上方的有限多个仿射开集 $\Spec(A_i) \subset X$，使得每个
$x_j$ 的像都包含在并集 $\bigcup \Spec(A_i)$ 中。因此，用各
$\Spec(B_j)$ 自身的标准仿射 Zariski 覆盖替换它之后，可以假设有标准
仿射 fppf 覆盖 $\{\Spec(B_i) \to \Spec(B)\}$，使得每个
$\Spec(B_i) \to Y$ 都通过位于 $\Spec(B)$ 上方的一个仿射开集
$\Spec(A_i) \subset X$ 分解。换言之，对每个 $i$，都有环同态
$C \to B \to A_i \to B_i$。注意，还可以考虑
$$
C \to B \to A = \prod A_i \to B' = \prod B_i
$$
，并且环同态 $B \to \prod B_i$ 忠实平坦且有限表示。

\medskip\noindent
$P = flat$ 的情形。此时已知 $C \to A$ 平坦，须证明 $C \to B$ 平坦。
设 $N \to N' \to N''$ 为 $C$-模的正合列。我们要证明
$N \otimes_C B \to N' \otimes_C B \to N'' \otimes_C B$ 正合。令 $H$
为其上同调，并令 $H'$ 为
$N \otimes_C B' \to N' \otimes_C B' \to N'' \otimes_C B'$ 的上同调。
由于 $B \to B'$ 平坦，有 $H' = H \otimes_B B'$。另一方面，
$N \otimes_C A \to N' \otimes_C A \to N'' \otimes_C A$ 正合，因而上同调
为零。所以态射 $H \to H'$ 为零（因为它通过零模分解）。故
$H' = 0$。由于 $B \to B'$ 忠实平坦，可知 $H = 0$，正如所需。

\medskip\noindent
$P = locally\ finite\ type$ 的情形。此时已知 $C \to A$ 有限型，须证明
$C \to B$ 有限型。由于 $B \to B'$ 有限表示（因而有限型），可知
$A \to B'$ 有限型；见 Algebra, Lemma
\ref{algebra-lemma-compose-finite-type}。所以 $C \to B'$ 有限型，结论由
Lemma \ref{descent-lemma-finite-type-local-source-fppf-algebra} 得出。

\medskip\noindent
$P = locally\ finite\ presentation$ 的情形。此时已知 $C \to A$ 有限表示，
须证明 $C \to B$ 有限表示。由于 $B \to B'$ 有限表示，而 $B \to A$
有限型，可知 $A \to B'$ 有限表示；见 Algebra, Lemma
\ref{algebra-lemma-compose-finite-type}。所以 $C \to B'$ 有限表示，结论
由 Lemma \ref{descent-lemma-flat-finitely-presented-permanence-algebra} 得出。
\end{proof}










\section{概形的局部性质}
\label{descent-section-descending-properties}

\noindent
常常可以证明概形的一个覆盖的各成员具有某种性质。在许多情形下，这蕴含
该概形也具有此性质。例如，若 $S$ 是概形，$f : S' \to S$ 是满平坦态射，
且 $S'$ 是既约概形，则 $S$ 既约。考察局部环并使用 Algebra, Lemma
\ref{algebra-lemma-descent-reduced} 即可证明这一点。我们称既约这一性质
{\it 沿平坦满态射下降}。这类结果的一部分汇集于 Algebra, Section
\ref{algebra-section-descending-properties}；关于概形的结果见 Section
\ref{descent-section-variants}。关于态射性质下降的一些类似结果见 Section
\ref{descent-section-descent-finiteness-morphisms}。

\medskip\noindent
另一方面，也存在满平坦态射 $f : S' \to S$，其中 $S$ 既约而 $S'$ 不
既约；例如态射 $\Spec(k[x]/(x^2)) \to \Spec(k)$。因此既约这一性质不
{\it 沿平坦态射上升}。剩余域无限这一性质确实沿平坦态射上升（当然不
沿满平坦态射下降）。这类结果的一部分汇集于 Algebra, Section
\ref{algebra-section-ascending-properties}。

\medskip\noindent
最后，若一个性质沿平坦态射上升，并沿平坦满态射下降，则称它
{\it 关于平坦拓扑为局部的}。一个略显琐碎的例子是剩余域具有给定特征
这一性质。为给出更精确的表述，并把它与概形上的各种拓扑联系起来，作
如下正式定义。

\begin{definition}
\label{descent-definition-property-local}
设 $\mathcal{P}$ 为概形的一项性质。设
$\tau \in \{fpqc, \linebreak[0] fppf, \linebreak[0] syntomic, \linebreak[0]
smooth, \linebreak[0] \etale, \linebreak[0] Zariski\}$.
若对任意 $\tau$-覆盖 $\{S_i \to S\}_{i \in I}$（见 Topologies, Section
\ref{topologies-section-procedure}）都有
$$
S \text{ 具有 }\mathcal{P}
\Leftrightarrow
\text{每个 }S_i \text{ 具有 }\mathcal{P}.
$$
则称 $\mathcal{P}$ {\it 在 $\tau$-拓扑中为局部的}。
\end{definition}

\noindent
确切地说，由于同构总是覆盖，可知（或要求）$S$ 具有性质
$\mathcal{P}$，当且仅当任意与 $S$ 同构的概形 $S'$ 具有该性质。事实上，
若 $\tau = fpqc, \linebreak[0] fppf, \linebreak[0] syntomic,
\linebreak[0] smooth, \linebreak[0] \etale$ 或 $Zariski$，且 $S$ 具有
$\mathcal{P}$，而 $S' \to S$ 相应地为平坦、平坦且局部有限表示、
syntomic、smooth、\'etale 或开浸入，则 $S'$ 具有 $\mathcal{P}$。这是
因为总能把 $\{S' \to S\}$ 扩充为一个 $\tau$-覆盖。

\medskip\noindent
有如下蕴含关系：
$\mathcal{P}$ 在 fpqc 拓扑中为局部的
$\Rightarrow$
$\mathcal{P}$ 在 fppf 拓扑中为局部的
$\Rightarrow$
$\mathcal{P}$ 在 syntomic 拓扑中为局部的
$\Rightarrow$
$\mathcal{P}$ 在 smooth 拓扑中为局部的
$\Rightarrow$
$\mathcal{P}$ 在 \'etale 拓扑中为局部的
$\Rightarrow$
$\mathcal{P}$ 在 Zariski 拓扑中为局部的。这由 Topologies, Lemmas
\ref{topologies-lemma-zariski-etale},
\ref{topologies-lemma-zariski-etale-smooth},
\ref{topologies-lemma-zariski-etale-smooth-syntomic},
\ref{topologies-lemma-zariski-etale-smooth-syntomic-fppf}, and
\ref{topologies-lemma-zariski-etale-smooth-syntomic-fppf-fpqc}.
得出。

\begin{lemma}
\label{descent-lemma-descending-properties}
设 $\mathcal{P}$ 为概形的一项性质。设
$\tau \in \{fpqc, \linebreak[0] fppf, \linebreak[0]
\etale, \linebreak[0] smooth, \linebreak[0] syntomic\}$。假设
\begin{enumerate}
\item 该性质在 Zariski 拓扑中为局部的；
\item 对任意仿射概形态射 $S' \to S$，若按照 $\tau$ 为 fpqc、fppf、
\'etale、smooth 或 syntomic，该态射相应地为平坦、平坦且有限表示、
\'etale、smooth 或 syntomic，则只要 $S$ 具有性质 $\mathcal{P}$，$S'$
也具有性质 $\mathcal{P}$；
\item 对任意满的仿射概形态射 $S' \to S$，若按照 $\tau$ 为 fpqc、fppf、
\'etale、smooth 或 syntomic，该态射相应地为平坦、平坦且有限表示、
\'etale、smooth 或 syntomic，则只要 $S'$ 具有性质 $\mathcal{P}$，$S$
也具有性质 $\mathcal{P}$。
\end{enumerate}
则 $\mathcal{P}$ 在基底上是 $\tau$-局部的。
\end{lemma}

\begin{proof}
这几乎立即由 $\tau$-覆盖的定义得出；见 Topologies, Definition
\ref{topologies-definition-fpqc-covering}
\ref{topologies-definition-fppf-covering}
\ref{topologies-definition-etale-covering}
\ref{topologies-definition-smooth-covering}，或
\ref{topologies-definition-syntomic-covering}
以及 Topologies, Lemma
\ref{topologies-lemma-fpqc-affine},
\ref{topologies-lemma-fppf-affine},
\ref{topologies-lemma-etale-affine},
\ref{topologies-lemma-smooth-affine}，或
\ref{topologies-lemma-syntomic-affine}。省略细节。
\end{proof}

\begin{remark}
\label{descent-remark-descending-properties-standard}
在上面的 Lemma \ref{descent-lemma-descending-properties} 中，若
$\tau = smooth$，则在条件 (3) 中可以假设该态射是（满的）标准 smooth
态射。$\tau = syntomic$ 或 $\tau = \etale$ 时也类似。
\end{remark}





\section{在 fppf 拓扑中为局部的概形性质}
\label{descent-section-descending-properties-fppf}

\noindent
本节给出一些在 fppf 拓扑中于基底上为局部的概形性质。

\begin{lemma}
\label{descent-lemma-Noetherian-local-fppf}
性质 $\mathcal{P}(S) =$``$S$ 局部 Noether'' 在 fppf 拓扑中为局部的。
\end{lemma}

\begin{proof}
我们使用 Lemma \ref{descent-lemma-descending-properties}。首先注意，``局部
Noether'' 在 Zariski 拓扑中为局部的。这由定义即得；见 Properties,
Definition \ref{properties-definition-noetherian}。接着，若
$S' \to S$ 是仿射概形之间的平坦有限表示态射，而 $S$ 局部 Noether，
则 $S'$ 局部 Noether；这是 Morphisms, Lemma
\ref{morphisms-lemma-finite-type-noetherian}。最后，若 $S' \to S$ 是仿射
概形之间的满、平坦、有限表示态射，而 $S'$ 局部 Noether，则须证明
$S$ 局部 Noether。这由 Algebra, Lemma
\ref{algebra-lemma-descent-Noetherian} 得出。因此 Lemma
\ref{descent-lemma-descending-properties} 的 (1)、(2)、(3) 成立，命题得证。
\end{proof}

\begin{lemma}
\label{descent-lemma-Jacobson-local-fppf}
性质 $\mathcal{P}(S) =$``$S$ 是 Jacobson 概形'' 在 fppf 拓扑中为局部的。
\end{lemma}

\begin{proof}
我们使用 Lemma \ref{descent-lemma-descending-properties}。首先注意，``Jacobson''
在 Zariski 拓扑中为局部的；这是 Properties, Lemma
\ref{properties-lemma-locally-jacobson}。接着，若 $S' \to S$ 是仿射概形
之间的平坦有限表示态射，而 $S$ 是 Jacobson 概形，则 $S'$ 是 Jacobson
概形；这是 Morphisms, Lemma
\ref{morphisms-lemma-Jacobson-universally-Jacobson}。最后，设
$f : S' \to S$ 是仿射概形之间的满、平坦、有限表示态射，且 $S'$ 是
Jacobson 概形；须证明 $S$ 是 Jacobson 概形。记 $S = \Spec(A)$、
$S' = \Spec(B)$，且 $S' \to S$ 由 $A \to B$ 给出。于是 $A \to B$
有限表示且忠实平坦。此外，环 $B$ 是 Jacobson 环；见 Properties,
Lemma \ref{properties-lemma-locally-jacobson}。

\medskip\noindent
由 Algebra, Lemma \ref{algebra-lemma-fppf-fpqf}，存在图表
$$
\xymatrix{
B \ar[rr] & & B' \\
& A \ar[ru] \ar[lu] &
}
$$
，其中 $A \to B'$ 有限表示、忠实平坦且拟有限。特别地，$B \to B'$
有限型，而由 Algebra, Proposition
\ref{algebra-proposition-Jacobson-permanence}，可知 $B'$ 是 Jacobson 环。
所以可以假设 $A \to B$ 不仅忠实平坦且有限表示，而且拟有限。

\medskip\noindent
反设 $A$ 不是 Jacobson 环以导出矛盾。由 Algebra, Lemma
\ref{algebra-lemma-characterize-jacobson}，存在非极大素理想
$\mathfrak p \subset A$ 以及元素 $f \in A$，$f \not \in \mathfrak p$，
使得 $V(\mathfrak p) \cap D(f) = \{\mathfrak p\}$。

\medskip\noindent
如下导出矛盾。先取 $A$ 的极大理想，使得
$\mathfrak p \subset \mathfrak m$。选取位于 $\mathfrak m$ 上方的素理想
$\mathfrak m' \subset B$（它存在，因为 $A \to B$ 忠实平坦；见 Algebra,
Lemma \ref{algebra-lemma-ff-rings}）。由于 $A \to B$ 平坦，由下降定理
（见 Algebra, Lemma \ref{algebra-lemma-flat-going-down}），可找到位于
$\mathfrak p$ 上方的素理想 $\mathfrak q \subset \mathfrak m'$。特别地，
$\mathfrak q$ 不是极大理想。因此再次由 Algebra, Lemma
\ref{algebra-lemma-characterize-jacobson}，集合
$V(\mathfrak q) \cap D(f)$ 是无限的（此处终于用到 $B$ 是 Jacobson 环）。
$V(\mathfrak q) \cap D(f)$ 的所有点都映到
$V(\mathfrak p) \cap D(f) = \{\mathfrak p\}$。因此
$\mathfrak p$ 上的纤维是无限的。这与 $A \to B$ 拟有限矛盾（见 Algebra,
Lemma \ref{algebra-lemma-quasi-finite}，或更明确地见 Morphisms, Lemma
\ref{morphisms-lemma-quasi-finite}）。引理得证。
\end{proof}

\begin{lemma}
\label{descent-lemma-locally-finite-nr-irred-local-fppf}
性质 $\mathcal{P}(S) =$``$S$ 的每个拟紧开集都只有有限多个不可约分支''
在 fppf 拓扑中为局部的。
\end{lemma}

\begin{proof}
我们使用 Lemma \ref{descent-lemma-descending-properties}。首先注意，
$\mathcal{P}$ 在 Zariski 拓扑中为局部的。接着，若 $T \to S$ 是仿射概形
之间的平坦有限表示态射，而 $S$ 只有有限多个不可约分支，则须证明
$T$ 也如此。事实上，由于 $T \to S$ 平坦，$T$ 的一般点映到 $S$ 的
一般点；见 Morphisms, Lemma
\ref{morphisms-lemma-generalizations-lift-flat}。所以只须证明，对
$s \in S$，纤维 $T_s$ 只有有限多个一般点。注意，$T_s$ 是
$\kappa(s)$ 上的有限型仿射概形；见 Morphisms, Lemma
\ref{morphisms-lemma-base-change-finite-type}。因此 $T_s$ 是 Noether
概形，并且只有有限多个不可约分支（Morphisms, Lemma
\ref{morphisms-lemma-finite-type-noetherian} 及 Properties, Lemma
\ref{properties-lemma-Noetherian-irreducible-components}）。最后，若
$T \to S$ 是仿射概形之间的满、平坦、有限表示态射，而 $T$ 只有有限
多个不可约分支，则须证明 $S$ 也如此。这由 Topology, Lemma
\ref{topology-lemma-surjective-continuous-irreducible-components} 得出。
因此 Lemma \ref{descent-lemma-descending-properties} 的 (1)、(2)、(3) 成立，
命题得证。
\end{proof}




\section{在 syntomic 拓扑中为局部的概形性质}
\label{descent-section-descending-properties-syntomic}

\noindent
本节给出一些在 syntomic 拓扑中于基底上为局部的概形性质。

\begin{lemma}
\label{descent-lemma-Sk-local-syntomic}
性质 $\mathcal{P}(S) =$``$S$ 局部 Noether 且满足 $(S_k)$'' 在 syntomic
拓扑中为局部的。
\end{lemma}

\begin{proof}
我们验证 Lemma \ref{descent-lemma-descending-properties} 的 (1)、(2)、(3)。由于
syntomic 态射平坦且有限表示（Morphisms, Lemmas
\ref{morphisms-lemma-syntomic-flat} 与
\ref{morphisms-lemma-syntomic-locally-finite-presentation}），在 Lemma
\ref{descent-lemma-Noetherian-local-fppf} 的证明中已经对 ``局部 Noether'' 验证
了这些条件；下文不再提及而直接使用。首先注意，$\mathcal{P}$ 在 Zariski
拓扑中为局部的。这由定义即得；见 Cohomology of Schemes, Definition
\ref{coherent-definition-depth}。接着，若 $S' \to S$ 是仿射概形之间的
syntomic 态射，而 $S$ 具有 $\mathcal{P}$，则 $S'$ 具有
$\mathcal{P}$。这是 Algebra, Lemma \ref{algebra-lemma-Sk-goes-up}
（使用 Morphisms, Lemma \ref{morphisms-lemma-syntomic-characterize}、
Algebra, Definition \ref{algebra-definition-lci} 以及 Lemma
\ref{algebra-lemma-lci-CM}）。最后，若 $S' \to S$ 是仿射概形之间的满
syntomic 态射，而 $S'$ 具有 $\mathcal{P}$，则 $S$ 具有
$\mathcal{P}$。这是 Algebra, Lemma \ref{algebra-lemma-descent-Sk}。
因此 Lemma \ref{descent-lemma-descending-properties} 的 (1)、(2)、(3) 成立，
命题得证。
\end{proof}

\begin{lemma}
\label{descent-lemma-CM-local-syntomic}
性质 $\mathcal{P}(S) =$``$S$ 是 Cohen--Macaulay 概形'' 在 syntomic 拓扑
中为局部的。
\end{lemma}

\begin{proof}
这由上面的 Lemma \ref{descent-lemma-Sk-local-syntomic} 立即得出，因为一个概形
是 Cohen--Macaulay 概形，当且仅当它局部 Noether，且对所有
$k \geq 0$ 都满足 $(S_k)$；见 Properties, Lemma
\ref{properties-lemma-scheme-CM-iff-all-Sk}。
\end{proof}







\section{在 smooth 拓扑中为局部的概形性质}
\label{descent-section-descending-properties-smooth}

\noindent
本节给出一些在 smooth 拓扑中于基底上为局部的概形性质。

\begin{lemma}
\label{descent-lemma-reduced-local-smooth}
性质 $\mathcal{P}(S) =$``$S$ 既约'' 在 smooth 拓扑中为局部的。
\end{lemma}

\begin{proof}
我们使用 Lemma \ref{descent-lemma-descending-properties}。首先注意，``既约'' 在
Zariski 拓扑中为局部的。这由定义即得；见 Schemes, Definition
\ref{schemes-definition-reduced}。接着，若 $S' \to S$ 是仿射概形之间的
smooth 态射，而 $S$ 既约，则 $S'$ 既约；这是 Algebra, Lemma
\ref{algebra-lemma-reduced-goes-up}。最后，若 $S' \to S$ 是仿射概形之间
的满 smooth 态射，而 $S'$ 既约，则 $S$ 既约；这是 Algebra, Lemma
\ref{algebra-lemma-descent-reduced}。因此 Lemma
\ref{descent-lemma-descending-properties} 的 (1)、(2)、(3) 成立，命题得证。
\end{proof}

\begin{lemma}
\label{descent-lemma-normal-local-smooth}
\begin{slogan}
正规性在 smooth 拓扑中为局部的。
\end{slogan}
性质 $\mathcal{P}(S) =$``$S$ 正规'' 在 smooth 拓扑中为局部的。
\end{lemma}

\begin{proof}
我们使用 Lemma \ref{descent-lemma-descending-properties}。首先证明 ``正规'' 在
Zariski 拓扑中为局部的。这由定义即得；见 Properties, Definition
\ref{properties-definition-normal}。接着，若 $S' \to S$ 是仿射概形之间的
smooth 态射，而 $S$ 正规，则 $S'$ 正规；这是 Algebra, Lemma
\ref{algebra-lemma-normal-goes-up}。最后，若 $S' \to S$ 是仿射概形之间的
满 smooth 态射，而 $S'$ 正规，则 $S$ 正规；这是 Algebra, Lemma
\ref{algebra-lemma-descent-normal}。因此 Lemma
\ref{descent-lemma-descending-properties} 的 (1)、(2)、(3) 成立，命题得证。
\end{proof}

\begin{lemma}
\label{descent-lemma-Rk-local-smooth}
性质 $\mathcal{P}(S) =$``$S$ 局部 Noether 且满足 $(R_k)$'' 在 smooth
拓扑中为局部的。
\end{lemma}

\begin{proof}
我们验证 Lemma \ref{descent-lemma-descending-properties} 的 (1)、(2)、(3)。由于
smooth 态射平坦且有限表示（Morphisms, Lemmas
\ref{morphisms-lemma-smooth-flat} 与
\ref{morphisms-lemma-smooth-locally-finite-presentation}），在 Lemma
\ref{descent-lemma-Noetherian-local-fppf} 的证明中已经对 ``局部 Noether'' 验证
了这些条件；下文不再提及而直接使用。首先注意，$\mathcal{P}$ 在 Zariski
拓扑中为局部的。这由定义即得；见 Properties, Definition
\ref{properties-definition-Rk}。接着，若 $S' \to S$ 是仿射概形之间的
smooth 态射，而 $S$ 具有 $\mathcal{P}$，则 $S'$ 具有
$\mathcal{P}$。这是 Algebra, Lemmas \ref{algebra-lemma-Rk-goes-up}
（使用 Morphisms, Lemma \ref{morphisms-lemma-smooth-characterize} 以及
Algebra, Lemmas \ref{algebra-lemma-base-change-smooth} 与
\ref{algebra-lemma-characterize-smooth-over-field}）。最后，若
$S' \to S$ 是仿射概形之间的满 smooth 态射，而 $S'$ 具有
$\mathcal{P}$，则 $S$ 具有 $\mathcal{P}$。这是 Algebra, Lemma
\ref{algebra-lemma-descent-Rk}。因此 Lemma
\ref{descent-lemma-descending-properties} 的 (1)、(2)、(3) 成立，命题得证。
\end{proof}

\begin{lemma}
\label{descent-lemma-regular-local-smooth}
性质 $\mathcal{P}(S) =$``$S$ 正则'' 在 smooth 拓扑中为局部的。
\end{lemma}

\begin{proof}
这由上面的 Lemma \ref{descent-lemma-Rk-local-smooth} 立即得出，因为局部 Noether
概形正则，当且仅当它局部 Noether，且对所有 $k \geq 0$ 都满足
$(R_k)$。
\end{proof}

\begin{lemma}
\label{descent-lemma-Nagata-local-smooth}
性质 $\mathcal{P}(S) =$``$S$ 是 Nagata 概形'' 在 smooth 拓扑中为局部的。
\end{lemma}

\begin{proof}
我们验证 Lemma \ref{descent-lemma-descending-properties} 的 (1)、(2)、(3)。首先
注意，Nagata 性在 Zariski 拓扑中为局部的；这是 Properties, Lemma
\ref{properties-lemma-locally-nagata}。接着，若 $S' \to S$ 是仿射概形
之间的 smooth 态射，而 $S$ 是 Nagata 概形，则 $S'$ 是 Nagata 概形；
这是 Morphisms, Lemma \ref{morphisms-lemma-finite-type-nagata}。最后，
若 $S' \to S$ 是仿射概形之间的满 smooth 态射，而 $S'$ 是 Nagata
概形，则 $S$ 是 Nagata 概形；这是 Algebra, Lemma
\ref{algebra-lemma-descent-nagata}。因此 Lemma
\ref{descent-lemma-descending-properties} 的 (1)、(2)、(3) 成立，命题得证。
\end{proof}





\section{下降性质的变体}
\label{descent-section-variants}

\noindent
有时可以下降并非局部的性质。本节汇集这类结果。关于光滑性等态射性质
的下降，另见 Section \ref{descent-section-descent-finiteness-morphisms}。

\begin{lemma}
\label{descent-lemma-descend-reduced}
若 $f : X \to Y$ 是概形的平坦满态射，且 $X$ 既约，则 $Y$ 既约。
\end{lemma}

\begin{proof}
考察局部环（Schemes, Definition \ref{schemes-definition-reduced}），并应用
Algebra, Lemma \ref{algebra-lemma-descent-reduced}，即得结论。
\end{proof}

\begin{lemma}
\label{descent-lemma-descend-regular}
设 $f : X \to Y$ 为代数空间态射。若 $f$ 局部有限表示、平坦且满，并且
$X$ 正则，则 $Y$ 正则。
\end{lemma}

\begin{proof}
本引理归约为下列代数陈述：若 $A \to B$ 是忠实平坦的有限表示环同态，
且 $B$ Noether 并正则，则 $A$ Noether 并正则。$A$ Noether 由 Algebra,
Lemma \ref{algebra-lemma-descent-Noetherian} 得出，而其正则性由 Algebra,
Lemma \ref{algebra-lemma-flat-under-regular} 得出。
\end{proof}









\section{概形芽}
\label{descent-section-germs}

\begin{definition}
\label{descent-definition-germs}
概形芽。
\begin{enumerate}
\item 由概形 $X$ 与点 $x \in X$ 组成的偶对 $(X, x)$ 称为
{\it $X$ 在 $x$ 处的芽}。
\item {\it 芽态射} $f : (X, x) \to (S, s)$ 是满足 $f(x) = s$ 的概形
态射 $f : U \to S$ 的等价类，其中 $U \subset X$ 是 $x$ 的开邻域。
两个这样的 $f$、$f'$ 等价，当且仅当 $f$ 与 $f'$ 在 $x$ 的某个开邻域
上一致。
\item 通过复合代表元定义{\it 芽态射的复合}（这是良定义的）。
\end{enumerate}
\end{definition}

\noindent
继续之前还需要一个定义。

\begin{definition}
\label{descent-definition-etale-morphism-germs}
设 $f : (X, x) \to (S, s)$ 为芽态射。若 $f$ 存在一个代表元
$f : U \to S$，它是概形的 \'etale 态射（相应地，smooth 态射），则称
$f$ {\it \'etale}（相应地，{\it smooth}）。
\end{definition}








\section{芽的局部性质}
\label{descent-section-properties-germs-local}

\begin{definition}
\label{descent-definition-local-at-point}
设 $\mathcal{P}$ 为概形芽的一项性质。若对任意 \'etale（相应地，smooth）
芽态射 $(U', u') \to (U, u)$，都有
$\mathcal{P}(U, u) \Leftrightarrow \mathcal{P}(U', u')$，则称
$\mathcal{P}$ {\it \'etale 局部}（相应地，{\it smooth 局部}）。
\end{definition}

\noindent
设 $(X, x)$ 为概形芽。$X$ 在 $x$ 处的维数，是 $x$ 在 $X$ 中各开邻域
维数的最小值；任意充分小的开邻域都具有这一维数。因此它是该芽同构类
的不变量，简记为 $\dim_x(X)$。下一个引理说明，断言
$\dim_x(X) = d$ 是芽的 \'etale 局部性质。

\begin{lemma}
\label{descent-lemma-dimension-at-point-local}
设 $f : U \to V$ 为概形的 \'etale 态射。设 $u \in U$ 且
$v = f(u)$。则 $\dim_u(U) = \dim_v(V)$。
\end{lemma}

\begin{proof}
陈述中的 $\dim_u(U)$ 是 Topology, Definition
\ref{topology-definition-Krull} 所定义的 $U$ 在 $u$ 处的维数，即 $u$
在 $U$ 中各开邻域的 Krull 维数的最小值。$\dim_v(V)$ 也类似。

\medskip\noindent
先证明 $\dim_v(V) \geq \dim_u(U)$。设 $V'$ 为 $v$ 在 $V$ 中的开邻域。
则存在 $u$ 在 $U$ 中的开邻域 $U'$，它包含于 $f^{-1}(V')$，且
$\dim_u(U) = \dim(U')$。设
$Z_0 \subset Z_1 \subset \ldots \subset Z_n$ 为 $U'$ 中不可约闭子概形的
一条链。若 $\xi_i \in Z_i$ 是一般点，则有特化
$\xi_n \leadsto \xi_{n - 1} \leadsto \ldots \leadsto \xi_0$.
这在 $V'$ 中给出特化
$f(\xi_n) \leadsto f(\xi_{n - 1}) \leadsto \ldots \leadsto f(\xi_0)$
。当 $i \not = j$ 时，$f(\xi_j) \not = f(\xi_i)$，因为 $f$ 的纤维离散
（见 Morphisms, Lemma \ref{morphisms-lemma-etale-over-field}）。所以
$\dim(V') \geq n$。形式地即可推出不等式
$\dim_v(V) \geq \dim_u(U)$。

\medskip\noindent
再证明 $\dim_u(U) \geq \dim_v(V)$。设 $U'$ 为 $u$ 在 $U$ 中的开邻域。
由 Morphisms, Lemma \ref{morphisms-lemma-fppf-open}，
$V' = f(U')$ 是 $v$ 的开邻域。因此 $\dim(V') \geq \dim_v(V)$。选取
$V'$ 中不可约闭子概形的一条链
$Z_0 \subset Z_1 \subset \ldots \subset Z_n$。令
$\xi_i \in Z_i$ 为一般点，于是有特化
$\xi_n \leadsto \xi_{n - 1} \leadsto \ldots \leadsto \xi_0$.
由于 $\xi_0 \in f(U')$，可找到一点 $\eta_0 \in U'$，满足
$f(\eta_0) = \xi_0$。考虑局部环同态
$$
\mathcal{O}_{V', \xi_0} \longrightarrow \mathcal{O}_{U', \eta_0}
$$
，它由 Morphisms, Lemma \ref{morphisms-lemma-etale-flat} 是平坦局部环同态。
注意，由 Schemes, Lemma \ref{schemes-lemma-specialize-points}，各点
$\xi_i$ 对应于左侧环的素理想。因此，对上示环同态应用下降定理（见
Algebra, Section \ref{algebra-section-going-up}），可以在 $U'$ 中找到一列
特化
$\eta_n \leadsto \eta_{n - 1} \leadsto \ldots \leadsto \eta_0$
，它们在 $f$ 下映到序列
$\xi_n \leadsto \xi_{n - 1} \leadsto \ldots \leadsto \xi_0$
。这蕴含 $\dim_u(U) \geq \dim_v(V)$。
\end{proof}

\noindent
设 $(X, x)$ 为概形芽。局部环 $\mathcal{O}_{X, x}$ 的同构类是该芽的
不变量。下一个引理说明，性质 $\dim(\mathcal{O}_{X, x}) = d$ 是芽的
\'etale 局部性质。

\begin{lemma}
\label{descent-lemma-dimension-local-ring-local}
设 $f : U \to V$ 为概形的 \'etale 态射。设 $u \in U$ 且
$v = f(u)$。则
$\dim(\mathcal{O}_{U, u}) = \dim(\mathcal{O}_{V, v})$.
\end{lemma}

\begin{proof}
须证明的代数陈述如下：若 $A \to B$ 是 \'etale 环同态，且
$\mathfrak q$ 是 $B$ 的一个素理想，位于 $\mathfrak p \subset A$ 上方，
则 $\dim(A_{\mathfrak p}) = \dim(B_{\mathfrak q})$。这是 More on Algebra,
Lemma \ref{more-algebra-lemma-dimension-etale-extension}。
\end{proof}

\noindent
设 $(X, x)$ 为概形芽。局部环 $\mathcal{O}_{X, x}$ 的同构类是该芽的
不变量。下一个引理说明，性质 ``$\mathcal{O}_{X, x}$ 正则'' 是芽的
\'etale 局部性质。

\begin{lemma}
\label{descent-lemma-regular-local-ring-local}
设 $f : U \to V$ 为概形的 \'etale 态射。设 $u \in U$ 且
$v = f(u)$。则 $\mathcal{O}_{U, u}$ 是正则局部环，当且仅当
$\mathcal{O}_{V, v}$ 是正则局部环。
\end{lemma}

\begin{proof}
须证明的代数陈述如下：若 $A \to B$ 是 \'etale 环同态，且
$\mathfrak q$ 是 $B$ 的一个素理想，位于 $\mathfrak p \subset A$ 上方，
则 $A_{\mathfrak p}$ 正则，当且仅当 $B_{\mathfrak q}$ 正则。这是 More on
Algebra, Lemma \ref{more-algebra-lemma-regular-etale-extension}。
\end{proof}







\section{在目标上为局部的态射性质}
\label{descent-section-descending-properties-morphisms}

\noindent
设 $f : X \to Y$ 为概形态射，设 $g : Y' \to Y$ 为概形态射，并设
$f' : X' \to Y'$ 为 $f$ 沿 $g$ 的基变换：
$$
\xymatrix{
X' \ar[d]_{f'} \ar[r]_{g'} & X \ar[d]^f \\
Y' \ar[r]^g & Y
}
$$
设 $\mathcal{P}$ 为概形态射的一项性质。可以问 (a)
$\mathcal{P}(f) \Rightarrow \mathcal{P}(f')$ 是否成立，也可以问其逆命题
(b) $\mathcal{P}(f') \Rightarrow \mathcal{P}(f)$ 是否成立。若每当 $g$
平坦时 (a) 都成立，则称 $\mathcal{P}$ 在平坦基变换下保持。若每当 $g$
满且平坦时 (b) 都成立，则称 $\mathcal{P}$ 沿满平坦基变换下降。若
$\mathcal{P}$ 在平坦基变换下保持，并沿满平坦基变换下降，则称
$\mathcal{P}$ 在目标上为平坦局部的。比较 Section
\ref{descent-section-descending-properties} 中的讨论。这是一个非常重要的概念，
正式定义如下。

\begin{definition}
\label{descent-definition-property-morphisms-local}
设 $\mathcal{P}$ 为基底上概形态射的一项性质，并设
$\tau \in \{fpqc, fppf, syntomic, smooth, \etale, Zariski\}$。若对任意
$\tau$-覆盖 $\{Y_i \to Y\}_{i \in I}$（见 Topologies, Section
\ref{topologies-section-procedure}）以及 $S$ 上的任意概形态射
$f : X \to Y$，都有
$$
f \text{ 具有 }\mathcal{P}
\Leftrightarrow
\text{每个 }Y_i \times_Y X \to Y_i\text{ 具有 }\mathcal{P}.
$$
则称 $\mathcal{P}$ {\it 在基底上为 $\tau$-局部的}，或
{\it 在目标上为 $\tau$-局部的}，或
{\it 关于 $\tau$-拓扑在基底上为局部的}。
\end{definition}

\noindent
确切地说，由于同构总是覆盖，可知（或要求）$X \to Y$ 具有性质
$\mathcal{P}$，当且仅当任意与 $X \to Y$ 同构的箭头 $X' \to Y'$ 具有
该性质。若一个性质在目标上为 $\tau$-局部的，则沿 $\tau$-覆盖中出现的
态射作基变换时，该性质保持。正式陈述如下。

\begin{lemma}
\label{descent-lemma-pullback-property-local-target}
设 $\tau \in \{fpqc, fppf, syntomic, smooth, \etale, Zariski\}$。设
$\mathcal{P}$ 为在目标上 $\tau$-局部的态射性质，并设 $f : X \to Y$
具有性质 $\mathcal{P}$。对任意态射 $Y' \to Y$，若它相应地为平坦、平坦
且局部有限表示、syntomic、\'etale、开浸入，则 $f$ 的基变换
% P05-ERR-0035
$f' : Y' \times_Y X \to Y'$ 具有性质 $\mathcal{P}$。
\end{lemma}

\begin{proof}
这是因为可以把 $Y' \to Y$ 加入一族构成 $\tau$-覆盖的态射中。
\end{proof}

\noindent
上述结果有一个经常使用的简单推论：若 $f : X \to Y$ 具有在目标上
$\tau$-局部的性质 $\mathcal{P}$，并且对某个开子概形 $V \subset Y$ 有
$f(X) \subset V$，则所诱导的态射 $X \to V$ 也具有 $\mathcal{P}$。
证明：沿 $V \to Y$ 对 $f$ 作基变换即得 $X \to V$。

\begin{lemma}
\label{descent-lemma-largest-open-of-the-base}
设 $\tau \in \{fppf, syntomic, smooth, \etale\}$，并设 $\mathcal{P}$ 为在
目标上 $\tau$-局部的态射性质。对任意概形态射 $f : X \to Y$，存在最大
开集 $W(f) \subset Y$，使得限制 $X_{W(f)} \to W(f)$ 具有
$\mathcal{P}$。此外，
\begin{enumerate}
\item 若 $g : Y' \to Y$ 平坦且局部有限表示、syntomic、smooth 或
\'etale，并且基变换 $f' : X_{Y'} \to Y'$ 具有 $\mathcal{P}$，则
$g(Y') \subset W(f)$；
\item 若 $g : Y' \to Y$ 平坦且局部有限表示、syntomic、smooth 或
\'etale，则 $W(f') = g^{-1}(W(f))$；
\item 若 $\{g_i : Y_i \to Y\}$ 是 $\tau$-覆盖，则
$g_i^{-1}(W(f)) = W(f_i)$，其中 $f_i$ 是 $f$ 沿 $Y_i \to Y$ 的基变换。
\end{enumerate}
\end{lemma}

\begin{proof}
考虑所有满足下列性质的态射 $g : Y' \to Y$ 的像 $g(Y') \subset Y$ 之并
$W$：
\begin{enumerate}
\item $g$ 平坦且局部有限表示、syntomic、smooth 或 \'etale；
\item 基变换 $Y' \times_{g, Y} X \to Y'$ 具有性质 $\mathcal{P}$。
\end{enumerate}
这类态射 $g$ 是开映射（见 Morphisms, Lemma
\ref{morphisms-lemma-fppf-open}），所以 $W \subset Y$ 是 $Y$ 的开子集。
由于 $\mathcal{P}$ 在 $\tau$-拓扑中为局部的，且给定了 $W$ 的覆盖
$\{Y' \to W\}$，使得各拉回具有 $\mathcal{P}$，故限制 $X_W \to W$
具有性质 $\mathcal{P}$。这证明了存在性，也证明了 $W(f)$ 具有性质 (1)。
为证明性质 (2)，注意 $W(f') \supset g^{-1}(W(f))$，因为在沿平坦且局部
有限表示、syntomic、smooth 或 \'etale 态射作基变换时，
$\mathcal{P}$ 保持；见 Lemma \ref{descent-lemma-pullback-property-local-target}。
另一方面，若 $Y'' \subset Y'$ 为开集，且 $X_{Y''} \to Y''$ 具有性质
$\mathcal{P}$，则由构造，$Y'' \to Y$ 通过 $W$ 分解；也就是说，
$Y'' \subset g^{-1}(W(f))$。这证明了 (2)。断言 (3) 由 (2) 得出，因为按
$\tau$-覆盖的定义，每个态射 $Y_i \to Y$ 都平坦且局部有限表示、
syntomic、smooth 或 \'etale。
\end{proof}

\begin{lemma}
\label{descent-lemma-descending-properties-morphisms}
设 $\mathcal{P}$ 为基底上概形态射的一项性质，并设
$\tau \in \{fpqc, fppf, \etale, smooth, syntomic\}$。假设
\begin{enumerate}
\item 按照 $\tau$ 为 fpqc、fppf、\'etale、smooth 或 syntomic，该性质
相应地在平坦、平坦且局部有限表示、\'etale、smooth 或 syntomic 基变换
下保持（比较 Schemes, Definition
\ref{schemes-definition-preserved-by-base-change}）；
\item 该性质在基底上为 Zariski 局部的；
\item 对任意满的仿射概形态射 $S' \to S$，若按照 $\tau$ 为 fpqc、fppf、
\'etale、smooth 或 syntomic，该态射相应地为平坦、平坦且有限表示、
\'etale、smooth 或 syntomic，并且对任意概形态射 $f : X \to S$，只要
基变换 $f' : X' = S' \times_S X \to S'$ 具有性质 $\mathcal{P}$，
$f$ 就具有性质 $\mathcal{P}$。
\end{enumerate}
则 $\mathcal{P}$ 在基底上为 $\tau$-局部的。
\end{lemma}

\begin{proof}
这几乎立即由 $\tau$-覆盖的定义得出；见 Topologies, Definition
\ref{topologies-definition-fpqc-covering}
\ref{topologies-definition-fppf-covering}
\ref{topologies-definition-etale-covering}
\ref{topologies-definition-smooth-covering}，或
\ref{topologies-definition-syntomic-covering}
以及 Topologies, Lemma
\ref{topologies-lemma-fpqc-affine},
\ref{topologies-lemma-fppf-affine},
\ref{topologies-lemma-etale-affine},
\ref{topologies-lemma-smooth-affine}，或
\ref{topologies-lemma-syntomic-affine}。省略细节。
\end{proof}

\begin{remark}
\label{descent-remark-descending-properties-morphisms-standard}
（这是对上面 Remark \ref{descent-remark-descending-properties-standard} 的重复。）
在上面的 Lemma \ref{descent-lemma-descending-properties-morphisms} 中，若
$\tau = smooth$，则在条件 (3) 中可以假设该态射是（满的）标准 smooth
态射。$\tau = syntomic$ 或 $\tau = \etale$ 时也类似。
\end{remark}




\section{在目标的 fpqc 拓扑中为局部的态射性质}
\label{descent-section-descending-properties-morphisms-fpqc}

\noindent
本节给出大量在 fpqc 拓扑中于基底上为局部的概形态射性质。与此相反，
我们将在 Examples, Section
\ref{examples-section-non-descending-property-projective} 中证明，即使在
Zariski 拓扑中，``射影'' 与 ``拟射影'' 也不在基底上为局部的。

\begin{lemma}
\label{descent-lemma-descending-property-quasi-compact}
性质 $\mathcal{P}(f) =$``$f$ 拟紧'' 在基底上为 fpqc 局部的。
\end{lemma}

\begin{proof}
拟紧态射的基变换拟紧；见 Schemes, Lemma
\ref{schemes-lemma-quasi-compact-preserved-base-change}。拟紧性在基底上为
Zariski 局部的；见 Schemes, Lemma
\ref{schemes-lemma-quasi-compact-affine}。最后，设 $S' \to S$ 为仿射概形
之间的平坦满态射，并设 $f : X \to S$ 为态射。假设基变换
$f' : X' \to S'$ 拟紧。则 $X'$ 拟紧，而 $X' \to X$ 满，所以 $X$
拟紧。这蕴含 $f$ 拟紧。因此应用 Lemma
\ref{descent-lemma-descending-properties-morphisms} 即得结论。
\end{proof}

\begin{lemma}
\label{descent-lemma-descending-property-quasi-separated}
性质 $\mathcal{P}(f) =$``$f$ 拟分离'' 在基底上为 fpqc 局部的。
\end{lemma}

\begin{proof}
拟分离态射的任意基变换拟分离；见 Schemes, Lemma
\ref{schemes-lemma-separated-permanence}。拟分离性在基底上为 Zariski
局部的（由定义，或由 Schemes, Lemma
\ref{schemes-lemma-characterize-quasi-separated}）。最后，设
$S' \to S$ 为仿射概形之间的平坦满态射，并设 $f : X \to S$ 为态射。
假设基变换 $f' : X' \to S'$ 拟分离。这意味着
$\Delta' : X' \to X'\times_{S'} X'$ 拟紧。注意，$\Delta'$ 是
$\Delta : X \to X \times_S X$ 沿 $S' \to S$ 的基变换。由 Lemma
\ref{descent-lemma-descending-property-quasi-compact}，这蕴含 $\Delta$ 拟紧，
所以 $f$ 拟分离。因此应用 Lemma
\ref{descent-lemma-descending-properties-morphisms} 即得结论。
\end{proof}

\begin{lemma}
\label{descent-lemma-descending-property-universally-closed}
性质 $\mathcal{P}(f) =$``$f$ 普遍闭'' 在基底上为 fpqc 局部的。
\end{lemma}

\begin{proof}
按定义，普遍闭态射的基变换普遍闭。普遍闭性在基底上为 Zariski 局部的
（由定义，或由 Morphisms, Lemma
\ref{morphisms-lemma-universally-closed-local-on-the-base}）。最后，设
$S' \to S$ 为仿射概形之间的平坦满态射，并设 $f : X \to S$ 为态射。
假设基变换 $f' : X' \to S'$ 普遍闭。设 $T \to S$ 为任意态射。考虑图表
$$
\xymatrix{
X' \ar[d] &
S' \times_S T \times_S X \ar[d] \ar[r] \ar[l] &
T \times_S X \ar[d] \\
S' &
S' \times_S T \ar[r] \ar[l] &
T
}
$$
，其中两个方块都是 Cartesian 的。因此假设蕴含中间的竖直箭头为闭
映射。右侧的水平箭头平坦、拟紧且满（因为它们是 $S' \to S$ 的基变换）。
所以 $T$ 的一个子集为闭集，当且仅当它在 $S' \times_S T$ 中的逆像为
闭集；见 Morphisms, Lemma
\ref{morphisms-lemma-fpqc-quotient-topology}。简单的图追踪表明右侧竖直
箭头也为闭映射，所以 $X \to S$ 普遍闭。因此应用 Lemma
\ref{descent-lemma-descending-properties-morphisms} 即得结论。
\end{proof}

\begin{lemma}
\label{descent-lemma-descending-property-universally-open}
性质 $\mathcal{P}(f) =$``$f$ 普遍开'' 在基底上为 fpqc 局部的。
\end{lemma}

\begin{proof}
证明与 Lemma \ref{descent-lemma-descending-property-universally-closed} 的证明相同。
\end{proof}

\begin{lemma}
\label{descent-lemma-descending-property-universally-submersive}
性质 $\mathcal{P}(f) =$``$f$ 普遍为商映射'' 在基底上为 fpqc 局部的。
\end{lemma}

\begin{proof}
证明与 Lemma \ref{descent-lemma-descending-property-universally-closed} 的证明相同；
其中使用 Morphisms, Lemma
\ref{morphisms-lemma-fpqc-quotient-topology} 所给出的事实：拟紧平坦满态射
普遍为商映射。
\end{proof}

\begin{lemma}
\label{descent-lemma-descending-property-separated}
性质 $\mathcal{P}(f) =$``$f$ 分离'' 在基底上为 fpqc 局部的。
\end{lemma}

\begin{proof}
分离态射的基变换分离；见 Schemes, Lemma
\ref{schemes-lemma-separated-permanence}。分离性在基底上为 Zariski 局部
的（由定义，或由 Schemes, Lemma
\ref{schemes-lemma-characterize-separated}）。最后，设 $S' \to S$ 为仿射
概形之间的平坦满态射，并设 $f : X \to S$ 为态射。假设基变换
$f' : X' \to S'$ 分离。这意味着
$\Delta' : X' \to X'\times_{S'} X'$ 是闭浸入，因而普遍闭。注意，
$\Delta'$ 是 $\Delta : X \to X \times_S X$ 沿 $S' \to S$ 的基变换。
由 Lemma \ref{descent-lemma-descending-property-universally-closed}，这蕴含
$\Delta$ 普遍闭。由于它是浸入（Schemes, Lemma
\ref{schemes-lemma-diagonal-immersion}），可知 $\Delta$ 是闭浸入。所以
$f$ 分离。因此应用 Lemma \ref{descent-lemma-descending-properties-morphisms}
即得结论。
\end{proof}

\begin{lemma}
\label{descent-lemma-descending-property-surjective}
性质 $\mathcal{P}(f) =$``$f$ 满'' 在基底上为 fpqc 局部的。
\end{lemma}

\begin{proof}
这是显然的。
\end{proof}

\begin{lemma}
\label{descent-lemma-descending-property-dominant}
性质 $\mathcal{P}(f) =$``$f$ 拟紧且支配'' 在基底上为 fpqc 局部的。
\end{lemma}

\begin{proof}
由 Morphisms, Lemma \ref{morphisms-lemma-base-change-dominant}，
拟紧支配态射在平坦拉回下得以保持。反方向较为容易。事实上，由 Lemma
\ref{descent-lemma-descending-property-quasi-compact}，拟紧性在基底上为 fpqc
局部的。支配性显然在基底上为 Zariski 局部的。最后，设 $S' \to S$
为概形的满态射（不必平坦），并设 $f : X \to S$ 为态射。假设基变换
$f' : X' \to S'$ 支配。于是 $X' \to S'\to S$ 是两个支配态射的复合，
因而支配。由于这也是复合 $X' \to X \to S$，可知 $X\to S$ 支配。
\end{proof}

\noindent
一般而言，支配态射在平坦拉回下未必保持，但下面这个特殊情形值得一提。
另见 Lemma \ref{descent-lemma-descending-fppf-property-dominant}。

\begin{lemma}
\label{descent-lemma-descending-property-dominant-over-field}
设 $E/k$ 为域扩张。则 $k$ 上的态射 $X \to Y$ 支配，当且仅当拉回
$X_E \to Y_E$ 支配。
\end{lemma}

\begin{proof}
由 Morphisms, Lemma
\ref{morphisms-lemma-scheme-over-field-universally-open}，态射
$\Spec(E) \to \Spec(k)$ 普遍开。因此 $Y_E \to Y$ 开。所以，若
$X \to Y$ 支配，则 Morphisms, Lemma
\ref{morphisms-lemma-open-base-change-dominant} 给出 $X_E \to Y_E$
支配。反之，假设 $X_E \to Y_E$ 支配。由 Morphisms, Lemma
\ref{morphisms-lemma-base-change-surjective}，态射 $Y_E \to Y$ 满，故复合
$X_E \to Y_E \to Y$ 支配。这也是复合 $X_E \to X \to Y$，所以
$X \to Y$ 支配。
\end{proof}

\begin{lemma}
\label{descent-lemma-descending-property-universally-injective}
性质 $\mathcal{P}(f) =$``$f$ 普遍单射'' 在基底上为 fpqc 局部的。
\end{lemma}

\begin{proof}
普遍单射态射的基变换普遍单射（这是形式的）。普遍单射性在基底上为
Zariski 局部的；这由定义显然可见。最后，设 $S' \to S$ 为仿射概形之间
的平坦满态射，并设 $f : X \to S$ 为态射。假设基变换
$f' : X' \to S'$ 普遍单射。设 $K$ 为域，并设
$a, b : \Spec(K) \to X$ 为两个满足 $f \circ a = f \circ b$ 的态射。
由于 $S' \to S$ 满，并由 Schemes, Section
\ref{schemes-section-points} 中的讨论，存在域扩张 $K'/K$ 和态射 $\Spec(K')
\to S'$，使下列实线图交换：
$$
\xymatrix{
\Spec(K') \ar[rrd] \ar@{-->}[rd]_{a', b'} \ar[dd] \\
 &
X' \ar[r] \ar[d] &
S' \ar[d] \\
\Spec(K) \ar[r]^{a, b} &
X \ar[r] &
S
}
$$
由于方块为 Cartesian 的，我们得到使图交换的两条虚线箭头 $a'$、$b'$。
因为 $X' \to S'$ 普遍单射，由 Morphisms, Lemma
\ref{morphisms-lemma-universally-injective} 可得 $a' = b'$。显然这迫使
$a = b$（由 Schemes, Section \ref{schemes-section-points} 中的讨论）。
因此应用 Lemma \ref{descent-lemma-descending-properties-morphisms} 即得结论。

\medskip\noindent
另一个证明可使用如下刻画：一个态射普遍单射，当且仅当其对角态射满；见
Morphisms, Lemma \ref{morphisms-lemma-universally-injective}。于是本引理由
满性在基底上为 fpqc 局部这一事实推出；见 Lemma
\ref{descent-lemma-descending-property-surjective}。（提示：使用对角态射的基变换
就是基变换的对角态射这一事实。）
\end{proof}

\begin{lemma}
\label{descent-lemma-descending-property-universal-homeomorphism}
性质 $\mathcal{P}(f) =$``$f$ 为普遍同胚'' 在基底上为 fpqc 局部的。
\end{lemma}

\begin{proof}
这可用与 Lemma \ref{descent-lemma-descending-property-universally-closed} 完全相同的
方式证明。或者，可以使用如下事实：拓扑空间的映射为同胚，当且仅当它
单射、满且开。因此，普遍同胚恰为满、普遍单射且普遍开的态射。于是本引理
由 Lemmas \ref{descent-lemma-descending-property-surjective}、
\ref{descent-lemma-descending-property-universally-injective} 和
\ref{descent-lemma-descending-property-universally-open} 推出。
\end{proof}

\begin{lemma}
\label{descent-lemma-descending-property-locally-finite-type}
性质 $\mathcal{P}(f) =$``$f$ 局部有限型'' 在基底上为 fpqc 局部的。
\end{lemma}

\begin{proof}
局部有限型性在基变换下保持；见 Morphisms, Lemma
\ref{morphisms-lemma-base-change-finite-type}。局部有限型性在基底上为
Zariski 局部的；见 Morphisms, Lemma
\ref{morphisms-lemma-locally-finite-type-characterize}。最后，设
$S' \to S$ 为仿射概形之间的平坦满态射，并设 $f : X \to S$ 为态射。
假设基变换 $f' : X' \to S'$ 局部有限型。设 $U \subset X$ 为仿射开集。
则 $U' = S' \times_S U$ 是 $S'$ 上的有限型仿射概形。写成
$S = \Spec(R)$,
$S' = \Spec(R')$,
$U = \Spec(A)$, and
$U' = \Spec(A')$.
我们知道 $R \to R'$ 忠实平坦，$A' = R' \otimes_R A$，并且
$R' \to A'$ 有限型。我们必须证明 $R \to A$ 有限型。这正是 Algebra,
Lemma \ref{algebra-lemma-finite-type-descends} 的结论。因此 $f$ 局部有限型。
所以应用 Lemma \ref{descent-lemma-descending-properties-morphisms} 即得结论。
\end{proof}

\begin{lemma}
\label{descent-lemma-descending-property-locally-finite-presentation}
性质 $\mathcal{P}(f) =$``$f$ 局部有限表示'' 在基底上为 fpqc 局部的。
\end{lemma}

\begin{proof}
局部有限表示性在基变换下保持；见 Morphisms, Lemma
\ref{morphisms-lemma-base-change-finite-presentation}。
% P05-ERR-0036
局部有限型性在基底上为 Zariski 局部的；见 Morphisms, Lemma
\ref{morphisms-lemma-locally-finite-presentation-characterize}。最后，设
$S' \to S$ 为仿射概形之间的平坦满态射，并设 $f : X \to S$ 为态射。
假设基变换 $f' : X' \to S'$ 局部有限表示。设 $U \subset X$ 为仿射开集。
则 $U' = S' \times_S U$ 是 $S'$ 上的有限型仿射概形。写成
$S = \Spec(R)$,
$S' = \Spec(R')$,
$U = \Spec(A)$, and
$U' = \Spec(A')$.
我们知道 $R \to R'$ 忠实平坦，$A' = R' \otimes_R A$，并且
$R' \to A'$ 有限表示。我们必须证明 $R \to A$ 有限表示。这正是
Algebra, Lemma \ref{algebra-lemma-finite-presentation-descends} 的结论。
因此 $f$ 局部有限表示。所以应用 Lemma
\ref{descent-lemma-descending-properties-morphisms} 即得结论。
\end{proof}

\begin{lemma}
\label{descent-lemma-descending-property-finite-type}
性质 $\mathcal{P}(f) =$``$f$ 有限型'' 在基底上为 fpqc 局部的。
\end{lemma}

\begin{proof}
合并 Lemmas \ref{descent-lemma-descending-property-quasi-compact} 和
\ref{descent-lemma-descending-property-locally-finite-type} 即得。
\end{proof}

\begin{lemma}
\label{descent-lemma-descending-property-finite-presentation}
性质 $\mathcal{P}(f) =$``$f$ 有限表示'' 在基底上为 fpqc 局部的。
\end{lemma}

\begin{proof}
合并 Lemmas \ref{descent-lemma-descending-property-quasi-compact}、
\ref{descent-lemma-descending-property-quasi-separated} 和
\ref{descent-lemma-descending-property-locally-finite-presentation} 即得。
\end{proof}

\begin{lemma}
\label{descent-lemma-descending-property-proper}
性质 $\mathcal{P}(f) =$``$f$ 固有'' 在基底上为 fpqc 局部的。
\end{lemma}

\begin{proof}
合并 Lemmas \ref{descent-lemma-descending-property-universally-closed}、
\ref{descent-lemma-descending-property-separated} 和
\ref{descent-lemma-descending-property-finite-type} 即得本引理。
\end{proof}

\begin{lemma}
\label{descent-lemma-descending-property-flat}
性质 $\mathcal{P}(f) =$``$f$ 平坦'' 在基底上为 fpqc 局部的。
\end{lemma}

\begin{proof}
平坦性在任意基变换下保持；见 Morphisms, Lemma
\ref{morphisms-lemma-base-change-flat}。由定义，平坦性在基底上为 Zariski
局部的。最后，设 $S' \to S$ 为仿射概形之间的平坦满态射，并设
$f : X \to S$ 为态射。假设基变换 $f' : X' \to S'$ 平坦。设
$U \subset X$ 为仿射开集，则 $U' = S' \times_S U$ 为仿射概形。写成
$S = \Spec(R)$,
$S' = \Spec(R')$,
$U = \Spec(A)$, and
$U' = \Spec(A')$.
我们知道 $R \to R'$ 忠实平坦，$A' = R' \otimes_R A$，并且
$R' \to A'$ 平坦。目标：证明 $R \to A$ 平坦。这立即由 Algebra,
Lemma \ref{algebra-lemma-flatness-descends} 推出。因此 $f$ 平坦。所以应用
Lemma \ref{descent-lemma-descending-properties-morphisms} 即得结论。
\end{proof}

\begin{lemma}
\label{descent-lemma-descending-property-open-immersion}
性质 $\mathcal{P}(f) =$``$f$ 为开浸入'' 在基底上为 fpqc 局部的。
\end{lemma}

\begin{proof}
开浸入性在基变换下稳定；见 Schemes, Lemma
\ref{schemes-lemma-base-change-immersion}。开浸入性在基底上为 Zariski
局部的（这是显然的）。

\medskip\noindent
设 $S' \to S$ 为仿射概形之间的平坦满态射，并设 $f : X \to S$ 为态射。
假设基变换 $f' : X' \to S'$ 为开浸入。我们断言 $f$ 为开浸入。此时
$f'$ 普遍开且普遍单射。因此由 Lemma
\ref{descent-lemma-descending-property-universally-open}，$f$ 普遍开；由 Lemma
\ref{descent-lemma-descending-property-universally-injective}，它普遍单射。特别地，
$f(X) \subset S$ 为开集。若对每个仿射开集 $U \subset f(X)$，我们都能证明
$f^{-1}(U) \to U$ 为同构，则 $f$ 为开浸入，结论即得。若以
$U' \subset S'$ 表示 $U$ 的逆像，则 $U' \to U$ 是仿射概形之间的忠实平坦
态射，且 $(f')^{-1}(U') \to U'$ 为同构（由于我们对 $U$ 的选取，
$f'(X')$ 包含 $U'$）。这样便归约到下一段讨论的情形。

\medskip\noindent
设 $S' \to S$ 为仿射概形之间的平坦满态射，设 $f : X \to S$ 为态射，并
假设基变换 $f' : X' \to S'$ 为同构。我们还必须证明 $f$ 为同构。显然，
$f$ 满、普遍单射且普遍开（后两者见上面的论证）。所以 $f$ 为双射，也即
$f$ 为同胚。因此由 Morphisms, Lemma
\ref{morphisms-lemma-homeomorphism-affine}，$f$ 为仿射态射。由于
$$
\mathcal{O}(S') \to
\mathcal{O}(X') =
\mathcal{O}(S') \otimes_{\mathcal{O}(S)} \mathcal{O}(X)
$$
为同构，并且 $\mathcal{O}(S) \to \mathcal{O}(S')$ 忠实平坦，这蕴含
$\mathcal{O}(S) \to \mathcal{O}(X)$ 为同构。因此 $f$ 为同构。这就完成了
上面断言的证明。所以应用 Lemma
\ref{descent-lemma-descending-properties-morphisms} 即得结论。
\end{proof}

\begin{lemma}
\label{descent-lemma-descending-property-isomorphism}
性质 $\mathcal{P}(f) =$``$f$ 为同构'' 在基底上为 fpqc 局部的。
\end{lemma}

\begin{proof}
合并 Lemmas \ref{descent-lemma-descending-property-surjective} 和
\ref{descent-lemma-descending-property-open-immersion} 即得。
\end{proof}

\begin{lemma}
\label{descent-lemma-descending-property-affine}
性质 $\mathcal{P}(f) =$``$f$ 为仿射态射'' 在基底上为 fpqc 局部的。
\end{lemma}

\begin{proof}
仿射态射的基变换为仿射态射；见 Morphisms, Lemma
\ref{morphisms-lemma-base-change-affine}。仿射性在基底上为 Zariski 局部的；
见 Morphisms, Lemma \ref{morphisms-lemma-characterize-affine}。最后，设
$g : S' \to S$ 为仿射概形之间的平坦满态射，并设 $f : X \to S$ 为态射。
假设基变换 $f' : X' \to S'$ 为仿射态射。换言之，$X'$ 为仿射概形，比如
$X' = \Spec(A')$。再写成 $S = \Spec(R)$ 和 $S' = \Spec(R')$。我们必须
证明 $X$ 为仿射概形。

\medskip\noindent
由 Lemmas \ref{descent-lemma-descending-property-quasi-compact} 和
\ref{descent-lemma-descending-property-separated}，$X \to S$ 分离且拟紧。因此
$f_*\mathcal{O}_X$ 是 $\mathcal{O}_S$-代数的拟凝聚层；见 Schemes,
Lemma \ref{schemes-lemma-push-forward-quasi-coherent}。所以对某个 $R$-代数
$A$，有 $f_*\mathcal{O}_X = \widetilde{A}$。当然，事实上
$A = \Gamma(X, \mathcal{O}_X)$。此外，由平坦基变换（例如见 Cohomology
of Schemes, Lemma \ref{coherent-lemma-flat-base-change-cohomology}），有
$g^*f_*\mathcal{O}_X = f'_*\mathcal{O}_{X'}$。换言之，
$A' = R' \otimes_R A$。考虑 Schemes, Lemma
\ref{schemes-lemma-morphism-into-affine} 中 $S$ 上的典范态射
$$
X \longrightarrow \Spec(A)
$$
。由上述结论，该态射到 $S'$ 的基变换为同构。因此由 Lemma
\ref{descent-lemma-descending-property-isomorphism}，它本身为同构。所以应用
Lemma \ref{descent-lemma-descending-properties-morphisms} 即得结论。
\end{proof}

\begin{lemma}
\label{descent-lemma-descending-property-closed-immersion}
性质 $\mathcal{P}(f) =$``$f$ 为闭浸入'' 在基底上为 fpqc 局部的。
\end{lemma}

\begin{proof}
设 $f : X \to Y$ 为概形态射，设 $\{Y_i \to Y\}$ 为 fpqc 覆盖。假设每个
$f_i : Y_i \times_Y X \to Y_i$ 都是闭浸入。这蕴含每个 $f_i$ 都是仿射
态射；见 Morphisms, Lemma
\ref{morphisms-lemma-closed-immersion-affine}。由 Lemma
\ref{descent-lemma-descending-property-affine}，可知 $f$ 为仿射态射。还需证明
$\mathcal{O}_Y \to f_*\mathcal{O}_X$ 满。对每个 $y \in Y$，存在 $i$ 和
映到 $y$ 的点 $y_i \in Y_i$。由 Cohomology of Schemes, Lemma
\ref{coherent-lemma-flat-base-change-cohomology}，层
$f_{i, *}(\mathcal{O}_{Y_i \times_Y X})$ 是 $f_*\mathcal{O}_X$ 的拉回。
由假设，它是 $\mathcal{O}_{Y_i}$ 的商。因此可见
$$
\Big(
\mathcal{O}_{Y, y} \longrightarrow (f_*\mathcal{O}_X)_y
\Big)
\otimes_{\mathcal{O}_{Y, y}} \mathcal{O}_{Y_i, y_i}
$$
满。由于 $\mathcal{O}_{Y_i, y_i}$ 在 $\mathcal{O}_{Y, y}$ 上忠实平坦，
这蕴含 $\mathcal{O}_{Y, y} \longrightarrow (f_*\mathcal{O}_X)_y$ 满，
正如所需。
\end{proof}

\begin{lemma}
\label{descent-lemma-descending-property-quasi-affine}
性质 $\mathcal{P}(f) =$``$f$ 为拟仿射态射'' 在基底上为 fpqc 局部的。
\end{lemma}

\begin{proof}
设 $f : X \to Y$ 为概形态射，设 $\{g_i : Y_i \to Y\}$ 为 fpqc 覆盖。
假设每个 $f_i : Y_i \times_Y X \to Y_i$ 都是拟仿射态射。这蕴含每个
$f_i$ 都拟紧且分离。由 Lemmas
\ref{descent-lemma-descending-property-quasi-compact} 和
\ref{descent-lemma-descending-property-separated}，这蕴含 $f$ 拟紧且分离。考虑
$\mathcal{O}_Y$-代数层 $\mathcal{A} = f_*\mathcal{O}_X$。由 Schemes,
Lemma \ref{schemes-lemma-push-forward-quasi-coherent}，它是拟凝聚
$\mathcal{O}_Y$-代数。考虑典范态射
$$
j : X \longrightarrow \underline{\Spec}_Y(\mathcal{A})
$$
；见 Constructions, Lemma \ref{constructions-lemma-canonical-morphism}。
由平坦基变换（例如见 Cohomology of Schemes, Lemma
\ref{coherent-lemma-flat-base-change-cohomology}），有
$g_i^*f_*\mathcal{O}_X = f_{i, *}\mathcal{O}_{X'}$，其中
$g_i : Y_i \to Y$ 是给定的平坦映射。因此，$j$ 沿 $g_i$ 的基变换 $j_i$
就是 Constructions, Lemma \ref{constructions-lemma-canonical-morphism}
中对应于态射 $f_i$ 的典范态射。由假设及 Morphisms, Lemma
\ref{morphisms-lemma-characterize-quasi-affine}，所有这些态射 $j_i$ 都是
拟紧开浸入。因此由 Lemmas \ref{descent-lemma-descending-property-quasi-compact}
和 \ref{descent-lemma-descending-property-open-immersion}，$j$ 为拟紧开浸入。再次
由 Morphisms, Lemma \ref{morphisms-lemma-characterize-quasi-affine}，可知
$f$ 为拟仿射态射。
\end{proof}

\begin{lemma}
\label{descent-lemma-descending-property-quasi-compact-immersion}
性质 $\mathcal{P}(f) =$``$f$ 为拟紧浸入'' 在基底上为 fpqc 局部的。
\end{lemma}

\begin{proof}
设 $f : X \to Y$ 为概形态射，设 $\{Y_i \to Y\}$ 为 fpqc 覆盖。写
$X_i = Y_i \times_Y X$，并以 $f_i : X_i \to Y_i$ 表示 $f$ 的基变换；
还以 $q_i : Y_i \to Y$ 表示给定的平坦态射。假设每个 $f_i$ 都是拟紧浸入。
由 Schemes, Lemma \ref{schemes-lemma-immersions-monomorphisms}，每个 $f_i$
都分离。由 Lemmas \ref{descent-lemma-descending-property-quasi-compact} 和
\ref{descent-lemma-descending-property-separated}，这蕴含 $f$ 拟紧且分离。令
$X \to Z \to Y$ 为 $f$ 经过其概形论像的分解。由于 $f$ 拟紧，由
Morphisms, Lemma
\ref{morphisms-lemma-quasi-compact-scheme-theoretic-image}，闭子概形
$Z \subset Y$ 由拟凝聚理想层
$\mathcal{I} = \Ker(\mathcal{O}_Y \to f_*\mathcal{O}_X)$ 截出。由平坦
基变换（例如见 Cohomology of Schemes, Lemma
\ref{coherent-lemma-flat-base-change-cohomology}；这里使用 $f$ 分离），
$f_{i, *}\mathcal{O}_{X_i}$ 是拉回 $q_i^*f_*\mathcal{O}_X$。因此，
$Y_i \times_Y Z$ 由拟凝聚理想层
$q_i^*\mathcal{I} =
\Ker(\mathcal{O}_{Y_i} \to f_{i, *}\mathcal{O}_{X_i})$ 截出。由
Morphisms, Lemma \ref{morphisms-lemma-quasi-compact-immersion}，态射
$X_i \to Y_i \times_Y Z$ 为开浸入。所以由 Lemma
\ref{descent-lemma-descending-property-open-immersion}，$X \to Z$ 为开浸入，
因而 $f$ 如所需为浸入（我们已知它拟紧）。
\end{proof}

\begin{lemma}
\label{descent-lemma-descending-property-integral}
性质 $\mathcal{P}(f) =$``$f$ 为整态射'' 在基底上为 fpqc 局部的。
\end{lemma}

\begin{proof}
整态射恰为仿射且普遍闭的态射；见 Morphisms, Lemma
\ref{morphisms-lemma-integral-universally-closed}。因此合并 Lemmas
\ref{descent-lemma-descending-property-universally-closed} 和
\ref{descent-lemma-descending-property-affine} 即得本引理。
\end{proof}

\begin{lemma}
\label{descent-lemma-descending-property-finite}
性质 $\mathcal{P}(f) =$``$f$ 为有限态射'' 在基底上为 fpqc 局部的。
\end{lemma}

\begin{proof}
有限态射恰为局部有限型的整态射；见 Morphisms, Lemma
\ref{morphisms-lemma-finite-integral}。因此合并 Lemmas
\ref{descent-lemma-descending-property-locally-finite-type} 和
\ref{descent-lemma-descending-property-integral} 即得本引理。
\end{proof}

\begin{lemma}
\label{descent-lemma-descending-property-quasi-finite}
性质 $\mathcal{P}(f) =$``$f$ 局部拟有限'' 和
$\mathcal{P}(f) =$``$f$ 拟有限'' 在基底上为 fpqc 局部的。
\end{lemma}

\begin{proof}
设 $f : X \to S$ 为概形态射，并设 $\{S_i \to S\}$ 为 fpqc 覆盖，使得每个
基变换 $f_i : X_i \to S_i$ 都局部拟有限。我们已经知道（Lemma
\ref{descent-lemma-descending-property-locally-finite-type}）``局部有限型'' 在基底上
为 fpqc 局部的，因而 $f$ 局部有限型。于是由 Morphisms, Lemma
\ref{morphisms-lemma-base-change-quasi-finite} 可知 $f$ 局部拟有限。拟有限的
情形也随之成立，因为我们已经知道``拟紧''在基底上为 fpqc 局部的（Lemma
\ref{descent-lemma-descending-property-quasi-compact}）。
\end{proof}

\begin{lemma}
\label{descent-lemma-descending-property-relative-dimension-d}
性质 $\mathcal{P}(f) =$``$f$ 局部有限型且相对维数为 $d$'' 在基底上为
fpqc 局部的。
\end{lemma}

\begin{proof}
这立即由局部有限型性在基底上为 fpqc 局部这一事实及 Morphisms, Lemma
\ref{morphisms-lemma-dimension-fibre-after-base-change} 推出。
\end{proof}

\begin{lemma}
\label{descent-lemma-descending-property-syntomic}
性质 $\mathcal{P}(f) =$``$f$ 为 syntomic 态射'' 在基底上为 fpqc 局部的。
\end{lemma}

\begin{proof}
一个态射为 syntomic 态射，当且仅当它局部有限表示、平坦，并且其纤维
为局部完全交。我们已经知道平坦性和局部有限表示性在基底上为 fpqc 局部
的（Lemmas \ref{descent-lemma-descending-property-flat} 和
\ref{descent-lemma-descending-property-locally-finite-presentation}）。因此 syntomic
的结论由 Morphisms, Lemma
\ref{morphisms-lemma-set-points-where-fibres-lci} 推出。
\end{proof}

\begin{lemma}
\label{descent-lemma-descending-property-smooth}
性质 $\mathcal{P}(f) =$``$f$ 光滑'' 在基底上为 fpqc 局部的。
\end{lemma}

\begin{proof}
一个态射光滑，当且仅当它局部有限表示、平坦，并且具有光滑纤维。我们已经
知道平坦性和局部有限表示性在基底上为 fpqc 局部的（Lemmas
\ref{descent-lemma-descending-property-flat} 和
\ref{descent-lemma-descending-property-locally-finite-presentation}）。因此光滑的
结论由 Morphisms, Lemma
\ref{morphisms-lemma-set-points-where-fibres-smooth} 推出。
\end{proof}

\begin{lemma}
\label{descent-lemma-descending-property-unramified}
性质 $\mathcal{P}(f) =$``$f$ 非分歧'' 在基底上为 fpqc 局部的。
性质 $\mathcal{P}(f) =$``$f$ 为 G-非分歧态射'' 在基底上为 fpqc 局部的。
\end{lemma}

\begin{proof}
一个态射非分歧（相应地，为 G-非分歧态射），当且仅当它局部有限型
（相应地，有限表示），并且其对角态射为开浸入（见 Morphisms, Lemma
\ref{morphisms-lemma-diagonal-unramified-morphism}）。我们已经知道，局部
有限型性（相应地，局部有限表示性）和开浸入性在基底上为 fpqc 局部的
（Lemmas \ref{descent-lemma-descending-property-locally-finite-presentation}、
\ref{descent-lemma-descending-property-locally-finite-type} 和
\ref{descent-lemma-descending-property-open-immersion}）。因此结论形式地成立。
\end{proof}

\begin{lemma}
\label{descent-lemma-descending-property-etale}
性质 $\mathcal{P}(f) =$``$f$ 为 \'etale 态射'' 在基底上为 fpqc 局部的。
\end{lemma}

\begin{proof}
一个态射为 \'etale 态射，当且仅当它平坦且为 G-非分歧态射；见
Morphisms, Lemma \ref{morphisms-lemma-flat-unramified-etale}。我们已经知道
平坦性和 G-非分歧性在基底上为 fpqc 局部的（Lemmas
\ref{descent-lemma-descending-property-flat} 和
\ref{descent-lemma-descending-property-unramified}）。故结论成立。
\end{proof}

\begin{lemma}
\label{descent-lemma-descending-property-finite-locally-free}
性质 $\mathcal{P}(f) =$``$f$ 有限局部自由'' 在基底上为 fpqc 局部的。
设 $d \geq 0$。性质 $\mathcal{P}(f) =$``$f$ 为次数 $d$ 的有限局部自由
态射'' 在基底上为 fpqc 局部的。
\end{lemma}

\begin{proof}
有限局部自由等价于有限、平坦且局部有限表示（Morphisms, Lemma
\ref{morphisms-lemma-finite-flat}）。因此这由 Lemmas
\ref{descent-lemma-descending-property-finite}、
\ref{descent-lemma-descending-property-flat} 和
\ref{descent-lemma-descending-property-locally-finite-presentation} 推出。若
$f : Z \to U$ 有限局部自由，且 $\{U_i \to U\}$ 是一族满的态射，使得每个
拉回 $Z \times_U U_i \to U_i$ 的次数均为 $d$，则 $Z \to U$ 的次数为
$d$；例如，这是因为在点 $u \in U$ 处，可以从纤维
$(f_*\mathcal{O}_Z)_u \otimes_{\mathcal{O}_{U, u}} \kappa(u)$ 读出次数。
\end{proof}

\begin{lemma}
\label{descent-lemma-descending-property-monomorphism}
性质 $\mathcal{P}(f) =$``$f$ 为单态射'' 在基底上为 fpqc 局部的。
\end{lemma}

\begin{proof}
设 $f : X \to S$ 为概形态射。设 $\{S_i \to S\}$ 为 fpqc 覆盖，并假设
$f$ 的每个基变换 $f_i : X_i \to S_i$ 都是单态射。设
$a, b : T \to X$ 为两个满足 $f \circ a = f \circ b$ 的态射。我们必须证明
$a = b$。由于 $f_i$ 为单态射，可知 $a_i = b_i$，其中
$a_i, b_i : S_i \times_S T \to X_i$ 为相应的基变换。特别地，复合
$S_i \times_S T \to T \to X$ 相等。由于
$\coprod S_i \times_S T \to T$ 为上态射（例如见 Lemma
\ref{descent-lemma-fpqc-universal-effective-epimorphisms}），故 $a = b$。
\end{proof}

\begin{lemma}
\label{descent-lemma-descending-property-regular-immersion}
下列性质
\begin{enumerate}
\item[] $\mathcal{P}(f) =$``$f$ 为 Koszul-正则浸入''，
\item[] $\mathcal{P}(f) =$``$f$ 为 $H_1$-正则浸入''，以及
\item[] $\mathcal{P}(f) =$``$f$ 为拟正则浸入''
\end{enumerate}
在基底上为 fpqc 局部的。
\end{lemma}

\begin{proof}
我们使用 Lemma \ref{descent-lemma-descending-properties-morphisms} 的判据来证明。
由 Divisors, Definition \ref{divisors-definition-regular-immersion}，
Koszul-正则（相应地，$H_1$-正则、拟正则）浸入性在基底上为 Zariski 局部
的。由 Divisors, Lemma
\ref{divisors-lemma-flat-base-change-regular-immersion}，Koszul-正则
（相应地，$H_1$-正则、拟正则）浸入性在平坦基变换下保持。Lemma
\ref{descent-lemma-descending-properties-morphisms} 的最后一个假设 (3) 转化为如下
代数陈述：设 $A \to B$ 为忠实平坦环映射，设 $I \subset A$ 为理想。若
$IB$ 在 $\Spec(B)$ 上局部地由 $B$ 中的 Koszul-正则（相应地，
$H_1$-正则、拟正则）序列生成，则 $I \subset A$ 在 $\Spec(A)$ 上局部地
由 $A$ 中的 Koszul-正则（相应地，$H_1$-正则、拟正则）序列生成。这正是
More on Algebra, Lemma
\ref{more-algebra-lemma-flat-descent-regular-ideal}。
\end{proof}





\section{在目标上的 fppf 拓扑中局部的态射性质}
\label{descent-section-descending-properties-morphisms-fppf}

\noindent
本节给出概形态射的一些性质：我们（尚）不能证明它们在 fpqc 拓扑中于基底
上为局部的，但它们确实在 fppf 拓扑中于基底上为局部的。

\begin{lemma}
\label{descent-lemma-descending-fppf-property-immersion}
性质 $\mathcal{P}(f) =$``$f$ 为浸入'' 在基底上为 fppf 局部的。
\end{lemma}

\begin{proof}
浸入性在基变换下稳定；见 Schemes, Lemma
\ref{schemes-lemma-base-change-immersion}。浸入性在基底上为 Zariski 局部
的。最后，设 $\pi : S' \to S$ 为仿射概形之间平坦且局部有限表示的满态射。
注意，由 Morphisms, Lemma \ref{morphisms-lemma-fppf-open}，
$\pi : S' \to S$ 为开映射。设 $f : X \to S$ 为态射。假设基变换
$f' : X' \to S'$ 为浸入。特别地，$f'(X') = \pi^{-1}(f(X))$ 局部闭。
因此由 Topology, Lemma \ref{topology-lemma-open-morphism-quotient-topology}，
$f(X) \subset S$ 局部闭。设 $Z \subset S$ 为闭子集
$Z = \overline{f(X)} \setminus f(X)$。再次由 Topology, Lemma
\ref{topology-lemma-open-morphism-quotient-topology}，$f'(X')$ 在
$S' \setminus Z'$ 中为闭集。因此可将 Lemma
\ref{descent-lemma-descending-property-closed-immersion} 应用于 fpqc 覆盖
$\{S' \setminus Z' \to S \setminus Z\}$，从而得到
$f : X \to S \setminus Z$ 为闭浸入。换言之，$f$ 为浸入。因此应用
Lemma \ref{descent-lemma-descending-properties-morphisms} 即得结论。
\end{proof}

\begin{lemma}
\label{descent-lemma-descending-fppf-property-dominant}
性质 $\mathcal{P}(f) =$``$f$ 支配'' 在基底上为 fppf 局部的。
\end{lemma}

\begin{proof}
由 Morphisms, Lemma \ref{morphisms-lemma-open-base-change-dominant}，支配态射
在沿开态射拉回时保持，因而在沿局部有限表示的平坦态射拉回时保持
（Morphisms, Lemma \ref{morphisms-lemma-fppf-open}）。反方向较为容易。
事实上，支配性显然在基底上为 Zariski 局部的。再设 $S' \to S$ 为概形的
满态射（不必平坦或局部有限表示），并设 $f : X \to S$ 为态射。假设基变换
$f' : X' \to S'$ 支配。则 $X' \to S'\to S$ 是两个支配态射的复合，故而
支配。由于这也是复合 $X' \to X \to S$，可知 $X\to S$ 支配。
\end{proof}






\section{态射性质的 fpqc 下降的应用}
\label{descent-section-application-descending-properties-morphisms}

\noindent
下面的引理看似有些轻巧，却是研究 \'etale 态射和非分歧态射的有用工具。

\begin{lemma}
\label{descent-lemma-flat-surjective-quasi-compact-monomorphism-isomorphism}
设 $f : X \to Y$ 为平坦、拟紧且满的单态射。则 f 为同构。
\end{lemma}

\begin{proof}
由于 $f$ 是平坦、拟紧且满的态射，$\{X \to Y\}$ 是 $Y$ 的 fpqc 覆盖。
对角态射 $\Delta : X \to X \times_Y X$ 为同构（Schemes, Lemma
\ref{schemes-lemma-monomorphism}）。这蕴含 $f$ 沿 $f$ 的基变换为同构。
因此由 Lemma \ref{descent-lemma-descending-property-isomorphism}，$f$ 为同构。
\end{proof}

\noindent
可以用这个引理证明下面的重要结果；我们还会给出一个不使用 fpqc 下降的
证明。我们将在 \'Etale Morphisms, Section
\ref{etale-section-topological-etale} 中更详细地讨论这个结果及相关结果。

\begin{lemma}
\label{descent-lemma-universally-injective-etale-open-immersion}
普遍单射的 \'etale 态射为开浸入。
\end{lemma}

\begin{proof}[第一种证明]
设 $f : X \to Y$ 为普遍单射的 \'etale 态射。则 $f$ 为开映射
（Morphisms, Lemma \ref{morphisms-lemma-etale-open}），所以可以用 $f(X)$
替换 $Y$，并假设 $f$ 满。于是 $f$ 双射且开，因而为同胚。所以 $f$ 拟紧。
因此由 Lemma
\ref{descent-lemma-flat-surjective-quasi-compact-monomorphism-isomorphism}，只需证明
$f$ 为单态射。由于 $X \to Y$ 为 \'etale 态射，由 Morphisms, Lemma
\ref{morphisms-lemma-diagonal-unramified-morphism}（以及 Morphisms, Lemma
\ref{morphisms-lemma-flat-unramified-etale}），态射
$\Delta_{X/Y} : X \to X \times_Y X$ 为开浸入。由于 $f$ 普遍单射，
$\Delta_{X/Y}$ 还满；见 Morphisms, Lemma
\ref{morphisms-lemma-universally-injective}。因此 $\Delta_{X/Y}$ 为同构，
也就是说，$X \to Y$ 为单态射。
\end{proof}

\begin{proof}[第二种证明]
设 $f : X \to Y$ 为普遍单射的 \'etale 态射。则 $f$ 为开映射（Morphisms,
Lemma \ref{morphisms-lemma-etale-open}），所以可以用 $f(X)$ 替换 $Y$，并
假设 $f$ 满。由于这些假设在任意基变换后仍然成立，可知 $f$ 为普遍同胚。
因此 $f$ 为整态射；见 Morphisms, Lemma
\ref{morphisms-lemma-universal-homeomorphism}。进而由 Morphisms, Lemma
\ref{morphisms-lemma-finite-integral}，$f$ 为有限态射；再由 Morphisms,
Lemma \ref{morphisms-lemma-finite-flat}，$f$ 有限局部自由。为完成证明，只需
证明 $f$ 是次数 $1$ 的有限局部自由态射（次数 $1$ 的有限局部自由态射为
同构）。存在 $Y$ 到开闭子概形 $V_d$ 的分解，使得
$f^{-1}(V_d) \to V_d$ 是次数 $d$ 的有限局部自由态射；见 Morphisms,
Lemma \ref{morphisms-lemma-finite-locally-free}。若 $V_d$ 非空，则可以选取
态射 $\Spec(k) \to V_d \subset Y$，其中 $k$ 为代数闭域（只需取 $V_d$
某点的剩余域的代数闭包）。于是由 Morphisms, Lemma
\ref{morphisms-lemma-etale-over-field} 及 $k$ 代数闭这一事实，
$\Spec(k) \times_Y X \to \Spec(k)$ 是若干份 $\Spec(k)$ 的不交并。然而，
由于 $f$ 普遍单射，只能有一份，因此如所需有 $d = 1$。
\end{proof}

\noindent
利用下面对平坦普遍单射态射的刻画，可以稍加改写上一个引理中的假设。

\begin{lemma}
\label{descent-lemma-flat-universally-injective}
设 $f : X \to Y$ 为概形态射。以 $X^0$ 表示 $X$ 的不可约分支的泛点集合。
若
\begin{enumerate}
\item $f$ 平坦且分离；
\item 对 $\xi \in X^0$，有 $\kappa(f(\xi)) = \kappa(\xi)$；并且
\item 若 $\xi, \xi' \in X^0$ 且 $\xi \not = \xi'$，则
$f(\xi) \not = f(\xi')$，
\end{enumerate}
则 $f$ 普遍单射。
\end{lemma}

\begin{proof}
我们必须证明 $\Delta : X \to X \times_Y X$ 满；见 Morphisms, Lemma
\ref{morphisms-lemma-universally-injective}。由于 $X \to Y$ 分离，$\Delta$
的像为闭集。所以若 $\Delta$ 不满，就可以找到 $X \times_S X$ 的某个不可约
分支的泛点
% P05-ERR-0037
$\eta \in X \times_S X$，它不在 $\Delta$ 的像中。投影
$\text{pr}_1 : X \times_Y X \to X$ 是平坦态射 $X \to Y$ 的基变换，因而
平坦；见 Morphisms, Lemma \ref{morphisms-lemma-base-change-flat}。因此沿
$\text{pr}_1$ 可以提升泛化；见 Morphisms, Lemma
\ref{morphisms-lemma-generalizations-lift-flat}。可知
$\xi = \text{pr}_1(\eta) \in X^0$。然而，假设 (2) 和 (3) 保证，对每个
$\xi \in X^0$，概形 $(X \times_Y X)_{f(\xi)}$ 至多有一个点。换言之，
$\Delta(\xi) = \eta$，矛盾。
\end{proof}

\noindent
因此可将 Lemma \ref{descent-lemma-universally-injective-etale-open-immersion} 改述如下。

\begin{lemma}
\label{descent-lemma-characterize-open-immersion}
设 $f : X \to Y$ 为概形态射。以 $X^0$ 表示 $X$ 的不可约分支的泛点集合。
若
\begin{enumerate}
\item $f$ 为 \'etale 态射且分离；
\item 对 $\xi \in X^0$，有 $\kappa(f(\xi)) = \kappa(\xi)$；并且
\item 若 $\xi, \xi' \in X^0$ 且 $\xi \not = \xi'$，则
$f(\xi) \not = f(\xi')$，
\end{enumerate}
则 $f$ 为开浸入。
\end{lemma}

\begin{proof}
立即由 Lemmas \ref{descent-lemma-flat-universally-injective} 和
\ref{descent-lemma-universally-injective-etale-open-immersion} 推出。
\end{proof}

\begin{lemma}
\label{descent-lemma-descending-property-proper-over-base}
设 $f : X \to Y$ 为局部有限型的概形态射，设 $Z$ 为 $X$ 的闭子集。若存在
fpqc 覆盖 $\{Y_i \to Y\}$，使得逆像 $Z_i \subset Y_i \times_Y X$ 在
$Y_i$ 上固有（Cohomology of Schemes, Definition
\ref{coherent-definition-proper-over-base}），则 $Z$ 在 $Y$ 上固有。
\end{lemma}

\begin{proof}
赋予 $Z$ 诱导的约化闭子概形结构；见 Schemes, Definition
\ref{schemes-definition-reduced-induced-scheme}。对每个 $i$，基变换
$Y_i \times_Y Z$ 是 $Y_i \times_Y X$ 的闭子概形，其底层闭子集为 $Z_i$。
由定义（借助 Cohomology of Schemes, Lemma
\ref{coherent-lemma-closed-proper-over-base}），投影
$Y_i \times_Y Z \to Y_i$ 为固有态射。因此由 Lemma
\ref{descent-lemma-descending-property-proper}，$Z \to Y$ 为固有态射。故由定义，
$Z$ 在 $Y$ 上固有。
\end{proof}

\begin{lemma}
\label{descent-lemma-descending-property-ample}
设 $f : X \to S$ 为概形态射，设 $\mathcal{L}$ 为可逆
$\mathcal{O}_X$-模，设 $\{g_i : S_i \to S\}_{i \in I}$ 为 fpqc 覆盖。
以 $f_i : X_i \to S_i$ 表示 $f$ 的基变换，并以 $\mathcal{L}_i$ 表示
$\mathcal{L}$ 到 $X_i$ 的拉回。下列条件等价：
\begin{enumerate}
\item $\mathcal{L}$ 在 $X/S$ 上丰沛；
\item 对每个 $i \in I$，$\mathcal{L}_i$ 在 $X_i/S_i$ 上丰沛。
\end{enumerate}
\end{lemma}

\begin{proof}
蕴含 (1) $\Rightarrow$ (2) 由 Morphisms, Lemma
\ref{morphisms-lemma-ample-base-change} 给出。假设对每个 $i \in I$，
$\mathcal{L}_i$ 在 $X_i/S_i$ 上丰沛。由 Morphisms, Definition
\ref{morphisms-definition-relatively-ample}，这蕴含 $X_i \to S_i$ 拟紧；
再由 Morphisms, Lemma \ref{morphisms-lemma-relatively-ample-separated}，
% P05-ERR-0038
这蕴含 $X_i \to S$ 分离。因此由 Lemmas
\ref{descent-lemma-descending-property-quasi-compact} 和
\ref{descent-lemma-descending-property-separated}，$f$ 拟紧且分离。

\medskip\noindent
这意味着
$\mathcal{A} = \bigoplus_{d \geq 0} f_*\mathcal{L}^{\otimes d}$
是拟凝聚分次 $\mathcal{O}_S$-代数（Schemes, Lemma
\ref{schemes-lemma-push-forward-quasi-coherent}）。此外，由 Cohomology of
Schemes, Lemma \ref{coherent-lemma-flat-base-change-cohomology}，
$\mathcal{A}$ 的构造与平坦基变换交换。特别地，若置
$\mathcal{A}_i = \bigoplus_{d \geq 0} f_{i, *}\mathcal{L}_i^{\otimes d}$
，则有 $\mathcal{A}_i = g_i^*\mathcal{A}$。由此，自然映射
$\psi_d : f^*\mathcal{A}_d \to \mathcal{L}^{\otimes d}$
（作为 $\mathcal{O}_X$-模的映射）拉回为自然映射
$\psi_{i, d} : f_i^*(\mathcal{A}_i)_d \to \mathcal{L}_i^{\otimes d}$
（作为 $\mathcal{O}_{X_i}$-模的映射）。由于 $\mathcal{L}_i$ 在 $X_i/S_i$
上丰沛，可知对任意点 $x_i \in X_i$，存在 $d \geq 1$，使得
$f_i^*(\mathcal{A}_i)_d \to \mathcal{L}_i^{\otimes d}$ 在 $x_i$ 处的茎上
为满射。这或者直接由相对丰沛模的定义推出，或者由 Morphisms, Lemma
\ref{morphisms-lemma-characterize-relatively-ample} 推出。若 $x \in X$，
则可以选取 $i$ 和映到 $x$ 的 $x_i \in X_i$。由于
$\mathcal{O}_{X, x} \to \mathcal{O}_{X_i, x_i}$ 平坦，因而忠实平坦，
可知对每个 $x \in X$，存在 $d \geq 1$，使得
$f^*\mathcal{A}_d \to \mathcal{L}^{\otimes d}$
在 $x$ 处的茎上为满射。这蕴含 Constructions, Lemma
\ref{constructions-lemma-invertible-map-into-relative-proj}
中对应于分次 $\mathcal{O}_X$-代数映射
$\psi : f^*\mathcal{A} \to \bigoplus_{d \geq 0} \mathcal{L}^{\otimes d}$
的开子集 $U(\psi) \subset X$ 等于 $X$。考虑对应的态射
$$
r_{\mathcal{L}, \psi} : X \longrightarrow \underline{\text{Proj}}_S(\mathcal{A})
$$
。由上述内容显然可知，$r_{\mathcal{L}, \psi}$ 到 $S_i$ 的基变换是态射
$r_{\mathcal{L}_i, \psi_i}$；由 Morphisms, Lemma
\ref{morphisms-lemma-characterize-relatively-ample}，后者为开浸入。因此由
Lemma \ref{descent-lemma-descending-property-open-immersion}，
$r_{\mathcal{L}, \psi}$ 为开浸入；再由 Morphisms, Lemma
\ref{morphisms-lemma-characterize-relatively-ample}，可知 $\mathcal{L}$ 在
$X/S$ 上丰沛。
\end{proof}















\section{在源上局部的态射性质}
\label{descent-section-properties-morphisms-local-source}

\noindent
经常会出现这样的情形：在与另一个态射预复合之后，可以证明某个态射具有
某一性质。在许多情况下，这也蕴含原态射具有该性质。我们用下面的定义将其
形式化。

\begin{definition}
\label{descent-definition-property-morphisms-local-source}
设 $\mathcal{P}$ 为概形态射的一个性质。设
$\tau \in \{Zariski, \linebreak[0] fpqc, \linebreak[0] fppf, \linebreak[0]
\etale, \linebreak[0] smooth, \linebreak[0] syntomic\}$。若对 $S$ 上任意概形
态射 $f : X \to Y$ 及任意 $\tau$-覆盖 $\{X_i \to X\}_{i \in I}$，都有
$$
f \text{ 具有 }\mathcal{P}
\Leftrightarrow
\text{每个 }X_i \to Y\text{ 具有 }\mathcal{P}.
$$
，则称 $\mathcal{P}$ 是{\it 在源上 $\tau$ 局部的}，或称它{\it 对
$\tau$-拓扑在源上局部}。
\end{definition}

\noindent
需要明确的是，由于同构总是覆盖，我们可见（或者要求）性质
$\mathcal{P}$ 对 $X \to Y$ 成立，当且仅当它对任意与 $X \to Y$ 同构的
箭头 $X' \to Y'$ 成立。若一个性质在源上为 $\tau$-局部的，则与出现在
$\tau$-覆盖中的态射预复合会保持该性质。下面是形式陈述。

\begin{lemma}
\label{descent-lemma-precompose-property-local-source}
设 $\tau \in \{fpqc, fppf, syntomic, smooth, \etale, Zariski\}$。设
$\mathcal{P}$ 为在源上 $\tau$-局部的态射性质。设 $f : X \to Y$ 具有性质
$\mathcal{P}$。对任意态射 $a : X' \to X$，若它分别为平坦态射、平坦且
局部有限表示的态射、syntomic 态射、
% P05-ERR-0039
\'etale 态射、开浸入，则复合 $f \circ a : X' \to Y$ 具有性质
$\mathcal{P}$。
\end{lemma}

\begin{proof}
这是因为可以把 $X' \to X$ 放入构成 $\tau$-覆盖的一族态射中。
\end{proof}

\begin{lemma}
\label{descent-lemma-largest-open-of-the-source}
设 $\tau \in \{fppf, syntomic, smooth, \etale\}$。设 $\mathcal{P}$ 为在源上
$\tau$-局部的态射性质。对任意概形态射 $f : X \to Y$，存在最大的开集
$W(f) \subset X$，使限制 $f|_{W(f)} : W(f) \to Y$ 具有 $\mathcal{P}$。
此外，若 $g : X' \to X$ 平坦且局部有限表示、为 syntomic 态射、光滑或为
\'etale 态射，并且 $f' = f \circ g : X' \to Y$，则
$g^{-1}(W(f)) = W(f')$。
\end{lemma}

\begin{proof}
考虑所有具有下列性质的态射 $g : X' \to X$ 的像 $g(X') \subset X$ 之并
$W$：
\begin{enumerate}
\item $g$ 平坦且局部有限表示、为 syntomic 态射、光滑或为 \'etale 态射；
\item 复合 $X' \to X \to Y$ 具有性质 $\mathcal{P}$。
\end{enumerate}
由于这样的态射 $g$ 为开映射（见 Morphisms, Lemma
\ref{morphisms-lemma-fppf-open}），$W \subset X$ 是 $X$ 的开子集。由于
$\mathcal{P}$ 在 $\tau$ 拓扑中为局部的，限制 $f|_W : W \to Y$ 具有性质
$\mathcal{P}$：事实上，我们得到 $W$ 的一个 $\tau$ 覆盖
$\{X' \to W\}$，其各拉回均具有 $\mathcal{P}$。这证明了 $W(f)$ 的存在性。
最后一句所述的相容性立即由 $W(f)$ 的构造推出。
\end{proof}

\begin{lemma}
\label{descent-lemma-properties-morphisms-local-source}
设 $\mathcal{P}$ 为概形态射的一个性质。设
$\tau \in \{fpqc, \linebreak[0] fppf, \linebreak[0]
\etale, \linebreak[0] smooth, \linebreak[0] syntomic\}$。假设：
\begin{enumerate}
\item 依照 $\tau$ 分别为 fpqc、fppf、\'etale、smooth 或 syntomic，该性质
在与平坦态射、平坦且局部有限表示的态射、\'etale 态射、光滑态射或
syntomic 态射预复合时保持；
\item 该性质在源上为 Zariski 局部的；
\item 该性质在目标上为 Zariski 局部的；
\item 对任意仿射概形态射 $f : X \to Y$ 及仿射概形的任意满态射
$X' \to X$，若依照 $\tau$ 分别为 fpqc、fppf、\'etale、smooth 或
syntomic，该满态射相应地为平坦、平坦有限表示、\'etale、光滑或
syntomic 态射，并且性质 $\mathcal{P}$ 对复合 $f' : X' \to Y$ 成立，
则性质 $\mathcal{P}$ 对 $f$ 成立。
\end{enumerate}
则 $\mathcal{P}$ 在源上为 $\tau$ 局部的。
\end{lemma}

\begin{proof}
这几乎立即由 $\tau$-覆盖的定义推出；见 Topologies, Definition
\ref{topologies-definition-fpqc-covering}
\ref{topologies-definition-fppf-covering}
\ref{topologies-definition-etale-covering}
\ref{topologies-definition-smooth-covering}, or
\ref{topologies-definition-syntomic-covering}
以及 Topologies, Lemma
\ref{topologies-lemma-fpqc-affine},
\ref{topologies-lemma-fppf-affine},
\ref{topologies-lemma-etale-affine},
\ref{topologies-lemma-smooth-affine}, or
\ref{topologies-lemma-syntomic-affine}.
省略细节。（提示：使用在源和目标上的局部性，把性质 $\mathcal{P}$ 的验证
归约到仿射概形之间的态射。然后应用 (1) 和 (4)。）
\end{proof}

\begin{remark}
\label{descent-remark-properties-morphisms-local-source-standard}
（这是对上面 Remarks \ref{descent-remark-descending-properties-standard} 和
\ref{descent-remark-descending-properties-morphisms-standard} 的重复。）在上述 Lemma
\ref{descent-lemma-properties-morphisms-local-source} 中，若 $\tau = smooth$，则在
条件 (4) 中可以假设该态射为（满的）标准光滑态射。$\tau = syntomic$ 或
$\tau = \etale$ 时也类似。
\end{remark}



\section{在源上的 fpqc 拓扑中局部的态射性质}
\label{descent-section-fpqc-local-source}

\noindent
下面给出一些在源上为 fpqc 局部的态射性质。

\begin{lemma}
\label{descent-lemma-flat-fpqc-local-source}
性质 $\mathcal{P}(f)=$``$f$ 平坦'' 在源上为 fpqc 局部的。
\end{lemma}

\begin{proof}
由于平坦性以局部环的映射来定义（Morphisms, Definition
\ref{morphisms-definition-flat}），必须证明如下代数事实：设
$A \to B \to C$ 为局部环的局部同态，并假设 $B \to C$ 平坦。则
$A \to B$ 平坦，当且仅当 $A \to C$ 平坦。若 $A \to B$ 平坦，则由
Algebra, Lemma \ref{algebra-lemma-composition-flat}，$A \to C$ 平坦。
反之，假设 $A \to C$ 平坦。注意 $B \to C$ 忠实平坦；见 Algebra,
Lemma \ref{algebra-lemma-local-flat-ff}。因此由 Algebra, Lemma
\ref{algebra-lemma-flat-permanence}，$A \to B$ 平坦。（另见 Morphisms,
Lemma \ref{morphisms-lemma-flat-permanence} 的直接证明。）
\end{proof}

\begin{lemma}
\label{descent-lemma-injective-local-rings-fpqc-local-source}
性质 $\mathcal{P}(f : X \to Y)=$``对每个 $x \in X$，局部环映射
$\mathcal{O}_{Y, f(x)} \to \mathcal{O}_{X, x}$ 单射'' 在源上为 fpqc 局部的。
\end{lemma}

\begin{proof}
省略。这只是一次（或许不太恰当的）幽默尝试。
\end{proof}





\section{在源上的 fppf 拓扑中局部的态射性质}
\label{descent-section-fppf-local-source}

\noindent
下面给出一些在源上为 fppf 局部的态射性质。

\begin{lemma}
\label{descent-lemma-locally-finite-presentation-fppf-local-source}
性质 $\mathcal{P}(f)=$``$f$ 局部有限表示'' 在源上为 fppf 局部的。
\end{lemma}

\begin{proof}
局部有限表示性在源和目标上都为 Zariski 局部的；见 Morphisms, Lemma
\ref{morphisms-lemma-locally-finite-presentation-characterize}。该性质在复合
下保持；见 Morphisms, Lemma
\ref{morphisms-lemma-composition-finite-presentation}。这证明了 Lemma
\ref{descent-lemma-properties-morphisms-local-source} 的 (1)、(2) 和 (3)。最后的
条件 (4) 是 Lemma \ref{descent-lemma-flat-finitely-presented-permanence-algebra}。
故结论成立。
\end{proof}

\begin{lemma}
\label{descent-lemma-locally-finite-type-fppf-local-source}
性质 $\mathcal{P}(f)=$``$f$ 局部有限型'' 在源上为 fppf 局部的。
\end{lemma}

\begin{proof}
局部有限型性在源和目标上都为 Zariski 局部的；见 Morphisms, Lemma
\ref{morphisms-lemma-locally-finite-type-characterize}。该性质在复合下保持；
见 Morphisms, Lemma \ref{morphisms-lemma-composition-finite-type}；而局部有限
表示的平坦态射局部有限型；见 Morphisms, Lemma
\ref{morphisms-lemma-finite-presentation-finite-type}。这证明了 Lemma
\ref{descent-lemma-properties-morphisms-local-source} 的 (1)、(2) 和 (3)。最后的
条件 (4) 是 Lemma \ref{descent-lemma-finite-type-local-source-fppf-algebra}。
故结论成立。
\end{proof}

\begin{lemma}
\label{descent-lemma-open-fppf-local-source}
性质 $\mathcal{P}(f)=$``$f$ 为开映射'' 在源上为 fppf 局部的。
\end{lemma}

\begin{proof}
开态射性显然在源和目标上都为 Zariski 局部的。该性质在复合下保持；见
Morphisms, Lemma \ref{morphisms-lemma-composition-open}；而有限表示的平坦
态射为开映射；见 Morphisms, Lemma \ref{morphisms-lemma-fppf-open}。这证明
了 Lemma \ref{descent-lemma-properties-morphisms-local-source} 的 (1)、(2) 和 (3)。
最后的条件 (4) 由 Morphisms, Lemma
\ref{morphisms-lemma-fpqc-quotient-topology} 推出。故结论成立。
\end{proof}

\begin{lemma}
\label{descent-lemma-universally-open-fppf-local-source}
性质 $\mathcal{P}(f)=$``$f$ 普遍开'' 在源上为 fppf 局部的。
\end{lemma}

\begin{proof}
设 $f : X \to Y$ 为概形态射，设 $\{X_i \to X\}_{i \in I}$ 为 fppf 覆盖。
% P05-ERR-0040
以 $f_i : X_i \to X$ 表示这些复合。我们必须证明 $f$ 普遍开，当且仅当每个
$f_i$ 都普遍开。若 $f$ 普遍开，则每个 $f_i$ 也普遍开，因为映射
$X_i \to X$ 普遍开，且普遍开态射的复合普遍开（Morphisms, Lemmas
\ref{morphisms-lemma-fppf-open} 和
\ref{morphisms-lemma-composition-open}）。反之，假设每个 $f_i$ 都普遍开。
设 $Y' \to Y$ 为概形态射，记 $X' = Y' \times_Y X$ 及
$X'_i = Y' \times_Y X_i$。注意，$\{X_i' \to X'\}_{i \in I}$ 也是 fppf
覆盖。由假设，态射 $f'_i : X_i' \to Y'$ 为开映射。因此由上面的 Lemma
\ref{descent-lemma-open-fppf-local-source}，$f' : X' \to Y'$ 为开映射，正如所需。
\end{proof}



\section{在源上的 syntomic 拓扑中局部的态射性质}
\label{descent-section-syntomic-local-source}

\noindent
下面给出一些在源上为 syntomic 局部的态射性质。

\begin{lemma}
\label{descent-lemma-syntomic-syntomic-local-source}
性质 $\mathcal{P}(f)=$``$f$ 为 syntomic 态射'' 在源上为 syntomic 局部的。
\end{lemma}

\begin{proof}
合并 Lemma \ref{descent-lemma-properties-morphisms-local-source}、Morphisms, Lemma
\ref{morphisms-lemma-syntomic-characterize}（在源和目标上为 Zariski 局部
的）、Morphisms, Lemma \ref{morphisms-lemma-composition-syntomic}
（预复合）以及 Lemma \ref{descent-lemma-syntomic-smooth-etale-permanence}
（第 (4) 部分）即得。
\end{proof}




\section{在源上的光滑拓扑中局部的态射性质}
\label{descent-section-smooth-local-source}

\noindent
下面给出一些在源上为光滑局部的态射性质。还请注意经由平坦态射下降光滑性
的结果 Lemma \ref{descent-lemma-smooth-permanence}；它在某些方面更强。

\begin{lemma}
\label{descent-lemma-smooth-smooth-local-source}
性质 $\mathcal{P}(f)=$``$f$ 光滑'' 在源上为光滑局部的。
\end{lemma}

\begin{proof}
合并 Lemma \ref{descent-lemma-properties-morphisms-local-source}、Morphisms, Lemma
\ref{morphisms-lemma-smooth-characterize}（在源和目标上为 Zariski 局部
的）、Morphisms, Lemma \ref{morphisms-lemma-composition-smooth}（预复合）
以及 Lemma \ref{descent-lemma-syntomic-smooth-etale-permanence}（第 (4) 部分）即得。
\end{proof}



\section{在源上的 \'etale 拓扑中局部的态射性质}
\label{descent-section-etale-local-source}

\noindent
下面给出一些在源上为 \'etale 局部的态射性质。

\begin{lemma}
\label{descent-lemma-etale-etale-local-source}
性质 $\mathcal{P}(f)=$``$f$ 为 \'etale 态射'' 在源上为 \'etale 局部的。
\end{lemma}

\begin{proof}
合并 Lemma \ref{descent-lemma-properties-morphisms-local-source}、Morphisms, Lemma
\ref{morphisms-lemma-etale-characterize}（在源和目标上为 Zariski 局部的）、
Morphisms, Lemma \ref{morphisms-lemma-composition-etale}（预复合）以及
Lemma \ref{descent-lemma-syntomic-smooth-etale-permanence}（第 (4) 部分）即得。
\end{proof}

\begin{lemma}
\label{descent-lemma-locally-quasi-finite-etale-local-source}
性质 $\mathcal{P}(f)=$``$f$ 局部拟有限'' 在源上为 \'etale 局部的。
\end{lemma}

\begin{proof}
我们使用 Lemma \ref{descent-lemma-properties-morphisms-local-source}。由 Morphisms,
Lemma \ref{morphisms-lemma-locally-quasi-finite-characterize}，局部拟有限性
在源和目标上为 Zariski 局部的。由 Morphisms, Lemmas
\ref{morphisms-lemma-composition-quasi-finite} 和
\ref{morphisms-lemma-etale-locally-quasi-finite}，局部拟有限态射与 \'etale
态射预复合后仍局部拟有限。最后，假设 $X \to Y$ 为仿射概形态射，且
$X' \to X$ 为仿射概形的满 \'etale 态射，并且 $X' \to Y$ 局部拟有限。
则 $X' \to Y$ 有限型；由 Lemma
\ref{descent-lemma-finite-type-local-source-fppf-algebra}，$X \to Y$ 也有限型。此外，
由假设 $X' \to Y$ 的纤维有限，因而 $X \to Y$ 的纤维也有限。由
Morphisms, Lemma \ref{morphisms-lemma-quasi-finite}，可知 $X \to Y$
拟有限。这证明了 Lemma \ref{descent-lemma-properties-morphisms-local-source} 的
最后一个假设，从而完成证明。
\end{proof}

\begin{lemma}
\label{descent-lemma-unramified-etale-local-source}
性质 $\mathcal{P}(f)=$``$f$ 非分歧'' 在源上为 \'etale 局部的。
性质 $\mathcal{P}(f)=$``$f$ 为 G-非分歧态射'' 在源上为 \'etale 局部的。
\end{lemma}

\begin{proof}
我们使用 Lemma \ref{descent-lemma-properties-morphisms-local-source}。由 Morphisms,
Lemma \ref{morphisms-lemma-unramified-characterize}，非分歧性（相应地，
G-非分歧性）在源和目标上为 Zariski 局部的。由 Morphisms, Lemmas
\ref{morphisms-lemma-composition-unramified} 和
\ref{morphisms-lemma-etale-smooth-unramified}，非分歧（相应地，G-非分歧）
态射与 \'etale 态射预复合后仍非分歧（相应地，G-非分歧）。最后，假设
$X \to Y$ 为仿射概形态射，且 $f : X' \to X$ 为仿射概形的满 \'etale
态射，并且 $X' \to Y$ 非分歧（相应地，G-非分歧）。则 $X' \to Y$
有限型（相应地，有限表示）；由 Lemma
\ref{descent-lemma-finite-type-local-source-fppf-algebra}（相应地，Lemma
\ref{descent-lemma-flat-finitely-presented-permanence-algebra}），$X \to Y$ 也有限型
（相应地，有限表示）。由 Morphisms, Lemma
\ref{morphisms-lemma-triangle-differentials-smooth}，有短正合列
$$
0 \to f^*\Omega_{X/Y} \to \Omega_{X'/Y} \to \Omega_{X'/X} \to 0.
$$
。由于 $X' \to Y$ 非分歧，中间项为零。又因 $f$ 忠实平坦，可知
$\Omega_{X/Y} = 0$。因此 $X \to Y$ 非分歧（相应地，G-非分歧）；见
Morphisms, Lemma \ref{morphisms-lemma-unramified-omega-zero}。这证明了
Lemma \ref{descent-lemma-properties-morphisms-local-source} 的最后一个假设，从而
完成证明。
\end{proof}




\section{在源与目标上 \'etale 局部的态射性质}
\label{descent-section-properties-etale-local-source-target}

\noindent
设 $\mathcal{P}$ 为概形态射的一个性质。``$\mathcal{P}$ 在源和目标上为
\'etale 局部的''这句话有一种直观含义。然而，这个概念并不等同于要求
$\mathcal{P}$ 既在源上为 \'etale 局部、又在目标上为 \'etale 局部。在
进一步讨论之前，先给出两个简单的例子。

\begin{example}
\label{descent-example-silly-one}
考虑按规则 $\mathcal{P}(X \to Y) = $``$Y$ 局部 Noether'' 定义的概形态射
性质 $\mathcal{P}$。读者可以验证，它在源上为 \'etale 局部的，也在目标上
为 \'etale 局部的（省略；见 Lemma \ref{descent-lemma-Noetherian-local-fppf}）。
但是，下列陈述{\bf 不}成立：若 $f : X \to Y$ 具有 $\mathcal{P}$，并且
$g : Y \to Z$ 为 \'etale 态射，则 $g \circ f$ 具有 $\mathcal{P}$。事实上，
$f$ 可以是 $Y$ 上的恒等态射，而 $g$ 可以是局部 Noether 概形 $Y$ 到非局部
Noether 概形 $Z$ 的开浸入。
\end{example}

\noindent
下面的例子在某种意义上更糟。

\begin{example}
\label{descent-example-silly-two}
考虑按下列规则定义的概形态射性质 $\mathcal{P}$：
$\mathcal{P}(f : X \to Y) = $``对每个 $y \in Y$，若它是某个 $f(x)$ 的特化，
其中 $x \in X$，则局部环 $\mathcal{O}_{Y, y}$ 为 Noether 环''。我们来验证
它在源上和目标上均为 \'etale 局部的。下面将自由使用 Schemes, Lemma
\ref{schemes-lemma-specialize-points}。

\medskip\noindent
在目标上局部：设 $\{g_i : Y_i \to Y\}$ 为 \'etale 覆盖。以
$f_i : X_i \to Y_i$ 表示 $f$ 的基变换，并以 $h_i : X_i \to X$ 表示投影。
假设 $\mathcal{P}(f)$。设
% P05-ERR-0041
$f(x_i) \leadsto y_i$ 为一个特化。则 $f(h_i(x_i)) \leadsto g_i(y_i)$，所以
$\mathcal{P}(f)$ 蕴含 $\mathcal{O}_{Y, g_i(y_i)}$ 为 Noether 环。此外，
$\mathcal{O}_{Y, g_i(y_i)} \to \mathcal{O}_{Y_i, y_i}$ 是某个 \'etale 环映射
的局部化。因此由 Algebra, Lemma \ref{algebra-lemma-Noetherian-permanence}，
$\mathcal{O}_{Y_i, y_i}$ 为 Noether 环。反之，假设对所有 $i$，
$\mathcal{P}(f_i)$ 均成立。设 $f(x) \leadsto y$ 为一个特化。选取 $i$ 和映到
$y$ 的 $y_i \in Y_i$。由于 $x$ 可视为
$\Spec(\mathcal{O}_{Y, y}) \times_Y X$ 的一个点，并且
$\mathcal{O}_{Y, y} \to \mathcal{O}_{Y_i, y_i}$ 忠实平坦，存在映到 $x$ 的点
$x_i \in \Spec(\mathcal{O}_{Y_i, y_i}) \times_Y X$。于是 $x_i \in X_i$，
并且 $f_i(x_i)$ 特化到 $y_i$。由 $\mathcal{P}(f_i)$ 可知
$\mathcal{O}_{Y_i, y_i}$ 为 Noether 环；再由 Algebra, Lemma
\ref{algebra-lemma-descent-Noetherian}，这蕴含 $\mathcal{O}_{Y, y}$ 为
Noether 环。

\medskip\noindent
在源上局部：设 $\{h_i : X_i \to X\}$ 为 \'etale 覆盖。以
$f_i : X_i \to Y$ 表示复合 $f \circ h_i$。假设 $\mathcal{P}(f)$。设
% P05-ERR-0042
$f(x_i) \leadsto y$ 为一个特化。则 $f(h_i(x_i)) \leadsto y$，所以
$\mathcal{P}(f)$ 蕴含 $\mathcal{O}_{Y, y}$ 为 Noether 环。因此
$\mathcal{P}(f_i)$ 成立。反之，假设对所有 $i$，$\mathcal{P}(f_i)$ 均成立。
设 $f(x) \leadsto y$ 为一个特化。选取 $i$ 和映到 $x$ 的 $x_i \in X_i$。
则 $y$ 是 $f_i(x_i) = f(x)$ 的特化。因此 $\mathcal{P}(f_i)$ 蕴含
$\mathcal{O}_{Y, y}$ 为 Noether 环，正如所需。

\medskip\noindent
我们断言，存在交换图
$$
\xymatrix{
U \ar[d]_a \ar[r]_h & V \ar[d]^b \\
X \ar[r]^f & Y
}
$$
，其中竖直箭头是满的 \'etale 态射，$h$ 具有 $\mathcal{P}$，而 $f$ 不具有
$\mathcal{P}$。事实上，令
$$
Y =
\Spec\Big(
\mathbf{C}[x_n; n \in \mathbf{Z}]/(x_n x_m; n \not = m)
\Big)
$$
，并设 $X \subset Y$ 为如下开子概形：它是所有坐标都满足 $x_n = 0$ 的点
的补集。令 $U = X$，令 $V = X \amalg Y$，令 $a, b$ 为显然的映射，并令
$h : U \to V$ 为 $U = X$ 到 $V$ 的第一个直和项的包含。上述断言成立，
因为 $U$ 局部 Noether，而 $Y$ 不是。
\end{example}

\noindent
那么，在源与目标上为 \'etale 局部的性质，其正确概念应是什么？我们认为，
类比 Morphisms, Definition \ref{morphisms-definition-property-local}，应取
下面的定义。

\begin{definition}
\label{descent-definition-local-source-target}
设 $\mathcal{P}$ 为概形态射的一个性质。若满足下列条件，则称
$\mathcal{P}$ {\it 在源与目标上为 \'etale 局部的}：
\begin{enumerate}
\item （在与 \'etale 映射预复合时稳定）若 $f : X \to Y$ 为 \'etale 态射，
且 $g : Y \to Z$ 具有 $\mathcal{P}$，则 $g \circ f$ 具有 $\mathcal{P}$；
\item （在 \'etale 基变换下稳定）若 $f : X \to Y$ 具有 $\mathcal{P}$，且
$Y' \to Y$ 为 \'etale 态射，则基变换 $f' : Y' \times_Y X \to Y'$ 具有
$\mathcal{P}$；
\item （局部性）给定态射 $f : X \to Y$，下列条件等价：
\begin{enumerate}
\item $f$ 具有 $\mathcal{P}$；
\item 对每个 $x \in X$，存在交换图
$$
\xymatrix{
U \ar[d]_a \ar[r]_h & V \ar[d]^b \\
X \ar[r]^f & Y
}
$$
，其中竖直箭头为 \'etale 态射，且存在满足 $a(u) = x$ 的 $u \in U$，使
$h$ 具有 $\mathcal{P}$。
\end{enumerate}
\end{enumerate}
\end{definition}

\noindent
事实证明，这个定义排除了 Examples \ref{descent-example-silly-one} 和
\ref{descent-example-silly-two} 中所见的行为。我们将在 Remark
\ref{descent-remark-compare-definitions} 中把它与 Deligne 和 Mumford 的论文
\cite{DM} 中的定义作比较。此外，一个在源与目标上为 \'etale 局部的性质，
在源上和目标上分别都是 \'etale 局部的。最后，正如 Lemma
\ref{descent-lemma-etale-local-source-target} 将表明的，其逆命题几乎成立。

\begin{lemma}
\label{descent-lemma-local-source-target-implies}
设 $\mathcal{P}$ 为在源与目标上 \'etale 局部的概形态射性质。则
\begin{enumerate}
\item $\mathcal{P}$ 在源上为 \'etale 局部的；
\item $\mathcal{P}$ 在目标上为 \'etale 局部的；
\item $\mathcal{P}$ 在与 \'etale 态射后复合时稳定：若 $f : X \to Y$
具有 $\mathcal{P}$，且 $g : Y \to Z$ 为 \'etale 态射，则 $g \circ f$
具有 $\mathcal{P}$；
\item $\mathcal{P}$ 具有保持性：给定 $f : X \to Y$ 及 \'etale 态射
$g : Y \to Z$，若 $g \circ f$ 具有 $\mathcal{P}$，则 $f$ 具有
$\mathcal{P}$。
\end{enumerate}
\end{lemma}

\begin{proof}
我们完整写出所有细节。

\medskip\noindent
(1) 的证明。设 $f : X \to Y$ 为概形态射，设
$\{X_i \to X\}_{i \in I}$ 为 $X$ 的 \'etale 覆盖。若每个复合
$h_i : X_i \to Y$ 都具有 $\mathcal{P}$，则对每个 $x \in X$，可以找到
$i \in I$ 及映到 $x$ 的点 $x_i \in X_i$。于是
$(X_i, x_i) \to (X, x)$ 是芽的 \'etale 态射，
$\text{id}_Y : Y \to Y$ 为 \'etale 态射，而 $h_i$ 正如 Definition
\ref{descent-definition-local-source-target} 第 (3) 部分所述。因此 $f$ 具有
$\mathcal{P}$。反之，若 $f$ 具有 $\mathcal{P}$，则由 Definition
\ref{descent-definition-local-source-target} 第 (1) 部分，每个 $X_i \to Y$ 都具有
$\mathcal{P}$。

\medskip\noindent
(2) 的证明。设 $f : X \to Y$ 为概形态射，设
$\{Y_i \to Y\}_{i \in I}$ 为 $Y$ 的 \'etale 覆盖。写
$X_i = Y_i \times_Y X$，并以 $h_i : X_i \to Y_i$ 表示 $f$ 的基变换。若
每个 $h_i : X_i \to Y_i$ 都具有 $\mathcal{P}$，则对每个 $x \in X$，选取
$i \in I$ 及映到 $x$ 的点 $x_i \in X_i$。于是
$(X_i, x_i) \to (X, x)$ 是芽的 \'etale 态射，$Y_i \to Y$ 为 \'etale
态射，而 $h_i$ 正如 Definition \ref{descent-definition-local-source-target}
第 (3) 部分所述。因此 $f$ 具有 $\mathcal{P}$。反之，若 $f$ 具有
$\mathcal{P}$，则由 Definition \ref{descent-definition-local-source-target}
第 (2) 部分，每个 $X_i \to Y_i$ 都具有 $\mathcal{P}$。

\medskip\noindent
(3) 的证明。假设 $f : X \to Y$ 具有 $\mathcal{P}$，且 $g : Y \to Z$ 为
\'etale 态射。对每个 $x \in X$，可将 $(X, x) \to (X, x)$ 视为芽的
\'etale 态射；$Y \to Z$ 为 \'etale 态射；而 $h = f$ 正如 Definition
\ref{descent-definition-local-source-target} 第 (3) 部分所述。因此 $g \circ f$
具有 $\mathcal{P}$。

\medskip\noindent
(4) 的证明。设 $f : X \to Y$ 为态射，$g : Y \to Z$ 为 \'etale 态射，
并且 $g \circ f$ 具有 $\mathcal{P}$。由 Definition
\ref{descent-definition-local-source-target} 第 (2) 部分，
$\text{pr}_Y : Y \times_Z X \to Y$ 具有 $\mathcal{P}$。另一方面，态射
$(f, 1) : X \to Y \times_Z X$ 是 \'etale 投影
$\text{pr}_X : Y \times_Z X \to X$ 的截面，因而为 \'etale 态射；见
Morphisms, Lemma \ref{morphisms-lemma-etale-permanence}。所以由 Definition
\ref{descent-definition-local-source-target} 第 (1) 部分，
$f = \text{pr}_Y \circ (f, 1)$ 具有 $\mathcal{P}$。
\end{proof}

\noindent
下面的引理是 Morphisms, Lemma
\ref{morphisms-lemma-locally-P-characterize} 的类似结果。

\begin{lemma}
\label{descent-lemma-local-source-target-characterize}
设 $\mathcal{P}$ 为在源与目标上 \'etale 局部的概形态射性质。设
$f : X \to Y$ 为概形态射。下列条件等价：
\begin{enumerate}
\item[(a)] $f$ 具有性质 $\mathcal{P}$；
\item[(b)] 对每个 $x \in X$，存在芽的 \'etale 态射
$a : (U, u) \to (X, x)$、\'etale 态射 $b : V \to Y$ 和态射
$h : U \to V$，使得 $f \circ a = b \circ h$，且 $h$ 具有
$\mathcal{P}$；
\item[(c)] 对任意交换图
$$
\xymatrix{
U \ar[d]_a \ar[r]_h & V \ar[d]^b \\
X \ar[r]^f & Y
}
$$
，若 $a$、$b$ 为 \'etale 态射，则态射 $h$ 具有 $\mathcal{P}$；
\item[(d)] 对某个如 (c) 的图，若 $a : U \to X$ 满，则 $h$ 具有
$\mathcal{P}$；
\item[(e)] 存在 \'etale 覆盖 $\{Y_i \to Y\}_{i \in I}$，使得每个基变换
$Y_i \times_Y X \to Y_i$ 都具有 $\mathcal{P}$；
\item[(f)] 存在 \'etale 覆盖 $\{X_i \to X\}_{i \in I}$，使得每个复合
$X_i \to Y$ 都具有 $\mathcal{P}$；
\item[(g)] 存在 \'etale 覆盖 $\{Y_i \to Y\}_{i \in I}$，并且对每个
$i \in I$，存在 \'etale 覆盖
$\{X_{ij} \to Y_i \times_Y X\}_{j \in J_i}$，使得每个态射
$X_{ij} \to Y_i$ 都具有 $\mathcal{P}$。
\end{enumerate}
\end{lemma}

\begin{proof}
(a) 与 (b) 的等价性是 Definition \ref{descent-definition-local-source-target} 的
一部分。(a) 与 (e) 的等价性是 Lemma
\ref{descent-lemma-local-source-target-implies} 的第 (2) 部分。(a) 与 (f) 的等价性
是 Lemma \ref{descent-lemma-local-source-target-implies} 的第 (1) 部分。由于现在
(a) 与 (e)、(f) 等价，可知 (a) 与 (g) 等价。

\medskip\noindent
显然 (c) 蕴含 (a)。若 (a) 成立，则对任意如 (c) 的图，由 Definition
\ref{descent-definition-local-source-target} 第 (1) 部分，态射 $f \circ a$ 具有
$\mathcal{P}$；进而由 Lemma \ref{descent-lemma-local-source-target-implies}
第 (4) 部分，$h$ 具有 $\mathcal{P}$。所以 (a) 与 (c) 等价。显然 (c)
蕴含 (d)。为证明 (d) 蕴含 (a)，假设有如 (c) 的图，其中
$a : U \to X$ 满，且 $h$ 具有 $\mathcal{P}$。由 Lemma
\ref{descent-lemma-local-source-target-implies} 第 (3) 部分，$b \circ h$ 具有
$\mathcal{P}$。由于 $\{a : U \to X\}$ 是 \'etale 覆盖，由 Lemma
\ref{descent-lemma-local-source-target-implies} 第 (1) 部分，可知 $f$ 具有
$\mathcal{P}$。
\end{proof}

\noindent
下面这个引理的结论似乎并非形式的；也就是说，它确实用到了 \'etale 态射
几何中的某些内容。

\begin{lemma}
\label{descent-lemma-etale-local-source-target}
设 $\mathcal{P}$ 为概形态射的一个性质。假设
\begin{enumerate}
\item $\mathcal{P}$ 在源上为 \'etale 局部的；
\item $\mathcal{P}$ 在目标上为 \'etale 局部的；
\item $\mathcal{P}$ 在与开浸入后复合时稳定：若 $f : X \to Y$ 具有
$\mathcal{P}$，且 $Y \subset Z$ 为开子概形，则 $X \to Z$ 具有
$\mathcal{P}$。
\end{enumerate}
则 $\mathcal{P}$ 在源与目标上为 \'etale 局部的。
\end{lemma}

\begin{proof}
设 $\mathcal{P}$ 为满足本引理条件 (1)、(2)、(3) 的概形态射性质。由 Lemma
\ref{descent-lemma-precompose-property-local-source}，$\mathcal{P}$ 在与 \'etale
态射预复合时稳定。由 Lemma \ref{descent-lemma-pullback-property-local-target}，
$\mathcal{P}$ 在 \'etale 基变换下稳定。因此只需证明 Definition
\ref{descent-definition-local-source-target} 的第 (3) 部分成立。

\medskip\noindent
更准确地，假设概形态射 $f : X \to Y$ 满足 Definition
\ref{descent-definition-local-source-target} 第 (3)(b) 部分。换言之，对每个
$x \in X$，存在 \'etale 态射 $a_x : U_x \to X$、映到 $x$ 的点
$u_x \in U_x$、\'etale 态射 $b_x : V_x \to Y$ 及态射
$h_x : U_x \to V_x$，使得 $f \circ a_x = b_x \circ h_x$，且 $h_x$ 具有
$\mathcal{P}$。一旦证明 $f$ 具有 $\mathcal{P}$，本引理的证明就完成了。
置 $U = \coprod U_x$、$a = \coprod a_x$、$V = \coprod V_x$、
$b = \coprod b_x$ 及 $h = \coprod h_x$。得到交换图
$$
\xymatrix{
U \ar[d]_a \ar[r]_h & V \ar[d]^b \\
X \ar[r]^f & Y
}
$$
，其中 $a$、$b$ 为 \'etale 态射，且 $a$ 满。注意，$h$ 具有
$\mathcal{P}$，因为每个 $h_x$ 都如此，且 $\mathcal{P}$ 在目标上为
\'etale 局部的。由于 $a$ 满，而 $\mathcal{P}$ 在源上为 \'etale 局部的，
只需证明 $b \circ h$ 具有 $\mathcal{P}$。这样，本引理归约为证明
$\mathcal{P}$ 在与 \'etale 态射后复合时稳定。

\medskip\noindent
在余下的证明中，设 $f : X \to Y$ 为具有性质 $\mathcal{P}$ 的态射，且
$g : Y \to Z$ 为 \'etale 态射。考虑下列陈述：
\begin{enumerate}
\item[(-)] 不加额外假设时，$g \circ f$ 具有性质 $\mathcal{P}$。
\item[(A)] 只要 $Z$ 为仿射概形，$g \circ f$ 就具有性质 $\mathcal{P}$。
\item[(AA)] 只要 $X$ 和 $Z$ 为仿射概形，$g \circ f$ 就具有性质
$\mathcal{P}$。
\item[(AAA)] 只要 $X$、$Y$ 和 $Z$ 均为仿射概形，$g \circ f$ 就具有性质
$\mathcal{P}$。
\end{enumerate}
一旦证明 (-)，本引理的证明就完成了。

\medskip\noindent
断言 1：(AAA) $\Rightarrow$ (AA)。事实上，设 $f : X \to Y$、
$g : Y \to Z$ 如上，且 $X$、$Z$ 为仿射概形。由于 $X$ 仿射，因而拟紧，
可以找到有限多个仿射开集 $Y_i \subset Y$，$i = 1, \ldots, n$，使得
$X = \bigcup_{i = 1, \ldots, n} f^{-1}(Y_i)$。置 $X_i = f^{-1}(Y_i)$。
由 Lemma \ref{descent-lemma-pullback-property-local-target}，每个态射
$X_i \to Y_i$ 都具有 $\mathcal{P}$。由于 $\mathcal{P}$ 在目标上为
\'etale 局部的，
$\coprod_{i = 1, \ldots, n} X_i \to \coprod_{i = 1, \ldots, n} Y_i$
具有 $\mathcal{P}$。将 (AAA) 应用于
$\coprod_{i = 1, \ldots, n} X_i \to \coprod_{i = 1, \ldots, n} Y_i$
及 \'etale 态射
$\coprod_{i = 1, \ldots, n} Y_i \to Z$，可知
$\coprod_{i = 1, \ldots, n} X_i \to Z$ 具有 $\mathcal{P}$。现在
$\{\coprod_{i = 1, \ldots, n} X_i \to X\}$ 是 \'etale 覆盖；由于
$\mathcal{P}$ 在源上为 \'etale 局部的，可知 $X \to Z$ 具有
$\mathcal{P}$，正如所需。

\medskip\noindent
断言 2：(AAA) $\Rightarrow$ (A)。事实上，设 $f : X \to Y$、
$g : Y \to Z$ 如上，且 $Z$ 为仿射概形。选取仿射开覆盖
$X = \bigcup X_i$。由于 $\mathcal{P}$ 在源上为 \'etale 局部的，每个
$f|_{X_i} : X_i \to Y$ 都具有 $\mathcal{P}$。由断言 1 从 (AAA) 推出的
(AA)，对每个 $i$，$X_i \to Z$ 都具有 $\mathcal{P}$。由于
$\{X_i \to X\}$ 是 \'etale 覆盖，而 $\mathcal{P}$ 在源上为 \'etale
局部的，可知 $X \to Z$ 具有 $\mathcal{P}$。

\medskip\noindent
断言 3：(AAA) $\Rightarrow$ (-)。事实上，设 $f : X \to Y$、
$g : Y \to Z$ 如上。选取仿射开覆盖 $Z = \bigcup Z_i$。置
$Y_i = g^{-1}(Z_i)$ 及 $X_i = f^{-1}(Y_i)$。由 Lemma
\ref{descent-lemma-pullback-property-local-target}，每个态射 $X_i \to Y_i$ 都具有
$\mathcal{P}$。由断言 2 从 (AAA) 推出的 (A)，对每个 $i$，
$X_i \to Z_i$ 都具有 $\mathcal{P}$。由于 $\mathcal{P}$ 在目标上为局部的，
且 $X_i = (g \circ f)^{-1}(Z_i)$，可知 $X \to Z$ 具有 $\mathcal{P}$。

\medskip\noindent
因此为证明本引理，只需证明 (AAA)。设 $f : X \to Y$ 和 $g : Y \to Z$
如上，且 $X, Y, Z$ 均为仿射概形。注意，仿射概形的 \'etale 态射具有
普遍有界的纤维；见 Morphisms, Lemma
\ref{morphisms-lemma-etale-locally-quasi-finite} 和 Lemma
\ref{morphisms-lemma-locally-quasi-finite-qc-source-universally-bounded}。
因此可以对界定 $Y \to Z$ 的纤维度数的整数 $n$ 作归纳。在 \'etale 态射
的情形，该整数的描述见 Morphisms, Lemma
\ref{morphisms-lemma-etale-universally-bounded}。若 $n = 1$，则
$Y \to Z$ 为开浸入；见 Lemma
\ref{descent-lemma-universally-injective-etale-open-immersion}；结论由本引理的假设
(3) 推出。假设 $n > 1$。

\medskip\noindent
考虑下列交换图：
$$
\xymatrix{
X \times_Z Y \ar[d] \ar[r]_{f_Y} &
Y \times_Z Y \ar[d] \ar[r]_-{\text{pr}} &
Y \ar[d] \\
X \ar[r]^f &
Y \ar[r]^g &
Z
}
$$
注意，我们有开闭子概形分解
$Y \times_Z Y = \Delta_{Y/Z}(Y) \amalg Y'$；见 Morphisms, Lemma
\ref{morphisms-lemma-diagonal-unramified-morphism}。作为基变换，第二投影
$\text{pr} : Y \times_Z Y \to Y$ 的纤维度数以 $n$ 为界；见 Morphisms,
Lemma \ref{morphisms-lemma-base-change-universally-bounded}。另一方面，
$\text{pr}|_{\Delta(Y)} : \Delta(Y) \to Y$ 为同构，且每个纤维恰有一个点。
因此应用 Morphisms, Lemma
\ref{morphisms-lemma-etale-universally-bounded}，可知限制
$\text{pr}|_{Y'} : Y' \to Y$ 的纤维度数以 $n - 1$ 为界。置
$X' = f_Y^{-1}(Y')$。图示如下：
$$
\xymatrix{
X \amalg X' \ar@{=}[d] \ar[r]_-{f \amalg f'} &
\Delta(Y) \amalg Y' \ar@{=}[d] \ar[r] &
Y \ar@{=}[d] \\
X \times_Z Y \ar[r]^{f_Y} &
Y \times_Z Y \ar[r]^-{\text{pr}} &
Y
}
$$
由于 $\mathcal{P}$ 在目标上为 \'etale 局部的，因而在 \'etale 基变换下
稳定（见 Lemma \ref{descent-lemma-pullback-property-local-target}），可知 $f_Y$
具有 $\mathcal{P}$。又因 $\mathcal{P}$ 在源上为 \'etale 局部的，
$f' = f_Y|_{X'}$ 具有 $\mathcal{P}$。由归纳假设，$X' \to Y$ 具有
$\mathcal{P}$。由于 $\mathcal{P}$ 在源上为局部的，并且
% P05-ERR-0043
$\{X \to X \times_Z Y, X' \to X \times_Y Z\}$ 是 \'etale 覆盖，可知
$\text{pr} \circ f_Y$ 具有 $\mathcal{P}$。注意，可将 $g \circ f$ 视为态射
$g \circ f : X \to g(Y)$。由于 $\text{pr} \circ f_Y$ 是态射
$g \circ f : X \to g(Y)$ 沿 \'etale 覆盖 $\{Y \to g(Y)\}$ 的拉回，且
$\mathcal{P}$ 在目标上为 \'etale 局部的，可知
$g \circ f : X \to g(Y)$ 具有性质 $\mathcal{P}$。最后，再次应用本引理的
假设 (3)，可知 $g \circ f : X \to Z$ 具有性质 $\mathcal{P}$。
\end{proof}

\begin{remark}
\label{descent-remark-list-local-source-target}
利用 Lemma \ref{descent-lemma-etale-local-source-target} 及本章前面各节完成的工作，
很容易列出在源与目标上为 \'etale 局部的态射类型。在每种情形中，我们
分别列出蕴含该性质在源上为 \'etale 局部的引理，以及蕴含它在目标上为
\'etale 局部的引理。在每种情形中，Lemma
\ref{descent-lemma-etale-local-source-target} 的第三个假设都显然成立，故予省略。
列表如下：
\begin{enumerate}
\item 平坦；见 Lemmas \ref{descent-lemma-flat-fpqc-local-source} 和
\ref{descent-lemma-descending-property-flat}；
\item 局部有限表示；见 Lemmas
\ref{descent-lemma-locally-finite-presentation-fppf-local-source} 和
\ref{descent-lemma-descending-property-locally-finite-presentation}；
\item 局部有限型；见 Lemmas \ref{descent-lemma-locally-finite-type-fppf-local-source}
和 \ref{descent-lemma-descending-property-locally-finite-type}；
\item 普遍开；见 Lemmas \ref{descent-lemma-universally-open-fppf-local-source} 和
\ref{descent-lemma-descending-property-universally-open}；
\item syntomic；见 Lemmas \ref{descent-lemma-syntomic-syntomic-local-source} 和
\ref{descent-lemma-descending-property-syntomic}；
\item 光滑；见 Lemmas \ref{descent-lemma-smooth-smooth-local-source} 和
\ref{descent-lemma-descending-property-smooth}；
\item \'etale；见 Lemmas \ref{descent-lemma-etale-etale-local-source} 和
\ref{descent-lemma-descending-property-etale}；
\item 局部拟有限；见 Lemmas
\ref{descent-lemma-locally-quasi-finite-etale-local-source} 和
\ref{descent-lemma-descending-property-quasi-finite}；
\item 非分歧；见 Lemmas \ref{descent-lemma-unramified-etale-local-source} 和
\ref{descent-lemma-descending-property-unramified}；
\item G-非分歧；见 Lemmas \ref{descent-lemma-unramified-etale-local-source} 和
\ref{descent-lemma-descending-property-unramified}；
\item 可按需在此添加更多性质。
\end{enumerate}
\end{remark}

\begin{remark}
\label{descent-remark-compare-definitions}
至此，对于态射性质 $\mathcal{P}$ ``在源和目标上为 \'etale 局部''的含义，
我们有三种可能的定义：
\begin{enumerate}
\item[(ST)] $\mathcal{P}$ 在源上为 \'etale 局部的，且 $\mathcal{P}$ 在
目标上为 \'etale 局部的；
\item[(DM)] （Deligne 和 Mumford 的论文 \cite[Page 100]{DM} 中的定义）
对每个图
$$
\xymatrix{
U \ar[d]_a \ar[r]_h & V \ar[d]^b \\
X \ar[r]^f & Y
}
$$
，若竖直箭头为满的 \'etale 态射，则
$\mathcal{P}(h) \Leftrightarrow \mathcal{P}(f)$；
\item[(SP)] $\mathcal{P}$ 在源与目标上为 \'etale 局部的。
\end{enumerate}
本节已经看到 (SP) $\Rightarrow$ (DM) $\Rightarrow$ (ST)。Examples
\ref{descent-example-silly-one} 和 \ref{descent-example-silly-two} 表明，两个蕴含均不能
反转。最后，Lemma \ref{descent-lemma-etale-local-source-target} 表明，对在与开浸入
后复合时稳定的态射性质，这一区别会消失；实践中总是这种情形。
\end{remark}

\begin{lemma}
\label{descent-lemma-etale-etale-local-source-target}
设 $\mathcal{P}$ 为在源与目标上 \'etale 局部的概形态射性质。给定概形的交换图
$$
\vcenter{
\xymatrix{
X' \ar[d]_{g'} \ar[r]_{f'} & Y' \ar[d]^g \\
X \ar[r]^f & Y
}
}
\quad\text{带标记点}\quad
\vcenter{
\xymatrix{
x' \ar[d] \ar[r] & y' \ar[d] \\
x \ar[r] & y
}
}
$$
，若 $g'$ 在 $x'$ 处为 \'etale 态射，且 $g$ 在 $y'$ 处为 \'etale 态射，
则
$x \in W(f) \Leftrightarrow x' \in W(f')$
，其中 $W(-)$ 如 Lemma \ref{descent-lemma-largest-open-of-the-source} 中所定义。
\end{lemma}

\begin{proof}
Lemma \ref{descent-lemma-largest-open-of-the-source} 适用，因为由 Lemma
\ref{descent-lemma-local-source-target-implies}，$\mathcal{P}$ 在源上为 \'etale
局部的。

\medskip\noindent
假设 $x \in W(f)$。设 $U' \subset X'$ 和 $V' \subset Y'$ 分别为 $x'$ 和
$y'$ 的开邻域，使得 $f'(U') \subset V'$、$g'(U') \subset W(f)$，且
$g'|_{U'}$ 与 $g|_{V'}$ 为 \'etale 态射。由 Definition
\ref{descent-definition-local-source-target} 的性质 (1)，
$f \circ g'|_{U'} = g \circ f'|_{U'}$ 具有 $\mathcal{P}$。再由 Lemma
\ref{descent-lemma-local-source-target-implies} 的 (4)，
$f'|_{U'} : U' \to V'$ 具有性质 $\mathcal{P}$。进而由 Lemma
\ref{descent-lemma-local-source-target-implies} 的 (3)，
% P05-ERR-0044
$f'_{U'} : U' \to Y'$ 具有 $\mathcal{P}$。所以由定义
$U' \subset W(f')$，进而 $x' \in W(f')$。

\medskip\noindent
假设 $x' \in W(f')$。设 $U' \subset X'$ 和 $V' \subset Y'$ 分别为 $x'$ 和
$y'$ 的开邻域，使得 $f'(U') \subset V'$、$U' \subset W(f')$，且
$g'|_{U'}$ 与 $g|_{V'}$ 为 \'etale 态射。由 $W(f')$ 的定义，
$U' \to Y'$ 具有 $\mathcal{P}$。由 Lemma
\ref{descent-lemma-local-source-target-implies} 的 (4)，$U' \to V'$ 具有
$\mathcal{P}$。再由 Lemma \ref{descent-lemma-local-source-target-implies} 的 (3)，
$U' \to Y$ 具有 $\mathcal{P}$。令
$U \subset X$ 为 \'etale（因而开）态射
% P05-ERR-0045
$g'|_U' : U' \to X$ 的像。则 $\{U' \to U\}$ 为 \'etale 覆盖；由 Lemma
\ref{descent-lemma-local-source-target-implies} 的 (1)，可知 $U \to Y$ 具有
$\mathcal{P}$。所以由定义 $U \subset W(f)$，进而 $x \in W(f)$。
\end{proof}

\begin{lemma}
\label{descent-lemma-orbits}
设 $k$ 为域，设 $n \geq 2$。对满足 $1 \leq i, j \leq n$、$i \not = j$ 和
$d \geq 0$ 的 i、j、d，以 $T_{i, j, d}$ 表示 $\mathbf{A}^n_k$ 的自同构；
它在坐标中由下式给出：
$$
(x_1, \ldots, x_n) \longmapsto
(x_1, \ldots, x_{i - 1}, x_i + x_j^d, x_{i + 1}, \ldots, x_n)
$$
设 $W \subset \mathbf{A}^n_k$ 为非空开子概形，并且对上述所有 $i, j, d$，
都有 $T_{i, j, d}(W) = W$。则或者 $W = \mathbf{A}^n_k$，或者 $k$ 的特征
为 $p > 0$，且 $\mathbf{A}^n_k \setminus W$ 是由闭点组成的有限集，这些点
的坐标在 $\mathbf{F}_p$ 上代数。
\end{lemma}

\begin{proof}
为证明本结论，可以用任意扩域替换 $k$。设 $Z$ 为
$\mathbf{A}^n_k \setminus W$ 的一个不可约分支。为得到矛盾，假设
$\dim(Z) \geq 1$。则存在域扩张 $k'/k$ 及
$Z_{k'} \subset \mathbf{A}^n_{k'}$ 的 $k'$-值点
$\xi = (\xi_1, \ldots, \xi_n) \in (k')^n$，使得 $x_1, \ldots, x_n$
中至少一个在素域上超越。断言：$\xi$ 在由变换 $T_{i, j, d}$ 生成的群作用
下的轨道在 $\mathbf{A}^n_{k'}$ 中 Zariski 稠密。该断言将给出所需的矛盾。

\medskip\noindent
若 $k'$ 的特征为零，则仅使用算子 $T_{i, j, 0}$ 就已经足够，因为这些算子
将 $\xi$ 变换为诸点
$$
(\xi_1 + a_1, \ldots, \xi_n + a_n)
$$
，其中 $(a_1, \ldots, a_n) \in \mathbf{Z}_{\geq 0}^n$ 任意。若特征为
$p > 0$，重新编号后可以假设 $\xi_n$ 在 $\mathbf{F}_p$ 上超越。对
$i < n$ 依次应用算子 $T_{i, n, d}$，可见 $\xi$ 的轨道包含诸元素
$$
(\xi_1 + P_1(\xi_n), \ldots, \xi_{n - 1} + P_{n - 1}(\xi_n), \xi_n)
$$
，其中 $(P_1, \ldots, P_{n - 1}) \in \mathbf{F}_p[t]$ 任意。因此轨道的
Zariski 闭包包含坐标超平面 $x_n = \xi_n$。对另一个坐标重复论证，可知
对任意 $P \in \mathbf{F}_p[t]$，只要 $\xi_i + P(\xi_n)$ 在
$\mathbf{F}_p$ 上超越，该 Zariski 闭包就包含
$x_i = \xi_i + P(\xi_n)$。这样的 $P$ 有无穷多个，故断言成立。

\medskip\noindent
当然，若 $Z = \{z\}$ 的维数为 $0$，且 $z$ 在 $\kappa(z)$ 中的坐标并非
都在 $\mathbf{F}_p$ 上代数，上一段的论证也同样适用。引理得证。
\end{proof}

\begin{lemma}
\label{descent-lemma-etale-tau-local-source-target}
设 $\mathcal{P}$ 为概形态射的一个性质。假设
\begin{enumerate}
\item $\mathcal{P}$ 在源上为 \'etale 局部的；
\item $\mathcal{P}$ 在目标上为光滑局部的；
\item $\mathcal{P}$ 在与开浸入后复合时稳定：若 $f : X \to Y$ 具有
$\mathcal{P}$，且 $Y \subset Z$ 为开子概形，则 $X \to Z$ 具有
$\mathcal{P}$。
\end{enumerate}
给定概形的交换图
$$
\vcenter{
\xymatrix{
X' \ar[d]_{g'} \ar[r]_{f'} & Y' \ar[d]^g \\
X \ar[r]^f & Y
}
}
\quad\text{带标记点}\quad
\vcenter{
\xymatrix{
x' \ar[d] \ar[r] & y' \ar[d] \\
x \ar[r] & y
}
}
$$
，使得
% P05-ERR-0046
$g$ 在 $y'$ 处光滑，并且 $X' \to X \times_Y Y'$ 在 $x'$ 处为 \'etale
态射，则 $x \in W(f) \Leftrightarrow x' \in W(f')$，其中 $W(-)$ 如
Lemma \ref{descent-lemma-largest-open-of-the-source} 中所定义。
\end{lemma}

\begin{proof}
由于 $\mathcal{P}$ 在源上为 \'etale 局部的，可知 $x \in W(f)$ 当且仅当
$x$ 在 $X \times_Y Y'$ 中的像属于 $W(X \times_Y Y' \to Y')$。因此可以
假设本引理中的图是 Cartesian 的。

\medskip\noindent
假设 $x \in W(f)$。由于 $\mathcal{P}$ 在目标上为光滑局部的，
$(g')^{-1}W(f) = W(f) \times_Y Y' \to Y'$ 具有 $\mathcal{P}$。所以
$(g')^{-1}W(f) \subset W(f')$，从而 $x' \in W(f')$。

\medskip\noindent
假设 $x' \in W(f')$。对 $y'$ 的任意开邻域 $V' \subset Y'$，可以用 $V'$
替换 $Y'$，并用 $U' = (f')^{-1}V'$ 替换 $X'$：这是因为
$V' \to Y'$ 光滑，因而 $W(f') \to Y'$ 的基变换
$W(f') \cap U' \to V'$ 具有性质 $\mathcal{P}$。所以可以假设存在 $Y$ 上的
\'etale 态射 $Y' \to \mathbf{A}^n_Y$；见 Morphisms, Lemma
\ref{morphisms-lemma-smooth-etale-over-affine-space}。图示如下：
$$
\xymatrix{
X' \ar[r] \ar[d] & Y' \ar[d] \\
\mathbf{A}^n_X \ar[r]_{f_n} \ar[d] & \mathbf{A}^n_Y \ar[d] \\
X \ar[r]^f & Y
}
$$
由 Lemma \ref{descent-lemma-etale-local-source-target}（并利用 \'etale 覆盖也是光滑
覆盖），$\mathcal{P}$ 在源与目标上为 \'etale 局部的。由 Lemma
\ref{descent-lemma-etale-etale-local-source-target}，$W(f')$ 是开集
$W(f_n) \subset \mathbf{A}^n_X$ 的逆像。特别地，$W(f_n)$ 包含一个位于
$x$ 上方的点。用 $W(f_n)$ 的像（它是开集）替换 $X$ 后，可以假设
$W(f_n) \to X$ 满。断言：$W(f_n) = \mathbf{A}^n_X$。该断言蕴含 $f$ 具有
$\mathcal{P}$，因为 $\mathcal{P}$ 在光滑拓扑中为局部的，并且
$\{\mathbf{A}^n_Y \to Y\}$ 是光滑覆盖。

\medskip\noindent
本质上，该断言源于以下事实：$W(f_n) \subset \mathbf{A}^n_X$ 是一个与每个
$\mathbf{A}^n_X \to X$ 的纤维相交的``平移不变''开集。不过，若要按这一路径
给出论证，就必须在 $Y$ 上作 \'etale 局部化以产生足够多的平移，这会有些
麻烦。我们改用 Lemma \ref{descent-lemma-orbits} 的自同构及仿射空间的 \'etale
态射。可以假设 $n \geq 2$。事实上，若 $n = 0$，结论已经成立。若
$n = 1$，考虑图
$$
\xymatrix{
\mathbf{A}^2_X \ar[r]_{f_2} \ar[d]_p & \mathbf{A}^2_Y \ar[d] \\
\mathbf{A}^1_X \ar[r]^{f_1} & \mathbf{A}^1_Y
}
$$
由本证明第一段，有 $p^{-1}(W(f_1)) \subset W(f_2)$。所以
$W(f_2) \to X$ 仍然满，我们可以改用 $f_2$。假设 $n \geq 2$。

\medskip\noindent
对满足 $1 \leq i, j \leq n$、$i \not = j$ 和 $d \geq 0$ 的任意
i、j、d，以 $T_{i, j, d}$ 表示 Lemma \ref{descent-lemma-orbits} 中定义的
$\mathbf{A}^n$ 的自同构。于是得到交换图
$$
\xymatrix{
\mathbf{A}^n_X \ar[r]_{f_n} \ar[d]_{T_{i, j, d}} &
\mathbf{A}^n_Y \ar[d]^{T_{i, j, d}} \\
\mathbf{A}^n_X \ar[r]^{f_n} & \mathbf{A}^n_Y
}
$$
，其竖直箭头为同构。因此
$T_{i, j, d}(W(f_n)) = W(f_n)$。应用 Lemma \ref{descent-lemma-orbits}，可知对任意
$x \in X$，纤维 $W(f_n)_x \subset \mathbf{A}^n_x$ 或者等于
$\mathbf{A}^n_x$（这正是我们要证明的），或者 $\kappa(x)$ 的特征为
$p > 0$，且 $W(f_n)_x$ 是由闭点组成的有限集
$Z_x \subset \mathbf{A}^n_x$ 的补集。第二种可能不成立。事实上，考虑
由下式给出的态射 $T_p : \mathbf{A}^n \to \mathbf{A}^n$：
$$
(x_1, \ldots, x_n) \mapsto (x_1 - x_1^p, \ldots, x_n - x_n^p)
$$
如上，我们得到交换图
$$
\xymatrix{
\mathbf{A}^n_X \ar[r]_{f_n} \ar[d]_{T_p} &
\mathbf{A}^n_Y \ar[d]^{T_p} \\
\mathbf{A}^n_X \ar[r]^{f_n} & \mathbf{A}^n_Y
}
$$
态射 $T_p : \mathbf{A}^n_X \to \mathbf{A}^n_X$ 在位于 $x$ 上方的每个点
处为 \'etale 态射，而态射 $T_p : \mathbf{A}^n_Y \to \mathbf{A}^n_Y$
在位于 $x$ 在 $Y$ 中的像上方的每个点处为 \'etale 态射。（省略细节；
提示：计算导数。）因此
% P05-ERR-0047
$$
T_p^{-1}(W) \cap \mathbf{A}^n_x = W \cap \mathbf{A}^n_x
$$
，这是由 Lemma \ref{descent-lemma-etale-etale-local-source-target} 得到的（我们已经
知道 $\mathcal{P}$ 在源与目标上为 \'etale 局部的）。由于
$T_p : \mathbf{A}^n_x \to \mathbf{A}^n_x$ 是次数 $p^n > 1$ 的有限
\'etale 态射，若 $Z_x$ 非空，则它包含更大的集合 $T_p^{-1}(Z_x)$。
这个矛盾完成了证明。
\end{proof}





\section{在源与目标上局部的芽态射性质}
\label{descent-section-local-source-target-at-point}

\noindent
本节讨论 Section \ref{descent-section-properties-etale-local-source-target} 的内容
对于概形芽态射的类似版本。

\begin{definition}
\label{descent-definition-local-source-target-at-point}
设 $\mathcal{Q}$ 为概形芽态射的一个性质。若对任意交换图
$$
\xymatrix{
(U', u') \ar[d]_a \ar[r]_{h'} & (V', v') \ar[d]^b \\
(U, u) \ar[r]^h & (V, v)
}
$$
，其中竖直箭头为芽的 \'etale 态射，都有
$\mathcal{Q}(h) \Leftrightarrow \mathcal{Q}(h')$，则称 $\mathcal{Q}$
{\it 在源与目标上为 \'etale 局部的}。
\end{definition}

\begin{lemma}
\label{descent-lemma-local-source-target-global-implies-local}
设 $\mathcal{P}$ 为在源与目标上 \'etale 局部的概形态射性质。考虑按下列规则
定义的芽态射性质 $\mathcal{Q}$：
$$
\mathcal{Q}((X, x) \to (S, s))
\Leftrightarrow
\text{存在一个代表元 }U \to S
\text{ 其具有 }\mathcal{P}
$$
则 $\mathcal{Q}$ 在源与目标上为 \'etale 局部的，含义如 Definition
\ref{descent-definition-local-source-target-at-point}。
\end{lemma}

\begin{proof}
若芽态射 $(X, x) \to (S, s)$ 具有 $\mathcal{Q}$，则存在任意小的 $x$ 的邻域
$U \subset X$ 和 $s$ 的邻域 $V \subset S$，使得
$(X, x) \to (S, s)$ 的代表 $U \to V$ 具有 $\mathcal{P}$。这由 Lemma
\ref{descent-lemma-local-source-target-implies} 推出。设
$$
\xymatrix{
(U', u') \ar[r]_{h'} \ar[d]_a & (V', v') \ar[d]^b \\
(U, u) \ar[r]^h & (V, v)
}
$$
如 Definition \ref{descent-definition-local-source-target-at-point} 所述。选取
$U_1 \subset U$ 及 $h$ 的代表 $h_1 : U_1 \to V$。选取
$V'_1 \subset V'$ 及 $b$ 的 \'etale 代表 $b_1 : V'_1 \to V$（Definition
\ref{descent-definition-etale-morphism-germs}）。选取 $U'_1 \subset U'$ 及 $a$、
$h'$ 的代表 $a_1 : U'_1 \to U_1$、$h'_1 : U'_1 \to V'_1$，其中 $a_1$
为 \'etale 态射。缩小 $U'_1$ 后，可以假设
$h_1 \circ a_1 = b_1 \circ h'_1$。由证明开头的说明，我们要证明
$u' \in W(h'_1) \Leftrightarrow u \in W(h_1)$，其中 $W(-)$ 如 Lemma
\ref{descent-lemma-largest-open-of-the-source} 中所定义。因此本引理由 Lemma
\ref{descent-lemma-etale-etale-local-source-target} 推出。
\end{proof}

\begin{lemma}
\label{descent-lemma-local-source-target-local-implies-global}
设 $\mathcal{P}$ 为在源与目标上 \'etale 局部的概形态射性质。令
% P05-ERR-0048
$Q$ 为相伴的芽态射性质；见 Lemma
\ref{descent-lemma-local-source-target-global-implies-local}。设 $f : X \to Y$ 为
概形态射。下列条件等价：
\begin{enumerate}
\item $f$ 具有性质 $\mathcal{P}$；
\item 对每个 $x \in X$，芽态射 $(X, x) \to (Y, f(x))$ 具有性质
$\mathcal{Q}$。
\end{enumerate}
\end{lemma}

\begin{proof}
蕴含 (1) $\Rightarrow$ (2) 直接由定义给出。蕴含 (2) $\Rightarrow$ (1)
也由 Definition \ref{descent-definition-local-source-target} 的第 (3) 部分推出。
\end{proof}

\noindent
芽态射 $(X, x) \to (S, s)$ 决定一个良定义的局部环映射。因此下面的引理
是有意义的。

\begin{lemma}
\label{descent-lemma-flat-at-point}
芽态射的性质
$$
\mathcal{P}((X, x) \to (S, s)) =
\mathcal{O}_{S, s} \to \mathcal{O}_{X, x}\text{ 平坦}
$$
在源与目标上为 \'etale 局部的。
\end{lemma}

\begin{proof}
给定如 Definition \ref{descent-definition-local-source-target-at-point} 中的图，得到
如下局部环的局部同态图：
$$
\xymatrix{
\mathcal{O}_{U', u'} & \mathcal{O}_{V', v'} \ar[l] \\
\mathcal{O}_{U, u} \ar[u] & \mathcal{O}_{V, v} \ar[l] \ar[u]
}
$$
注意，竖直箭头为 \'etale 环映射的局部化；特别地，它们本质有限表示、
平坦且非分歧（见 Algebra, Section \ref{algebra-section-etale}）。特别地，
竖直映射忠实平坦；见 Algebra, Lemma
\ref{algebra-lemma-local-flat-ff}。现在，若上方水平箭头平坦，则将 Algebra,
Lemma \ref{algebra-lemma-flat-permanence} 应用于
$R = \mathcal{O}_{V, v}$、$S = \mathcal{O}_{U, u}$ 和
$M = \mathcal{O}_{U', u'}$，可知下方水平箭头平坦。若下方水平箭头平坦，
则环映射
$$
\mathcal{O}_{V', v'} \otimes_{\mathcal{O}_{V, v}} \mathcal{O}_{U, u}
\longleftarrow
\mathcal{O}_{V', v'}
$$
由 Algebra, Lemma \ref{algebra-lemma-flat-base-change} 可知平坦。而环映射
$$
\mathcal{O}_{U', u'}
\longleftarrow
\mathcal{O}_{V', v'} \otimes_{\mathcal{O}_{V, v}} \mathcal{O}_{U, u}
$$
是 $\mathcal{O}_{U, u}$ 的两个 \'etale 环扩张之间映射的局部化，因此由
Algebra, Lemma \ref{algebra-lemma-map-between-etale} 可知平坦。
\end{proof}

\begin{lemma}
\label{descent-lemma-etale-on-fiber}
考虑概形态射的交换图
$$
\xymatrix{
U' \ar[r] \ar[d] & V' \ar[d] \\
U \ar[r] & V
}
$$
，其中竖直箭头为 \'etale 态射，且点 $v' \in V'$ 映到 $v \in V$。则纤维
的态射 $U'_{v'} \to U_v$ 为 \'etale 态射。
\end{lemma}

\begin{proof}
注意，$U'_v \to U_v$ 是 \'etale 态射 $U' \to U$ 的基变换，因而为
\'etale 态射。概形 $U'_v$ 是 $V'_v$ 上的概形。由 Morphisms, Lemma
\ref{morphisms-lemma-etale-over-field}，概形 $V'_v$ 是 $\kappa(v)$ 的有限
可分域扩张的谱的不交并。其中一个是 $v' = \Spec(\kappa(v'))$。所以
$U'_{v'}$ 是 $U'_v$ 的开闭子概形，从而 $U'_{v'} \to U'_v \to U_v$
为 \'etale 态射（它是开浸入与 \'etale 态射的复合；见 Morphisms,
Section \ref{morphisms-section-etale}）。
\end{proof}

\noindent
给定概形的芽态射 $(X, x) \to (S, s)$，可以把{\it 纤维}定义为芽
$(U_s, x)$ 的同构类，其中 $U \to S$ 为任意代表。我们常常略微滥用记号，
直接写成 $(X_s, x)$。

\begin{lemma}
\label{descent-lemma-dimension-local-ring-fibre}
设 $d \in \{0, 1, 2, \ldots, \infty\}$。芽态射的性质
$$
\mathcal{P}_d((X, x) \to (S, s)) =
\text{局部环 }
\mathcal{O}_{X_s, x}
\text{ 的纤维维数为 }d
$$
在源与目标上为 \'etale 局部的。
\end{lemma}

\begin{proof}
给定如 Definition \ref{descent-definition-local-source-target-at-point} 中的图，得到
把 $u'$ 映到 $u$ 的纤维的 \'etale 态射 $U'_{v'} \to U_v$；见 Lemma
\ref{descent-lemma-etale-on-fiber}。因此结论由 Lemma
\ref{descent-lemma-dimension-local-ring-local} 推出。
\end{proof}

\begin{lemma}
\label{descent-lemma-transcendence-degree-at-point}
设 $r \in \{0, 1, 2, \ldots, \infty\}$。芽态射的性质
$$
\mathcal{P}_r((X, x) \to (S, s))
\Leftrightarrow
\text{trdeg}_{\kappa(s)} \kappa(x) = r
$$
在源与目标上为 \'etale 局部的。
\end{lemma}

\begin{proof}
给定如 Definition \ref{descent-definition-local-source-target-at-point} 中的图，得到
如下局部环的局部同态图：
$$
\xymatrix{
\mathcal{O}_{U', u'} & \mathcal{O}_{V', v'} \ar[l] \\
\mathcal{O}_{U, u} \ar[u] & \mathcal{O}_{V, v} \ar[l] \ar[u]
}
$$
注意，竖直箭头为 \'etale 环映射的局部化；特别地，它们非分歧（见
Algebra, Section \ref{algebra-section-etale}）。所以
$\kappa(u')/\kappa(u)$ 和 $\kappa(v')/\kappa(v)$ 是有限可分域扩张。因此
% P05-ERR-0049
$\text{trdeg}_{\kappa(v)} \kappa(u) = \text{trdeg}_{\kappa(v')} \kappa(u)$
，这证明了本引理。
\end{proof}

\noindent
设 $(X, x)$ 为概形的芽。$X$ 在 $x$ 处的维数，是 $x$ 在 $X$ 中各开邻域的
维数之最小值；任意足够小的开邻域都具有此维数。因此，这是芽的同构类的
不变量。我们把它简记为 $\dim_x(X)$。

\begin{lemma}
\label{descent-lemma-dimension-at-point}
设 $d \in \{0, 1, 2, \ldots, \infty\}$。芽态射的性质
$$
\mathcal{P}_d((X, x) \to (S, s))
\Leftrightarrow
\dim_x (X_s) = d
$$
在源与目标上为 \'etale 局部的。
\end{lemma}

\begin{proof}
给定如 Definition \ref{descent-definition-local-source-target-at-point} 中的图，得到
把 $u'$ 映到 $u$ 的纤维的 \'etale 态射 $U'_{v'} \to U_v$；见 Lemma
\ref{descent-lemma-etale-on-fiber}。于是等式
$\dim_u(U_v) = \dim_{u'}(U'_{v'})$ 由 Lemma
\ref{descent-lemma-dimension-at-point-local} 推出。
\end{proof}







\section{概形上概形的下降数据}
\label{descent-section-descent-datum}

\noindent
本节的大部分论证都是形式的，仅依赖于下降数据的定义。在 Simplicial
Spaces, Section \ref{spaces-simplicial-section-simplicial-descent} 中，我们
将考察它与单纯概形的关系，这会在一定程度上澄清情形。

\begin{definition}
\label{descent-definition-descent-datum}
设 $f : X \to S$ 为概形态射。
\begin{enumerate}
\item 设 $V \to X$ 为 $X$ 上的概形。一个{\it $V/X/S$ 的下降数据}是
$X \times_S X$ 上的概形同构
$\varphi : V \times_S X \to X \times_S V$，它满足{\it 上循环条件}：图
$$
\xymatrix{
V \times_S X \times_S X \ar[rd]^{\varphi_{01}} \ar[rr]_{\varphi_{02}} &
&
X \times_S X \times_S V\\
&
X \times_S V \times_S X \ar[ru]^{\varphi_{12}}
}
$$
交换（记号含义显然）。
\item 也称对 $(V/X, \varphi)$ 是一个{\it 相对于 $X \to S$ 的下降数据}。
\item 相对于 $X \to S$ 的下降数据之间的一个{\it 态射
$g : (V/X, \varphi) \to (V'/X, \varphi')$}，是使下图交换的 $X$ 上概形
态射 $g : V \to V'$：
$$
\xymatrix{
V \times_S X \ar[r]_{\varphi} \ar[d]_{g \times \text{id}_X} &
X \times_S V \ar[d]^{\text{id}_X \times g} \\
V' \times_S X \ar[r]^{\varphi'} & X \times_S V'
}
$$
。
\end{enumerate}
\end{definition}

\noindent
上述定义会产生各种``奇妙的''恒等式。例如，$\varphi$ 沿对角态射
$\Delta : X \to X \times_S X$ 的拉回可视为态射
$\Delta^*\varphi : V \to V$。这是因为
$X \times_{\Delta, X \times_S X} (V \times_S X) = V$，并且
$X \times_{\Delta, X \times_S X} (X \times_S V) = V$。事实上，
$\Delta^*\varphi$ 等于恒等态射。若你不熟悉这些材料，这是一个很好的练习。

\begin{remark}
\label{descent-remark-easier}
设 $X \to S$ 为概形态射，设 $(V/X, \varphi)$ 为相对于 $X \to S$ 的下降
数据。可以把同构 $\varphi$ 看作同构
$$
(X \times_S X) \times_{\text{pr}_0, X} V
\longrightarrow
(X \times_S X) \times_{\text{pr}_1, X} V
$$
（作为 $X \times_S X$ 上的概形）。因此粗略地说，可以把 $\varphi$ 看作映射
$\varphi : \text{pr}_0^*V \to \text{pr}_1^*V$\footnote{遗憾的是，我们在这里
为箭头选择了``错误的''方向。根据``函数''与``空间''对偶的一般原则，
Definitions \ref{descent-definition-descent-datum} 和
\ref{descent-definition-descent-datum-for-family-of-morphisms} 中的方向，应当与
Definition \ref{descent-definition-descent-datum-quasi-coherent} 中所取的方向相反。}。
此时上循环条件断言
$\text{pr}_{02}^*\varphi =
\text{pr}_{12}^*\varphi \circ \text{pr}_{01}^*\varphi$.
从这个角度看，它与拟凝聚层上的下降数据情形十分相似。
\end{remark}

\noindent
下面给出具有固定目标的一族态射的情形下的定义。

\begin{definition}
\label{descent-definition-descent-datum-for-family-of-morphisms}
设 $S$ 为概形，设 $\{X_i \to S\}_{i \in I}$ 为以 $S$ 为目标的一族态射。
\begin{enumerate}
\item 一个{\it 相对于态射族 $\{X_i \to S\}$ 的下降数据
$(V_i, \varphi_{ij})$}由下列内容给出：对每个 $i \in I$，一个 $X_i$ 上的
概形 $V_i$；对每对 $(i, j) \in I^2$，一个 $X_i \times_S X_j$ 上的概形同构
$\varphi_{ij} : V_i \times_S X_j \to X_i \times_S V_j$
，使得对每个三元指标组 $(i, j, k) \in I^3$，图
$$
\xymatrix{
V_i \times_S X_j \times_S X_k
\ar[rd]^{\text{pr}_{01}^*\varphi_{ij}}
\ar[rr]_{\text{pr}_{02}^*\varphi_{ik}} &
&
X_i \times_S X_j \times_S V_k\\
&
X_i \times_S V_j \times_S X_k
\ar[ru]^{\text{pr}_{12}^*\varphi_{jk}}
}
$$
（作为 $X_i \times_S X_j \times_S X_k$ 上的概形）交换（记号含义显然）。
\item 下降数据之间的一个{\it 态射
$\psi : (V_i, \varphi_{ij}) \to (V'_i, \varphi'_{ij})$}由一族 $X_i$-概形
态射 $\psi_i : V_i \to V'_i$ 给出，记为
$\psi = (\psi_i)_{i \in I}$，并要求所有图
$$
\xymatrix{
V_i \times_S X_j \ar[r]_{\varphi_{ij}} \ar[d]_{\psi_i \times \text{id}} &
X_i \times_S V_j \ar[d]^{\text{id} \times \psi_j} \\
V'_i \times_S X_j \ar[r]^{\varphi'_{ij}} & X_i \times_S V'_j
}
$$
交换。
\end{enumerate}
\end{definition}

\noindent
例如，当询问相对曲线的纤维化范畴在 fpqc 拓扑中是否为叠时，这个概念会
自然出现（它不是——至少在坚持只使用概形时不是）。

\begin{remark}
\label{descent-remark-easier-family}
设 $S$ 为概形，设 $\{X_i \to S\}_{i \in I}$ 为以 $S$ 为目标的一族态射。
设 $(V_i, \varphi_{ij})$ 为相对于 $\{X_i \to S\}$ 的下降数据。可以把同构
$\varphi_{ij}$ 看作同构
$$
(X_i \times_S X_j) \times_{\text{pr}_0, X_i} V_i
\longrightarrow
(X_i \times_S X_j) \times_{\text{pr}_1, X_j} V_j
$$
（作为 $X_i \times_S X_j$ 上的概形）。因此粗略地说，可以把
$\varphi_{ij}$ 看作 $X_i \times_S X_j$ 上的同构
$\text{pr}_0^*V_i \to \text{pr}_1^*V_j$。此时上循环条件断言
$\text{pr}_{02}^*\varphi_{ik} =
\text{pr}_{12}^*\varphi_{jk} \circ \text{pr}_{01}^*\varphi_{ij}$.
从这个角度看，它与拟凝聚层上的下降数据情形十分相似。
\end{remark}

\noindent
我们通常使用仅由单个态射组成的族这一版本，理由是下面的引理。

\begin{lemma}
\label{descent-lemma-family-is-one}
设 $S$ 为概形，设 $\{X_i \to S\}_{i \in I}$ 为以 $S$ 为目标的一族态射。
置 $X = \coprod_{i \in I} X_i$，并将其视为 $S$-概形。存在范畴的典范等价
$$
\begin{matrix}
\text{下降数据范畴 } \\
\text{相对于该族 } \{X_i \to S\}_{i \in I}
\end{matrix}
\longrightarrow
\begin{matrix}
\text{ 下降数据范畴} \\
\text{ 相对于 } X/S
\end{matrix}
$$
，它把 $(V_i, \varphi_{ij})$ 映到 $(V, \varphi)$，其中
$V = \coprod_{i\in I} V_i$ 且 $\varphi = \coprod \varphi_{ij}$。
\end{lemma}

\begin{proof}
注意，$X \times_S X = \coprod_{ij} X_i \times_S X_j$，更高纤维积的情形
也类似。给出态射 $V \to X$ 恰好等同于给出一族态射 $V_i \to X_i$。
而给出下降数据 $\varphi$ 恰好等同于给出一族 $\varphi_{ij}$。
\end{proof}

\begin{lemma}
\label{descent-lemma-pullback}
概形上概形下降数据的拉回如下。
\begin{enumerate}
\item 设
$$
\xymatrix{
X' \ar[r]_f \ar[d]_{a'} & X \ar[d]^a \\
S' \ar[r]^h & S
}
$$
为概形态射的交换图。构造
$$
(V \to X, \varphi) \longmapsto f^*(V \to X, \varphi) = (V' \to X', \varphi')
$$
，其中 $V' = X' \times_X V$，且 $\varphi'$ 定义为复合
$$
\xymatrix{
V' \times_{S'} X' \ar@{=}[r] &
(X' \times_X V) \times_{S'} X' \ar@{=}[r] &
(X' \times_{S'} X') \times_{X \times_S X} (V \times_S X)
\ar[d]^{\text{id} \times \varphi} \\
X' \times_{S'} V' \ar@{=}[r] &
X' \times_{S'} (X' \times_X V) &
(X' \times_{S'} X') \times_{X \times_S X} (X \times_S V) \ar@{=}[l]
}
$$
，它定义了从相对于 $X \to S$ 的下降数据范畴到相对于 $X' \to S'$ 的
下降数据范畴的函子。
\item 给定两个使该图交换的态射 $f_i : X' \to X$，$i = 0, 1$，函子
$f_0^*$ 与 $f_1^*$ 典范同构。
\end{enumerate}
\end{lemma}

\begin{proof}
省略 (1) 的证明，但指出在 Remark \ref{descent-remark-easier} 引入的记号下，态射
$\varphi'$ 就是态射 $(f \times f)^*\varphi$。对 (2)，我们说明给出该函子
同构的是哪个态射 $f_0^*V \to f_1^*V$。事实上，由于 $f_0$ 和 $f_1$ 都
嵌入上述交换图，存在唯一的态射 $r : X' \to X \times_S X$，使得
$f_i = \text{pr}_i \circ r$。于是取
\begin{eqnarray*}
f_0^*V & = &
X' \times_{f_0, X} V \\
& = &
X' \times_{\text{pr}_0 \circ r, X} V \\
& = &
X' \times_{r, X \times_S X} (X \times_S X) \times_{\text{pr}_0, X} V \\
& \xrightarrow{\varphi} &
X' \times_{r, X \times_S X} (X \times_S X) \times_{\text{pr}_1, X} V \\
& = &
X' \times_{\text{pr}_1 \circ r, X} V \\
& = &
X' \times_{f_1, X} V \\
& = & f_1^*V
\end{eqnarray*}
省略对此构造确实有效的验证。
\end{proof}

\begin{definition}
\label{descent-definition-pullback-functor}
设 $S, S', X, X', f, a, a', h$ 如 Lemma \ref{descent-lemma-pullback} 中所述。该引理
构造的函子
$$
(V, \varphi) \longmapsto f^*(V, \varphi)
$$
称为下降数据上的{\it 拉回函子}。
\end{definition}

\begin{lemma}[概形族上概形下降数据的拉回]
\label{descent-lemma-pullback-family}
设 $\mathcal{U} = \{U_i \to S'\}_{i \in I}$ 和
$\mathcal{V} = \{V_j \to S\}_{j \in J}$ 为具有固定目标的态射族。设
$\alpha : I \to J$、$h : S' \to S$ 和 $g_i : U_i \to V_{\alpha(i)}$
构成具有固定目标的映射族之间的态射；见 Sites, Definition
\ref{sites-definition-morphism-coverings}。
\begin{enumerate}
\item 设 $(Y_j, \varphi_{jj'})$ 为相对于态射族
% P05-ERR-0050
$\{V_j \to S'\}$ 的下降数据。系统
$$
\left(
g_i^*Y_{\alpha(i)},
(g_i \times g_{i'})^*\varphi_{\alpha(i)\alpha(i')}
\right)
$$
（记号如 Remark \ref{descent-remark-easier-family}）
% P05-ERR-0051
是相对于 $\mathcal{V}$ 的下降数据。
\item 此构造定义了相对于 $\mathcal{U}$ 的下降数据与相对于
$\mathcal{V}$ 的下降数据之间的函子。
\item 再给定 $\alpha' : I \to J$、$h' : S' \to S$ 及具有固定目标的
映射族之间的态射 $g'_i : U_i \to V_{\alpha'(i)}$；若 $h = h'$，则所得到的
两个下降数据函子典范同构。
\item 经由 Lemma \ref{descent-lemma-family-is-one}，这些函子与 Lemma
\ref{descent-lemma-pullback} 中构造的拉回函子一致。
\end{enumerate}
\end{lemma}

\begin{proof}
这由 Lemma \ref{descent-lemma-pullback} 及 Lemma
\ref{descent-lemma-family-is-one} 的对应关系推出。
\end{proof}

\begin{definition}
\label{descent-definition-pullback-functor-family}
设 $\mathcal{U} = \{U_i \to S'\}_{i \in I}$、
$\mathcal{V} = \{V_j \to S\}_{j \in J}$、$\alpha : I \to J$、
$h : S' \to S$ 和 $g_i : U_i \to V_{\alpha(i)}$ 如 Lemma
\ref{descent-lemma-pullback-family} 中所述。该引理构造的函子
$$
(Y_j, \varphi_{jj'}) \longmapsto
(g_i^*Y_{\alpha(i)}, (g_i \times g_{i'})^*\varphi_{\alpha(i)\alpha(i')})
$$
称为下降数据上的{\it 拉回函子}。
\end{definition}

\noindent
若 $\mathcal{U}$ 和 $\mathcal{V}$ 具有同一目标 $S$，且 $\mathcal{U}$ 加细
$\mathcal{V}$（见 Sites, Definition
\ref{sites-definition-morphism-coverings}），但没有给出明确的偶
$(\alpha, g_i)$，我们仍可谈论拉回函子，因为 Lemma
\ref{descent-lemma-pullback-family} 已表明该偶的选择无关紧要（相差典范同构）。


\begin{definition}
\label{descent-definition-effective}
设 $S$ 为概形，设 $f : X \to S$ 为概形态射。
\begin{enumerate}
\item 给定 $S$ 上的概形 $U$，有 $U$ 相对于 $\text{id} : S \to S$ 的
{\it 平凡下降数据}，即 $U$ 上的恒等态射。
\item 由 Lemma \ref{descent-lemma-pullback}，沿 $f$ 拉回平凡下降数据，可在
$X \times_S U$ 上得到一个相对于 $X \to S$ 的{\it 典范下降数据}。常把
这个下降数据记为 $(X \times_S U, can)$。
\item 若相对于 $X/S$ 的下降数据 $(V, \varphi)$ 与某个 $S$ 上概形 $U$
所对应的典范下降数据 $(X \times_S U, can)$ 同构，则称 $(V, \varphi)$
{\it 有效}。
\end{enumerate}
\end{definition}

\noindent
因此，有效意味着存在 $S$ 上的概形 $U$ 及 $X$-概形同构
$\psi : V \to X \times_S U$，使得 $\varphi$ 等于复合
$$
V \times_S X \xrightarrow{\psi \times \text{id}_X}
X \times_S U \times_S X =
X \times_S X \times_S U
\xrightarrow{\text{id}_X \times \psi^{-1}}
X \times_S V
$$

\begin{definition}
\label{descent-definition-effective-family}
设 $S$ 为概形，设 $\{X_i \to S\}$ 为以 $S$ 为目标的一族态射。
\begin{enumerate}
\item 给定 $S$ 上的概形 $U$，拉回 $U$ 相对于
$\{\text{id} : S \to S\}$ 的平凡下降数据，即在概形族
$X_i \times_S U$ 上得到一个{\it 典范下降数据}。将其记为
$(X_i \times_S U, can)$。
\item 若存在 $S$ 上的概形 $U$，使相对于 $\{X_i \to S\}$ 的下降数据
$(V_i, \varphi_{ij})$ 与 $(X_i \times_S U, can)$ 同构，则称
$(V_i, \varphi_{ij})$ 为{\it 有效}。
\end{enumerate}
\end{definition}












\section{拉回函子的全忠实性}
\label{descent-section-fully-faithful}

\noindent
事实证明，fpqc 覆盖的下降数据之间的拉回函子是全忠实的。换言之，概形
态射满足 fpqc 下降。本节的目标是证明这一点；不过我们鼓励读者自行证明。
关键是使用 Lemma \ref{descent-lemma-fpqc-universal-effective-epimorphisms}。

\begin{lemma}
\label{descent-lemma-surjective-flat-epi}
满且平坦的态射在概形范畴中是上态射。
\end{lemma}

\begin{proof}
假设 $h : X' \to X$ 满且平坦，并有态射 $a, b : X \to Y$，满足
$a \circ h = b \circ h$。由于 $h$ 满，$a$ 与 $b$ 在底层拓扑空间上相同。
选取 $x' \in X'$，置 $x = h(x')$ 和 $y = a(x) = b(x)$。考虑局部环映射
$$
a^\sharp_x, b^\sharp_x : \mathcal{O}_{Y, y} \to \mathcal{O}_{X, x}
$$
。它们与平坦局部同态
$h^\sharp_{x'} : \mathcal{O}_{X, x} \to \mathcal{O}_{X', x'}$ 复合后相等。
平坦局部同态忠实平坦（Algebra, Lemma
\ref{algebra-lemma-local-flat-ff}），所以 $h^\sharp_{x'}$ 单射。因而
$a^\sharp_x = b^\sharp_x$，这蕴含 $a = b$，正如所需。
\end{proof}

\begin{lemma}
\label{descent-lemma-ff-base-change-faithful}
设 $h : S' \to S$ 为满平坦概形态射。基变换函子
$$
\Sch/S \longrightarrow \Sch/S', \quad
X \longmapsto S' \times_S X
$$
是忠实的。
\end{lemma}

\begin{proof}
设 $X_1$、$X_2$ 为 $S$ 上的概形，设
$\alpha, \beta : X_2 \to X_1$ 为 $S$ 上的态射。若 $\alpha$、$\beta$
基变换为同一态射，则得到如下交换图：
$$
\xymatrix{
X_2 \ar[d]^\alpha &
S' \times_S X_2 \ar[l] \ar[d] \ar[r] &
X_2 \ar[d]^\beta \\
X_1 &
S' \times_S X_1 \ar[l] \ar[r] &
X_1
}
$$
所以只需证明 $S' \times_S X_2 \to X_2$ 为上态射。作为满平坦态射的
基变换，它满且平坦（见 Morphisms, Lemmas
\ref{morphisms-lemma-base-change-surjective} 和
\ref{morphisms-lemma-base-change-flat}）。因此本引理由 Lemma
\ref{descent-lemma-surjective-flat-epi} 推出。
\end{proof}

\begin{lemma}
\label{descent-lemma-faithful}
在 Lemma \ref{descent-lemma-pullback} 的情形下，假设 $f : X' \to X$ 满且平坦。
则拉回函子忠实。
\end{lemma}

\begin{proof}
设 $(V_i, \varphi_i)$，$i = 1, 2$，为 $X \to S$ 的下降数据。设
$\alpha, \beta : V_1 \to V_2$ 为下降数据的态射。假设
$f^*\alpha = f^*\beta$。我们的任务是证明 $\alpha = \beta$。注意，
$\alpha$、$\beta$ 是 $X$ 上概形的态射，而 $f^*\alpha$、$f^*\beta$ 只是
$\alpha$、$\beta$ 到 $X'$ 上态射的基变换。因此本引理由 Lemma
\ref{descent-lemma-ff-base-change-faithful} 推出。
\end{proof}

\noindent
下面是本节的关键引理。

\begin{lemma}
\label{descent-lemma-fully-faithful}
在 Lemma \ref{descent-lemma-pullback} 的情形下，假设
\begin{enumerate}
\item $\{f : X' \to X\}$ 是 fpqc 覆盖（例如，若 $f$ 满、平坦且拟紧）；
\item $S = S'$.
\end{enumerate}
则拉回函子全忠实。
\end{lemma}

\begin{proof}
假设 (1) 蕴含 $f$ 满且平坦。因此由 Lemma \ref{descent-lemma-faithful}，拉回函子
忠实。设 $(V, \varphi)$ 和 $(W, \psi)$ 为相对于 $X \to S$ 的两个下降数据。
置 $(V', \varphi') = f^*(V, \varphi)$ 及
$(W', \psi') = f^*(W, \psi)$。设 $\alpha' : V' \to W'$ 为 $S$ 上 $X'$
的下降数据态射。我们必须证明存在 $S$ 上 $X$ 的下降数据态射
$\alpha : V \to W$，其拉回为 $\alpha'$。

\medskip\noindent
回忆 $V'$ 是 $V$ 沿 $f$ 的基变换，而 $\varphi'$ 是 $\varphi$ 沿
$f \times f$ 的基变换（见 Remark \ref{descent-remark-easier}）。由假设，图
$$
\xymatrix{
V' \times_S X' \ar[r]_{\varphi'} \ar[d]_{\alpha' \times \text{id}} &
X' \times_S V' \ar[d]^{\text{id} \times \alpha'} \\
W' \times_S X' \ar[r]^{\psi'} &
X' \times_S W'
}
$$
交换。我们断言两个复合
$$
\xymatrix{
V' \times_V V' \ar[r]^-{\text{pr}_i} &
V' \ar[r]^{\alpha'} &
W' \ar[r] &
W
}
, \quad i = 0, 1
$$
相等。建议读者自行证明，而不要阅读本段余下部分。（若你找到漂亮简洁的
论证，请发邮件。）设 $v_0, v_1$ 为 $V'$ 中映到同一点 $v \in V$ 的点。
设 $x_i \in X'$ 为 $v_i$ 的像，并设 $x$ 为 $X$ 中的点，即 $v$ 在 $X$ 中的像。换言之，
在 $V' = X' \times_X V$ 中，$v_i = (x_i, v)$。对 $V$ 的某个点 $v'$，写
$\varphi(v, x) = (x, v')$。这是可能的，因为 $\varphi$ 是
$X \times_S X$ 上的态射。记 $v_i' = (x_i, v')$，它是 $V'$ 的点。计算
（使用 $\varphi'$ 的定义）表明 $\varphi'(v_i, x_j) = (x_i, v'_j)$。记
$w_i = \alpha'(v_i)$ 及 $w'_i = \alpha'(v_i')$。现在可以把
$w_i = (x_i, u_i)$ 写成这种形式，其中 $u_i$ 是 $W$ 的某个点；类似地，
$w_i' = (x_i, u'_i)$，其中 $u_i'$ 是 $W$ 的某个点。该断言等价于
$u_0 = u_1$。使用 $\psi'$ 的定义作形式计算（见 Lemma
\ref{descent-lemma-pullback}），可知上述图的交换性意味着
$$
((x_i, x_j), \psi(u_i, x)) = ((x_i, x_j), (x, u'_j))
$$
（作为
$(X' \times_S X') \times_{X \times_S X} (X \times_S W)$
的点）对所有 $i, j \in \{0, 1\}$ 成立。这表明
$\psi(u_0, x) = \psi(u_1, x)$，取 $\psi^{-1}$ 即得 $u_0 = u_1$。上述论证
是形式的，并且可以取概形点（换言之，可以取
$(v_0, v_1) = \text{id}_{V' \times_V V'}$），故这证明了断言。

\medskip\noindent
此时可以使用 Lemma \ref{descent-lemma-fpqc-universal-effective-epimorphisms}。
事实上，$\{V' \to V\}$ 是态射 $f : X' \to X$ 的基变换，因而为 fpqc
覆盖。所以由 Lemma \ref{descent-lemma-fpqc-universal-effective-epimorphisms}，态射
$\alpha' : V' \to W' \to W$ 唯一地经过一个态射 $\alpha : V \to W$
分解，而后者的基变换必为 $\alpha'$。最后，图
$$
\xymatrix{
V \times_S X \ar[r]_{\varphi} \ar[d]_{\alpha \times \text{id}} &
X \times_S V \ar[d]^{\text{id} \times \alpha} \\
W \times_S X \ar[r]^{\psi} & X \times_S W
}
$$
交换，因为它到 $X' \times_S X'$ 的基变换交换，并且态射
$X' \times_S X' \to X \times_S X$ 满且平坦（使用 Lemma
\ref{descent-lemma-ff-base-change-faithful}）。所以 $\alpha$ 是下降数据态射
$(V, \varphi) \to (W, \psi)$，正如所需。
\end{proof}

\noindent
前面材料经过改进后的叙述已经使下面两个引理过时。但它们仍然正确！

\begin{lemma}
\label{descent-lemma-pullback-selfmap}
设 $X \to S$ 为概形态射，设 $f : X \to X$ 为 $X$ 在 $S$ 上的自映射。
此时，沿 $f$ 的拉回同构于相对于 $X \to S$ 的下降数据范畴上的恒等函子。
\end{lemma}

\begin{proof}
这由 Lemma \ref{descent-lemma-pullback} 显然可见，因为它给出
$f^* \cong \text{id}^*$。
\end{proof}

\begin{lemma}
\label{descent-lemma-morphism-with-section-equivalence}
设 $f : X' \to X$ 为基概形 $S$ 上的概形态射。假设存在 $S$ 上的态射
$g : X \to X'$；例如，$f$ 有截面时即如此。则 Lemma
\ref{descent-lemma-pullback} 的拉回函子定义了相对于 $X/S$ 与相对于 $X'/S$ 的
下降数据范畴之间的等价。
\end{lemma}

\begin{proof}
设 $g : X \to X'$ 为 $S$ 上的态射。上述 Lemma
\ref{descent-lemma-pullback-selfmap} 表明，函子
$f^* \circ g^* = (g \circ f)^*$ 和 $g^* \circ f^* = (f \circ g)^*$
分别同构于相应的恒等函子，正如所需。
\end{proof}

\begin{lemma}
\label{descent-lemma-morphism-source-faithfully-flat}
设 $f : X \to X'$ 为基概形 $S$ 上的概形态射。假设 $X \to S$ 满且平坦。
则 Lemma \ref{descent-lemma-pullback} 的拉回函子，是从相对于 $X'/S$ 的下降数据
范畴到相对于 $X/S$ 的下降数据范畴的忠实函子。
\end{lemma}

\begin{proof}
可以把 $X \to X'$ 分解为 $X \to X \times_S X' \to X'$。第一个态射有
截面，所以由 Lemma \ref{descent-lemma-morphism-with-section-equivalence}，它诱导
下降数据范畴的等价。第二个态射满且平坦，所以由 Lemma
\ref{descent-lemma-faithful}，它诱导忠实函子。
\end{proof}

\begin{lemma}
\label{descent-lemma-morphism-source-fpqc-covering}
设 $f : X \to X'$ 为基概形 $S$ 上的概形态射。假设 $\{X \to S\}$ 为 fpqc
覆盖（例如，若 $f$ 满、平坦且拟紧）。则 Lemma \ref{descent-lemma-pullback} 的
拉回函子，是从相对于 $X'/S$ 的下降数据范畴到相对于 $X/S$ 的下降数据
范畴的全忠实函子。
\end{lemma}

\begin{proof}
可以把 $X \to X'$ 分解为 $X \to X \times_S X' \to X'$。第一个态射有
截面，所以由 Lemma \ref{descent-lemma-morphism-with-section-equivalence}，它诱导
下降数据范畴的等价。第二个态射是 fpqc 覆盖，所以由 Lemma
\ref{descent-lemma-fully-faithful}，它诱导全忠实函子。
\end{proof}

\begin{lemma}
\label{descent-lemma-fpqc-refinement-coverings-fully-faithful}
设 $S$ 为概形。设 $\mathcal{U} = \{U_i \to S\}_{i \in I}$ 和
$\mathcal{V} = \{V_j \to S\}_{j \in J}$ 为以 $S$ 为目标的态射族。设
$\alpha : I \to J$、$\text{id} : S \to S$ 和
$g_i : U_i \to V_{\alpha(i)}$ 构成具有固定目标的映射族之间的态射；见
Sites, Definition \ref{sites-definition-morphism-coverings}。假设对每个
$j \in J$，态射族 $\{g_i : U_i \to V_j\}_{\alpha(i) = j}$ 是 $V_j$ 的
fpqc 覆盖。则 Lemma \ref{descent-lemma-pullback-family} 的拉回函子
$$
\text{相对下降数据 }
\mathcal{V}
\longrightarrow
\text{相对下降数据 }
\mathcal{U}
$$
是全忠实的。
\end{lemma}

\begin{proof}
考虑 $S$ 上的概形态射
$$
g :
X = \coprod\nolimits_{i \in I} U_i
\longrightarrow
Y = \coprod\nolimits_{j \in J} V_j
$$
；它在第 $i$ 个分支上经由态射
% P05-ERR-0053
$g_{\alpha(i)}$ 映入第 $\alpha(i)$ 个分支。我们断言
$\{g : X \to Y\}$ 是概形的 fpqc 覆盖。事实上，由 Topologies, Lemma
\ref{topologies-lemma-disjoint-union-is-fpqc-covering}，对每个 $j$，态射
$\{\coprod_{\alpha(i) = j} U_i \to V_j\}$ 是 fpqc 覆盖。因此，对每个仿射
开集 $V \subset V_j$（可把它视为 $Y$ 的仿射开集），可以找到有限多个
仿射开集 $W_1, \ldots, W_n \subset \coprod_{\alpha(i) = j} U_i$（可把它们
视为 $X$ 的仿射开集），使得
$V = \bigcup_{i = 1, \ldots, n} g(W_i)$。这给出了足够多的 $Y$ 的仿射开集，
它们可由有限多个 $X$ 的仿射开集覆盖，因而 Topologies, Lemma
\ref{topologies-lemma-recognize-fpqc-covering} 的第 (3) 部分适用，断言成立。
分别以 $DD(X/S)$ 和 $DD(\mathcal{U})$ 表示相对于 $X/S$ 和相对于
$\mathcal{U}$ 的下降数据范畴；对 $Y/S$ 和 $\mathcal{V}$ 也采用类似记号。
考虑图
$$
\xymatrix{
DD(Y/S) \ar[r] & DD(X/S) \\
DD(\mathcal{V}) \ar[u]^{\text{引理 }\ref{descent-lemma-family-is-one}} \ar[r] &
DD(\mathcal{U}) \ar[u]_{\text{引理 }\ref{descent-lemma-family-is-one}}
}
$$
该图交换；见 Lemma \ref{descent-lemma-pullback-family} 的证明。竖直箭头为等价。
因此本引理由 Lemma \ref{descent-lemma-fully-faithful} 推出，后者表明图中上方水平
箭头全忠实。
\end{proof}

\noindent
下一个引理表明，为了检验有效性，总可以对给定的以 $S$ 为目标的态射族作
Zariski 加细。

\begin{lemma}
\label{descent-lemma-Zariski-refinement-coverings-equivalence}
设 $S$ 为概形。设 $\mathcal{U} = \{U_i \to S\}_{i \in I}$ 和
$\mathcal{V} = \{V_j \to S\}_{j \in J}$ 为以 $S$ 为目标的态射族。设
$\alpha : I \to J$、$\text{id} : S \to S$ 和
$g_i : U_i \to V_{\alpha(i)}$ 构成具有固定目标的映射族之间的态射；见
Sites, Definition \ref{sites-definition-morphism-coverings}。假设对每个
$j \in J$，态射族 $\{g_i : U_i \to V_j\}_{\alpha(i) = j}$ 是 $V_j$ 的
Zariski 覆盖（见 Topologies, Definition
\ref{topologies-definition-zariski-covering}）。则 Lemma
\ref{descent-lemma-pullback-family} 的拉回函子
$$
\text{相对下降数据 }
\mathcal{V}
\longrightarrow
\text{相对下降数据 }
\mathcal{U}
$$
是范畴等价。特别地，$S$ 上概形的范畴等价于相对于 $S$ 的任意 Zariski
覆盖的下降数据范畴。
\end{lemma}

\begin{proof}
由 Lemma \ref{descent-lemma-fpqc-refinement-coverings-fully-faithful}，该函子忠实且
全忠实。下面说明如何证明它本质满。设 $(X_i, \varphi_{ii'})$ 为相对于
$\mathcal{U}$ 的下降数据。固定 $j \in J$，并置
$I_j = \{i \in I \mid \alpha(i) = j\}$。对 $i, i' \in I_j$，注意存在
典范态射
$$
c_{ii'} : U_i \times_{g_i, V_j, g_{i'}} U_{i'} \to U_i \times_S U_{i'}.
$$
。因此可沿此态射拉回 $\varphi_{ii'}$；对 $i, i' \in I_j$，置
$\psi_{ii'} = c_{ii'}^*\varphi_{ii'}$。这样得到相对于 Zariski 覆盖
$\{g_i : U_i \to V_j\}_{i \in I_j}$ 的下降数据
$(X_i, \psi_{ii'})$。注意，$\psi_{ii'}$ 是从 $X_i$ 的开集
$X_{i, U_i \times_{V_j} U_{i'}}$ 到 $X_{i'}$ 的相应开集的同构。由 Schemes,
Section \ref{schemes-section-glueing-schemes}，可以把
$(X_i, \psi_{ii'})$ 黏合成 $V_j$ 上的概形 $Y_j$。此外，对 $i \in I_j$ 和
$i' \in I_{j'}$，态射 $\varphi_{ii'}$ 黏合成满足上循环条件的态射
$\varphi_{jj'} : Y_j \times_S V_{j'} \to V_j \times_S Y_{j'}$（省略细节）。
因此得到所需的相对于 $\mathcal{V}$ 的下降数据
$(Y_j, \varphi_{jj'})$。
\end{proof}

\begin{lemma}
\label{descent-lemma-refine-coverings-fully-faithful}
设 $S$ 为概形。设 $\mathcal{U} = \{U_i \to S\}_{i \in I}$ 和
$\mathcal{V} = \{V_j \to S\}_{j \in J}$ 为 $S$ 的 fpqc 覆盖。若
$\mathcal{U}$ 是 $\mathcal{V}$ 的加细，则拉回函子
$$
\text{相对下降数据 }
\mathcal{V}
\longrightarrow
\text{相对下降数据 }
\mathcal{U}
$$
全忠实。特别地，$S$ 上概形的范畴可认同为相对于 $S$ 的任意 fpqc 覆盖的
下降数据范畴的一个全子范畴。
\end{lemma}

\begin{proof}
考虑 $S$ 的 fpqc 覆盖
$\mathcal{W} = \{U_i \times_S V_j \to S\}_{(i, j) \in I \times J}$。它既是
$\mathcal{U}$ 的加细，也是 $\mathcal{V}$ 的加细。因此有函子与范畴的
$2$-交换图
$$
\xymatrix{
DD(\mathcal{V}) \ar[rd] \ar[rr] & & DD(\mathcal{U}) \ar[ld] \\
& DD(\mathcal{W}) &
}
$$
记号如 Lemma \ref{descent-lemma-fpqc-refinement-coverings-fully-faithful} 的证明；交换
性由 Lemma \ref{descent-lemma-pullback-family} 第 (3) 部分给出。因此显然只需证明
函子 $DD(\mathcal{V}) \to DD(\mathcal{W})$ 和
$DD(\mathcal{U}) \to DD(\mathcal{W})$ 全忠实。这由 Lemma
\ref{descent-lemma-fpqc-refinement-coverings-fully-faithful} 推出，正如所需。
\end{proof}

\begin{remark}
\label{descent-remark-morphisms-of-schemes-satisfy-fpqc-descent}
Lemma \ref{descent-lemma-refine-coverings-fully-faithful} 断言概形态射满足 fpqc
下降。换言之，给定概形 $S$ 及 $S$ 上的概形 $X$、$Y$，函子
$$
(\Sch/S)^{opp} \longrightarrow \textit{Sets},
\quad
T \longmapsto \Mor_T(X_T, Y_T)
$$
满足 fpqc 拓扑的层条件。最简单的情形如下。假设 $T \to S$ 是仿射概形
之间的满平坦态射。设 $\psi_0 : X_T \to Y_T$ 为 $T$ 上与典范下降数据相容
的概形态射。则存在唯一的态射 $\psi : X \to Y$，其到 $T$ 的基变换为
$\psi_0$。事实上，这个特殊情形直接由 Lemma
\ref{descent-lemma-fully-faithful} 推出。而该引理又是下列两个事实的形式推论：
(a) 沿忠实平坦态射的基变换函子忠实，见 Lemma
\ref{descent-lemma-ff-base-change-faithful}；(b) 概形满足 fpqc 拓扑的层条件，见
Lemma \ref{descent-lemma-fpqc-universal-effective-epimorphisms}。
\end{remark}

\begin{lemma}
\label{descent-lemma-effective-for-fpqc-is-local-upstairs}
设 $X \to S$ 为满、拟紧、平坦的概形态射。设 $(V, \varphi)$ 为相对于
$X/S$ 的下降数据。假设对所有 $v \in V$，都存在开子概形
$v \in W \subset V$，使得 $\varphi(W \times_S X) \subset X \times_S W$，
且下降数据 $(W, \varphi|_{W \times_S X})$ 有效。则 $(V, \varphi)$ 有效。
\end{lemma}

\begin{proof}
设 $V = \bigcup W_i$ 为开覆盖，满足
$\varphi(W_i \times_S X) \subset X \times_S W_i$，且下降数据
$(W_i, \varphi|_{W_i \times_S X})$ 有效。设 $U_i \to S$ 为概形，并设
$\alpha_i : (X \times_S U_i, can) \to (W_i, \varphi|_{W_i \times_S X})$
为下降数据同构。对每对指标 $(i, j)$，考虑开集
$\alpha_i^{-1}(W_i \cap W_j) \subset X \times_S U_i$。由于所有内容都与
下降数据相容，且 $\{X \to S\}$ 为 fpqc 覆盖，可以应用 Lemma
\ref{descent-lemma-open-fpqc-covering}，找到开集 $U_{ij} \subset U_i$，使得
$\alpha_i^{-1}(W_i \cap W_j) = X \times_S U_{ij}$。现在，
$W_i \cap W_j$ 上的恒等态射与下降数据相容，因而来自唯一的 $S$ 上态射
$\varphi_{ij} : U_{ij} \to U_{ji}$（见 Remark
\ref{descent-remark-morphisms-of-schemes-satisfy-fpqc-descent}）。于是
$(U_i, U_{ij}, \varphi_{ij})$ 是 Schemes, Section
\ref{schemes-section-glueing-schemes} 中的黏合数据（省略证明）。因此可以
假设存在 $S$ 上的概形 $U$，使得 $U_i \subset U$ 为开集、
$U_{ij} = U_i \cap U_j$ 且
$\varphi_{ij} = \text{id}_{U_i \cap U_j}$；见 Schemes, Lemma
\ref{schemes-lemma-glue}。拉回到 $X$ 后，可以利用 $\alpha_i$ 得到所需
同构 $\alpha : X \times_S U \to V$。
\end{proof}







\section{态射类型的下降}
\label{descent-section-descending-types-morphisms}

\noindent
下面研究概形相对于 fpqc 覆盖的下降数据是否有效。首先要指出，这并非总是
成立；我们将在 Algebraic Spaces, Example
\ref{spaces-example-non-representable-descent} 中看到这一点。甚至投射态射也
不总是对 fpqc 覆盖满足下降，见 Examples, Lemma
\ref{examples-lemma-non-effective-descent-projective}。

\medskip\noindent
另一方面，如果试图下降的概形特别简单，那么对整个一类概形，下降数据有时
确实有效。这里引入抽象描述这一现象的术语，尽管以后若使用不当，它可能造成
混淆。

\begin{definition}
\label{descent-definition-descending-types-morphisms}
设 $\mathcal{P}$ 为基底上概形态射的一条性质，并设
$\tau \in \{Zariski, fpqc, fppf, \etale, smooth, syntomic\}$。如果对任意
$\tau$-覆盖 $\mathcal{U} : \{U_i \to S\}_{i \in I}$（见 Topologies,
Section \ref{topologies-section-procedure}），每个相对于 $\mathcal{U}$ 的
下降数据 $(X_i, \varphi_{ij})$，只要每个态射 $X_i \to U_i$ 具有性质
$\mathcal{P}$，就都是有效的，则称{\it $\mathcal{P}$ 型态射对
$\tau$-覆盖满足下降}。
\end{definition}

\noindent
注意，在每种情形下我们已经看到，从 $S$ 上概形到 $\mathcal{U}$ 上下降
数据的函子都是全忠实的（由 Lemma
\ref{descent-lemma-refine-coverings-fully-faithful} 以及 Topologies 中任意
$\tau$-覆盖也是 fpqc 覆盖的结果可知）。我们还看到，相对于 Zariski 覆盖，
下降数据总是有效的（Lemma
\ref{descent-lemma-Zariski-refinement-coverings-equivalence}）。为避免错误论证
（即在下降进行到一半时依赖 $\mathcal{P}$），谨慎的做法是：仅当
$\mathcal{P}$ 在任意基变换下稳定，或至少在 $\tau$-拓扑中对基底为局部性质
时，才研究刚引入的概念（见 Definition
\ref{descent-definition-property-morphisms-local}）。

\medskip\noindent
下面这个必备引理把问题约化到由仿射概形之间单个态射给出的覆盖。

\begin{lemma}
\label{descent-lemma-descending-types-morphisms}
设 $\mathcal{P}$ 为基底上概形态射的一条性质，并设
$\tau \in \{fpqc, fppf, \etale, smooth, syntomic\}$。假设：
\begin{enumerate}
\item $\mathcal{P}$ 在任意基变换下稳定（见 Schemes, Definition
\ref{schemes-definition-preserved-by-base-change}）；
\item 如果 $Y_j \to V_j$（$j = 1, \ldots, m$）具有 $\mathcal{P}$，则
$\coprod Y_j \to \coprod V_j$ 也具有该性质；并且
\item 对任意仿射概形之间的满射态射 $X \to S$，按照 $\tau$ 分别为 fpqc、
fppf、\'etale、smooth 或 syntomic，假设该态射分别是平坦、有限表现平坦、
\'etale、光滑或 syntomic 的；则相对于 $S$ 上 $X$ 的任意下降数据
$(V, \varphi)$，只要 $V \to X$ 具有 $\mathcal{P}$，就都是有效的。
\end{enumerate}
那么 $\mathcal{P}$ 型态射对 $\tau$-覆盖满足下降。
\end{lemma}

\begin{proof}
设 $S$ 为概形，设
$\mathcal{U} = \{\varphi_i : U_i \to S\}_{i \in I}$ 为 $S$ 的
$\tau$-覆盖。设 $(X_i, \varphi_{ii'})$ 为相对于 $\mathcal{U}$ 的下降
数据，并假设每个态射 $X_i \to U_i$ 具有性质 $\mathcal{P}$。我们必须证明
存在概形 $X \to S$，使得
$(X_i, \varphi_{ii'}) \cong (U_i \times_S X, can)$。

\medskip\noindent
在正式开始证明前先指出：对任意态射族
$\mathcal{V} : \{V_j \to S\}$ 及任意族态射
$\mathcal{V} \to \mathcal{U}$，如果把下降数据
$(X_i, \varphi_{ii'})$ 拉回成 $\mathcal{V}$ 上的下降数据
$(Y_j, \varphi_{jj'})$，那么每个态射 $Y_j \to V_j$ 也具有性质
$\mathcal{P}$。这是因为假设 (1) 说 $\mathcal{P}$ 在任意基变换下稳定，
并结合拉回的定义（见 Definition
\ref{descent-definition-pullback-functor-family}）。下文将不再说明地使用这一点。

\medskip\noindent
先在 $S$ 仿射时证明本引理。由 Topologies, Lemma
\ref{topologies-lemma-fpqc-affine},
\ref{topologies-lemma-fppf-affine},
\ref{topologies-lemma-etale-affine},
\ref{topologies-lemma-smooth-affine}, or
\ref{topologies-lemma-syntomic-affine}
，存在加细 $\mathcal{U}$ 的标准 $\tau$-覆盖
$\mathcal{V} : \{V_j \to S\}_{j = 1, \ldots, m}$。下降数据范畴之间的拉回
函子
$DD(\mathcal{U}) \to DD(\mathcal{V})$
由 Lemma \ref{descent-lemma-refine-coverings-fully-faithful} 可知是全忠实的。因此只需
证明标准 $\tau$-覆盖 $\mathcal{V}$ 上的下降数据有效。由假设 (2)，
$\coprod Y_j \to \coprod V_j$ 具有性质 $\mathcal{P}$。由 Lemma
\ref{descent-lemma-family-is-one}，问题约化到覆盖
$\{\coprod_{j = 1, \ldots, m} V_j \to S\}$；对此，引理的假设 (3) 已经给出
所需结果。因此，当 $S$ 仿射时引理成立。

\medskip\noindent
现设 $S$ 为一般概形，并设 $V \subset S$ 为仿射开集。由位点的性质，族
$\mathcal{U}_V = \{V \times_S U_i \to V\}_{i \in I}$ 是 $V$ 的
$\tau$-覆盖。用 $(X_i, \varphi_{ii'})_V$ 表示给定下降数据在
$\mathcal{U}_V$ 上的限制（或拉回）。由刚才的结论，得到 $V$ 上的概形
$X_V$，其相对于 $\mathcal{U}_V$ 的典范下降数据与
$(X_i, \varphi_{ii'})_V$ 同构。假设 $V' \subset V$ 是 $V$ 的仿射开集。
那么 $X_{V'}$ 与 $V' \times_V X_V$ 的典范下降数据都同构于
$(X_i, \varphi_{ii'})_{V'}$。再次由 Lemma
\ref{descent-lemma-refine-coverings-fully-faithful}，得到 $S$ 上的典范态射
$\rho^V_{V'} : X_{V'} \to X_V$，它把 $X_{V'}$ 等同于 $V'$ 在 $X_V$
中的逆像。略去如下验证：对 $S$ 的仿射开集
$V'' \subset V' \subset V$，有
$\rho^V_{V''} = \rho^V_{V'} \circ \rho^{V'}_{V''}$。

\medskip\noindent
由 Constructions, Lemma \ref{constructions-lemma-relative-glueing}，数据
$(X_V, \rho^V_{V'})$ 拼接成概形 $X \to S$。此外，我们有恢复态射
$\rho^V_{V'}$ 的同构 $V \times_S X \to X_V$。展开概形 $X_V$ 的构造，
得到同构
$$
V \times_S U_i \times_S X
\longrightarrow
V \times_S X_i
$$
，它们与态射 $\varphi_{ii'}$ 相容，也与限制到 $X$ 中更小的仿射开集相容。
这说明 $U_i \times_S X$ 上的典范下降数据同构于给定下降数据，证毕。
\end{proof}










\section{仿射态射的下降}
\label{descent-section-affine}

\noindent
本节证明“仿射态射对 fpqc 覆盖满足下降”。正式陈述如下。

\begin{lemma}
\label{descent-lemma-affine}
设 $S$ 为概形。设 $\{X_i \to S\}_{i\in I}$ 为 fpqc 覆盖，见
Topologies, Definition \ref{topologies-definition-fpqc-covering}。设
$(V_i/X_i, \varphi_{ij})$ 为相对于 $\{X_i \to S\}$ 的下降数据。如果每个
态射 $V_i \to X_i$ 都是仿射的，那么该下降数据有效。
\end{lemma}

\begin{proof}
仿射性是概形态射的一条性质，它在基底上为局部性质并在任意基变换下保持，
见 Morphisms, Lemmas \ref{morphisms-lemma-characterize-affine} 和
\ref{morphisms-lemma-base-change-affine}。因此可应用 Lemma
\ref{descent-lemma-descending-types-morphisms}，只需在 fpqc 覆盖由仿射概形之间的
单个平坦满射 $\{X \to S\}$ 给出时证明本引理。记
$X = \Spec(A)$ 及 $S = \Spec(R)$，于是 $R \to A$ 是忠实平坦环映射。设
$(V, \varphi)$ 为相对于 $S$ 上 $X$ 的下降数据，并假设 $V \to X$ 仿射。
由于 $V \to X$ 仿射，有 $V = \Spec(B)$，其中 $B$ 是某个 $A$-代数
（见 Morphisms, Definition \ref{morphisms-definition-affine}）。同构
$\varphi$ 对应于环同构
$$
\varphi^\sharp :
B \otimes_R A \longleftarrow A \otimes_R B
$$
，它是 $A \otimes_R A$-代数同构。$\varphi$ 的余圈条件说明
$$
\xymatrix{
B \otimes_R A \otimes_R A & &
A \otimes_R A \otimes_R B \ar[ll] \ar[ld]\\
& A \otimes_R B \otimes_R A \ar[lu] &
}
$$
交换。将这些箭头反向，便得到 Definition
\ref{descent-definition-descent-datum-modules} 意义下相对于 $R \to A$ 的模下降
数据。因此可应用 Proposition \ref{descent-proposition-descent-module}，得到
$R$-模
$C = \Ker(B \to A \otimes_R B)$
以及保持下降数据的同构 $A \otimes_R C \cong B$。对任意
$c, c' \in C$，因为态射 $\varphi$ 是代数同态，$B$ 中的乘积 $cc'$ 属于
$C$。因此 $C$ 是 $R$-代数，其到 $A$ 的基变换以与下降数据相容的方式
同构于 $B$。应用 $\Spec$，得到 $S$ 上的概形 $U$，使
$(V, \varphi) \cong (X \times_S U, can)$，正合所需。
\end{proof}

\begin{lemma}
\label{descent-lemma-closed-immersion}
设 $S$ 为概形。设 $\{X_i \to S\}_{i\in I}$ 为 fpqc 覆盖，见
Topologies, Definition \ref{topologies-definition-fpqc-covering}。设
$(V_i/X_i, \varphi_{ij})$ 为相对于 $\{X_i \to S\}$ 的下降数据。如果每个
态射 $V_i \to X_i$ 都是闭浸入，那么该下降数据有效。
\end{lemma}

\begin{proof}
这是因为闭浸入是仿射态射（Morphisms, Lemma
\ref{morphisms-lemma-closed-immersion-affine}），因而可应用 Lemma
\ref{descent-lemma-affine}。
\end{proof}


\section{拟仿射态射的下降}
\label{descent-section-quasi-affine}

\noindent
本节证明“拟仿射态射对 fpqc 覆盖满足下降”。正式陈述如下。

\begin{lemma}
\label{descent-lemma-quasi-affine}
设 $S$ 为概形。设 $\{X_i \to S\}_{i\in I}$ 为 fpqc 覆盖，见
Topologies, Definition \ref{topologies-definition-fpqc-covering}。设
$(V_i/X_i, \varphi_{ij})$ 为相对于 $\{X_i \to S\}$ 的下降数据。如果每个
态射 $V_i \to X_i$ 都是拟仿射的，那么该下降数据有效。
\end{lemma}

\begin{proof}
拟仿射性是概形态射的一条性质，它在任意基变换下保持，见 Morphisms,
Lemmas \ref{morphisms-lemma-characterize-quasi-affine} 和
\ref{morphisms-lemma-base-change-quasi-affine}。因此可应用 Lemma
\ref{descent-lemma-descending-types-morphisms}，只需在 fpqc 覆盖由仿射概形之间的
单个平坦满射 $\{X \to S\}$ 给出时证明本引理。记 $X = \Spec(A)$ 及
$S = \Spec(R)$，于是 $R \to A$ 是忠实平坦环映射。设 $(V, \varphi)$ 为
相对于 $S$ 上 $X$ 的下降数据，并假设 $\pi : V \to X$ 拟仿射。

\medskip\noindent
由 Morphisms, Lemma \ref{morphisms-lemma-characterize-quasi-affine}，这表示
$$
V \longrightarrow \underline{\Spec}_X(\pi_*\mathcal{O}_V) = W
$$
是 $X$ 上概形的拟紧开浸入。投影
$\text{pr}_i : X \times_S X \to X$ 是平坦的，因而由平坦基变换
（Cohomology of Schemes, Lemma
\ref{coherent-lemma-flat-base-change-cohomology}），有
$$
\text{pr}_0^*\pi_*\mathcal{O}_V =
(\pi \times \text{id}_X)_*\mathcal{O}_{V \times_S X}, \quad
\text{pr}_1^*\pi_*\mathcal{O}_V =
(\text{id}_X \times \pi)_*\mathcal{O}_{X \times_S V}
$$
。因此同构 $\varphi : V \times_S X \to X \times_S V$（它是
$X \times_S X$ 上的同构）诱导 $X \times_S X$ 上拟凝聚代数层的同构
$$
\varphi^\sharp :
\text{pr}_0^*\pi_*\mathcal{O}_V
\longrightarrow
\text{pr}_1^*\pi_*\mathcal{O}_V
$$
。$\varphi$ 的余圈条件蕴含 $\varphi^\sharp$ 的余圈条件。换言之，它在仿射
概形 $W$ 上产生相对于 $S$ 上 $X$ 的下降数据 $\varphi'$，并且态射
$V \to W$ 是下降数据的态射。因此由 Lemma \ref{descent-lemma-affine}（或由拟凝聚
代数下降的有效性），得到概形 $U' \to S$ 以及下降数据同构
$(W, \varphi') \cong (X \times_S U', can)$。顺便指出，由 Lemma
\ref{descent-lemma-descending-property-affine}，$U'$ 是仿射的。

\medskip\noindent
现在可把 $V$ 看成一个（拟紧）开集 $V \subset X \times_S U'$，它在下降
数据
$$
can : X \times_S U' \times_S X \to X \times_S X \times_S U',
(x_0, u', x_1) \mapsto (x_0, x_1, u').
$$
下稳定。换言之，对映到 $S$ 中同一点的任意 $x_0, x_1, u'$，有
$(x_0, u') \in V \Rightarrow (x_1, u') \in V$。因为 $X \to S$ 满射，
立即可知 $V$ 是某个子集 $U \subset U'$ 在态射
$X \times_S U' \to U'$ 下的逆像。因为 $X \to S$ 拟紧、平坦且满射，
$X \times_S U' \to U'$ 也拟紧、平坦且满射。因此由 Morphisms, Lemma
\ref{morphisms-lemma-fpqc-quotient-topology}，子集 $U \subset U'$ 是开集，
证毕。
\end{proof}












\section{以层表述下降数据}
\label{descent-section-descent-data-sheaves}


\noindent
对位点上的覆盖，下面给出理解下降数据的另一种方式。

\begin{lemma}
\label{descent-lemma-descent-data-sheaves}
设 $\tau \in \{Zariski, fppf, \etale, smooth, syntomic\}$\footnote{这里没有
fpqc 并非笔误。见 Topologies, Section
\ref{topologies-section-fpqc} 中的讨论。}。设 $\Sch_\tau$ 为大
$\tau$-位点，设 $S \in \Ob(\Sch_\tau)$。设
$\{S_i \to S\}_{i \in I}$ 为位点 $(\Sch/S)_\tau$ 中的覆盖。存在范畴等价
$$
\left\{
\begin{matrix}
\text{下降数据 }(X_i, \varphi_{ii'})\text{ 使得}\\
\text{每个 }X_i \in \Ob((\Sch/S)_\tau)
\end{matrix}
\right\}
\leftrightarrow
\left\{
\begin{matrix}
\text{层 }F\text{ 定义于 }(\Sch/S)_\tau\text{ 使得}\\
\text{每个 }h_{S_i} \times F\text{ 可表}
\end{matrix}
\right\}.
$$
此外：
\begin{enumerate}
\item 右侧表示 $h_{S_i} \times F$ 的对象对应于左侧的概形 $X_i$；并且
\item 层 $F$ 可表示，当且仅当相应的下降数据
$(X_i, \varphi_{ii'})$ 有效。
\end{enumerate}
\end{lemma}

\begin{proof}
我们已在 Section \ref{descent-section-fpqc-universal-effective-epimorphisms} 中看到，
可表示预层是位点 $(\Sch/S)_\tau$ 上的层。此外，米田引理（Categories,
Lemma \ref{categories-lemma-yoneda}）保证可表示层之间的映射与表示对象
之间的映射一一对应。在证明中将不再说明地使用这些事实。

\medskip\noindent
先构造从右到左的函子。设 $F$ 为 $(\Sch/S)_\tau$ 上的层，并且每个
$h_{S_i} \times F$ 都可表示。在这种情形下，令 $X_i$ 为
$(\Sch/S)_\tau$ 中的一个表示对象。它带有态射 $X_i \to S_i$。于是
$X_i \times_S S_{i'}$ 与 $S_i \times_S X_{i'}$ 都表示层
$h_{S_i} \times F \times h_{S_{i'}}$，因而得到同构
$$
\varphi_{ii'} : X_i \times_S S_{i'} \to S_i \times_S X_{i'}
$$
容易看出，态射 $\varphi_{ii'}$ 是 $S_i \times_S S_{i'}$ 上的态射并满足
余圈条件。从右到左的函子由构造
$F \mapsto (X_i, \varphi_{ii'})$ 给出。

\medskip\noindent
再构造从左到右的函子。对每个 $i$，用 $F_i$ 表示层 $h_{X_i}$。同构
$\varphi_{ii'}$ 给出同构
$$
\varphi_{ii'} :
F_i \times h_{S_{i'}}
\longrightarrow
h_{S_i} \times F_{i'}
$$
，它位于 $h_{S_i} \times h_{S_{i'}}$ 上。令 $F$ 等于下图中的余等化子
$$
\xymatrix{
\coprod_{i, i'} F_i \times h_{S_{i'}}
\ar@<1ex>[rr]^-{\text{pr}_0}
\ar@<-1ex>[rr]_-{\text{pr}_1 \circ \varphi_{ii'}}
& &
\coprod_i F_i \ar[r]
&
F
}
$$
余圈条件保证 $h_{S_i} \times F$ 同构于 $F_i$，因而可表示。从左到右的
函子由构造 $(X_i, \varphi_{ii'}) \mapsto F$ 给出。

\medskip\noindent
略去这些构造互为拟逆函子的验证。最后的陈述 (1) 与 (2) 由这些构造推出。
\end{proof}

\begin{remark}
\label{descent-remark-what-product-means}
在 Lemma \ref{descent-lemma-descent-data-sheaves} 的陈述中，
$h_{S_i} \times F$ 可表示这一条件，等价于 $F$ 在
$(\Sch/S_i)_\tau$ 上的限制可表示。
\end{remark}
